\documentclass{article}

% --------------------------------------------------------------------------
% Language definition
% --------------------------------------------------------------------------

\usepackage[utf8]{inputenc}
%\usepackage[czech]{babel}

% --------------------------------------------------------------------------
% Package definition and folder with images
% --------------------------------------------------------------------------

\usepackage{graphicx}
\graphicspath{ {./Obrazky/} }

% --------------------------------------------------------------------------
% Package difinition
% --------------------------------------------------------------------------

\usepackage{geometry}
 \geometry{
 a4paper,
 total={170mm,257mm},
 left=20mm,
 top=20mm,
 }
\usepackage{fancyhdr}
\usepackage{xcolor}
\usepackage{titling}
\usepackage{array}
\usepackage{todonotes}
\usepackage{tabularx} % Pro vytvoření tabulky, která bude mít full page width, musím zkontrolovat, zda tento package někde v práci něco nerozhodil
\usepackage{enumitem}
\usepackage{pdfpages}

% --------------------------------------------------------------------------
% Base header definition
% --------------------------------------------------------------------------

\setlength\headheight{26pt}
\lhead{\includegraphics[width=3cm,height=\dimexpr \headheight-\dp\strutbox]{newcevro}}
\rhead{\small{\leftmark}}

% --------------------------------------------------------------------------
% Bibliography definition
% --------------------------------------------------------------------------

\usepackage[backend=biber,defernumbers=true]{biblatex}
\addbibresource{Zdroje.bib}

% --------------------------------------------------------------------------
% Přidán backend biber pro sorting, nakonec jsem se ale rozhodl skrýt čísla v bibliografii pomocí defernumbers a omitnumbers
% Deklarace, že před bibliografií nic nebude, od kraje nebude žádný prostor a mezi bibligorafickými záznami bude větší mezera
% --------------------------------------------------------------------------

\DeclareFieldFormat{labelnumberwidth}{}
\setlength{\biblabelsep}{0pt}
\setlength\bibitemsep{1.5\itemsep}

% --------------------------------------------------------------------------
% Footnote package definition
% --------------------------------------------------------------------------

\usepackage{footnote} %Package použitý k tomu, aby byli citace uvnitř float elementů pod čarou
\makesavenoteenv{figure} %Nastavení toho, aby byli citace uvnitř figure pod čarou

% --------------------------------------------------------------------------
% Package definition
% --------------------------------------------------------------------------

\usepackage{xcolor}

% --------------------------------------------------------------------------
% Different language for contents and figures
% --------------------------------------------------------------------------

\renewcommand\contentsname{Obsah}
\renewcommand\listfigurename{Seznam obrázků}
\renewcommand\figurename{Obrázek}

% --------------------------------------------------------------------------
% For printing roman numerals
% --------------------------------------------------------------------------

\makeatletter
\newcommand*{\rom}[1]{\expandafter\@slowromancap\romannumeral #1@}
\makeatother

% --------------------------------------------------------------------------
% For hiding part of the work if needed
% Usage
% \ifcomment
% Part I want to hide
% \fi
% Then
% |           Command          | Result |
% | -------------------------- | ------ |
% | \commenttrue               | Shown  |
% | <percentsight>\commenttrue | Hidden |
% --------------------------------------------------------------------------

\newif\ifcomment
%\commenttrue

% --------------------------------------------------------------------------
% Special header style definition for sections with long text in header
% This is however useful only if page width would be too small to fit the whole section name on one line
% --------------------------------------------------------------------------

\fancypagestyle{smallertextinheader}{ % Dokončit styl s menším textem v hlavičce, již je dokončeno
   \fancyhf{}
   \fancyhead[L]{\includegraphics[width=3cm, height=\dimexpr \headheight-\dp\strutbox]{newcevro}}
   \fancyhead[R]{%
   \parbox[b]{\dimexpr \textwidth-3cm-\columnsep}%
   {\small\uppercase\leftmark}}%
   \fancyfoot[C]{\thepage}
}

% --------------------------------------------------------------------------
% Special header for contents only
% --------------------------------------------------------------------------

\fancypagestyle{Contents}{ % Dokončit styl s menším textem v hlavičce, již je dokončeno
   \fancyhf{}
   \fancyhead[LE,LO]{\includegraphics[width=3cm, height=\dimexpr \headheight-\dp\strutbox]{newcevro}}
   \fancyhead[RE,RO]{\small{\uppercase{\rightmark}}}
}

% --------------------------------------------------------------------------
% Basic page style definition
% --------------------------------------------------------------------------

\pagestyle{fancy}

% --------------------------------------------------------------------------
% Pretitle definition
% --------------------------------------------------------------------------

\pretitle{
	\begin{center}
	\LARGE
	\includegraphics[width=10cm,height=3cm,keepaspectratio]{newcevro}\\
}
\posttitle{\end{center}}

% --------------------------------------------------------------------------
% Document body --------------
% --------------------------------------------------------------------------

\begin{document}

% --------------------------------------------------------------------------
% Title definition for second page
% --------------------------------------------------------------------------

\title{Svěřenský fond jako nástroj transferu majetku\\ mezi generacemi}
\author{Dominik Bálint}
\date{01.12.2019}

% --------------------------------------------------------------------------
% First page of my bachelors work
% --------------------------------------------------------------------------

\pagenumbering{gobble}
  \thispagestyle{empty}
  \begin{center}
  \includegraphics[width=10cm,height=3cm,keepaspectratio]{newcevro} \\
  \end{center}
  \vspace{15mm}
  \begin{center}
  {\Large Vysoká škola Cevro Institut} \\
  \vspace{15mm}
  {\Large \textbf{Svěřenský fond jako mezigenerační nástroj převodu majetku}} \\
  \vspace{15mm}
  {\Large Dominik Bálint} \\
  \vspace{15mm}
  {\Large Bakalářská práce} \\
  \vspace{49mm}
  {\Large \textbf{Praha 2019}} \\
  \end{center}

% --------------------------------------------------------------------------  
% Second page information 
% --------------------------------------------------------------------------

\newpage
  \thispagestyle{empty}
  \maketitle
  \begin{center}
  {\Large Katedra veřejného práva} \\	
  \vspace{15mm}
  {\Large \textbf{Studijní program:} Veřejné právo} \\
  {\Large \textbf{Studijní obor:} Právo v obchodních vztazích} \\
  {\Large \textbf{Jméno vedoucího diplomové práce:} } \\
  {\Large Mgr. Ing.  Střeleček  Tomáš LL.M.} \\
  \end{center}

% --------------------------------------------------------------------------  
% Čestné prohlášení 
% --------------------------------------------------------------------------
 
\newpage
  \thispagestyle{empty}
  \vspace*{\fill}

\noindent \textbf{Čestné prohlášení} \\

Prohlašuji, že jsem předkládanou práci zpracoval samostatně, uvedl
v ní všechny použité prameny a zdroje, které jsou uvedeny v seznamu použité
literatury, a v textu řádně vyznačil jejich použití. \\

\noindent V Praze dne \today\\%{\selectlanguage{czech}\today} \\ %Změněno na české datum, možná budu muset odebrat, babel package způsobuje z nějakého důvodu error
\vspace{10mm} \\
\begin{tabular}{p{11.5cm}c}
& ................................................. \\
& Dominik Bálint
\end{tabular}

% --------------------------------------------------------------------------
% Poděkování
% --------------------------------------------------------------------------

\newpage

\thispagestyle{empty}

\vspace*{\fill}
\noindent \textbf{Poděkování} \\

	Na tomto místě bych rád poděkoval svému vedoucímu práce panu Mgr. Ing. 
Tomáši Střelečkovi LL.M. za podporu, trpělivost, spolupráci, podněty ke zlepšení a 
také za čas, který mi věnoval při vedení práce.

\vspace*{\fill}

% --------------------------------------------------------------------------
% Page with resumé
% --------------------------------------------------------------------------

\newpage
	\thispagestyle{empty}
	\begin{center}
		\textbf{\Large Resumé}\\
			\vspace*{5mm}
			Práce je zaměřena na problematiku svěřenských fondů. Zabývá se historickými souvislostmi vedoucími ke společenské potřebě vytvořit instituty podobné svěřenskému fondu a následným vývojem těchto institutů završených přijetím právní úpravy quebecké fiducie do českého právního řádu ve formě svěřenského fondu. Práce dále klade důraz na definování tohoto institutu, jeho právní kodifikaci a charakteristiku ve vztahu k jeho užití pro mezigenerační převod majetku. Práce je členěna na tři části. První část obsahuje teoretický úvod do problematiky a historii. Část druhá obsahuje popis právní kodifikace svěřenského fondu a načrtává spíše demonstrativní výčet výhod svěřenského fondu pro použití k mezigeneračnímu převodu majetku. Poslední kapitoly jsou věnovány problémům spojených se svěřenským fondem, přičemž výzkumná část obsahuje komparativní bádání zaměřující se především na problematiku kolize právní úpravy svěřenského fondu a právní úpravy dědictví společně s, pro český řád, novým pojetím bezsubjektového vlastnictví.\\
			\vspace*{5mm}
			Cílem bakalářské práce bylo zjistit, zda je svěřenský fond použitelný k mezigeneračním převodům majetků společně se zodpovězením otázky, zda je nepominutelný dědic chráněn před možností obejití jeho práva na povinný díl s použitím svěřenského fondu.\\
			\vspace*{5mm}
			Z výzkumu jednoznačně vyplývá, že svěřenský fond přispívá k extenzi individuální sféry vlastnické svobody a je tak velice dobře použitelný k mezigeneračním převodům majetku. Ve vztahu k otázce ochrany nepominutelného dědice však bádání neposkytlo jednoznačnou odpověď, zda je právo nepominutelného dědice chráněno, nebo nikoli.\\
		\vspace*{10mm}
		\textbf{\Large Summary}\\
			\vspace*{5mm}
			The work is focused on the issue of trusts. It deals with the historical context leading to the social need to create institutes similar to trusts and the subsequent development of these institutes culminating in the adoption of the Quebec fiducie into the Czech legal system in the form of a svěřenský fond. The work also focuses on the definition of this institute, its legal codification and characteristics in relation to its use for intergenerational transfer of property. The work is divided into three parts. The first part contains a theoretical introduction to the issue of trusts and its history. The second part contains a description of the legal codification of a svěřenský fond and outlines a rather illustrative list of its benefits with regard to its usage in the intergenerational transfer of assets. The last chapters are devoted to the problems associated with the svěřenský fond, while the research part contains comparative research focusing mainly on the issue of conflict between the legal regulations of the svěřenský fond and the legal regulations of inheritance together with the new concept of non-subject ownership.\\
			\vspace*{5mm}
			The aim of the thesis was to find out whether the svěřenský fond is usable for intergenerational transfers of property together with answering the question whether the forced heir is protected from the possibility of circumventing his right to a mandatory share using a svěřenský fond.\\
			\vspace*{5mm}
			The research clearly indicates that the svěřenský fond contributes to the extension of the individual right of exercising freedom of property dispositions and is thus very well applicable to intergenerational transfers of assets. However, in relation to the question of protecting the forced heir, the research did not provide a clear answer as to whether or not the right of a forced heir is protected.\\
		\vspace*{10mm}
		\textbf{\Large Klíčová slova}\\
			\vspace*{5mm}
			Trust, svěřenský fond, svěřenství, \textit{inter vivos}, \textit{mortis causa}, dědictví, dědic, nepominutelný dědic, smlouva, pořízení pro případ smrti, vlastnictví, fideikomis, Řím, římské právo, kontinentální právo, angloamerické právo, odkaz, Česká republika, Anglie, Kanada, Lichtenštejnsko, Německo, Itálie, Francie, Quebec, CCQ, občanský zákoník, evidence, zakladatelské právní jednání, statut, vznik, novelizace, právní subjektivita, právnická osoba, bezsubjektové vlastnictví, mezigenerační převod majetku, darování, dědická sukcese, daňové zatížení, věřitelé.
	\end{center}


% --------------------------------------------------------------------------
% Page with contents 
% --------------------------------------------------------------------------

\newpage
  \thispagestyle{Contents}
  \tableofcontents
  
 \newpage
  \listoffigures

% --------------------------------------------------------------------------
% Part of the page I list appendixes list on
% --------------------------------------------------------------------------

\vspace*{0.5cm}
  
\noindent{\Large{\textbf{Přílohy}}}\\

\normalsize{1\hspace*{1.8em}Vzor smlouvy o založení svěřenského fondu}\\
\indent\normalsize{2\hspace*{1.8em}Vzor listiny přání zakladatele o určení obmyšlených svěřenského fondu}

% --------------------------------------------------------------------------
% Strana se seznamem použitých zkratek  
% --------------------------------------------------------------------------

\newpage
 \pagenumbering{arabic}
 
\begin{center}
\section{Seznam použitých zkratek}
\end{center}

\vspace{5 mm}

\textbf{Právní předpisy}

\vspace{5 mm}

\begin{tabular}{p{3cm}p{8cm}}
\textbf{OZ} & zákon č. 89/2012 Sb., občanský zákoník	 \\
\textbf{NOZ} & zákon č. 89/2012 Sb., občanský zákoník	 \\
\textbf{CCQ} & Civil code of Quebec, neboli Quebecký občanský zákoník \\
\textbf{ABGB} & Všeobecný zákoník občanský	 \\
\textbf{ZŘS} & zákon č. 292/2013 Sb., zákon o zvláštních řízeních soudních	 \\
\textbf{ZMPS} & zákon č. 91/2012 Sb., zákon o mezinárodním právu soukromém	 \\
\end{tabular}

\vspace{5 mm}

\textbf{Zkratky zemí}

\vspace{5 mm}

\begin{tabular}{p{3cm}p{8cm}}
\textbf{JAR} & Jihoafrická republika	 \\
\textbf{US} & United States \\
\end{tabular}

\vspace{5 mm}

\textbf{Ostatní zkratky}

\vspace{5 mm}

\begin{tabular}{p{3cm}p{8cm}}
\textbf{ZČÚ} & Západočeská univerzita v Plzni \\
\end{tabular}

% --------------------------------------------------------------------------
% Samotný text práce začíná zde  
% --------------------------------------------------------------------------

\newpage
  
\section{Úvod}

Práce si klade za cíl v historickém a právním kontextu představit institut svěřenského fondu a převzetí jeho právní úpravy z právního řádu provincie Quebec.\\
\linebreak

\indent Dále se práce bude zabývat představením současné právní úpravy svěřens\-kého fondu, jeho možnostmi a výhodami, ukázanými na příkladu praktického využití – jako nástroj, který může substituovat, nebo být komplementem klasickému dědění a dalším způsobům převodů majetku a jeho porovnání s těmito způsoby. Dále se budu zabývat i potenciálními možnostmi zneužití a jak, a zda, jsou odstraněny obavy o zneužití svěřenských fondů v souvislosti s novelizací občanského zákoníku z roku 2018.\\
\linebreak

\indent Dalším bodem práce bude faktické porovnání právní úpravy dědického práva s právní úpravou svěřenského fondu. Těžištěm této časti je snaha potvrdit, nebo vyvrátit tvrzení, že právní úprava dědického práva a svěřenského fondu je vyvážená a svěřenským fondem se nedají obcházet práva dědiců. Pokud bych tedy hypotézu refrázoval jako otázku, zněla by následovně: "Dá se svěřenský fond použít k obcházení práv nepominutelných dědiců?"\\
\linebreak

\indent V tomto rámci se budu zabývat zejména zhodnocením právní úpravy svěřens\-kých fondů a její analýzou v souvislosti s ochranou dědiců, vyložení potenciálních problémů současné právní úpravy v rámci českého právního řádu, tzn. potenciální zneužití, v souvislosti s právem nepominutelných dědiců, které mohou vzniknout v případě souběhu dědického řízení a založení svěřenského fondu a dalšími spatřenými nedostatky.\\
\linebreak

\indent Tato část je nesmírně důležitá, neboť věřím, že s ohledem na nastupující novou generaci, která se již narodila do svobodného kapitalistického tržně ekonomického systému, čeká svěřenský fond období, ve kterém o něj bude velký zájem a začne se hromadně využívat, s čím souvisí i potenciální možnost rozšíření důvodů a účelů založení svěřenského fondu a jeho postupné přibližování angloamerické právní úpravě. Další faktorem mnou předvídaného budoucího rostoucího zájmu o tento institut je stárnoucí vrstva populace, která po pádu soci\-alismu naakumulovala určitý majetek a otázka jeho předání nastupující generaci je pro ně důležitá. Je tedy nezbytné ujasnit si určité otázky vyplývající z recepce\ tohoto právního institutu. \\
\linebreak

\indent V závěru práce budou vyloženy systémové nedostatky, založené na historickém kontextu, praktické části této práce a komparaci s právní úpravou tohoto institutu v jiných zemích, v případě, že je naleznu, a následovat bude návrh řešení nalezených problémů. Přílohu práce tvoří mnou navrhnutá smlouva o zřízení svěřenského fondu za účelem předání rodinného majetku.

% --------------------------------------------------------------------------
%{\color{green}\noindent===================Comment===================} \\
%Přidat hypotézu.\\
%Dále by bylo vhodné poznamenat do úvodu, že věřím, že s ohledem na nastupující novou generaci, která již se narodila do svobodného kapitalistického tržněekonomického systému, čeká svěřenský fond období, ve kterém o něj bude velký zájem a začně se hromadně využívat. Je tedy třeba ujasnit si určité otázky vyplývající z recepce tohoto právního institutu. Dále si i myslím, že by se mohly rozšířit důvody a účely založení trustu respektive svěřenského fondu s ohledem na to, jak o něj poroste zájem.\\
%Nelze převzít institut ale právní úpravu, tedy Komparace České právní úpravy s úpravou ostatních zemí
%Je možné komparovat čskou a quebeckou právní úpravu
%Zjistit úpravu zadání co se týče porovnání anglického Trustu s českou právní úpravou
%Musím komparovat právní úpravu ne institut

%V komparaci již představím právní úpravy, tedy nemusím dále představit momentální právní úpravu svěřenského fondu

%Porovnání dědického práva.

%V tomato rámci se budu zabývat zejména tím, že porovnám právní úpravu s právní úpravou svěřenského fondu, budu se zabývat její analýzou v souvisloti s ochranou dědiců

%Hypotéza: Pokud není jako otázka - ANO nebo NE, pokud bych dal jako otázku, tak bych dopředu nepresumoval jak to je, ale dospěl bych k tomu

%V závěru může krátce popsat rozdííy mezi Anglickou právní úpravou a Českou právní úpravou

%V závěru bych mohl komparovat pro případy návrhů jak zamezit zneužití
%Ještě se podívat na poznámku ohledně bakalářské práce v Notes
%{\color{green}\noindent============================================} \\

%Kombinace mortis causa a unter vivis pujde podle me jako pohledavka k pozustalosti
%Dle paragrafu 311 je dle meho sverensky fond ucastnikem dedickeho rizeni - nepominutelni dedicove tedy budou chraneni paragrafem 1452 a budou se moci domahat sveho dedickeho naroku nejspise zalobou na vydani, toto musim overit, podivam se na jine zaloby uplatnitelne v ramci dedickeho prava
%Doplnit tu Haagskou úmluvu o uznávání svěřenských fondů, myslím, že mi tento bod vyplynul když jsem četl o možnostech zneužití svěřenských fondů v Itálii
%Napíšu, že ustanovení § 300 něco týkající se nadací, na které odkazuje svěřenský fond je platné, tzn majetek vyčleněný do svěřenského fondu je součástí dědického řízení to vyplývá i ze zákona o zvláštních řízeních soudních, ale není součástí pozůstalosti nicméně stále bude platit to, že vyčlenění majetku do svěřenského fondu nesmí omezit povinný díl nepominutelného dědice, pokusit se vnést i pohled, že přesto, že svěřenský fond nemá právní subjektivitu, je na něj částečně nahlíženo například pro fungování daní jako na právnickou osobu, toto by mohlo posloužit analogicky i pro dědické řízení, nicméně účastníci asi nebudou dle http://www.bulletin-advokacie.cz/funkcni-kategorizace-sverenskych-fondu
%Co je ale vlastně právní subjektivita, jaké požadavky musí splňovat dědic, fond může podnikat, takže může uzavírat smlouvy, má tedy částečně právní subjektivitu, proč by tedy nemohl být součátí dědického řízení
% --------------------------------------------------------------------------

\newpage

\section{Svěřenský fond - obecně}

\subsection{Pojem svěřenství, svěřenský fond a struktura práce}

% --------------------------------------------------------------------------
% Zamyslet se nad tím, zda jinak nepojmenovat tuto část. Přejmenováno, zkonzultovat, zda to takto dává smysl.
% --------------------------------------------------------------------------

\indent Dne 22.3.2012 vstoupil v platnost nový občanský zákoník, který nabyl účinnosti prvním dnem roku 2014. Nový občanský zákoník s číslem 89/2012 Sb. rekodifikoval civilní právo a zavedl v českém právu několik nových institutů. Rekodifikace civilního práva se dotkla i právní úpravy správy cizího majetku, v jejichž mezích bylo zavedeno několik nových právních institutů. Mezi tyto instituty můžeme zařadit také institut svěřenství, respektive svěřenského nástupnictví, a svěřenského náhradnictví, který předchozí občanský zákoník, až na jeden paragraf\footfullcite[V 40/1964, §859 je upraven zánik všech omezení vyplývajících ze svěřenského náhradnictví]{noauthor_zakon_1964}, neupravoval a také svěřenský fond, na který se práce bude především zaměřovat. \\

\indent Nově zavedená podoba institutu svěřenského fondu a obecné správy cizího majetku je pro nás sice nová, nicméně jedná se o podobný instrument, který v českém právu v minulosti v určitých formách existoval až do 1.1.1951, do zrušení svěřens\-kého náhradnictví, zvané též jako fideikomisární substituce\footfullcite[§565]{noauthor_zakon_nodate}. Omezení vyplývající z institutu svěřenského náhradnictví nicméně trvala až do roku 1964, kdy vstoupil v platnost, a zanedlouho i v účinnost, nový federální občanský zákoník pod číslem 40/1960 Sb., který úplně zrušil svěřenské náhradnictví\footnote{Viz citace 4.}. Mezi roky 1964 a 2014 lze tedy hovořit o faktickém vakuu právní úpravy svěřen\-ských institutů v Československu, České a Slovenské republice a následně v České republice, respektive zákazu používání těchto institutů.\\

\indent Práce je zaměřena na toto období, na širší historii a vznik těchto institutů i na současnou právní úpravu a možnosti z ní vycházejících. Z tohoto důvodu se tedy sestává ze tří dílčích částí, z nichž každá saturuje informační potřebu nezbytnou pro další část práce. První část práce je čistě teoretická, pojednává o historii a pojmových znacích svěřenského fondu a jemu předcházejících institutů a obecnému úvodu do dědického práva, které je nezbytné pro zodpovězení stanovené hypotézy. Druhá část práce je zaměřena na praktické zkou\-mání možností svěřenského fondu zejména s ohledem na možnosti mezigeneračního převodu majetku s ním spojeného a zodpovězení nejasností z této možnosti vyplývajících, především s ohledem na institut nepominutelného dědice. Ve třetí, a tedy konečné, části práce bude následovat můj pohled na tyto problémy a mnou navrhnuté potenciální řešení těchto problémů. \\

% --------------------------------------------------------------------------
%Považuji za vhodné tyto jednotlivé části více popsat a činím proto tak na další straně.\\

%\newpage
% --------------------------------------------------------------------------

{\Large Teoretická část}\\

S ohledem na výše uvedený fakt, je v první řadě nezbytné obecné představení tohoto nově zavedeného institutu v teoretické rovině, které je detailněji popsáno v následující podkapitole a bude sloužit jako solidní odrazový můstek pro vylož\-ení možností, které nám svěřenský fond přinesl, stejně tak, jako jeho potenciálních problémů a nejasností. \\

Kapitola a jednotlivé podkapitoly následující rovněž ve zkratce představí institut svěřenství od dob Římské říše, následovat budou popis vývoje svěřenství během období monarchie a po vzniku samostatného Československého státu, až po zrušení svěřenství v roce 1964. V poslední části bude popsáno znovuzavedení institutu svěřenství a svěřenského fondu novým občanským zá\-koníkem v roce 2014, s důrazem na převzetí právní úpravy svěřenského fondu z právního řádu provincie Quebec, kterému bude, pro lepší pochopení české právní úpravy, předcházet popis a historie trustu, s důrazem na současnou úpravu v právním řádu provincie Quebec. Tato kapitola bude vycházet především z Komentovaného znění části nového občanského zákoníku upravující problematiku správy cizího majetku od Jaroslava Svejkovského a kolektivu. Takto vyložené informace jsou, s ohledem na jejich přímou návaznost na institut svěřenského fondu, zásadní pro pochopení historie a funkčnosti svěřenských fondů. Prioritně bude práce klást důraz na svěřenský fond jako na institut, který je možné použít k mezigeneračnímu převedení majetku. V návaznosti na znovuzavedení institutu svěřenství bude následovat základní popis novelizace právní úpravy svěřenských fondů z roku 2018, který bude dále podrobněji objasněn v druhé, tedy praktické, části práce. V neposlední řadě bude součástí této části i stručný úvod do českého dědického práva, který je důležitý pro pochopení kolizí, které mohou mezi právní úpravou dědictví a právní úpravou svěřenských fondů vzniknout.\\

% --------------------------------------------------------------------------
%které je možné používat pro velké množství účelů.
% --------------------------------------------------------------------------

\newpage

{\Large Praktická část}\\

Po historickém úvodu do problematiky bude následovat plynulý přechod k současné právní úpravě svěřenských fondů a bližšímu vysvětlení novelizace Občanského zákoníku z roku 2018, který mimo jiné přinesl povinnou evidenci svěřenských fondů. Poté se posunu k praktickému využití institutu svěřenského fondu pro výše zmíněné účely mezigeneračního převod majetku. V této části také zhodnotím výhody, respektive cesty, oproti jiným možnostem převodu majetků za stejným účelem. V tomto rámci se také naplno projeví jakási nedotaženost právní úpravy svěřenského fondu v našem právním řádu, pramenící v první řadě z faktu, že svěřenský fond je obdobou institutu trustu, který je především používaný v systémech common law, ve kterých neexistuje mnoho právních institutů vlastní našemu systému civil law. Přes fakt, že institut svěřenského fondu byl přejat z quebeckého občanského zákoníku, který si zachoval charakter systému civil law, tak i tato norma neupravuje mnoho institutů nám známým, stejně tak i na mnoho institutů nahlíží jinak. Z tohoto vyplývá, že i po šesti letech účinnosti nového občanského zákoníku, a tedy i možnosti zakládat svěřenské fondy, přetrvávají nejasnosti ohledně kolize úpravy svěřenského fondu s jinými právními předpisy. Jednu z těchto kolizí představuje i nejasnost mezi založením svěřenského fondu pro případ smrti a právem nepominutelných dědiců. Na tento rozkol existuje mnoho zásadně odlišných názorů a k dnešnímu dni neexistuje dokonce ani žádné precedenční soudní rozhodnutí, které by tuto otázku zodpovědělo. Na toto a další nejasnosti, které vyplývají z recepce institutu vlastního systémům common law do našeho právního řádu, bude klást důraz část následující po části popisující současnou právní úpravu - Systémové nedostatky a návrhy na řešení. Nejasnosti se pokusím, v historickém kontextu, s ohledem na cizí právní úpravu a zhodnocení odborníků, vyložit a nalézt na ně vhodnou odpověď nebo řešení, které bude představeno v následující, tedy závěrečné, části. \\

{\Large Závěrečná část - Návrhy na změny České právní úpravy a návrh smlouvy}\\

% --------------------------------------------------------------------------
%Návrh smlouvy - pokud nebude čas, tak tuto část odstraním
% --------------------------------------------------------------------------

Po analýze a nalezení nedostatků v rámci české právní úpravy je vhodné se nad těmito problémy hlouběji zamyslet a navrhnout jistá řešení. Daná řešení bych chtěl obsáhnout právě v této části práce. Tedy po zanalyzování a popsání jednotlivých nedostatků a vyložení jednotlivých názorových proudů na tyto nedostatky z předchozí části. Přílohou k práci bude v samém závěru mnou navrhnutá smlouva o zřízení svěřenského fondu, který bude zřízen za účelem předání rodinného majetku. Smlouva bude sepsána s ohledem na problémy v této práci vyložené tak, aby pro vznik svěřenského fondu podle této smlouvy nepředstavovaly žádnou překážku.

\newpage

\subsection{Úvod do historických počátků svěřenství a svěřenských fondů}

\indent Aby bylo možné dále hovořit o institutu svěřenství, respektive svěřenského fondu a jeho předlohy ve formě trustu v CCQ, je třeba si nejprve vymezit co institut svěřenství znamená, představuje a k jakému účelu slouží, tedy definovat pojmové znaky a jeho konstrukci.\\

Samotnou definici českým zákonodárcem přijatého trustu dobře shrnuje Bar\-bora Bednaříková, ta institut svěřenství definuje jako: "...takový vztah, kdy v principu jedna osoba svěří svůj určitý majetek druhé osobě ve prospěch osoby třetí."\footfullcite[Kapitola Úvod, str. IX]{bednarikova_barbora_sverenske_2014} K této definici bych dále doplnil, že takto svěřený majetek nemusí sloužit pouze třetí osobě, ale i zakladatelem určenému veřejně prospěšnému účelu.\\

Tato definice se do značné míry ztotožňuje s definicí čistě common law trustu z knihy Underhill and Hayton, Law of Trusts and Trustees, která trust definuje následovně:\\ 

% --------------------------------------------------------------------------
% což je, s ohledem na fakt, že trust sloužil jako inspirace, logické. V daném díle je definován následovně: \\
% --------------------------------------------------------------------------

\textit{"A trust is an equitable obligation, binding a person (called a trustee) to deal with property (called trust property) owned by him as a separate fund, distinct from his own private property, for the benefit of persons (called beneficiaries or, in old cases, cestuis que trust), of whom he may himself be one, and any one of whom may enforce the obligations."}\footfullcite{underhill_a_hayton_d_law_2010} \\

% --------------------------------------------------------------------------
%Doplnit zdroj a pod čáru dát poznámku, že jsem daný zdroj objevil v jiné Diplomové práci, překlad udělám sám, protože nevím koho bych měl ozdrojovat.
% --------------------------------------------------------------------------

\textit{"Trust je spravedlivý závazek, zavazující osobu (trustee - obdoba svěřenského správce), aby se starala o majetek, kterýžto vlastní jako o oddělené jmění, odlišné od jmění jím vlastněné, za účelem prospěchu osob (beneficient), kteroužto může rovněž býti i on sám, nebo kdokoliv další kdo má právo dané závazky vymáhat."}\footnote{Inspirováno překladem z práce Martina Skuhrovce, zdroj překladu není patrný.} \textsuperscript{,}  \footfullcite{skuhrovec_michal_sverensky_2018} \\

Tato podobnost není nikterak zvláštní, neboť, mimo samotné litery zákona a obou definicí, i sama důvodová zpráva k NOZ uvádí, že institut svěřenského fondu vychází z trustu, přičemž se přímo inspiroval quebeckou fiducií, tedy specifickou formou trustu.\\

% --------------------------------------------------------------------------
%Z litery zákona, definice České autorky i zahraničních autorů jasně vyplývá to, co stanovuje i důvodová zpráva, a to jasná inspirace trustem, tedy trustem upraveným v rámci CCQ. \\
% --------------------------------------------------------------------------

Účel takovéhoto svěření, jak je výše popsáno, ať již se bavíme o první, či druhé definici, je tedy zajištění určitého prospěchu pro třetí osobu/y. V historickém kontextu tento prospěch spočíval především v poskytování užitků z vyčleněného majetku a/nebo jeho převedení ve prospěch třetí osoby, nicméně v rámci současné právní úpravy lze hovořit o mnoho dalších účelech svěřenství, ke kterým se v průběhu práce také blíže dostanu. Jedná se jak o svěřenské fondy založené za veřejně prospěšným, tak i za soukromým účelem.\\

\indent V kontinentálním právu se instituty, které by se svým obsahem shodovaly s výše popsaným vymezením, začaly používat již v antickém Římě ve formě tak\-zvaného fideikomisu (latinsky \textit{fideicommissum}, pochází ze slov \textit{fīdē} v překladu "důvěra" a \textit{commissum} v překladu "svěřené", dohromady též jako "tvé důvěře svěřuji"\footfullcite{noauthor_fideicommissum_nodate}). V právu anglosaském, které se vyvíjelo odděleně od dob konce římské nadvlády a především od dob heptarchie v 7. století našeho letopočtu\footfullcite{jan_kuklik_dejiny_2011}, poté ve formě trustů (anglicky \textit{trust}, v překladu "důvěra", již zde lze i z čistě lingvistického hlediska pozorovat podobnost mezi trustem a fideicommissem, významy těchto slov jsou stejné: "důvěře svěřuj" a "důvěra"), které jsou ale obdobně ve své podstatě potomky římských fideicommissů, nicméně ke vzniku trustu jako takového došlo v 16.století, jako důsledek Zákona o závětích\footnote{Dějiny angloamerického práva - strana 86.}. Ve většině zemí západního světa, používajících obou právních sys\-témů, je tato forma správy cizího majetku dlouhodobě používána k uchování, převodu, správě, zachování celistvosti, nebo použití k dalším soukromým, či veřejným účelům\footfullcite[Kapitola Úvod, str. IX]{bednarikova_barbora_sverenske_2014}, v různých formách. Některé země tradičně využívají institutu trustu, například Irsko, jiné země právní úpravu trustu známou především ze zemí Common Law\footnote{Země používající anglosaský systém práva.} přizpůsobili kontinentálnímu právnímu myšlení. Mezi tyto země lze zařadit Lichtenštejnsko (Treuhandverhältnis) a Lucembursko (Fiducie). Další země se rozhodly funkci trustu nahradit jiným institutem - Německo (Treuhand) a Rakousko (Privatstftungsgesetz) \footfullcite{trust_2014}. \\

\indent V současné právní úpravě, tedy v Občanském zákoníku, jsou v paragrafech 1448 a 1449 fakticky vymezeny účely svěřenského fondu, těmito účely je účel soukromý a účel veřejný\footfullcite{noauthor_zakon_2012}. Soukromý svěřenský fond slouží ku prospěchu určité osoby, nebo na její památku, z tohoto ustanovení je možné dovodit, že svěřenský fond založený za soukromým účelem slouží k ochraně, uchování, rozmnožení, správě a převodu vyčleněného majetku, práce tedy bude primárně zaměřena na svěřenské fondy založené za tímto účelem, protože právě tyto svěřenské fondy se dají použít k mezigeneračnímu převodu majetku. Co se veřejného účelu týče, hlavním účelem nemůže být dosahování zisku nebo provozování závodu, svěřenský fond založený za veřejným účelem tedy slouží k veřejnému, především socioekonomickému, prospěchu, do kterého důvodová správa k § 1448 - 1452 řadí účely kulturní, vzdělávací, vědecké, náboženské a další obdobné účely.\footfullcite{jaroslav_svejkovsky_sprava_2015} \\

Stejně jako o základních účelech svěřenských fondů, tedy veřejném a soukro\-mém, lze hovořit i o jednotlivých typech svěřenských fondů používaných v praxi. Pro lepší pochopení je vhodné začít s trustem v common law. Trust jakožto velice flexibilní právní institut by měl mít systém, který dovoluje jisté setřídění individuálních trustů dle jejich účelu a dle potřeb, ke kterým jsou tyto jednotlivé skupiny využívané. Tento systém umožňuje rozdílný přístup k individuálním skupinám lidí jak ze strany úřadů, tak i v rámci samotné funkčnosti trustu. Toto rozdělení trustů je také vhodné pro účely zakladatele (settlor), správce (trustee) a nakonec i obmyšleného (beneficiary), neboť jim poskytuje jistou strukturu, která zjednodušuje rozhodovací proces související s volbou vhodného trustu.\\

Inspiraci je vhodné čerpat například v Anglii, kde trusty mají sofistikovaný systém, jehož jednotlivé části jsou vhodné k různým potřebám a kde slouží jako jeden z pěti možných způsobů jak držet majetek\footfullcite{andrew_burrows_english_2013}.\\

\begin{figure}[h]
\centering
\includegraphics[width=15cm,height=5cm,keepaspectratio]{English_UK_Structure.png}
\caption{Struktura trust systému v Anglii}
\label{fig:struktura}
\end{figure}

Krom výše načrtnuté struktury je možné trusty dělit dle dalších kategorií, living, nebo testamentary, revocable, či irrevocable, funded, či unfunded a dále taky je možné rozdělení dle typů trustů, které jsou používány v praxi. Zde se může jednat například o Insurance Trust, A Spendthrift Trust, Blind Trust a podobně. Vzhledem k tématu práce není třeba jít v tomto směru do detailu, představení tohoto rozdělení v rámci anglického trustu je důležité k představení typů svěřenských fondů v českém právu.\\

% --------------------------------------------------------------------------
%\newpage
% --------------------------------------------------------------------------

Za zmínku stojí pouze rozdíl mezi Express a Implied trustem. Charakteristickým rysem express trustu je dle Judith Bray to, že přímo zakladatel (settlor) stanoví účel trustu a vyčlení do něj majetek\footfullcite{judith_bray_students_2012}. Naopak Implied Trust není tvořen z rozhodnutí zakladatele, ale ze zákona.\footnote{Tamtéž.}\\

Český zákonodárce nešel v rámci rozdělení trustu takto do hloubky, nýbrž zavedl pouze 4 základní kategorie svěřenských fondů, které jsou již představeny výše. Jedná se tedy o svěřenský fond založený \textit{inter vivos} a \textit{mortis causa}. Druhé členění pak spočívá v účelu, pro který byl svěřenský fond založen, tedy soukromý, nebo veřejný. S dalším rozlišením, mimo zákonné dělení, se setkáv\-áme v praxi, kdy jsou svěřenské fondy členěny podle způsobu použití nebo podle typu\footnote{Vhodný inspirační zdroj pro výčet individuálních typů svěřenského fondu poskytují jednotlivé společnosti zaměřující se na poskytování služeb související se zakládáním a správou svěřenských fondů.}. Základními typy svěřenských fondů mohou být osobní, rodinný, obchodní, charitativní, investiční. Pro účely této práce se budu zaměřovat především na svěřenský fond rodinný, testamentární a ochranný\footnote{Bližší rozdělení poskytuji v kapitole Typy svěřenských fondů.}.\\

Následující kapitoly možná budou působit vzhledem k názvu a zaměření práce poněkud extenzivně, ale považuji za důležité zaměřit se v mé práci i na historii svěřenských institutů a správy cizího majetku, neboť k pochopení, vysvětlení a popsání mnoha sporných otázek, které si dnes mnozí právníci pokládají, je nutné znát historické souvislosti a kontext vývoje právní vědy na našem území, neboť současná právní úprava se těmito tradičními právními předpisy v značné míře inspirovala\footfullcite[str. XIX-XXI]{svestka_j_obcansky_2014}, proto je třeba se k zodpovězení těchto otázek občas vydat až k samotným kořenům tohoto institutu. Toto bude důležité především v částech zaměřujících se na rozdíly v pojetí vlastnictví a v otázkách dědického práva.\\

% --------------------------------------------------------------------------
%s ohledem na svěřenské fondy
% --------------------------------------------------------------------------

\newpage

% --------------------------------------------------------------------------
%s Začátek ifcomment, toto potenciálně odeberu
% --------------------------------------------------------------------------

%\ifcomment

\section{Správa cizího majetku - historie}

Počátky svěřenských institutů, které dnes známe v podobě trustů, svěřenských fondů, či jiných obdobných institutů\footnote{Především instituty jiných evropských zemí, které pojali institut trustu odlišným způsobem.}, je možné pozorovat již za dob antického Říma. K pochopení, co vedlo ke vzniku těchto institutů a s nimi spojených problémů, jejichž následky můžeme pozorovat i v dnešní době, bude následující kapitola obsahovat popis systému práva a dědického práva antického Říma, následovat bude představení svěřenských institutů, respektive fiduciárních s\-mluv, vzniklých v této době a následně bude popsán jejich vývoj s ohledem na změny právního systému Říma.

\subsection{Antický Řím}

\textit{Ius est ars boni et aequi}, tedy právo je umění dobra a slušnosti. Tímto Ulpiá\-novým výrokem by se dalo charakterizovat římské právo\footfullcite[Strana 23]{blaho_peter_haramia_ivan_a_zidlicka_michaela_zaklady_1997}. Na tomto principu stojí i určitý právní vývoj, který proběhl v římském právu ve vztahu k funkci dědictví a postupně dal římským úředníkům podnět a hlavně důvod vytvořit i alternativní způsoby převodu majetku, do kterých by se daly zařadit právě fiduciární smlouvy, které jsme od Římanů s jistými úpravami zdědili ve formě trustů a obdobných trust-like institutů.% K procesu tohoto vývoje je však nejdříve nezbytné jisté obecné představení římského práva a specificky římské dědické právo.

\subsubsection{Praetorské právo a Ius Civile}

Oním důvodem, který dal vzniknou fiduciárním institutům byla svazující pravid\-la římského dědického práva stojící na úpravě občanského práva, která jsou, pokud je porovnáme, v přímém rozporu s principy, na kterých stojí novodobý Občanský zákoník, ten se snaží ponechat co nejširší prostor svobodné iniciativě jednotlivce. Tato pravidla byla značně omezující ve vztahu k vůli zakladatele a stála na teorii projevu, která byla značně formalizovaná\footfullcite{bartosek_milan_encyklopedie_1981}. Postupem času však k vývoji římského dědického práva přispívala činnost praetorů, která jej velmi posunula směrem k našemu modernímu pojetí možností projevu vůle zůstavitele a povinností a možností vztahujících se nejen na něj, ale i na jeho dědice, či další osoby a vtělila do něj dříve právně nevymahatelné instituty jako například \textit{fideicommissum}. Kvůli činnosti praetorů a z jejich činnosti plynoucího \textit{ius honorario}, které společně s {ius civile} ovlivňovalo dennodenní život každého Římana je možné na římské právo, a specificky pro účely této práce dědické právo, nahlížet jako na dvojkolejné.\\

K pochopení dvoukolejnosti římského dědického práva je v první řadě třeba vymezit rozdíly mezi \textit{ius civile} a \textit{ius honorarium}, které se v rámci dědického práva lišily způsoby stanovování dědické posloupnosti a přístupu k dědickému řízení. Pro celkové pochopení struktury římského práva soukromého je vhodné zmínit a popsat i další části, kterými jsou \textit{ius gentium} a \textit{ius naturale}. \\

Soukromé právo se sestávalo ze 4 částí, pro jejich vysvětlení je vhodné zkombinovat práci dvou významných římských právníků Ulpiána a Gaiuse. Můžeme přitom vycházet z \textit{Corpus Juris Civilis}, jakožto souhrnného občanského zákoníku vydaného ve Východořímské říši za vlády císaře Justiniána a specificky ze sbírky děl římských právníků s názvem \textit{Digesta} vydaného roku 530 našeho letopočtu. \\

% --------------------------------------------------------------------------
%\newpage
% --------------------------------------------------------------------------

Ulpián píše následující: \textit{Ius naturale est, quod natura omnia animalia docuit: nam ius istud non humani generis proprium, sed omnium animalium, quae in terra, quae in mari nascuntur, avium quoque commune est}, což lze volně přeložit jako: "Přírodní řád je řádem, kterému příroda učí všem živým tvorům, tento řád nejen že není člověku cizí, ale ovlivňuje všechny tvory ať již pochází ze země, moře, či jsou ptáky."\footfullcite[Dig. 1.1.1.3, Ulpianus 1 inst., vlastní překlad]{noauthor_digest_nodate}\textsuperscript{,}\footfullcite[Použito k překladu]{noauthor_digest_book1} \\

Ulpián dále říká: \textit{privatum ius tripertitum est: collectum etenim est ex naturalibus praeceptis aut gentium aut civilibus}, což lze volně přeložit jako: "Soukromé právo skládá se ze tří částí, neboť odvozeno jest buď z příkazů přírodních, příkazů národů, nebo těch vycházejíc z práva civilního."\footfullcite[Dig. 1.1.1.2, Ulpianus 1 inst., vlastní překlad]{noauthor_digest_nodate}\textsuperscript{,}\footfullcite[Použito k překladu]{noauthor_digest_nodate} \\

Gaius v knize první (Dig. 1.1.9., Gaius 1 inst.) soukromé právo rozděluje obdobě na právo civilní a právo národů: \textit{Omnes populi, qui legibus et moribus reguntur, partim suo proprio, partim communi omnium hominum.}, což lze volně přeložit jako: "Všechny národy, které se řídí zvyky a právem částečně používají práva svého a práva národů", tedy \textit{ius civile} a \textit{ius gentium}. \\

Poslední součást římského soukromého práva tvoří tzv. \textit{ius honorarium}, tedy právo tvořené praetory, které bylo pro rozvoj římského práva zásadní. \\

% --------------------------------------------------------------------------
%\newpage
% --------------------------------------------------------------------------

Nejobecněji řečeno, se tedy římské soukromé právo sestávalo ze 4 základních prvků, respektive práv:
\vspace{5 mm}

\begin{itemize}
\item \textit{ius civile},
\item \textit{ius honorarium},
\item \textit{ius gentium},
\item \textit{ius naturale}.
\end{itemize}

\vspace{5 mm}

% --------------------------------------------------------------------------
%\newpage
% --------------------------------------------------------------------------

Při ponechání \textit{Ius naturale} jako základního práva všeho stranou, by se struktura římského soukromého práva, s ohledem na popis v této kapitole, dala vizuálně vyobrazit následovně, pro popis fungování dědického práva je důležité především \textit{ius civile} a \textit{ius honorarium}:

\begin{figure}[h]
\centering
\includegraphics[width=15cm,height=5cm,keepaspectratio]{rimskepravostruktura.jpeg}
\caption{Struktura soukromého římského práva}
\label{fig:struktura}
\end{figure}

\underline{\textbf{\textit{Ius civile}}} je právo, jehož subjekty jsou římští občané, původně vycházelo z obyčejů. Nejstarší \textit{ius scriptum}\footnote{Psané právo.}, vykládající ius civile, který je nám znám je \textbf{Zákon \MakeUppercase{\romannumeral 12} desek} (Latinsky \textit{"Lex duodecim tabularum"}) z 5. století před naším letopočtem. Dle několika dochovaných informací potenciálně existovala psaná úprava již na samém počátku republiky ve formě sbírky \textit{leges regiae}\footnote{Zákony vydávané Římskými králi.}, která byla kodifikována do takzvaného \textit{ius civile Papirianum}\footfullcite[str.31]{mousourakis_roman_2014}. \textit{Ius civile bylo} tvořeno pontifikálními interpretacemi\footnote{Zákony tvořeny Pontifiky, kteří vykládali právo.} a komiciálními zákony\footnote{Zákony tvořeny lidovým shromážděním.} a v době císařství císařskými konstitucemi\footnote{Zákony tvořeny císařem.}. Problém \textit{ius civile} však spočíval v jeho nepružnosti, která se projevovala špatným přizpůsobováním změnám ve společnosti a vývoji doby, a jednoduchého a striktního určení některých zásad\footfullcite[kapitola 1, str.3]{bednarikova_barbora_sverenske_2014}. \\

% --------------------------------------------------------------------------
%Dle několika informací, doplnit zdroj

%\newpage
% --------------------------------------------------------------------------

\underline{\textbf{\textit{Ius honorario}}}, neboli praetorské právo, bylo tvořeno vysokým římským úředníkem zvaným \textit{Praetor}. Úřad praetora byl vytvořen roku 367 před naším letopočtem Liciniovým zákonem\footfullcite[str.35]{blaho_peter_haramia_ivan_a_zidlicka_michaela_zaklady_1997}. Praetor řídil soudní procesy a za pomoci \textit{aequitas}\footnote{Spravedlnost a slušnost.} a \textit{bona fide}\footnote{Dobrá víra, dobrý úmysl.} mohl ovlivňovat a řídit soudní řízení a tím i vytvářet, respektive nalézat, \textit{ius honorarium}, jež rozvíjelo neobratné a zastaralé principy civilního práva, ze kterého vycházel a jehož principy vztahoval na skutkové vztahy, které původní \textit{ius civile} neupravovalo. V situacích, kdy se pro svůj nedostatek formy a obsahu z \textit{ius civile} vycházet nedalo, mohl praetor vytvořit hypotézu (\textit{Formula in factum concepta}) skutkové podstaty, která když je naplněna, má být něco vykonáno, například má být žalovaný odsouzen, tyto hypotézy nemusely být založeny na žádných existujících normách. Praetor zároveň v řízení účastníkům umožnil podávat námitku (\textit{exceptio}), na kterou dle civilního práva nebyl brán zřetel, nicméně pokud se v rámci \textit{ius honorario} daná námitka prokázala, brala se v úvahu\footnote{Tamtéž.}. Praetor zároveň vydával na začátku svého jednoročního funkčního období takzvaný edikt, který fakticky sloužil jako soupis honorárního práva a tímto byla odstraněna potřeba sepisovat \textit{Ius honorarium} ve formě zákoníků. \\

%Doplnit další informace. \\

\underline{\textbf{\textit{Ius gentium}}} je obligační majetkové právo\footnote{Například právo rodinné a dědické zůstávalo vyhrazeno zásadně civilnímu, tedy národnímu, právu a právu tvořeného praetory.}, které bylo tvořeno cizineckým praetorem za účelem poskytování právní ochrany při obchodních stycích Římanů s cizinci, případně cizinců s cizinci na území Říma. V překladu právo národů kombinovalo \textit{Ius civile}, zahraniční, především řecké, právo a usance obchodníků. Ze své podstaty se jednalo o právní předpis jednoduchý, který nekladl na právní subjekty zásadnější formální požadavky. Ve svém důsledku tedy, obdobně jako \textit{Ius honorarium} přispěl i rozvoji \textit{Ius civile}, kterému se stalo podnětem k určitému zjednodušování norem, které byli rigidní, zastaralé a svojí existencí překážely\footfullcite{kincl_j_rimske_1995}.\\

\underline{\textbf{\textit{Ius naturale}}} není v rámci římské jurisprudence vždy vykládáno stejně. Dle mého názoru nejosobitější stanovisko zastává Ulpianus, který \textit{Ius naturale} vysvětluje jako jakési nejvyšší právo působící na všechny tvory, tedy jak lidi, tak i zvířata, které jim bylo vnuknuto samotnou přírodou. V uvozovkách se tedy jedná o právo propůjčující možnost svazků mezi opačným pohlavím, plozením potomstva, či jeho výchovu. Pro téma této práce není nikterak důležité, z toho důvodu jsem ho vynechal i ve výše načrtnutém obrázku přibližujícím strukturu římského soukromého práva a dále se jím již ani nebudu zajímat, uvedl jsem ho pouze pro úplnost\footnote{Tamtéž.}.\\

% --------------------------------------------------------------------------
%\newpage
% --------------------------------------------------------------------------

\subsubsection{Dědické právo v dobách antického Říma}

Po vyjasnění struktury římského práva je tedy možné přesunout se přímo k popisu způsobu fungování Římského dědického práva. Tato část je důležitá pro historické pochopení vzniku svěřenských institutů ve starověkém Říme a pokračování dalšího vývoje těchto institutů, které mají svůj prvopočátek právě v Římské říši. \\

% --------------------------------------------------------------------------
%Římské dědické právo bylo prvně zakotveno v rámci v Zákonu desek, jež dávalo římským občanům právo se domáhat svého dědického práva žalobou\footfullcite{kincl_j_rimske_1995}\textsuperscript{,}\footnote{Obdobně Barbora Bednaříková Svěřenské fondy Institut pro uchovaní a převody rodinného majetku}.
% --------------------------------------------------------------------------

Římské dědické právo bylo prvně zakotveno v rámci \textit{Ius civile} v Zákonu \MakeUppercase{\romannumeral 12} desek, jež dávalo římským občanům právo domáhat se svého dědického práva žalobou\footfullcite{kincl_j_rimske_1995}\textsuperscript{,}\footnote{Obdobně Barbora Bednaříková: Svěřenské fondy. Institut pro uchování a převody rodinného majetku.}. Přes takovouto, dnes by se dalo říci, vyspělost tehdejšího římského práva však tento rigidní zákon kladl před zůstavitele v rámci jeho volnosti pořizovat o svém majetku i značné překážky a obdobně značně omezoval i samotné dědice. Z důvodu této rigidnosti římského civilního práva se začalo do popředí dostávat dědické právo, tak jak na něj nahlížel praetor, to ve svém důsledku zapříčinilo jeho dvojkolejnost\footnote{Ta byla zakotvená právě z důvodu výše popsaného systému práva a jejich, pro vývoj dědického práva, dvou nejdůležitějších složek, \textit{Ius civile} a \textit{Ius honorario}.}.\\

Dědické právo v Říme, vycházelo z principu univerzální sukcese\footnote{Do dědictví spadá jmění, tedy jak majetek, tak i závazky.}\textsuperscript{,} \footfullcite{blaho_peter_haramia_ivan_a_zidlicka_michaela_zaklady_1997}, což sehrálo spolu s pasivní dědickou legitimací významnou roli při tvorbě alternativních způsobů odkazu majetku zůstavitele. \\

Do doby samotného převzetí pozůstalosti byl pozůstalostí majetek ležící pozůstalostí. Ten se mohl v průběhu své existence zvětšovat, či zmenšovat. Později byla ležící pozůstalost chápána jako právnická osoba, existující po ome\-zenou dobu, dokud nebyla převzata povolanými dědici. V případě odúmrti byl majetek \textit{bona vacantia} a připadl \textit{fiscusu}, tedy státní pokladně, ta ručila za pohledávky ručila do výše nabytého dědictví.\\

Právní důsledky přijetí dědictví tedy pro dědice znamenaly, že se stal subjektem všech práv a závazků zůstavitele. Přijetím dědictví došlo ke konfuzi (splynutí) majetku zůstavitele a dědice. Dědic tak plně ručil za dluhy zůstavitele a věřitel se tak mohl dožadovat splnění takovýchto dluhů i v případě, že bylo dědictví předlužené\footnote{Tamtéž.}. Toto mohlo mít negativní dopad i na věřitele dědice. Praetor proto dle situace uděloval \textit{beneficium inventarii}, nebo \textit{beneficium separationis}.\\

V případě určení více dědiců, přičemž někteří se dědici nestali, se uvolněné podíly poměrně přirůstaly ostatním dědicům i s právním břemenem v případě jeho existence. Tato akresence uvolněného podílu proběhla i v případě, že si to dotčený dědic nepřál.\\

% --------------------------------------------------------------------------
%\newpage
% --------------------------------------------------------------------------

Po zůstavitelově smrti šlo dědit na základě zákona, nebo na základě testamentu, toto zároveň určovalo i dědickou posloupnost. \\

\vspace{5 mm}

Římské právo tedy rozlišovalo dvě posloupnosti:
\begin{enumerate}
\item \textit{hereditas testamentaria} - dědická posloupnost na základě testamentu,
\item \textit{hereditas legitima} - dědická posloupnost na základě zákonu.
\end{enumerate}

\vspace{5 mm}

% --------------------------------------------------------------------------
%\newpage
% --------------------------------------------------------------------------

\underline{\textbf{\textit{Dědická posloupnost testamentární}}}

\vspace*{5 mm}

Posloupnost zakládající se na testamentu, tedy jednostranného právního dokumentu ustanovujícího zůstavitelovi dědice, měla přednost před posloupností zákonnou a římské dědické právo stálo na principu \textit{Nemo pro parte testatus, pro parte intestatus decedere potest}. Nevyčerpání celého dědictví ze strany hereditas testamentaria, tedy nemohlo založit právní důvod pro delaci\footnote{Povolání.} intestátních dědiců\footnote{Dědicové dle hereditas legitima.}, v takovém případě dědil i zbytek pozůstalosti dědic povolaný na základě testamentu. \\

Přes fakt, že dědění na základě testamentu mělo přednost, znalo římské právo i institut nepominutelného dědice, jednalo se o osoby v blízkém příbuzen\-ství se zůstavitelem. Právo na část pozůstalosti se ve svých prvopočátcích vymáhalo takzvanou kverelou, tedy žalobou před soudem na zrušení testamentu. Až v pozdější době císařské měli nepominutelní dědici při jejich opomenutí zůstavitelem v závěti právo na dorovnání svého zákonného podílu bez nutnosti rušit testament\footfullcite[3. vydání z roku 2012, str.230]{marek_karel_redakcni_rada_casopis_2012}. Povinnost zohlednit nepominutelného dědice se mohl zůstavitel zprostit vyděděním daného dědice. Z počátku nemusel zůstavitel uvádět důvod, toto bylo změněno vydáním Justiniánských reforem, které taxativně vymezili důvod pro vydědění nepominutelného dědice \footfullcite[str.147]{blaho_peter_haramia_ivan_a_zidlicka_michaela_zaklady_1997} . \\

Z hlediska budoucího vývoje římského dědického práva a v návaznosti na to vývoje fiduciárních smluv je důležitá pouze posloupnost testamentární. Zde bylo kladeno před testátora, i přes krok vpřed, který představoval zákon dvanácti desek a obecně i pozdější iterace \textit{Ius civile}, mnoho překážek, které omezovali jeho testovací volnost. Tyto se projevovali například ve smyslu omezení volnosti určení dědiců v testamentu vycházejícího z rigidnosti římského civilního práva, například co do osob blíže neurčených, aneb \textit{personae incertae} a zapříčinilo jejich dědickou nezpůsobilost. Takovéto osoby jsou určeny jen popisem, namísto jejich jména, tedy například \textit{"dědicem budiž ten, kdo přijde první na můj pohřeb}, takovéto vymezení bylo neplatné. Původně se toto omezení vztahovalo obdobně i na pohrobky a právnické osoby\footnote{Původně mohl dědit pouze \textit{fiscus}, tedy císařská pokladna, a až později města a obce.}, tento přístup se nicméně později obrátil v jejich prospěch\footfullcite{blaho_peter_haramia_ivan_a_zidlicka_michaela_zaklady_1997}.\\

Značně formalizované a omezující římské dědické právo vycházející ve svých počátcích výhradně z \textit{Ius civile} značně omezovalo zůstavitele a stavělo jiné občany do pozice pasivní dědické nezpůsobilost (latinsky \textit{testamenti factio pasiva}) i v mnoha dalších ohledech týkajících se například pohlaví, majetku, rodinného stavu nebo typu příbuzenství\footnote{Tamtéž.}.\\

Jisté řešení skýtalo použití jedné ze tří variant dědické substituce. Tato však sloužila pouze k ustanovení náhradního dědice v případě, že se dědic nemohl, nebo nechtěl dědictví ujmout a neposkytovala tak řešení dědické nezpůsobilosti\footfullcite{kincl_j_rimske_1995}.\\

Zde se naplno projevila důležitost práva praetorského a jeho důležitost na vývoj římského práva a v návaznosti na to i dalších právních systémů vzniknuv\-ších po úpadku a následném zániku Římské říše a s tím i spojených fiduciárních institutů.\\

\newpage

\underline{\textbf{\textit{Dědická posloupnost intestátní}}}

\vspace*{5 mm}

V případě úmrtí zůstavitele, který nezanechal testament, popřípadě byl zanechaný testament neplatný, nastupovalo dědictví \textit{ex lege}.\\

% --------------------------------------------------------------------------
%\newpage
% --------------------------------------------------------------------------

Římské právo rozlišovalo právní postavení dědiců ve smyslu povinnosti dědit na dvě posloupnosti \footfullcite[str.141]{blaho_peter_haramia_ivan_a_zidlicka_michaela_zaklady_1997}:

\vspace{5 mm}

% --------------------------------------------------------------------------
%Římské právo tedy rozlišovalo dvě posloupnosti \footfullcite[str.141]{blaho_peter_haramia_ivan_a_zidlicka_michaela_zaklady_1997}:
% --------------------------------------------------------------------------

\begin{enumerate}
\item \textit{heredes voluntarii} - dobrovolní dědicové, jsou povoláni k dědictví projevem vůle nebo zachováním se ve smyslu přijetí - tedy například uhrazení pohledávky zůstavitele,
\item \textit{heredes necessarii} - nutní dědicové, jsou povoláni k dědictví smrtí zůstav\-itele bez možnosti odmítnutí dědictví, tato skupina se dále dělí na další dvě podskupiny:
\begin{itemize}
\item \textit{heredes sui et necessarii} - dědici nutní a vlastní, římská rodina byla silně patriarchálního charakteru a vztahy se rozlišovaly na takzvané agnátské a kognátské. Paterfamilias, tedy otec rodiny měl moc nad členy jeho rodiny. Tito členové, kteří v době zůstavitelovi smrti spadali pod jeho moc se stali členy této podskupiny. Do vzniku \textit{beneficium abstinendi} v rámci praetorského práva neměli možnost dědictví odmítnout.
\item \textit{heredes necessarii} - jedná se o otroky, kteří byli v testamentu zároveň osvobozeni, neměli možnost dědictví odmítnout. Pokud byl nicméně v testamentu jako dědic označen otrok v moci jiného římského občana, stal se dědicem onen římský občan.
\end{itemize}
\end{enumerate}

\vspace{5 mm}

% --------------------------------------------------------------------------
%Toto ještě ověřit.

%Pokud byl nicméně v testamentu jako dědic označen otrok v moci jiného římského občana, stal se dědicem onen římský občan. \todo{Dohledat zdroj a najít vhodné místo k doplnění v textu} \\
% --------------------------------------------------------------------------

V souvislosti s dalším vývojem v rámci praetorského práva vyhradil dědicům praetor takzvané beneficium abstinenti, které jim umožňovalo vzdát se dědictví a přesto, že stále byli dědici dle ius civile, zabraňovalo také možnosti podávání žaloby zůstavitelových věřitelů proti těmto dědicům\footnote{Tamtéž.}.\\

S dalším vývojem římského práva se rozvíjela i intestátní dědická posloupnost, kterou vnímalo jinak \textit{ius civile}, pro které bylo rozhodující příbuzenství agnátské a \textit{ius honorarium}, pro které bylo rozhodující příbuzenství kognátské. Rozdíly mezi oběma soustavami byly smazány Justiniánovými novelami, ty založili novou posloupnost, která byla velice podobná naší moderní úpravě dědění ze zákona a jednotlivým třídám.\\

% --------------------------------------------------------------------------
%\newpage
% --------------------------------------------------------------------------

\subsubsection{Jiná pořízení pro případ smrti}

V rámci římského práva a zvyklostí bylo možné přenechat určitou část majetku někomu, kdo nebyl dědicem, formou takzvaného legata (v překladu odkazu), jedná se o takzvanou singulární sukcesi. Povinnost převést majetek formou odkazu mohla být uložena pouze testamentárnímu dědici. Tento způsob převodu majetku byl u Římanů velice oblíbený. Smyslem legatu bylo, aby dědic po smrti zůstavitele poskytl nějakou část pozůstalosti jiné osobě. Často se však stávalo, že legatum tvořilo celou pozůstalost, což z hlediska nastavení římského práva nic nepřinášelo dědici, který nejenže neměl žádný prospěch z dědictví, ale stal se fakticky pouhým vykonavatelem závěti s z toho plynoucími obligacemi. Takováto přetěžovaná dědictví pak byla často dědici odmítána\footfullcite{blaho_peter_haramia_ivan_a_zidlicka_michaela_zaklady_1997}. Testament bez dědice však nemohl v římském dědickém právu existovat, takovýto testament se tedy stal nulitním. Další důsledek odmítnutí dědictví zahrnoval užití intestátní posloupnost a neplatnost všech dispozic v počátečním testamentu. Aby se předešlo těmto důsledkům, které byly škodlivé pro všechny zúčastněné strany, začali Římané proti tendenci ukládat dědici příliš mnoho odkazů vydávat zákony\footfullcite{stagl_lex_2018}.\\

% --------------------------------------------------------------------------
%Následně tedy s ohledem na pojetí římského dědického práva, kdy testament nemohl platně existovat bez dědice, se stal nulitním.\\ 
% --------------------------------------------------------------------------

Dále byl tento způsob zatížen formálními požadavky na jeho formu, z tohoto důvodu se začaly vyčleňovat i jiné postupy převodu majetku, které nebyly zatíženy takovou formálností\footfullcite{noauthor_fideicommissum_nodate}. Jeden z těchto postupů byl již výše zmíněný fideikomis, jehož splnění bylo ve svých prvopočátcích svěřeno dědici čistě na základě důvěry. Faktem však nadále zůstávalo, že fideikomis nebyl právně vymahatelný a tím pádem nebylo možné splnění zůstavitelova přání požadovat formou žaloby\footnote{Tamtéž.}.\\

% --------------------------------------------------------------------------
%Doplnit další zdroje dle předmětu římské právo
% --------------------------------------------------------------------------

\subsubsection{Fideikomis v antickém Římě}

V důsledku faktických omezení kladených na \textit{legatum} a obecně na římské dědické právo, ať již dle \textit{Ius civile}, nebo \textit{Ius honorario}, se tedy, již za dob Římské republiky, začal vyvíjet i jiný, z počátku neformální, institut k účelu převodu majetku po smrti zůstavitele, který byl potřebný pro specifické, \textit{Ius civile} přehlížené, situace v rámci dědění ve smyslu dědické nezpůsobilosti, kladené přísnosti na formu testamentu, či situací, kdy chtěl zůstavitel vyjádřit určité jeho přání ohledně nakládání s majetkem\footnote{Takovouto možnost neposkytovalo ani civilní ani honorární právo, mohlo se jednat například o předání majetku nezpůsobilým dědicům, tedy například cizincům, či použití majetku k nějakému opakujícímu se úkonu, Gaius II, 246.}.\\

 Tímto institutem byly takzvané \textit{fideikomisy}, respektive fiduciární smlouvy. Podstata \textit{fideikomisu}, na který lze nahlížet jako otce současně používaného trustu a obdobných fiduciárních institutů používaných jako suplementy nebo substituty klasického dědění, respektive převodu majetku, spočívala v tom, že zůstavitel v testamentu žádal jím určenou osobu (především univerzálního dědice), aby po jeho smrti naplnil určité zůstavitelovo přání. Toto přání spočívalo nejčastěji v převodu části zůstavitelova majetku na jinou osobu, než na určené dědice. Jak již je poznamenáno v úvodu, slovo \textit{fideicommissum} se dá volně přeložit jako "tvé důvěře svěřuji", toto označení bylo výstižné, protože tato žádost se opravdu zakládala čistě na důvěře mezi zůstavitelem a osobou jím určenou, šlo tedy pouze o formální obligaci. Fideikomisář (rozuměj osoba v jejíž prospěch má být plněno) se tedy nemohl žalobou domáhat svého práva. Tato žádost se obvykle připojovala ke kodicilu, jež tvořil právě s odkazem a fideikomisem ustanovení pro případ smrti\footfullcite{michal_skrejpek_rimske_2011}\textsuperscript{,}\footfullcite{leopold_heyrovsky_dejiny_1910}. \\

Hlavní výhody tohoto institutu spočívaly právě v převedení majetku i oso\-bám nezpůsobilým dědit podle \textit{Ius civile} a dále možnost zřídit ho bez složitějších formálních požadavků\footfullcite{bednarikova_barbora_sverenske_2014}.\\

% --------------------------------------------------------------------------
%Za vlády císaře Augustina byla \textit{fideicomisům} přiznána právní závaznost, Claudius později pro fideikomisi zřídil magistráta a Justinianus ve své podstatě sloučil legatum s fideikomisem\footfullcite{blaho_peter_haramia_ivan_a_zidlicka_michaela_zaklady_1997}.\\
% --------------------------------------------------------------------------

Běžné byly univerzální fideikomisy, kterými se měl převádět celý zůstavitelův majetek. Touto formou šlo, především díky fakticky neexistující regulaci a malým nárokům na formu, obejít principy dědického práva, dědic si při tom nemohl zpočátku strhnout falcidiální quartu tak, jako tomu bylo u odkazu. Tato forma převodu majetku však nebyla možná (uskutečnitelná) bez součinnosti dědice, ten musel jako univerzální dědic dědictví přijmout, což bylo v případě, kdyby měl majetek celý převést (obdobně jako u odkazu), pro dědice nevýhodné a opět by se stal pouhým vykonavatelem závěti ale se všemi obligacemi vycháze\-jícími z toho, že dědictví přijmul, dědic tedy ručil za dluhy zůstavitele, přičemž krytí ze strany fideikomisáře bylo pouze obligačního rázu. Je však pravdou, že po dlouhou dobu, po kterou nebyl fideikomis právně zakotven a nebyla mu přiznána právní závaznost\footnote{Ta mu byla přiznána až za vlády Císaře Augustiana.} mohl dědic dědictví v zásadě přijmout a neřídit se zůstavitelovým přáním vyjádřeným v rámci fideikomisu.\\

% --------------------------------------------------------------------------
%\newpage
% --------------------------------------------------------------------------

Dle Barbory Bednaříkové je možné hledat prameny takovéto konstrukce v římských \textit{commendationes}, neboli mandátech. Ty spočívaly v předání jisté věci osobě blízké na smrtelném lóži k jejímu uchování obdobně se však jednalo pouze o nevymahatelnou morální obligaci\footfullcite{johnston_d_roman_1988}.\\

% --------------------------------------------------------------------------
%Je však pravda, že později bylo možné v rámci fideikomisu předat i určitý úkol, který měla daná osoba vykonat ve vztahu k tomuto majetku. Je zde tedy určitá paralela mezi fideikomisem a mandátem.
% --------------------------------------------------------------------------

A právě \textit{commendationes} představovali inspiraci ke zrození konstrukce i pozdějšího fideikomisu. Ten se však od mandátu liší v tom, že umožňuje krom předání určitého majetku i předání určitého úkolu, týkajícího se tohoto majetku, fiduciáři u smrtelného lóže, který měl daný úkol vykonat vůči třetí osobě, tedy fideikomisáři. Fiduciář tedy takto označený majetek držel a spravoval a při splnění či nesplnění určité zůstavitelem uložené podmínky měl fiduciář daný nerozdělený\footnote{Snaha uchovat nerozdělený majetek a předat ho další generaci hrála důležitou roli v celém vývoji svěřenských institutů, jak bude dále vidno při popisu historického vývoje svěřenství na našem území.} majetek předat zůstavitelem určené osobě\footfullcite{bednarikova_barbora_sverenske_2014}.\\

Tento typ fideikomisu se však stále velmi lišil od našeho českého svěřenského fondu, či trustu v Québecu, tak jak jej známe dnes. Zde narážím na, dle mnohých zásadní problém přijímání trustu do evropských právních řádů, tím je pojetí vlastnictví (viz kapitola 7.9 Rozdílné pojetí vlastnictví). V době římské se totiž stále nejednalo o účelově určený majetek dle Lepaulla tak jak jej známe u nás, nebo v Québecu. Tehdejší institut by se dal mnohem více připodobnit k úpravě trustu ve Skotsku, americkém státě Louisiana, nebo v Jihoafrické republice, v těch se fiducář stává vlastníkem, avšak se značným omezením, jedná se tedy v zásadě o vlastnictví \textit{sui generis}\footfullcite{lionel_d_smith_re-imagining_2012}.\\

% --------------------------------------------------------------------------
%Fiduciář tedy takto označený majetek držel a spravoval a při splnění či nesplnění určité zůstavitelem uložené podmínky měl fiduciář daný nerozdělený majetek předat zůstavitelem určené osobě - HOTOVO. 
% --------------------------------------------------------------------------

\subsubsection{Legislativní činnost směřující ke změnám v odkazech\\ a fideikomisech}

Používání odkazů bylo u Římanů velice rozšířené a oblíbené. Problémem však bylo, že jimi přetěžovali dědictví, která pak byla odmítána. Bylo proto žádoucí, aby se zřizování odkazů určitým způsobem omezilo. To s ohledem na budoucí vývoj, který směřoval ke sloučení odkazů a fideikomisů ovlivnilo i samotný fideikomis.\\

\noindent\textbf{\textit{Lex Falcidia}}\\

Zákon přijatý v roce 40 př. Kr. směřoval k omezení odkazů. Nově musel dědic obdržet minimálně jednu čtvrtinu aktivní pozůstalosti ve formě takzvané falcidiánské kvarty. Zůstavitel tak mohl za pomoci odkazu odkázat maximálně 3/4 svého majetku. Tento zákon měl vést k tomu, aby dědicové byli motivováni přijímat dědictví právě z důvodu obdržení minimálně jedné čtvrtiny z celého dědictví a zůstavitelova vůle, která byla vyjádřena ve formě odkazu, byla respektována\footfullcite{blaho_peter_haramia_ivan_a_zidlicka_michaela_zaklady_1997}.\\

\noindent\textbf{\textit{SC Trebellianum}}\\

Jak již jsem však popsal výše tento institut byl stále právem neupravený. Toto se však s jeho vývojem zanedlouho změnilo neboť za dob vlády císaře Augusta byla fideikomisům přiznána právní závaznost. Následně za vlády Nera s vydáním normativního právního aktu s názvem Senatus Consultum Trebe\-llianum senátem z roku 56 našeho letopočtu, tedy již z dob císařství, bylo právně upraveno fungování fideikomisu. Tato právní úprava stanovila, že po vydání dědictví fideikomisáři mají žaloby z aktiv a žaloby věřitelů směřovat výhradně vůči němu\footfullcite[Str. 185]{frydek_mirosla_rimske_2009}. Fiduciář tedy získal exception restitutae hereditatis, která právě toto zaručovala\footfullcite[Str. 460]{noauthor_encyclopedic_1968}. Univerzální fideikomisář nesl tíži odkazu, nemohl si však srazit čtvrtinu pro sebe.\\ 

\noindent\textbf{\textit{SC Pegasianum}}\\

% --------------------------------------------------------------------------
%https://books.google.cz/books?id=oR0LAAAAIAAJ&pg=PA460&lpg=PA460&dq=exceptio+restitutae+hereditatis&source=bl&ots=sUQwV-jwfF&sig=ACfU3U2RoXeRycRvU6f90NbZiPpO7aq6rQ&hl=en&sa=X&ved=2ahUKEwijuezSyoHpAhXxt3EKHbaRAj8Q6AEwAHoECAYQAQ#v=onepage&q=exceptio%20restitutae%20hereditatis&f=false - HOTOVO, ozdrojoval jsem výše. 
% --------------------------------------------------------------------------

Stále však v rámci fideikomisu existoval zásadní problém, který spočíval v tom, že dědic dědictví častokrát odmítl, to mu však koneckonců zaručovala samo římské právo.\\ 

Toto časté odmítání, a problém ve smyslu toho, že tíhu odkazu nesl právě dědic, přijetí dědictví ze strany fiducáře vyřešila až Senatus Consultum Pegasiana. Ta rozšiřuje na fideikomisy ustanovení Lex Falcidia, dědic měl tedy nyní možnost strhnout si z dědictví falcidiánskou kvartu, tedy čtvrtinu převáděného majetku\footfullcite{michal_skrejpek_rimske_2011}, jelikož tvůrci chtěli, aby dědic přijímal dědictví dobrovolně. Tvůrci SC tedy zavádějí také povinou 1/4 dědictví pro dědice, s tím, že pokud není zaručena zůstavitelem je možné, aby si ji dědic sám podržel pro sebe. V případě toho, že by se zdráhal dědictví přijmou, je možné k přijetí a následnému vydání celého dědictví donutit, bez nebezpečí placení pozůstalostních dluhů.\\

Fideikomis s takto vymezenými možnostmi začal být v průběhu času v praxi využíván pro rozšíření možností projevení vůle ze strany zůstavitele. Takto vyčleněný majetek se nazýval fideikomisární substituce jeho specifickým případ\-em byl právě rodinný fideikomis, kterýžto je obdobou současného typu rodinného svěřenského fondu v českém Občanském zákoníku, na nějž bude s ohledem na téma práce kladen největší důraz. Ten hrál zásadní roli v případech ochrany majetku zámožných římských rodin. Takovýto nerozdělený majetek je spravován výhradním dědicem, kterého zůstavitel vybral, a nemůže být rozprodám ani rozdělen. Jeho účelem je hmotné zabezpečení římské rodiny pro další generace\footfullcite[Strana 714 až 715]{bonfante_pietro_instituce_1932}.\\

V římském právu však existovaly i další instituty pro převody majetku, výkladem výše jsem tedy nevyčerpal všechny římským právem umožněné mož\-nosti, jednou z těchto možností představuje například darování pro případ smrti, které se také standardně mezi římskými občany používalo. Částečně bych tedy tuto část uzavřel s tím, že i přes dlouhou cestu, kterou musel fideikomis ještě ujít, aby se z něj vyvinuly trust-like instituty dnešní podoby, je zřejmé, že právě zde se začal vyvíjet tento v dnešní době, minimálně v jiných koutech světa, tak často využívaný institut.\\

% --------------------------------------------------------------------------
%Předělat to, dát tu dolní část před tím a říct, že byli další možnosti, ale neni již moc přínosné o nich mluvit - HOTOVO.

%Obsahem tedy bylo takový přesně označený majetek držet a spravovat, případně ho při splnění či naopak nesplnění určité podmínky předat další osobě, nerozdělený (na toto slovo je potřeba klást zvláštní důraz, neb co do dalšího vývoje v budoucnu bude hrát naprosto stěžejní roli v evoluci institutů fideikomisního typu). Gaius pak ve svých Institucích taktéž poskytuje určitý návod, jakým způsobem fideikomis zřídit, kdy za vhodná slova označuje "peto, rogo, volo, fidei commito", která bychom mohli přeložit jako "žádám, prosím, přeji si, věrnosti svěřuji". Takové doporučení pak jasně podtrhuje charakter fideikomisu jakožto GAIUS Učebnice práva ve čtyřech knihách pouhé morální obligace, kterou v této době skutečně byla - HOTOVO, napsáno o stěžejní roli nerozdělování majetku. 


%Toto psát asi nebudu
%Zásadním zlom pro právní vymahatelnost a stavební kámen celého dalšího vývoje představuje případ prokonzula Lucia Lentula, sloužícího v Africe, který si zhruba roku 4 n. l. za svého fiduciáře zvolil samotného císaře Augusta. Císař povinnost svého svěřenského správcovství splnil a tento případ tak nabyl až precedenčního charakteru. Koneckonců, kam sám císař vedl, kdo jemu nižší mohl takovou povinnost nesplnit? Citace 28 Augustus dále rozhodl, že v případech, kdy fiduciář selže a nesplní svou fiduciární povinnost, mohou se osoby takovým ommitere dotčené domáhat spravedlnosti u konzulů. 

%TOTO DOPLNIT
%Kodicil tvořil v této době společně s legatem a fideikomisem 

%\newpage
% --------------------------------------------------------------------------

\subsection{Historie svěřenských fondů u nás a ve světě}

Není pochyb, že Římská říše a následnické státy (Západořímská říše a Byzantská říše) měli zásadní vliv na rozvoj dědického práva, fideikomisu a práva jako celku, nejen na svém území, ale i ve státních útvarech zbytku Evropy, do které se římské právo jako poměrně vyspělý konstrukt šířilo, a potažmo i celého světa.\\

 Právní vývoj mnoha zemí byl tedy ovlivněn římským právem přičemž právní řády mnoha z těchto zemí dnes stojí na kontinentálním právním řádu, pro který položilo základ právě římské právo.\\

 Česká republika se řadí do skupiny zemí používající, obdobně jako celý evropský kontinent, kontinentální systém práva, přijetí institutu trustu tak představovalo jistou výzvu. Je tedy otázkou proč se český zákonodárce nezaměřil spíše na přijetí institutu například Německého Treuhandu, jež v rámci Německého práva řeší jak institut nepominutelného dědice, tak i problematiku odděleného vlastnictví. Věřím nicméně, že polemika nad co by kdyby nemá pro zbytek práce nikterak zásadní význam, můj pohled na tuto otázku by se tedy dal shrnout následovně: Ve skrze si troufám tvrdit, že důvod pro který se Český zákonodárce rozhodl přijmout trust vycházející z Common Law místo například Německého Treuhandu, jehož historický vývoj se mnohem více blíží fideikomisu\footfullcite[Piotr Stec Strana 46]{cooke_modern_2003} souvisí v zásadě s Globalizací nesoucí s sebou významnou měrou mezinárodní obchod\footfullcite[Strana 61]{mcauley_m_talpis_j_quebec_1992}\textsuperscript{,}\footfullcite[Strana 256]{becker_fiducie_2007} a s tím související jakousi soutěž mezinárodních právních řádů.\\
 
 Na první pohled by se mohlo zdát, že německý treuhand je s naším právním řádem kompatibilnější než trust, není však koneckonců tajemství, že Německo používá, stejně jako my, kontinentální systém práva, tato zjevná kompatibilita tak není překvapující. Německé právo s ohledem na treuhand řeší jak samotný pojem vlastnictví, tak i chrání v případě dalších institucí jako nadace, nepominutelného dědice.\\
 
 Je nepochybné, že hospodářsky je trust potřebný, to koneckonců stanovuje sám zákonodárce v důvodové zprávě, neboť nadnárodní propojení obchodů a finančních transakcí vytváří na přijímání těchto institutů, analogických k anglosaskému pojetí trustu, jakýsi tlak. Profesor Rainer Kulms s tímto zmiňuje problematiku německého přístupu k trustu. To dle něho pomíjí koncepci právního vztahu vytvořeného zakladatelem inter vivos či pro případ smrti, v němž je jmění svěřeno pod kontrolu svěřenského správce ve prospěch obmyšleného či za určitým účelem. Německé soudy pak přistupují k otázce odděleného vlastnictví, k němuž má titul svěřenský správce, velmi zdrženlivě\footfullcite[Německo mezi trustem a treuhandem]{trust_2014}. Čl. 1 odst. 2 písm. h Nařízení o právu rozhodném pro smluvní závazkové vztahy (Řím I) dokládá, že trusty a vztah zakladatele, svěřenského správce a obmyšlených nelze jednoduše zařadit do závazkového práva\footnote{Tamtéž.}\textsuperscript{,}\footfullcite{noauthor_narizeni_nodate}. Další argumenty pro recepci trustů do civil law systémů poskytuje Irina Gvelesiani ve své práci na téma "Treuhand" and "fiducie" (Terminological Problematics), mezi nimiž například:\\
 
 \begin{itemize}
 \item \textit{"the globalization of the legal practice has exhibited an overall tendency of unification and harmonization of legal systems of the world contries",}	
 \item \textit{"the dominance of U.S. and English law firms in complex transactions including non-trust jurisdictions...",}
 \item \textit{"all mature legal systems have catered in various ways to the same needs that have promoted and informed the developement of the common law trust: division between the economic value of assets and their holding and management..."}\footfullcite{gvelesiani__2016}
 \end{itemize}

Je tedy zřejmé, že s ohledem na soutěž právních řádů a použití k oběma institutům vychází trust lépe a nezbývá nic jiného, než pochválit českého zákonodárce, že při jeho přijímání hleděl směrem k trustu. Obdobně by se dalo hovořit o faktoru ekonomického přínosu. Samotný inspirační zdroj ve formě Quebecké fiducie už pak představoval jasné završení procesu jeho vybírání pro rekodifikační komise, o čem ostatně budu hovořit v kapitole Recepce právní úpravy Svěřenského fondu z právního řádu provincie Quebec. Ohledně kompatibility obou institutů obdobně nemá treuhand nikterak navrch viz kapitola Nejasný výklad ochrany nepominutelných dědiců.\\

Samotný historický vývoj fiduciárních institutů na našem území je v obdobích dalece předcházející snaze české rekodifikační komise o nalezení vhodného inspiračního zdroje právní úpravy trustu a jeho přijetí také velice důležitý a svědčí o tom, že instituce obdobné svěřenskému fondu nejsou v Česku ničím zas tak převratně novým. Pojďme si ho tedy ve zkratce představit.

% --------------------------------------------------------------------------
%Samotný hostorický vývoj fiduciárních institutů na našem území v jehož průběhu jsme se dostali až k dnešní uzákoněné formě svěřenského fonduPo krátkém objasnění proč český zákonodárce pohlédl zrovna jako na trust jako inspirační zdroj pro českou právní úpravu
 
 %Napsat že jsem kontinentální řádem a proto je dobré podrobit analýze respektive osvětlit historický vývoj na našem území, aby bylo jasné proč jsme přijali v podstatě trust, který je vlastní právním řádům common law a místo něho jsme nepřijali spíše nějaký systém, který vychází přímo z římského práva a tedy potažmo i z římského fideicomisu.
 
 %Popsat to, že římské právo mělo například vyřešen problém, který nás momentálně s ohledem na svěřenský fond trápí v Českém právu a to vyplacení povinného dílu nepominutelnému dědici. Dědic si totiž mohl strhnout falcidiánskou quartu. 
 
 % a jsou založeny na základě kontinentálního systému právapřičemž mnoho z těchto zemí stojí na kontinentálním právním řádu, pro který položilo základ právě římské právo.\\

%Za vlády císaře Justiniána kdy bylo římské právo a římská právní kultura na svém vrcholu, na naše území začaly přicházet první slovanské kmeny, které poté dali v 9. a 10. století vzniknout základu Českého státu. Celkový historický vývoj na našem území byl, dalo by se říci až do pádu socialismu a několika následujích let, poměrně turbulentní a s tím šel ruku v ruce i vývoj právní kultury a práva jako takového. \\

%Co se našeho území týče, tak

%Co se historie svěřenský fondů týče, tak \footnote[4]{Zde bude nějaká informace}\\
%---Koment\\
%Doplnit úpravu za dob království/monarchie a první republiky a vývoj svěřenských institutů jinde ve světě, držet se struktury popsané v druhém paragrafu v kapitole druhé, to znamená popsat i vznik trustu.
%Dále popsat zrušení svěřenských institutů v období první republiky a úplné zrušení po druhé světové válce.\\
%---\\
% --------------------------------------------------------------------------

\subsubsection{Svět - Anglie}

Anglie jakožto kolébka angloamerického\footnote{Rovněž zvaný jako anglosaský systém práva, nebo common law.} právního systému dala vzniknout i samotnému institutu trustu. Trust anglického původu byl s jistými koncepčními změnami přejat právě do právního řádu provincie Quebec a ačkoliv je mladším z dvojice, položil společně s římským fideikomisem základy současného trustu a jiných obdobných institutů stojících na jeho základu.\\

% --------------------------------------------------------------------------
%Otázku, kterou je vhodné si položit i přes nediskutovatelný fakt, že římské právo ovlivnilo mnohé oblasti a právní vývoj, v Anglii nevyjímaje,

%Základní exkurz do historického vývoje trustu v Anglii začíná v 11. století našeho letopočtu při počátcích vývoje obecného práva, v angličtině common law.

%Institut trustu, jehož podobu, ačkoli v Quebeckém pojetí, jsme přijali, má rodiště v Anglii, respektive na Britských ostrovech. Anglie je zároveň rodištěm systému common law\footnote{Rovněž zvaný jako anglosaský, angloamerický systém práva.}, přijde mi proto vhodné, abych strušně uvedl i tento právní systém.\\
% --------------------------------------------------------------------------

Common law, neboli obecné právo, se v Anglii vyvinulo mezi 11. a 12. století za vlády králů Viléma I. Dobyvatele a Jindřicha II.\\

Před institucionalizací common law Jindřichem II. se velká část právních sporů řešila při různých lidových soudech, městských soudech, panských soudech, či církevních soudech na základě nesouvislého a rozdrobeného práva jednotlivých oblastí\footfullcite[Strana 4]{langbein_john_h_lerner_renee_lettow_smith_bruce_p_history_2009}. Účelem common law bylo právo sjednotit a stát se všeobecně běžně používaným právem celé Anglie.\\

Takto konstruované právo bylo tvořeno spojením původních anglosaských a novějších normanských obyčejů se soudními rozhodnutími královských soudů. Common law se tak stalo základním pramenem soudcovského práva\footfullcite[Strana 3]{noauthor_pravni_2005}.\\

% --------------------------------------------------------------------------
%Možná strana 100
% --------------------------------------------------------------------------

Kvůli zvýšenému počtu žádostí o vyřešení sporů králem, například v případech, kdy se strany sporu domnívali, že byli poškozeni kvůli systémem obecného práva se vyvinul systém equity, rozhodovaný lordem kancléřem, který se stal druhým pramenem soudcovského práva. Kvůli samotné povaze common law a equity byly tyto složky práva často v konfliktu. Soudní spory se tak často protahovaly na celé roky, v 17. století bylo proto ustanoveno, že rozhodnutí podle equity má přednost\footfullcite[Strana 34]{salmond_john_william_jurisprudence_1907}.\\

Soudy common law a equity byly sjednoceny na základě The Judicature Acts z roku 1873 a 1875. V případě konfliktu mělo přednost právo podle equity\footfullcite[strana 96]{viktor_knapp_teorie_1995}. Tato duality lze spatřovat například v common law pojetí vlastnictví v trustu.\\

Za první nezpochybnitelného předchůdce dnešního trustu na území Anglie lze bezesporu považovat právo \textit{use of land} či \textit{usus}\footfullcite[Strana 104]{bednarikova_barbora_sverenske_2014}, jež zahrnovalo převod majetku na jinou osobu, za určitým účelem, s ohledem na nějž s ním daná osoba mohla nakládat, ve prospěch osoby třetí a \textit{entail}, jež se svojí konstrukcí a užitím dal připodobnit k, u nás používanému, nedílu\footnote{Tamtéž, strana 105.}. \\

Institutu trustu dala vzniknout později poměrně přirozená, avšak poněkud na barbarských základech stojící, potřeba anglických rytířů vyrážejících na křížové výpravy či do jiných válečných konfliktů postarat se o svůj doma v Anglii zanechaný majetek. Trust jim, za pomocí převedení tohoto majetku na jinou osobou, jež rytíře zastupovala při veškeré správě daného majetku, přesně toto umožnil. Zároveň se daná osoba krom o převedený majetek starala i o doma zanechanou rodinu. V případě toho, že se rytíř z výpravy vrátil, majetek mu byl převeden zpět, v případě, že se nevrátil, převedl by pak správce majetek na jeho dospělého syna\footfullcite[Strana 15]{david_j_hayton_law_2003}. Základní znaky definice trustu podle Haytona\footnote{Viz kapitola 3.2 Úvod do historických počátků svěřenství a svěřenských fondů.}, tedy aby se správce staral o majetek, ve prospěch jiné osoby jsou tak kromě požadavku na oddělené jmění splněny. Dle common law totiž došlo k reálnému převodu majetku, potenciální spory se tak rozhodovaly na základě equity králem.\\

I přes další vývoj trustu adresující například výše zmíněnou situaci se však samotná konstrukce trustu moc nezměnila. Současná úprava trustu tak předpokládá, že zakladatel zřídí za života, nebo pro případ smrti ve prospěch obmyšleného trust, přičemž majetek převede na správce, který ho bude v rámci trustu spravovat\footnote{Dle pokynů zakladatele, s péčí řádného hospodáře, ve prospěch obmyšlených, ti pak samozřejmě nemusí být určeni například u charitable trustů.}. Obmyšlení, stejně jako správce mohou být jak právnická, tak fyzická osoba. Zakladatel si může ponechat právo aktivní kontroly a možnost rozhodovat sám o věcech zásadních pro fungování trustu, stejně tak může být správce či obmyšleným.\\

Aby však nedocházelo k perpetuitám, bylo zavedeno \textit{Rule Agains Perpetuities}. Toto pravidlo zabraňuje, aby zakladatel za použití zákonných instrumentů vykonával kontrolu nad vlastnictvím soukromého majetku v trustu po dobu delší než trvání životů lidí žijících v čas, kdy byl trust založen. Respektive stanovuje, že trust může existovat maximálně 21 po smrti osoby, jež byla na živu v době založení trustu(zahrnut je i nasciturus). Pravidlo v zásadě brání osobě v aplikaci stanovených pravidel, nebo kritérií, která by i nadále ovlivňovala vlastnictví majetku dlouho poté, co daná osoba zemřela a navíc ho vyděloval z jakéhokoliv hospodářského dění a kontroly státu, což se často označuje jako kontrola „mrtvou rukou“ nebo "mortmain". Po uplynutí této doby je trust automaticky rozpuštěn a majetek je rozdělen do osobního vlastnictví jednotlivých beneficientů\footfullcite[Strana 109]{bednarikova_barbora_sverenske_2014}\textsuperscript{,}\footnote{Ve Velké Británii momentálně perpetuitu upravuje Perpetuities and Accumulations Act 2009, ten v sekci 5 stanovuje maximální dobu založení trustu na 125 let, vztahuje se však pouze na trusty založené po 6.4.2010, text dostupný zde: http://www.legislation.gov.uk/ukpga/2009/18/section/5.}.\\

% --------------------------------------------------------------------------
%., zatímco jeden soud proti druhému čelil [105], i když bylo v 17. století prokázáno, že spravedlnost by měla převažovat.

%Rostoucí počet sporů zapříčinil v důsledku omezeného přístupu k soudům a jejich adresování králi

%Možná strana 100 - nerelevantní, nerad jsem to posunul.
% --------------------------------------------------------------------------

Otázku, kterou je vhodné si položit i přes nediskutovatelný fakt, že římské právo ovlivnilo mnohé oblasti a právní vývoj, v Anglii nevyjímaje, je, zda-li je trust opravdu čistě common law institutem, nebo zda je pokračovatelem fideicomissu, jenž se zpět na kontinent vrátil velkou územní i časovou oklikou do České republiky ve formě svěřenského fondu a současný trust tak stojí na základech \textit{civil law}.\\

Z historického exkurzu do historie trustu v Anglii je nezpochybnitelná značná podobnost k institutu \textit{fideicommissu}. Vzhledem k tomu, že unikátní vývoj trustu v systému \textit{common law} a z tohoto důvodu i jeho konstrukce je často považována za nekompatibilní s právními systémy \textit{civil law}, mnoho autorů v této souvislosti poukazuje právě na existenci \textit{fideicommissu} jakožto trust-like institutu v římském právu\footfullcite[Strana 26]{barriere_francois_reception_2004}. Jsou slyšet i hlasy tvrdící, že právě z fideikomisu se vyvinul trust, a že tak tento čistě common law institut má své pevné kořeny v civilním právu\footnote{Tamtéž, strana 92.}. \textit{Fideicommissum} je v mnoha aspektech podobné anglickému institutu \textit{usus}, který je bezesporu považován jako předchůdce institutu trustu\footfullcite[Strana 6]{waters_donovan_wm_gillen_mark_smith_lionel_waters_2012}.\\

Podobnost mezi oboumi instituty je značná, strany jsou podobné a oba zahrnují převod majetku podléhajícího trust-like závazkům vůči jiné osobě, je však otázkou, zda oba dva instituty vznikly kvůli rigidnosti a neflexibilitě tehdejší právní úpravy a kvůli potřebám tehdejších lidí, nebo zda trust opravdu stojí na římských základech. Nelze tak se stoprocentní jistotou tvrdit, že anglický trust by byl přímým pokračovatelem římského fideikomisu, mnoho tomu ale nasvědčuje\footfullcite[Strana 1]{johnston_d_roman_1988}.

% --------------------------------------------------------------------------
%---Další informace dostupné zde:
%https://books.google.cz/books?id=p-p6DwAAQBAJ&pg=PA245&lpg=PA245&dq=does+trust+come+from+fideicommissum?&source=bl&ots=NCbZLKv1XZ&sig=ACfU3U2ARMUwpOmamr5rR5i623dmccUbiA&hl=en&sa=X&ved=2ahUKEwiv_6bUhJLqAhWL16QKHUuQB0QQ6AEwC3oECAYQAQ#v=onepage&q=does%20trust%20come%20from%20fideicommissum%3F&f=false
% --------------------------------------------------------------------------

\subsubsection{České území}

Za vlády císaře Justiniána kdy bylo římské právo a římská právní kultura na svém vrcholu, na naše území začaly přicházet první slovanské kmeny, které poté daly v 9. a 10. století vzniknout základu Českého státu. Celkový historický vývoj na našem území byl, dalo by se říci až do pádu socialismu a několika následujících let, poměrně turbulentní a s tím šel ruku v ruce i vývoj právní kultury a práva jako takového a s ohledem na to i svěřenských institutů na našem území. \\

Toto období lze také charakterizovat nízkým zájmem o zachování rodinného, respektive rodového, majetku, samotné vlastnictví majetku a jeho udržení vycházelo zásadně z rodových tradic. První zásadnější institut, který přineslo období feudalismu, a který lze přirovnat ke svěřenskému fondu, představoval takzvaný rodinný nedíl, který byl využíván šlechtou a představoval společné a nedílné rodinné vlastnictví. Princip rodinného nedílu spočíval v tom, že po smrti zůstavitele byl majetek v rodinném nedílu převeden na nejstaršího potomka. Vlastník majetku v nedílu měl volné právo s majetkem disponovat, nemohl o něm však pořídit závěť. Jak poukazuje Barbora Bednaříková, lze si tak představit, že tento institut vedl k uchování rodového majetku, takto tomu však nebylo. Kvůli tehdejší úpravě dědického práva, se mohl jakýkoliv člen z nedílu oddělit a měl právo na část majetku v tomto nedílu, samotná existence nedílu tak vedla naopak k rozdrobování majetku na místo toho, aby ho udržovala po hromadě. Koneckonců toto oddělení samo předpokládala tehdejší dědické právo a samo ho aktivně podporovalo\footnote{Z pohledu historiků zřejmě jako úmysl krále oslabit svého hlavního kompetitora, kterýmžto byla šlechta.}. V případě, že nebylo nejstaršího syna, nebo dcery, připad majetek králi.\footfullcite[Strana 24]{bednarikova_barbora_sverenske_2014}.\\

V době pohusitské si pak Česká a Moravská šlechta postupně na králi vydobyla zrušení odúmrti, právo pořídit testament však měla stále pouze s jeho souhlasem. Zrušení odúmrti zapříčinilo to, že zůstala zachována přednost nedílných příbuzných při převodu nedílu. Namísto odúmrti tedy mohli dědictví přijmout dříve oddělení příbuzní. Toto částečně zabránilo většímu dělení majetku, bylo však stále možné se z nedílu oddělit\footnote{Tamtéž strana 26}.\\

Prvním skutečným nástrojem, který by se dal připodobnit k modernímu svěřenskému fondu, v tomto případě rodinnému, byl konstituován pro rod Rožmberků roku 1493. Jeho účelem bylo zabránění problémového drobení majetku tím, že celý majetek přecházel na dědice na základě primogenitury\footnote{Právo založené na senioritě, tedy na pořadí narození.}\textsuperscript{,}\footfullcite[Valentin Urfus Rodinný fideikomis v Čechách, z Sveřenské fondy, Barbora Bednaříková]{vaclav_vojtisek_sbornik_nodate} a ostatní členové měli pouze právo na požitky z tohoto majetku, vyměřené jim podle vůle hlavy rodu\footnote{Tamtéž strana 27}.\\

I rodinný statut rodiny Rožmberků však nebyl bezchybný a existovala oprávněná obava, že by se mohli rodinní členové takovémuto faktickému vydědění vzepřít. Proto šlechta prosadila roku 1498 stavovské usnesení, které bylo zapsáno jako zákon do zemských desk. Tento zákon umožňoval vznik rodinného zřízení, které pokud je stvrzeno králem ve formě rodinného statutu, neumožňuje, mohl být povolán za dědice pouze nejstarší člen rodiny. Takovýto majetek byl nezcizitelný a platil zákaz jeho zadlužování. Zároveň bylo povinností dědice majetek držet a bylo mu znemožněno, aby se proti němu dědic mohl jakkoliv zcizením bránit. Takto se ve své podstatě proměnil rodinný nedíl v 15. a 16. století v institut svěřenství\footnote{Přičemž termín svěřenství byl již za těchto dob nezřídka kdy používán.}, který byl i kodifikován v zemských deskách\footnote{Tamtéž strana 28}.\\ 

Přestože římskoprávní základ není jasně zřetelný lze shrnout, že samotná funkce rodinného nedílu a následně svěřenství přesto vychází z konstrukce a myšlenky římského fideikomisu, tedy že je majetek odevzdán jisté osobě, aby ho spravovala a ta ho v určitý čas nezkráceně předala osobě třetí. Což je jasná paralela k funkci římského fideikomisu.\\

Doba pobělohorská přinesla v rámci Obnoveného zřízení zemského zásadní změny do právního uspořádání svěřenských institutů. Zavedla možnost volného pořizování závětí v rámci dědického práva. Potřeba institutu, který by zajistil rodové udržení majetku však stále byla významná. Tato potřeba se postupně vyvinula v průběhu 17. století a dala vzniknout rodinnému fideikomisu. Rodinný fideikomis šlo zřídit pouze se souhlasem krále. V počátcích bylo možné konstruovat perpetuální rodinné fideikomisy, novela 159 však omezila jejich trvání maximálně na čtyři generace. Majetek vyčleněný do rodinného fideikomisu přecházel na předem omezenou dobu do správy zakladatelem určené osoby, která vykonávala správu v souladu s instrukcemi zakladatele, pro tento vyčleněný majetek byla zaručena nedělitelnost a nezcizitelnost. Rodinný fideikomis sloužil šlechtě, měšťani mohli využít fideikomisární substituce\footfullcite[Strana 31]{bednarikova_barbora_sverenske_2014}. Pokračující vývojové peripetie rodinného fideikomisu jsou však za hranicí tématu mé bakalářské práce, dosti dopodrobna jde ve svém díle například Barbora Bednaříková, která v souvislosti s následným vývojem rodinného fideikomisu po Obnoveném zřízení zemském nashromáždila obrovské množství zdrojů.\\

% --------------------------------------------------------------------------
%Ještě napsat, že se začalo pohlížet na povinný dědický podíl, stránka 49 z Barbory Bednaříkové.
% --------------------------------------------------------------------------

Další vývojový krok přišel s vydáním Všeobecného občanského zákoníku, takzvaného ABGB. Ten obsahoval instituty, které se víceméně blíží současnému pojení svěřenských institutů v Novém občanském zákoníku. Ten v paragrafech 604 až 646 obsahoval ustanovení upravující obecnou substituci (v NOZ je jeho obdobou svěřenské náhradnictví), svěřenské náhradnictví (v NOZ je jeho obdobou svěřenské nástupnictví) a rodinný fideikomis (v NOZ je jeho obdobou svěřenský fond)\footfullcite{noauthor_9461811_nodate}.\\

Rodinný fideikomis v ABGB byl konstruován jako vyčlenění určitého majetku jako nezcizitelného pro budoucí potomky. ABGB rozlišovalo tři druhy nástupnictví, \textit{primogenituru}, \textit{seniorát} a \textit{majorát}\footnote{Prvorozenost, nejstarší člen rodiny, nejstarší příbuzný.}.\\

Založení rodinného fideikomisu vyžadovalo dle ABGB souhlas zákonodárné moci\footfullcite[§ 627]{noauthor_9461811_nodate}. Dle ustanovení § 629 ABGB pak bylo založeno i dělené vlastnictví mezi čekateli a dočasným majitelem svěřenství, tedy správcem. Vlastnictví bylo děleno jako vrchní, které náleželo čekatelům, a užitkové, které náleželo správci, je zde tedy možné pozorovat paralelu mezi anglickým trustem, kdy je vlastnictví děleno dle equity a common law\footnote{Jedná se zde tak o typický prvek systému práva common law, a ukazuje, že jistou dualitu vlastnictví lze založit i v kontinentálním právním řádu, v moderním pojetí tak především vlastnictvím a věcným právem k věci cizí. Nicméně nelze se inspirovat systémem práva common law co do rozdělení titulu k vlastnictví dle zvláštního souboru norem spravedlnosti a obecného práva, srovnání Svejkovský, Marek a kol., Správa cizího majetku v novém občanském zákoníku.}.\\

Po vzniku samostatného Československého státu se začalo debatovat nad šlechtictvím a jeho zrušením, jakožto dalším odstřihnutím se od struktur Rakouska-Uherska z dob minulých. Tohoto bylo docíleno zákonem č. 61/1918 Sb., jímž se zrušují šlechtictví, řády a tituly.\\

Nicméně na svěřenství bylo stále nahlíženo jako na šlechtické privilegium, proto bylo pro postmonarchistickou společnost nežádoucí. Roku 1924 tak v souvislosti s přijetím zákona č. 179/1924 Sb., o zrušení svěřenství bylo zrušeno i to.\\

Podle § 1 se zrušilo svěřenství (rodinné fideikomisy) a nová nemohla být zřizována. Rovněž pozbyla platnost ustanovení zřizovacích listin, rodinných usnesení a jiných opatření o svěřenství. Zanikly také substituční nároky. Jednotlivé poměry měly být vypořádány dohodou. V případě nedosažení dohody do 6 měsíců od vydání tohoto zákona připadl majetek do vlastnictví správce, jako přední dědic přičemž svěřenské nástupnictví měl první čekatel.\\

Svěřenské náhradnictví pak bylo zrušeno Občanským zákoníkem z roku 1950 a stalo se tak posledním hřebíkem do rakve svěřenských institutů v České republice. Ty byly znovu zavedeny až přijetím a účinností nového Občanského zákoníku\footnote{Jistou obdobu před přijetím nového Občanského zákoníku také můžeme spatřovat ve fondech kolektivního investování v již zrušeném Zákoně o kolektivním investování z roku 2004.}.\\

Na přípravě nového Občanského zákoníku pracovala z pověření ministra spravedlnosti rekodifikační komise vedená prof. Dr. JUDr. Karlem Eliášem z právnické fakulty ZČU v Plzni, ti svojí prací završili dlouholetou práci významných českých právníků na rekodifikaci civilního práva tak, aby odráželo nové společenské podmínky, které se v České republice vytvořily po roce 1989\footfullcite[Úvod, strana IX]{bednarikova_barbora_sverenske_2014}. Tato příprava vyústila až v nový Občanský zákoník, který je s několika novelizacemi k současnému datu účinný.\\

% --------------------------------------------------------------------------
%Konstituován bych možná měl přepsat na strukturován

%která sama předpokládala vydělení

%\subsubsection{Vývoj v Anglii}

%Anglie jakožto kolébka angloamerického právního systému dala vzniknout i samotnému institutu trustu. Trust anglického původu byl s jistými koncepčními změnami přejat právě do právního řádu provincie Quebec a ačkoliv je mladším z dvojice, položil společně s římským fideikomisem základy současného trustu a jiných obdobných institutů stojících na jeho základu.\\

%Základní exkurz do historického vývoje trustu v Anglii začíná v 11. století našeho letopočtu při počátcích vývoje obecného práva, v angličtině common law.
% --------------------------------------------------------------------------

\newpage

\subsection{Recepce právní úpravy Svěřenského fondu z právního řádu provincie Quebec}

% --------------------------------------------------------------------------
%Možná napsat o rekodifikační komisy a tomu co dělala, ale spíš to nechám jenom takhle. Mohl bych se kdyžtak inspirovat začátkem knihy od Barbory Bednaříkové a napsat k tomu tak jeden odstaveček, ale spíš to nechám být.
% --------------------------------------------------------------------------

Trust si postupem času získal na mezinárodní scéně velkou oblibu. Mnoho autorů v souvislosti s tímto tvrdí, že přijetí tohoto fiduciárního institutu je nezbytné. Stejné stanovisko zaujímá i zákonodárce v důvodové zprávě, kdy hned na začátku říká následující: \textit{"Generální úprava správy cizího majetku je hospodářsky, sociálně i právně velmi potřebná a v našem soukromém právu je její nedostatek citelný."}\footfullcite{Duvodovazprava}.\\

Obdobné stanovisko zastává i profesorka Madeleine Cantin Cumyn, která nicméně vidí benefit trust like institutů v rámci systémů civil law především v jejich business potenciálu. Společně s tím dodává, že implementace těchto institutů může v Evropských zemích představovat velkou výzvu a být méně atraktivnější, neboť, na rozdíl od zemí používajících systém common law, jsou svázané restrikcemi v rámci svého dědického práva\footfullcite[Madeleine Cantin Cumyn, strana 14]{lionel_d_smith_re-imagining_2012} a z toho důvodu by se nemuseli používat tak výrazně k převodům rodinných majetků\footnote{A to i přes to, že historicky jak trust, tak i kontinentální fiduciární instituce byli používány zejména k mezigeneračním převodům majetku.}, jako tomu bylo většinou v zemích používající systému common law a trustu.\\

% --------------------------------------------------------------------------
%Doplnit, že přes business vhodnost je historicky trust spajt hlavně s mezigeneračním předáním majetku a také uchováním majetku.
% --------------------------------------------------------------------------

I na toto by se dalo pohlížet jako na jeden z důvodů, proč zákonodárce přímo neupravil určité náležitosti ve vztahu svěřenských fondů k dědickému právu. Právě v tomto ohledu zůstalo nejvíce mezer, což je vzhledem k úpravám, kterými prošly ostatní právní předpisy s ohledem na přijetí trustu zarážející\footnote{Například daňové normy.}, nelze však tvrdit, že zákonné mezery existují pouze ve vztahu k dědickému právu, nicméně jsou zde největší. Lze se domnívat, že takováto nedostatečná úprava nahrává zůstavitelově testovací svobodě, neboť zákonodárce nechtěl takto omezit možnost zakladatele pořizovat o svém jmění a znevýhodnit tak typ svěřenského fondu \textit{mortis causa} a omezit obecné používání svěřenských fondů pouze k business účelům\footnote{Příkladem mohou být fondy kolektivního investování.}.\\

Samotný inspirační zdroj pro převzetí nové právní úpravy si český zákonodárce vybral v Občanském zákoníku provincie Quebec. Jak již napovídá sama důvodová zpráva, bylo tomu tak kvůli tomu, že si Quebecké právo zachovalo svůj výrazný charakter kontinentálního práva, jemuž tento institut common law funkčně přizpůsobilo\footfullcite{Duvodovazprava}.\\

S tímto přizpůsobením jde ruku v ruce i originalita, kterou si Quebecký právní řád přizpůsobil Trust. Zejména se jedná o typ přijatého trustu a způsob pojetí vlastnictví.\\

S přijetím Qubecké právní úpravy jsme zároveň do Českého právního řádu přejali i tuto originální formu pojetí vlastnictví. Qubecké právo, kde z hlediska snahy inkorporace common law institutu do systému práva civil law, docházelo k prolínání či přímo koalescenci mezi oběma právními systémy a tím pádem se jevil jako vhodnější inspirační zdroj, než právní řády zemí jako je Jihoafrická republika, Skotsko či Americký stát Louisiana, které aplikují problematičtější pojetí vlastnického práva správce k vyčleněnému majetku formou vlastnictví majetku správcem s mnoha souvisejícími omezeními, kdy se tak ve své podstatě nejedná o účelově určený majetek, který vlastní sám sebe (pojetí v CCQ a Čr, tzv. Lepaullovo vlastnictví)\footfullcite[Vít Zvánovec, Založení svěřenského fondu se zvláštním důrazem na oddělení majetku]{trust_2014}, ale jedná se o vlastnictví sui generis\footfullcite{lionel_d_smith_re-imagining_2012}.\\

% --------------------------------------------------------------------------
%Toto prolínání a úspěšná integrace common law institutu do systému práva civil law
% --------------------------------------------------------------------------

Do Quebeckého právního řádu byl trust přejat s mírou adaptace odpovídající rozdílné právní kultuře z klasické konstrukce trustu v anglosaském právu. Do Quebeckého právní řádu tak byl přejat v zásadě pouze Express trust, jiné druhy trustů jejichž vznik stojí zásadně na principech common law jako je Resulting (Výsledné)\footnote{Chybí-li podstatná náležitost zakládacího aktu, majetek je převeden bez právního důvodu.}, nebo Constructive (Sestrojené)\footnote{Bezdůvodné obohacení.} a subsequentně instituční\footnote{Vznik je postaven na základě deklaratorního rozsudku.} a nápravný\footnote{Vznik je postaven na základě konstitutivního rozsudku.} trust nebyli přijaty, neboť se v jejich případě jedná o implicitní vznik (viz. kapitola Úvod do historických počátků svěřenství a svěřenských fondů), který je v systémech civil law vhodně řešet pomocí obligačních institutů\footfullcite[Vít Zvánovec, Založení svěřenského fondu se zvláštním důrazem na oddělení majetku]{trust_2014}\textsuperscript{,}\footnote{Viz komentář k CCQ, k přijetí trustu, specificky strana 375, odstavec 5, obdobně stanovující, že CCQ jakožto součást civil law systému má vhodné nástroje tyto situace řešící.}.\\

Je pak pouze logické, že český zákonodárce vycházejíc z úpravy trustu v CCQ\footfullcite[Strana 375]{noauthor_report_1978} přijal pouze úpravu Express trustu a jen s malými úpravami ji včlenil do našeho právního řádu v podobě svěřenského fondu.\\

Tuto část je možné uzavřít se závěrem, že analogickým převzetím pasáží z Québeckého občanského zákoníku, upravující trust, došlo k implementaci institutu, který není svázán podmínkami uvedenými v dědickém právu. Toto ostatně uznávají i autoři v rámci komentářové literatury jako je například Miloš Kocí\footfullcite[Strana 1996 až 1999]{svestka_j_obcansky_2014-1}, ale najdou se i tací, kteří v takovéto transplantaci spatřují chybu a právní ustanovení svěřenského fondu vykládají ve vztahu k dědickému právu jinak, například Vlastimil Pihera\footfullcite[Strana 1186 až 1189]{spacil_j_a_kolektiv_obcansky_2013}. Nicméně ať se ztotožním s jakýmkoliv z těchto proudů, je nesporné, že svěřenský fond, hlavně tedy testamentární, otevřel spoustu nových dveří, do kterých bych rád v průběhu práce vstoupil a představil možnosti, které se v nich nacházejí primárně ve vztahu k mezigeneračnímu převodu majetku.\\

Jasná inspirace je trustem je pak zřejmá již z prvního ustanovení věnující se svěřenskému fondu v OZ a CCQ.\\

% --------------------------------------------------------------------------
%PŘEPSAT
%Quebecké právo si zachovalo svůj výrazný charakter kontinentálního práva, jemuž tento institut common law funkčně přizpůsobilo. Důvodová zpráva

%PŘEPSAT
%Je poměrně logické, že při pátrání po vhodném institutu zákonodárce sledoval 10
%stopy kontinentální právní kultury tam, kde se nejspíše střetávala s common law a docházelo tak k
%určité "amalgamaci". Nicméně další adeptské právní řády, Skotsko, Louisiana a JAR přinášejí
%poněkud problematičtější koncepci vlastnického práva (kdy se svěřenský správce se stává
%vlastníkem, ovšem s takovými omezeními, že se jeho vlastnické právo stává v zásadě vlastnictvím
%sui generis)11. A právě v tomto aspektu se projevuje originalita CCQ, která sahá po tzv. "Lepaullově
%vlastnictví". Smith

%---Koment\\
% Následně popsat znovuzavedení těchto institutů v roce 2014, převzetí z Quebecké právní úpravy a popsat i novelu z roku 2018 - na základě stížnosti ministerstev a vrchního státního zastupitelství, odkážu poté na další kapitolu věnující se této novelizaci.
 
 %Tyto instituty v mnohém navazují na starší instituty používané již za dob antického říma a později v Česku, Československu, jejich tradice a dále na to dobré, na co se ohlížíme.
 
 %Existují i právní úpravy, respektive země, které kontinuálně využívali tohoto institutu.\\
 
 %Pro další pochopení takto flexibilního institutu a možnosti z něho vycházejících pro převody rodiného majetku je nejprve nutné vysvětlit dvě věci. První z těchto věcí je samotná úprava svěřenského fondu, kterou obsahuje kapitola následující, věcí druhou je pak samotný základ institutu správy cizího majetku a trustu, tedy jeho struktura a základní premisy a historické souvislosti vedoucí až k podobě tohoto institutu tak, jak jí známe dnes, ať již se jedná o historické souvislosti sahající až k dobám Římské říše, nebo znovupoložení si otázek, které si kladli slavní právníci 20. století a snahu zodpovědět je.\\
 %---\\
 %Tento odstavec podle mě nedává smysl, bude nutné ho mírně upravit, chci v něm říct to, že jsem v předchozích kapitolách vyložil historické souvisloti a v další kapitole popíšu právní úpravu. S touto přípravou je možné zaměřit se na možnosti svěřenského fondu v rámci mezigeneračního majetku a i na mojí výzkumnou otázku.

%\newpage
%\thispagestyle{smallertextinheader}

%\section{Pojem svěřenský fond, úvod do problematiky}

%Text.

%S výše uvedeným je již jistě mnohem pochopitelnější kontrukce CCQ a z něho vycházejícího OZ:
% --------------------------------------------------------------------------

Tato základní podobnost svěřenského fondu vycházející z konstrukce trustu respektive fiducie v CCQ se dá s ohledem na českou transplantaci shrnout v následujících ustanoveních, z kterých je jasně patrné, že jak v České republice, tak i v Quebecu lze založit pouze express trust \textit{mortis causa}, nebo \textit{inter vivos}, obsažených v Pododdílu 1 části upravující Svěřenský fond s odkazem na odpovídající ustanovení CCQ. Již z § 1448 OZ a jeho protějšků v CCQ je podobnost na výsost zřejmá.\\

Připomeňme však pro úplné dokreslení definici svěřenského fondu od Barbory Bednaříkové a doplňme jí o definici odděleného vlastnictví, ta by tedy mohla znít takto: \textit{"Svěřenským fondem je \textbf{samostatně vyčleněný majetek (\underline{účelově určený majetek})}, který jedna osoba \textbf{svěří ke správě osobě druhé, tedy správci (\underline{vznik trustu})}, ve prospěch osoby třetí. Správce majetek spravuje dle vůle zakladatele, \textbf{stejně jako zakladatel a obmyšlený však není vlastníkem vyčleněného majetku, pouze vykonává vlastnická práva\\ (\underline{vlastnická práva})}."}\footnote{Jednotlivé níže uvedené ustanovení přísluší v citaci podtrženým slovům.}.\\

% --------------------------------------------------------------------------
%Pro vznik trusto oba zákoníky stanoví:\\

% a OZ z nichž je patrná zásadní podobnost obou institutů.
% --------------------------------------------------------------------------

---\\

Pro \textbf{\underline{vznik trustu}} oba zákoníky stanoví:\\

\textbf{§ 1260 CCQ}: \textit{"A trust results from an act whereby a person, the settlor, transfers property from his patrimony to another patrimony constituted by him which he appropriates to a particular purpose and which a trustee undertakes, by his acceptance, to hold and administer."}\\

\textbf{Překlad}: \textit{Trust vzniká na základě jednání osoby, zakladatele, který převádí majetek ze svého jmění do jiného jmění jím zřízeného, které vyčlení ke konkrétnímu účelu a kterého se akceptací ujme správce, aby jmění držel a spravoval}\footfullcite[Srovnávací tabulka]{jaroslav_svejkovsky_sprava_2015}.\\

\textbf{§ 1448 stanovuje}:
\textit{(1) Svěřenský fond se vytváří vyčleněním majetku z vlastnictví zakladatele tak, že ten svěří správci majetek k určitému účelu smlouvou nebo pořízením pro případ smrti a svěřenský správce se zaváže tento majetek držet a spravovat.}\\

---\\

K \textbf{\underline{účelově určenému majetku}} oba zákoníky stanoví:\\

\textbf{§ 1265 CCQ}: \textit{"Acceptance of the trust divests the settlor of the property, charges the trustee with seeing to the appropriation of the property and the administration of the trust patrimony and is sufficient to establish the right of the beneficiary with certainty."}\\

\textbf{Překlad}: \textit{Akceptací správcem zakladatel pozbývá majetek, správci vzniká povinnost respektovat přivlastnění majetku (účelu) a (povinnost) spravovat jmění trustu a (akceptací) je na jisto postaven vznik práva obmyšleného}\footnote{Tamtéž.}.\\

\textbf{§ 1448 (2)}:
\textit{(2) Vznikem svěřenského fondu vzniká oddělené a nezávislé vlastnictví vyčleněného majetku a svěřenský správce je povinen ujmout se tohoto majetku a jeho správy.}\\

---\\

K \textbf{\underline{vlastnickým právům}} k majetku oba zákoníky stanoví:\\

\textbf{§ 1261 CCQ}: \textit{"The trust patrimony, consisting of the property transferred in trust, constitutes a patrimony by appropriation, autonomous and distinct from that of the settlor, trustee or beneficiary and in which none of them has any real right."}\\

\textbf{Překlad}: \textit{Jmění trustu sestávající se z majetku převedeného do trustu, tvoří jmění přivlastněné účelu, autonomní a oddělené od jmění zakladatele, správce nebo obmyšleného a takové (jmění), ke kterému nemá žádný z nich jakékoliv věcné právo}\footnote{Tamtéž.}.\\

\textbf{§ 1448 (3)}:
\textit{(3) Vlastnická práva k majetku ve svěřenském fondu vykonává vlastním jménem na účet fondu svěřenský správce; majetek ve svěřenském fondu však není ani vlastnictvím správce, ani vlastnictvím zakladatele, ani vlastnictvím osoby, které má být ze svěřenského fondu plněno.} \\

---\\

Míra adaptace, respektive změn při přijímání institutu trustu je tak poměrně malá, avšak nelze hovořit o tom, že by český zákonodárce přijal Quebecký trust, respektive fiducii, v nepozměněné podobě. Určité unika české úpravy je tak možné pozorovat v rozdílu věcných práv, kdy Quebecký CCQ vylučuje ve vztahu k trustu i ty, český zákonodárce tak však explicitně nečiní\footnote{Srovnání, Svejkovský, Marek a kol. Správa cizího majetku v novém Občanském zákoníku, dualita vlastnických práv, strana 329.}, popřípadě v absenci statutu v CCQ\footnote{Srovnání, tamtéž, Statut svěřenského fondu, strana 428.}, dalším rozdílem je pak typizace fondů v CCQ a OZ.\\

V článcích 1266-1270 CCQ je pojednáno o typech, které je možné podle CCQ založit, já již jsem tento rozdíl ve zkratce zmínil v úvodu práce. Tímto rozdílem oproti CCQ je účel, za kterým lze založit v česku svěřenský fond, je jím účel veřejný a účel soukromý. V CCQ má však zakladatel možnost vybrat si ze tří účelů založení trustu, respektive fiducie.\\

\textbf{§ 1266 CCQ}: \textit{"Trusts are constituted for personal purposes or for purposes of private or social utility."}\\

\textbf{Překlad}: \textit{Trusty jsou zřizovány pro osobní účely nebo pro účely soukromé či sociální}\footnote{Tamtéž.}.\\

Článek 1266 CCQ stanoví, že v Quebecu je možné fiducii založit za 3 účely, a to za účelem osobním, soukromým a sociálním. Přičemž tyto se od sebe liší následovně:\\

\textbf{Osobní trust} je možné založit bez protiplnění za účelem poskytování prospěchu pro konkrétní či určitelnou osobu\footfullcite[Srovnávací tabulka, strana 526]{jaroslav_svejkovsky_sprava_2015}. Jedná se tedy o jakýsi protipól ke svěřenskému fondu založenému za soukromým účelem. Příkladem osobního trustu budiž trust testamentární, nebo \textit{inter vivos} family trust.\\

\textbf{Soukromý trust} je možné založit za účelem vztyčení, udržování nebo ochrany věci nebo k užití majetku přivlastněnému ke konkrétnímu užití a to buď jako nepřímou výhodu pro osobu, nebo na její památku, nebo pro jiný soukromý účel\footnote{Tamtéž.}. Rozdíl lze nalézt v šíři účelů, tento trust je tak možné konstituovat například za účelem retirement trustů, respektive korporátních penzijních fondů\footnote{RCA, RRSP, RPS například: https://www.canada.ca/en/revenue-agency/services/tax/businesses/topics/payroll/payroll-deductions-contributions/special-payments/retirement-compensation-arrangements.html.}\\

Trust zřízený s protiplněním (pro zakladatele), konkrétně takový (trust) zřízený s účelem dosažení zisku pomocí investování (za účelem) zajištění důchodu nebo (za účelem) obstarání prospěchu pro zakladatele nebo pro osoby, které (zakladatel) určí nebo pro společníky v korporaci, družstvu nebo pro zaměstnance či akcionáře, je také \textbf{soukromý trust}\footnote{Tamtéž.}, rozdíl tak lze shledat například i v možnosti protiplnění.\\

Detailnější popis rozdílů mezi osobním a soukromým fondem tak může být složitější a podrobnější rozbor není s ohledem na obsah práce vhodný, ostatně sám Smith říká následující: \textit{Distinguishing a personal from a private trust requires a more nuanced analysis. The C.C.Q. private trust encompasses various kinds of trusts which clearly constitute trusts for persons in the common law. Consequently, although  a trust to provide for the education of one's descendants incontrovertibly constitutes a trust for persons under the common law, it does not immediately follow that it constitutes a personal trust under the C.C.Q."}\footfullcite[Strana 349]{smith_worlds_2013} Základní diferenciaci je tak možné spatřovat v tom, že private trust neslouží určité osobě, ale určitému účelu, ten však nemusí být prospěšný pro širší veřejnost a nejedná se tak o social trust\footfullcite[Strana 127]{bednarikova_barbora_sverenske_2014}.\\

\textbf{Sociální trust} je možné založit za účelem veřejného zájmu jako např. kulturnímu, vzdělávacímu, filantropickému, náboženskému či vědeckému účelu, jeho hlavním účelem není dosahování zisku, nebo provoz závodu\footnote{Tamtéž.}.\\

Zajímavé je zmínit, že v článku 1273 CCQ se explicitně stanovuje, že trusty za soukromým a sociálním účelem mohou být perpetuity, mohou být tedy založeny na neomezeně dlouhou dobu, soukromý trust tedy může být v Quebecu založen na neomezeně dlouhou dobu, v České republice nikoli.\footnote{Neboť byl koncipován jako součást svěřenského fondu založeného za soukromými účely.}\\

Naopak český Občanský zákoník stanoví následující:\\

% --------------------------------------------------------------------------
%Dodělat část dle Michala skuhrovce a lépe popsat rozdíl mezi personal a private trustem dle Barbory Bednaříkové.
% --------------------------------------------------------------------------

\textbf{§ 1449 (1)}:
\textit{(1) Účel svěřenského fondu může být veřejně prospěšný, anebo soukromý.} Svěřenský fond lze tak založit pouze za soukromým a veřejným účelem.\\

Mohlo by se tedy zdát, že český zákonodárce přijal pouze personal a social trust, není tomu ovšem tak. Z definice § 1499 odst. 2 je zřejmé, že byl přejat i private trust, avšak český zákonodárce ho vtělil do personal trustu ve formě svěřenského fondu založeného za soukromým účelem\footnote{Srovnání textace § 1449 odst. 2 OZ "Svěřenský fond zřízený k soukromému účelu slouží k prospěchu určité osoby nebo na její památku." a článku 1268 CCQ.}\\

% --------------------------------------------------------------------------
%\textsuperscipt{,}\footnote{Potenciální problematiku neexistence možné perpetuity u private trustu pak shrnuje Barbora Bednaříková, srovnání strana,}.\\
% --------------------------------------------------------------------------

\newpage
\section{Dědické právo v České republice}

Svěřenský fond je třeba v rámci mezigeneračního převodu majetku, respektive práv a závazků, vnímat především, krom jiných, spíše okrajových forem, jako alternativu či suplement klasické možnosti převodu věci, (ať už se jedná o věc nemovitou, movitou, zastupitelnou, hromadnou etc.) tedy převodu v rámci dědického řízení, který potvrzuje soud\footfullcite[§ 1670]{noauthor_zakon_2012}.\\

Zde záleží na optice, kterou nahlédneme na svěřenský fond a na otázku, zda je svěřenský fond považován za dědice, účastní se dědického řízení a převod vyčleněné věci tedy ve svém důsledku potvrzuje soud, nebo stojí mimo toto řízení a věci vyčleněné do svěřenského fondu nejsou součástí pozůstalosti.\\

Takovéhoto klasického převodu lze ze strany zůstavitele dosáhnou různými způsoby, respektive instrumenty, které jsou upravené v rámci Dědického práva. Jedná s primárně o intestátní posloupnost ve chvíli, kdy zůstavitel o svém majetku nepořídil, popřípadě je majetek převeden na dědice v souladu s přáním, které vyjádří zůstavitel, zde se jedná o pořízení pro případ smrti. Na základě tohoto přání, nebo ze zákona, vzniká dědicům dědický titul, přičemž je možné dědit i na základě několika dědických titulů, tzn. kumulace dědických titulů. V dalším případě pak vzniká jiným osobám, v jejichž prospěch zřídil zůstavitel odkaz, pohledávka za dědici.\\

Za dědice je možné zvolit i právnickou osobu, například nadaci, ta je ale nadána právní subjektivitou a dle § 311 Občanského zákoníku se do ní vyčleňuje majetek povoláním za dědice, je tedy nepochybné, že ustanovení o nepominutelném dědici je zde plně uplatnitelné.\\

Se svěřenským fondem je to ovšem složitější a dle výše zmíněné optiky, kterou na svěřenský fond nahlížíme může, ale zároveň nemusí být účastníkem dědického řízení a s ohledem na to se na něj buď budou, nebo nebudou vztahovat jednotlivé ustanovení dědického práva.\\

Při založení svěřenského fondu \textit{mortis causa} zákonem povolenou cestou, tedy pořízením pro případ smrti je však nutné uvažovat nad právní úpravou vztahující se k tomuto právnímu jednání. To samo o sobě poskytuje v § 1492 ochranu nepominutelným dědicům a je tedy otázka, zda je toto bez dalšího na takovýmto způsobem vyčleněný majetek aplikovatelné.\\

% --------------------------------------------------------------------------
%Jak je to ale se svěřenským fonde? Na to se budu soustředit dále ve své práci.

%\newpage
% --------------------------------------------------------------------------

Aby tedy bylo možné hovořit o výhodách, které svěřenský fond oproti klasickému dědění\footnote{Tedy převodu majetku v rámci dědického řízení (jakožto instrumentu převodu majetkových práv v rámci dědického řízení), tohoto lze dosáhnou mnoha způsoby upravenými právě v rámci Dědického práva, především se jedná o pořízení pro případ smrti.} poskytuje a rozebrat jednotlivé výše nastíněné problémy, je třeba čtenáře stručně obeznámit s právní úpravou dědictví a jeho funkcí v rámci mezigeneračního převodu majetku a to zejména v bodech, ve kterých dochází k, právní obcí vnímaným, kolizím\footnote{Rozdílný pohled je představen v části Systémové nedostatky a návrhy na řešení, existují dva názorové proudy, přičemž jeden z nich, který zastává především Kocí v komentáři k Občanskému zákoníku kontrastuje s názory jiných autorů, mezi ty je možné zařadit Piheru, Svejkovského, Koláře a Horna.} obou právních úprav.\\

% --------------------------------------------------------------------------
%Doplnit či provést jakoukoliv komparaci.
% --------------------------------------------------------------------------

Tyto kolize spočívají především v úpravě pozůstalosti, respektive dědického řízení, vydědění a s ním související ochranou nepominutelného dědice.\\

% --------------------------------------------------------------------------
%Převod práv a závazků je bez pochyby nejčastější formou mezigeneračního převodu majetku. Dědické právo je obecně upraveno převážně v paragrafech 1475 až 1720, přičemž důležitou součást tvoří Díl 5, který upravuje institut nepominutelného dědice.\\

%, která je upravena v Hlavě III Občanského zákoníku v části Absolutní majetková práva v § 1475 a dále, tedy v rámci Dědického práva\footnote{Tedy převodu majetku v rámci dědického řízení (jakožto instrumentu převodu majetkových práv v rámci dědického řízení), tohoto lze dosáhnou mnoha způsoby upravenými právě v rámci Dědického práva, především se jedná o pořízení pro případ smrti.}. Tyto možnosti jsou bezpochyby nejčastější formou mezigeneračního převodu majetku a je tak nutné je stručně představit, aby bylo možné popsat nejen výhody svěřenského fondu oproti způsobům upraveným Dědickým právem, ale také právní obcí vnímané kolize\footnote{Rozdílný pohled je představen v části Systémové nedostatky a návrhy na řešení, existují dva názorové proudy, přičem jeden z nich, který zastává především Kocí v komentáři k Občanskému zákoníku kontrastuje s názory jiných autorú, mezi ty je možné zařadit Piheru, Svejkovského, Koláře a Horna.} mezi oboumi právními úpravami.\\

%Popsat dědický titul.

%Je svěřenský fond věcí hromadnou? Popsat rozdělení věcí podle občanského zákoníku.

%S ohledem na to, že svěřenský fond představuje alternativu či suplement k této možnosti převodu majetku, je třeba čtenáře stručne obeznámit s právní úpravou dědického práva a to především v bodech, ve kterých dochází ke kolizi obou právních úprav. Toto je důležité (nejen) především k definici a popsání těchto kolizí, ale i k vyložení výhod, které svěřenský fond poskytuje oproti formám upraveným v dědickém právu.

%Tyto kolize spočívají především v úpravě pozůstalosti, respektive dědického řízení, vydědění a s ním související ochranou nepominutelného dědice.

%Pro komparaci právní úpravy svěřenského fondu s právní úpravou dědického práva je v první řadě důležité popsat jakým způsobem v České republice funguje dědické právo. Dědické právo je obecně upraveno převážně v paragrafech 1475 až 1720, přičemž důležitou součást tvoří Díl 5, který upravuje institut nepominutelného dědice. \\
% --------------------------------------------------------------------------

\subsection{Stručný úvod do právní úpravy dědictví}

Převod práv a závazků v rámci dědictví je bez pochyby nejčastější formou mezigeneračního převodu majetku. Dědické právo je upraveno převážně v paragrafech 1475 až 1720, přičemž důležitou součást tvoří Díl 5, který upravuje institut nepominutelného dědice.\\

V rámci české právní úpravy vzniká dědické právo smrtí zůstavitele a je právem na pozůstalost nebo na poměrný podíl z ní, ten kdo zemře společně se zůstavitelem, nebo před ním, není způsobilý dědit\footfullcite[§ 1479]{noauthor_zakon_2012}. S ohledem na toto je vhodné poukázat na formální požadavky právní úpravy na zůstavitele a dědice. Nejobecnějším formálním požadavkem kladeným na zůstavitele, je pořizovací způsobilost, respektive nezpůsobilost. Obecně zákon neklade žádné překážky tomu, aby bylo jmění zůstavitele součástí dědického řízení, a aby tak mohlo přejít na jeho dědice. Je nicméně omezen způsob, který může zůstavitel použít k vyjádření své vůle o tom, jakým způsobem má toto převedení proběhnout. To znamená, že zůstavitel je omezen ve způsobech, kterými může pořizovat o svém majetku pokud byl nezletilcem, který nedovršil věku patnácti let, popřípadě pokud by byl omezen ve svéprávnosti\footfullcite[§ 1525 až § 1528]{noauthor_zakon_2012}. Tyto ustanovení dopadají pouze na pořízení pro případ smrti, dědění na základě dědického titulu ze zákona tímto není dotčeno. \\

Je vhodné stručně popsat všechny způsoby vyloučení dědice z dědického práva, neboť dle mnohých autorů je i samotný svěřenský fond způsob vyloučení dědice z dědického práva. Toto může a nemusí být pravda, pojďme tedy toto tvrzení podrobit analýze. K tomu je potřeba osvětlení základních principů českého dědického práva a jeho hmotných a procesněprávních náležitostí.\\

Množina požadavků kladených na jednotlivé dědice, jinak řečeno tedy náležitosti, které dědicové musí splňovat, aby mohli býti dědicky způsobilí, je z rozhodnutí zákonodárce širší. Ve chvíli, kdy tyto požadavky nejsou naplněny nebo je porušena nějaká povinnost kladená na dědice, není umožněno danému dědici nabýt celou pozůstalost, nebo její část. \\

Z tohoto vyplývá, že aby byl dědic způsobilý dědit, musí splňovat jak pasivní, tak i aktivní podmínky na něj zákonem kladené. Aktivní podmínky, které musí dědic splňovat jsou vymezeny v §1481-1482 Občanského zákoníku (Dědická nezpůsobilost), aby dědici nesvědčila dědická způsobilost, nesmí se tedy dopustit ani jednoho z vymezených jednání. Pokud se ho dopustí, může být ještě způsobilý dědit, pokud mu zůstavitel dané jednání výslovně odpustí\footfullcite[§ 1481]{noauthor_zakon_2012}. A podmínky pasivní, které ve své podstatě spočívají pouze v tom, že dědic musí být buď právnická, nebo fyzická osoba. Tato pasivní podmínka je nesmírně důležitá ve vztahu ke svěřenským fondům, neboť svěřenský fond není považován ani za jednu z těchto osob.\\

% --------------------------------------------------------------------------
%\newpage
% --------------------------------------------------------------------------

Jako jistou podskupinu k aktivním podmínkám bych dále zařadil i ostatní skutečnosti vylučující dědice z dědického práva, které mohou nastat na straně dědice, tedy již uvedená dědická nezpůsobilost, či další jednání na straně dědice, tak i na straně zůstavitele. Co se tedy tohoto vyloučení dědice z dědického práva v rámci aktivní podskupiny týče, Občanský zákoník taxativně vymezuje následující důvody k němu vedoucí:

\begin{figure}[h]
\centering
\includegraphics[width=12cm,height=4cm,keepaspectratio]{Dedicka_nezpusobilost.png}
\caption[Vyloučení dědice z dědického práva]{Vyloučení dědice z dědického práva\footfullcite{noauthor_informacni_nodate}}
\label{fig:komparace}
\end{figure}

Vyloučení z dědického práva se obecně dá rozdělit na důvody závislé na jednání dědice a nezávislé na jednání dědice. Do první skupiny patří ve své podstatě dědická nezpůsobilost, zřeknutí, odmítnutí, vzdání se a zcizení dědictví. Do druhé skupiny patří vydědění. \\

\noindent\textbf{Jednotlivé důvody vylučující dědice z dědického práva:} \\

Dědická nezpůsobilost je upravena v § 1481 až § 1483, jedná o případy, ve kterých se dědic dopustil úmyslného trestného činu proti zůstavitelovi, nebo jeho blízkým, popřípadě pokud se dopustil zavrženíhodného činu proti zůstavitelově poslední vůli, spočívající v lstivém svedení k projevu poslední vůle, popřípadě donucení k tomuto projevu, překažení projevu, zatajení, falšování a podvrhnutí nebo zničení pořízení. Dále pokud se dědic dopustil činu uvedeného v § 1482 OZ. Dědic se i přes toto stává způsobilým, pokud mu zůstavitel výslovně tento čin prominul. \\

Zřeknutí se dědického práva je upraveno v § 1484 a stanovuje, že se dědic může předem zříci dědického práva, nebo práva na povinný díl smlouvou se zůstavitelem. Tato smlouva může působit jak proti samotnému dědici, tak i proti jeho potomkům, jedná se o dispozitivní ustanovení občanského zákona a záleží tedy na vůli smluvních stran, aby si mezi sebou toto upravili. Smlouva musí být uzavřena formou veřejné listiny. \\

Odmítnutí dědictví je upraveno v § 1485 OZ, v případě, že to nevylučuje dědická smlouva, může dědic dědictví po smrti zůstavitele, výslovný prohlášením vůči soudu, odmítnout. Nepominutelný dědic může dědictví odmítnout s výhradou povinného dílu. V případě odmítnutí se pak na dědice pohlíží, jako by dědictví nikdy nenabyl.\\

Vzdání se dědictví je upraveno v § 1490, dědic se po smrti zůstavitele může vzdát dědictví ve prospěch jiného dědice. Pokud byl však tento dědic obtížen příkazem, nařízením odkazu, nebo jiným opatřením, které podle zůstavitelovi vůle může a má splnit osobně, pak se této povinnosti vzdáním se dědictví nezbavuje.\\

Pomocí prohlášení o vydědění lze vydědit nepominutelného dědice. Toto prohlášení způsobuje vyloučení dědice z dědění po zůstavitelovi. Takto vyděděný titul nenabývá práva založeného na dědickém titulu, ani nenastupuje po zůstavitelově smrti do jeho povinností kvůli ztrátě práva na svůj zákonem stanovený díl. Lze vydědit pouze v zákonech stanovených případech (viz kapitola Vydědění). S ohledem na toto omezení je možné nahlížet na svěřenský fond jako na možnost obejití tohoto omezení.\\

Dědic může po smrti zůstavitele své dědické právo rovněž bezúplatně, nebo úplatně zcizit dle § 1714 OZ. Nabyvatel tak vstupuje do práv a povinností náležejícím k pozůstalosti. Smlouva vyžaduje formu veřejné listiny.\\

\subsection{Nepominutelný dědic}

Dědicem jako takovým může být nejen člověk, ale i právnická osoba, pokud vznikne maximálně do jednoho roku od zůstavitelova úmrtí. Nepominutelným dědicem však může být pouze osoba označená v § 1643 OZ, tou je dítě zůstavitele, nebo jeho potomci.\\

V dědickém právu nová úprava směřuje především k poznání a respektování skutečné vůle a pohnutek zůstavitele, které jej vedly k určitému právnímu jednání, přes toto však zákonná úprava dědického práva staví před zůstavitele i určité podmínky a požadavky. Jednou z těchto podmínek je institut nepominutelného dědice.\\

Pominutí nepominutelného dědice v závěti však nezpůsobí její neplatnost, byť relativní, ale pouze založí právo nepominutelného dědice požadovat vyplacení svého dílu\footfullcite[§ 1654 odst. 1 OZ]{noauthor_zakon_2012}.\\

\subsection{Vydědění}

Nepominutelného dědice lze vydědit z jednoho ze čtyř stanovených důvodů\footfullcite[§ 1646 odst. 1]{noauthor_zakon_2012}. Zůstavitel může vydědit potomka pokud:

\begin{enumerate}
\item mu neposkytl potřebnou pomoc v nouzi,
\item o zůstavitele neprojevuje opravdový zájem, jaký by projevovat měl,
\item byl odsouzen pro trestný čin spáchaný za okolností svědčící o jeho zvrhlé povaze, nebo
\item vede trvale nezřízený život.
\end{enumerate}

Obdobně lze dědice vydědit v případě, že není dědicky způsobilý, další možnost faktického vydědění představuje zůstavení dílu dětem vyděděného dědice v případě, že by byl tak zadlužen nebo si počínal tak marnotratně, že by měl zůstavitel obavu, že se pro jeho potomky nezachová povinný díl\footnote{Tamtéž, § 1646, 1647.}.\\

Pokud zůstavitel neuvede v prohlášení o vydědění důvod vydědění (§ 1648 OZ), bude mít nepominutelný dědic právo na povinný díl. Toto právo však dědic pozbyde, prokáže-li se ale proti němu zákonný důvod vydědění.\\

Je tedy zřejmé, že založení svěřenského fondu nespadá ani pod jednu z těchto kategorií a zákon tedy explicitně nestanoví, že za jeho pomoci lze vydědit, nebo obejít právo nepominutelných dědiců.\\

V případě, že vyděděný s důvody vydědění nesouhlasí, může se bránit žalobou. Musí pak prokázat proč dle jeho názoru neměl zůstavitel pro vydědění zákonný důvod. Pokud zůstavitel neuvedl důvod vydědění, odkáže se k podání žaloby ten, kdo má dědit na jeho místě\footnote{Tamtéž, § 1673.}.\\

Další možnost pak představuje negativní závěť, kde může zůstavitel vyloučit dědice jiného, než nepominutelného, tomu však musí svědčit zákonná dědická posloupnost, zůstavitel tak může například vydědit svého manžela.

\subsection{Dědické tituly}

Výše jsem hovořil o dědickém titulu, pojďme si tedy ve zkratce představit i ty. \\

Samotný přechod jmění ze zůstavitele na dědice se neděje náhodně, k tomu aby došlo k přechodu jmění musí existovat jistý dědický titul, na němž bude samotný převod založen a na základě kteréhož se bude člověk považovat za dědice\footfullcite[§ 1476]{noauthor_zakon_2012}.\\

Dědickými tituly jsou:

\begin{enumerate}
\item dědická smlouva,
\item závěť,
\item zákon.	
\end{enumerate}

Dědická smlouva je co do právní síly nejsilnějším dědickým titulem. Předpoklady pro její uzavření jsou pořizovací a smluvní způsobilost. Musí mít dále formu veřejné listiny, v případě, že veřejnou listinou není, nahlíží se na ní jako na závěť. Nelze ja její pomoci však pořídit o celé pozůstalosti, čtvrtina pozůstalosti musí zůstat volná, o ní může zůstavitel pořídit pomocí zvlášť projevené vůle\footnote{Viz § 1582 OZ a dále.}.\\

% --------------------------------------------------------------------------
%Není toto taky důležité pro zakládání svěřenského fondu?
% --------------------------------------------------------------------------

Dědická smlouva je dvoustranným právním jednáním a nelze tak jednostranně zrušit, což hraje i důležitou roli pro zakládání svěřenského fondu \textit{mortis causa}. Dědická smlouva však nijak neomezuje zůstavitele na možnosti nakládat se svým majetkem\footnote{Mimo § 1588 odst. 2 OZ.}.\\

Závěť je odvolatelný projev vůle, má přednost před děděním ze zákona. Existuje několik typů závětí, jejich popis však není pro tuto práci důležitý. Zakladatel může v závěti též uvést podmínku, doložení času, nebo příkaz\footnote{Viz vedlejší doložky v závěti, § 1551 OZ.}, minimálně co do možnosti "výplaty" však nemá tolik možností jako u svěřenského fondu\footnote{Viz § 1494 OZ a dále.}.\\

V případě, že dědic o svém jmění nepořídí, přichází na řadu dědění ze zákona. Okruh osob, které mohou dědit ze zákona stanovuje Občanský zákoník v šesti dědických třídách\footfullcite[§ 1633 a dále]{noauthor_zakon_2012}.\\

V případě zůstavitelovi neaktivity tedy nastane dědění ze zákona, pokud zůstavitel aktivitu projeví, může tak učinit jedním z výše popsaných způsobů pomocí pořízení pro případ smrti (závět, dědická smlouva či dovětek, o tom viz níže).\\

Dědici vstupují do práv a povinností zůstavitele, jsou tedy odpovědni za dluhy zůstavitele v plném rozsahu. Dědic tedy musí uhradit dluhy zůstavitele i nad rozsah nabytého dědictví. Zákonodárce nicméně dědici v případě předlužené pozůstalosti přiznává ochranu ve formě výhrady soupisu pozůstalosti. V případě, že si dědic vyhradí soupis pozůstalosti, bude povinen k úhradě zůstavitelových dluhů jen do výše nabytého dědictví\footnote{Tamtéž, § 1647.}.

\subsection{Pozůstalost}

Laická veřejnost nezřídka kdy plete pojmy dědictví a pozůstalost. Rozdíl spočívá v tom, že jako pozůstalost je označováno veškeré jmění\footnote{Majetek představuje to, co osobě patří, přičemž jmění zahrnuje jak majetek, tak i dluhy, srovnání § 495 OZ.} zůstavitele k okamžiku jeho smrti, které je způsobilé přejít na právní nástupce. Do pozůstalosti nespadají práva a povinnosti, které jsou výhradně vázána na osobu zůstavitele, ledaže byla jako dluh uznána u orgánu veřejné moci, zde se jedná například o bolestné či satisfakci v penězích. Naproti tomu dědictví je skutečná část pozůstalosti, která přejde na zůstavitelovi dědice\footfullcite[§ 1475 odst. 2]{noauthor_zakon_2012}.\\

NOZ zachovalo koncepci dědického práva a nevrací se k římskému pojetí tzv. ležící pozůstalosti\footnote{Tak jako tomu bylo před rokem 1950.}.

\subsection{Odkaz a dovětek}

Nový Občanský zákoník navrátil možnost zřizování odkazů. Odkaz je výjimkou z univerzálního dědického nástupnictví, odkazovník tak neručí za dluhy zůstavitele. Za odkazovníka lze povolat osobu fyzickou nebo právnickou, ta však není dědicem. Zatímco nabytí dědictví potvrzuje soud, odkaz uplatňuje odkazovník proti dědici přímo bez soudní ingerence. Odkaz tak fakticky zřizuje odkazovníku pohledávku proti dědicům\footnote{Viz § 1594 a dále OZ.}.\\

Dovětek je společně se závětí a dědickou smlouvu dalším pořízením pro případ smrti. Dovětkem se však na rozdíl od závěti a dědické smlouvy neustanovuje dědic, ale zůstavitel v něm činí jiné opatření, tzn. může zřídit odkaz, zatížit dědice nebo odkazovníka příkazem a podobně. Na způsobilost zůstavitele pořídit dovětkem a také o formách dovětku se použijí stejná ustanovení jako na závěť. Dovětek může být pořízen současně se závětí, nebo samostatně\footfullcite[§ 1498]{noauthor_zakon_2012}.\\

\subsection{Procesní aspekty dědického práva}

Jak již jsem zmínil výše, účelem této části není podání komplexního přehledu o fungování dědického práva a jeho vysvětlení, ale pouze představení jeho nezbytných částí, které jsou podstatné pro následnou komparaci s právní úpravou svěřenského fondu. Z tohoto důvodu neshledávám jako důležité vyložit celý průběh dědického (nově pozůstalostního) řízení, ale pouze ve zkratce vyložit jeho existenci a funkci, neboť součástí analýzy kolizí dědického práva a práva svěřenských fondů bude otázka spojená s průběhem pozůstalostního řízení a v souvislosti s ním také krátký exkurz do Zákona o zvláštních řízeních soudních a do Občanského soudního řádu.\\

Řízení se standardně zahajuje ve chvíli, kdy matriční úřad oznámí soudu úmrtí zůstavitele. Řízení lze také zahájit na návrh dědice.\\

Průběh řízení se skládá ve své podstatě ze dvou fází, první začíná předběžným řízením u notáře a končí nařízením jednání u soudního komisaře, v oněž moment začíná druhá fáze pozůstalostního řízení. V první fázi se zjišťuje jmění zůstavitele, okruh dědiců a z Evidence o právních jednání pro případ smrti se zjišťuje, zda po sobě zůstavitel nezanechal listiny důležité v rámci řízení o pozůstalosti, tedy například závěť či dovětek.\\

V rámci jednání u soudního komisaře, jímž je soudem pověřený notář, se následně pozůstalost rozdělí mezi jednotlivé dědice, podle dohody dědiců, podle zákonných dědických podílů, nebo podle přání zůstavitele. Notář následně vydá usnesení, kterým se dědické řízení končí\footnote{Více například na stránkách Notářské komory, dostupné zde: https://www.nkcr.cz/sluzby/dedicke-pravo/projednani-pozustalosti.}.\\

% --------------------------------------------------------------------------
%Co když se pořízení pro případ smrti, nebo závěť učiněná před svědky neshoduje se závětí, dle které byl vyčleněn majetek do svěřenského fondu. Podle mě soud zhodnotí, kdo má slabší nárok respektive důkazy a toho pak odkáže, aby podal žalobu v občanskoprávním řízení.
% --------------------------------------------------------------------------

% --------------------------------------------------------------------------
% konec ifcomment, potenciálně odeberu
% --------------------------------------------------------------------------

%\fi

\newpage

% --------------------------------------------------------------------------
%{\pagestyle{smallertextinheader}
% --------------------------------------------------------------------------

% --------------------------------------------------------------------------
%\thispagestyle{smallertextinheader}
% --------------------------------------------------------------------------

\section[Svěřenský fond jako mezigenerační nástroj transferu majetku]{Svěřenský fond jako mezigenerační nástroj\\ transferu majetku}

V rámci důvodové zprávy k Novému Občanskému zákoníku uvedl zákonodárce velice zajímavou myšlenku: \textit{Návrh vychází z ideje, že funkční určení soukromého práva je sloužit člověku jako prostředek k prosazování jeho svobody. Účel občanského kodexu je umožnit i garantovat svobodné utváření soukromého života, a ponechat tedy co nejširší prostor svobodné iniciativě jednotlivce. Proto také osnova klade zásadní důraz na hledisko autonomie vůle.}\footfullcite[Strana 20]{Duvodovazprava}. Tato idea se zásadně projevila v oblasti správy majetku, kde představuje základní východisko pro nově zavedený institut svěřenského fondu, a dalších obdobných forem správy majetku, a staví na základním stavebním kamenu pro rekodifikaci soukromého práva. Tímto stavebním kamenem není nic jiného, než pohled na soukromé právo jako na právo, které má v soukromé sféře sloužit subjektům práva jako prostředek prosazování svobody a vytváření svobodného prostoru a iniciativy k vlastním řešením určitých problémů.\\

Velkou součást této části práva tvoří právo majetkové relativní a absolutní. Existence takovéto úpravy je naprosto zásadní pro lidi, kteří v průběhu svého života kumulovali majetek a otázka předání, udržení či rozmnožení, ať už vůči sobě, nebo jiným osobám, takto naakumulovaného majetku je pro ně tedy zásadní.\\

Je nepochybné, že v případě, kdy se vlastník tohoto majetku, respektive jmění, rozhodne svůj majetek svěřit jinému subjektu do správy, za jakýmkoliv účelem, tedy rozmnožením, zařízením převodu, správy a podobně, bude primárně vybírat z možností jemu představených zákonem. Těchto možností je mnoho, majetek bývá pravidelně spravován na základně smlouvy, standardně příkazního typu, osobami určenými v této smlouvě, standardně jimi může být investiční fond, záložna, družstvo, jiný k tomu pověřený správce či banka.\\

Nově lze krom výše zmíněných druhů správy majetku a jiných, pro běžné občany poněkud exotičtějších způsobů správy majetku, jako je vložení majetku do ústavu\footnote{V tomto případě zásadně za účelem sledujícím veřejný prospěch.}, obchodní korporace či fundace\footnote{Nadační fond a nadace, kdy musí být obdobně sledování společensky nebo hospodářsky prospěšný účel.} nebo využití i různých dalších nástrojů jako jsou různá spoření s investiční složkou či investičních fondů\footfullcite{noauthor_zakon_nodate_investicni_spol_a_investicni_fondy}, využít i jednotlivé formy velmi flexibilního institutu svěřenského fondu, jež nabízí obrovské množství účelů, za kterými může být založen.\\

S ohledem na tuto myšlenku se svěřenský fond i přes negativní ohlasy a kritiku jeví být jako velice užitečný a flexibilní institut, který se k uchování a převodům rodinného majetku na další generace velice hodí. Navíc v právních prostředích jiných zemí, především používajících anglosaský systém práva, jsou trust a obdobné trust-like instituty souvisle používány krom jiného běžně za účelem uchování a převodů rodinných majetků\footfullcite[Strana 152]{bednarikova_barbora_sverenske_2014}.\\

Zajímavé je zmínit i možnost využití zahraničních trust-like struktur, je ale nutné vzít na vědomí, že Česká republika není signatářem Haagské konvence o právu uplatnitelném na trusty a jejich uznávání a funkčnost těchto institutů, v praxi České právní úpravy tak může být, oproti zemím, které jsou signatáři a mají vyřešené kolizní normy, jako je například Itálie, problematický\footfullcite{noauthor_hcch_nodate}\textsuperscript{,}\footnote{Příklady dalších institutů uvádím v kapitole Úvod do historických počátků svěřenství a svěřenských fondů.}.\\

S ohledem na úpravu mezinárodního práva soukromého v Čr je však možné uznat zahraniční trust založený podle zahraničního práva obdobně jako podle Haagské úmluvy. Přestože Český republika není signatářem Haagské úmluvy výše zmíněné, samotná důvodová zpráva k ZMPS ho zmiňuje jako inspirační zdroj\footnote{K § 73 stanovuje následující: \textit{Návrh občanského zákoníku obsahuje ustanovení o svěřenském fondu. Ustanovení § 73 proto stanoví kolizní úpravu pro tento právní institut. Tato kolizní úprava vychází z úpravy obsažené v Haagské úmluvě o právu rozhodném pro trust a o jeho uznání z 1. 7. 1985. Úmluva není závazná pro Českou republiku, která není jejím smluvním státem. }.}\textsuperscript{,}\footfullcite[§ 73]{Duvodovazprava_ZMPS}. Svěřenský fond tak musí snést i porovnání s jinými trust-like instituty a potenciální zakladatel se tak může rozhodnout i pro založení trustu podle cizího práva, ten musí naplňovat obecné znaky trustu\footnote{O Tomto více v kapitole 6.10.1 Využití zahraničních trust-like institutů.}.\\

% --------------------------------------------------------------------------
%zajímavé dovození, že v Čr je možné uznat zahraniční trusty, které jsou založeny podle zahraničního práva. 

%Přestože Český republika není signatářem Haagské úmluvy výše zmíněné, samotná důvodová zpráva k ZMPS ho zmiňuje jako inspirační zdroj\footnote{K § 73 stanovuje následující: \textit{Návrh občanského zákoníku obsahuje ustanovení o svěřenském fondu. Ustanovení § 73 proto stanoví kolizní úpravu pro tento právní institut. Tato kolizní úprava vychází z úpravy obsažené v Haagské úmluvě o právu rozhodném pro trust a o jeho uznání z 1. 7. 1985. Úmluva není závazná pro Českou republiku, která není jejím smluvním státem. }.}\textsuperscript{,}\footfullcite[§ 73]{Duvodovazprava_ZMPS}. Svěřenský fond tak musí snést i porovnání s jinými trust-like instituty a potenciální zakladatel se tak může rozhodnout i pro založení trustu podle cizího práva, ten musí naplňovat obecné znaky trustu\footnote{O Tomto více v kapitole.}.\\

%--Koment\\
%Popsat jak je v momentální právní úpravě vyložen institut svěřenského fondu a jaké možnosti poskytuje. \\
%--\\
% --------------------------------------------------------------------------

Jak již jsem výše zmínil, česká právní úprava tedy zakládá mnoho možností, které lze použít k převedení majetku. Těmito mohou býti klasické aktivní způsoby běžně využívané na denní bázi všemi z nás, zahrnující například koupi, darování či výměnu. Lze nicméně použít i méně častých možností jako vydržení, nebo i zakládání právnických osob sloužících za účelem převodu majetku jako je například nadace. Troufám si však tvrdit, že nejčastější jsou způsoby upravené v rámci dědického práva počínajících od nabytí majetku na základě dědického titulu vyplývajícího ze zákona, tedy v případech kdy zůstavitel zůstal za svého života nečinný a neupravil způsob, jakým má být s jeho jměním naloženo po jeho smrti, nebo na základě dědických titulů vycházejících z pořízení pro případ smrti, tedy z právního jednání aktivně učiněného za života zůstavitele. \\

Způsoby, kterými lze dosáhnout převedení majetku jsou tedy dva, aktivní jednání zůstavitele, nebo ponechání převodu majetku na zákoně a na tom jak jsou dědicové schopni se mezi sebou dohodnout, tedy určitá pasivnost zůstavitele. Dá se říci, že s ohledem na převáděné jmění a na budoucí nakládání s takto převedeným jměním, je více než vhodné, aby zůstavitel upravil tyto záležitosti již za svého života. Může tím totiž i nadále v uvozovkách vykonávat správu nad svým jměním, respektive může určit jak s jeho jměním má být po jeho smrti naloženo, takovéto jednání je vhodné zejména s ohledem na mezilidské vtahy v rodině, ochránění jmění před neuvážlivým utrácením ze strany dědiců, ochraně před rozdrobením majetku a mnoha dalšími potenciálními nepříznivými dopady převedení majetku. S ohledem na toto primárně přichází v úvahu způsoby, které jsou upravené v rámci dědického práva, v očích běžného občana tímto bude především závěť, v očích těch znalejších to může být i dědická smlouva, dovětek, nebo odkaz. Tyto způsoby však mají i svá specifická omezení a ve svém důsledku spočívají v konečném převedení jmění na jednotlivé dědice, nebo odkazovníky. Zůstavitel také nemá tak širokou možnost stanovit svým dědicům různé podmínky\footnote{Například ve formě vedlejších doložek v závěti, viz kapitola 5.4 Dědické tituly.} a to i přes existenci dovětku. Z tohoto důvodu nemusí být vhodné jejich použití ve všech případech. \\

% --------------------------------------------------------------------------
%Porovnat a popsat možnosti, které má zůstavitel v dovětku a ve svěřenském fondu.

%\newpage
%\thispagestyle{smallertextinheader}
% --------------------------------------------------------------------------

Jako jistá alternativa k těmto klasickým způsobům byla do české právní úpravy transplantována úprava trustu z Quebeckého občanského zákoníku, tento poskytuje zůstavitelům další možnost, jak naložit se svým jměním již za svého života (založení svěřenského fondu \textit{inter vivos}), popřípadě po své smrti (založení svěřernského fondu \textit{mortis causa}). Institut svěřenského fondu přináší mnoho možností, které zůstavitel před jeho zavedením neměl, ale zároveň do českého práva vnáší i spoustu otázek, které je potřeba pro jeho správné používání a ochranu zúčastněných stran zodpovědět. V následujících kapitolách proto představím způsoby, kterými lze svěřenský fond použít pro mezigenerační převod majetku, představím jeho přínosy a benefity, ale zároveň se budu snažit zodpovědět otázky týkající se nejistého výkladu právní úpravy svěřenského fondu, jeho potenciálních problémů a možností zneužití. \\

% --------------------------------------------------------------------------
%\newpage

%\thispagestyle{smallertextinheader}
% --------------------------------------------------------------------------

\subsection{Svěřenský fond v České republice}

% --------------------------------------------------------------------------
%Za použití zcela shodných argumentů jako v kapitole 5.2.2 lze opět dojít k závěru, že majetek ve svěřenském fondu, jakožto autonomní vlastnictví, není součástí pozůstalosti a tedy ani není součástí dědického řízení. Nemůže k němu vznikat dědické právo a tím z něj nikdo ani nemůže být vyloučen. Pro takový případ nelze než doporučit, aby budoucí zůstavitelé-zakladatelé ponechali určitý prostor pro správce, aby ten mohl v případech podobných (jako předpokládá OZ u dědické nezpůsobilosti) rozhodnout o plnění obmyšleným podle obecných zásad spravedlnosti a dobrých mravů.

%Napsat i že správcem svěřenského fondu může být jak fyzická osoba, tak i osoba právnická, ta ale jenom stanoví li tak zákon. Tímto zákonem je pokud se nepletu zákon o obchodních a investičních společnostech.

%--- Toto asi odstraním\\
%Pro další pochopení svěřenského fondu a jeho možností, problémů, výhod, nevýhod a nejasností je v samém úvodu této kapitoly nutné čtenáře obeznámit s právní úpravou svěřenského fondu v rámci českého práva. \\
%---\\
% --------------------------------------------------------------------------

Pro správný výklad svěřenského fondu a jeho použití k mezigeneračnímu převodu majetku je prvně důležité popsat právní úpravu svěřenského fondu.\\

Definovat správně svěřenský fond, jakožto velice flexibilní právní institut, je poměrně složité a až doposud, tedy i relativně dlouhou dobu po jeho zavedení, existuje spousta protichůdných názorů (například v rámci lingvistického výkladu jednotlivých ustanovení, definicí majetku nebo koexistenci právního institutu přejatého z právního systému Common Law s instituty vlastní kontinentálním právním systémům, jako je například institut nepominutelného dědice) na tento institut, který v době svého přijetí do naší, tedy kontinentální, právní úpravy vyvolal poměrně mnoho emocí, od mnohých právních expertů se dočkal dílem chladného přijetí a od mnohých dílem přímé kritiky. Ať již nahlížíme na svěřenský fond jakkoliv, je zřejmé, že pro jednotlivé subjekty práva a pro jednotlivé účely, ke kterým je možné jej užít, skýtá i velké výhody. S ohledem na název bakalářské práce bude hlavní těžiště této práce v rámci praktické části spočívat především v představení těchto nových možností a výhod v procesu převádění rodinného majetku na další generace a jeho uchování jako alternativní, nebo rozšiřující prostředky ke stávajícím možnostem mezigeneračního převodu majetku, nebo převodu majetku obecně. \\

Podle právní úpravy Občanského zákoníku je svěřenský fond institutem pro uchování a správu majetku. Svěřenský fond je však specifickou právní entitou, neboť nedisponuje právní osobností. Svěřenský fond ale není ani věcí, ačkoli věci\footnote{I hromadné.} může pojímat\footfullcite[Komentář k § 1448]{spacil_j_a_kolektiv_obcansky_2013}, nýbrž věci ve svěřenském fondu náležejí svému účelu\footfullcite{svejkovsky_jaroslav_sverenske_2018}.\\

Občanský zákoník tak pojímá svěřenský fond jako entitu, do které lze vyčlenit majetek a do jisté míry  jmění, například v případě obchodního závodu, který spravuje v souladu se statutem svěřenský správce za určitým účelem vymezeným ve statutu zakladatelem fondu. Občanský zákoník předpokládá existenci tří osob, které k funkci svěřenského fondu neodmyslitelně náležejí. Těmi jsou:\\

\begin{itemize}
\item svěřenský správce, jež vykonává práva spojená s majetkem ve svěřenském fondu a vykonává správu svěřenského fondu, samotné vlastnické právo mu ale nenáleží, neboť se \textit{strictum sensum} jedná o věc ničí, nikoli však opuštěnou, vlastnické právo tak nenáleží ani zakladateli, ani obmyšlenému\footnote{Srovnání § 1448 odst. 3, dále viz kapitola 7.8 Rozdílné pojetí vlastnictví.},
\item zakladatel, jež zakládá svěřenský fond a vyčleňuje do něj část svého majetku, nebo všechen svůj majetek\footnote{Majetek může vyčlenit i jiná osoba, ta se ale ovšem pro účely funkce svěřenského fondu nepovažuje za zakladatele, viz § 1468 OZ.}, po vyčlenění ztrácí veškerá dispoziční oprávnění,
\item obmyšlený, jež má dle statutu, zakladatelského právního jednání, právo na určité plnění ze svěřenského fondu\footnote{V rámci majetku přivlastněného účelu tvůrce této teorie Pierre Lepaulle existencí obmyšleného nepodmiňuje existenci trustu, srovnání POPOVICI A. Trust québeckého a českého práva:autonomní vlastnictví? in: TICHÝ Luboš, RONOVSKÁ Kateřina, KOCÍ Miloš, Trust a srovnatelné instituty v Evropě. V případě veřejně prospěšného svěřenského fondu je toto částečně naplněno, neboť se neurčuje obmyšlený přímo. Z veřejně prospěšného svěřenského fondu však může být plněno obmyšlenému, který bude určen na základě podmínek stanovených statutem svěřenského fondu. Takovéto určení ale musí naplňovat znaky veřejné prospěšnosti.}.	
\end{itemize}

Majetek jež lze vyčlenit do svěřenského fondu může mít formu věci hmotné i nehmotné, přičemž se může jednat jak o jakékoliv věci movité, tedy například i finanční prostředky\footfullcite[Strana 380]{dvorak_j_svestka_j_zuklinova_m_a_kol_obcanske_2016}, tak i věci nemovité.\\ 

Zajímavá je rovněž myšlenka související s zákonodárcem přijatou bipartitou svěřenského fondu místo tripartity trustu upravené v CCQ\footnote{Viz kapitola 4.3 Recepce právní úpravy Svěřenského fondu z právního řádu provincie Quebec.}. Tou je otázka existence svěřenského fondu bez obmyšleného. V Quebecu je toto možné v rámci \textit{private trustu} dle ustanovení článku 1268 CCQ, jež může být zřízen ve prospěch věci. Občanský zákoník však o možnost zřízení svěřenského fondu bez obmyšleného mlčí, naopak stanovuje nutnost nastavit minimálně způsob určení obmyšlených\footfullcite[§ 1452 odst. 2 a § 1457 odst. 2]{noauthor_zakon_2012}, důvodová zpráva o možnosti zřízení svěřenského fondu bez určení obmyšleného rovněž mlčí. Dle ustanovení § 1449 odst. 2 může být soukromý svěřenský fond zřízen ku prospěchu určité osoby, nebo na její památku. Článek 1268 CCQ stanoví, že účelem \textit{private trustu} je udržování, ochrana, nebo zachování určité věci s přímým, či nepřímým užitkem pro určitou osobu, na její památku, nebo pro jiný soukromý účel.\\

Možnost zřízení trustu v Quebecu bez jakéhokoliv určení obmyšleného pak již jen potvrzuje článek 1289 odst. 2, jež stanovuje následující: \textit{"In a private trust of which no person, even determinable or future, may be a beneficiary, the rights granted to the beneficiary under this subsection may be exercised by the Public Curator."}.\\

Pihera pak k tomuto říká následující: \textit{"Je-li účel svěřenského fondu vymezen ve prospěch neosobního účelu, pak takový svěřenský fond nemá obmyšleného."}\footfullcite[Komentář k § 1457]{spacil_j_a_kolektiv_obcansky_2013}. Příkladem budiž dle Pihery svěřenský fond zřízení nikoli ku prospěchu určité osoby, ale k určitému účelu, tedy například k udržování určité věci\footnote{Tamtéž.}.\\

Typizace takovéhoto fondu je k delší diskuzi, nicméně jednalo by se o svěřenský fond, který by měl za účel prospěch něčeho, co není osobou\footnote{Zde je nutné mít na paměti ustanovení § 17 odst. 2 OZ, které stanoví, že zřídí-li někdo právo nebo uloží-li povinnost tomu, co osobou není, přičte se právo nebo povinnost osobě, které podle povahy právního případu náleží.}. Nad touto otázkou by šlo uvažovat ve vztahu k diskrečnímu svěřenskému fondu. Tomuto bodu se však dále nebudu s ohledem na téma práce věnovat, neb účel mezigeneračního přenosu majetku stojí na převodu majetku z jedné generace na generaci druhou, je tedy zřejmé, že v takovémto případě na straně obmyšlených osoba stát bude.\\

Za zmínku stojí rovněž absence lhůty svěřenského správce pro určení obmyšleného, pokud toto právo nevyužije zakladatel, důvodová zpráva k tomuto říká následující: \textit{"Označení obmyšlených nebo způsobu, jakým budou vybráni, není nezbytné, protože při absenci takové klauzule ve statutu obmyšleného vybere svěřenský správce."}\footfullcite[Komentář k § 1448 až 1452]{Duvodovazprava}. Zákon i komentář tak počítají s určením obmyšleného, ze zákonné dikce však neplyne sama lhůta, ve které by tak měl svěřenský správce učinit.\\

% --------------------------------------------------------------------------
%V rámci svěřenského fondu funguje správce, který

%Český svěřenský fond byl nicméně přejat ve formě express trustu, jež


% Dopsat někde de lege lata
% Dopsat alterace a modifikace testamentu
% Nevystupuje jako subjekt ale spíše jako objekt právních vztahů
% Jelikož nelze svěřenský fond zařadit mezi právnické osoby, není možné na něj aplikovat ustanovení o právnických osobách v procesních, daňových, trestních a dalších předpisech. Zákonodárce proto musel tyto právní předpisy novelizovat s ohledem na úpravu svěřenských fondů. Důsledky zařazení svěřenského fondu do českého právního řádu tak lze pozorovat v právních předpisech napříč celým právním systémem

%Nově lze využít i formy svěřenského fondu, jehož základ tvoří osamostatněný majetek (resp. jmění), vyčleněný jeho zakladatelem k určitému účelu, svěřený do správy svěřenskému správci.
% --------------------------------------------------------------------------

\subsubsection{Právní kodifikace a charakteristika}

Jednotlivá ustanovení upravující svěřenský fond jsou v rámci současné právní úpravy zařazeny v části třetí OZ nazvané Absolutní majetková práva, speciálně v hlavě  \MakeUppercase{{\romannumeral 2}} - Věcná práva, Díle 6, Správa cizího majetku, Oddílu 4, Svěřenský fond v § 1448 až § 1474. Oddíl 6 je pak členěn do jednotlivých pododdílů, které níže ve zkratce představím a okomentuji. \\

Toto zařazení například dle Lucie Joskové není samozřejmé a hůře odůvodnitelné kvůli konvergenci prvků institutů právnické osoby, závazkového vztahu a vlastnictví v jednom, tedy ve svěřenském fondu\footfullcite[Strana 109]{lucie_joskova_sprava_2017}, přičemž ten ale sám o sobě není ani jedním z nich. Lze soudit, že takto byla úprava svěřenského fondu jakožto úprava speciální ke správě cizího majetku zařazena s ohledem na převzetí z CCQ, kde je trust obdobně zařazen pod \textit{Property}\footfullcite[Komentář k § 1448 až 1452]{Duvodovazprava}. Osobně se však domnívám, že takovéto zařazení je správné, dispozitivní regulace vztahů uvnitř svěřenského fondu není primárním účelem svěřenského fondu v českém právním prostředí, neboť to ho definovalo v rámci kogentních ustanovení\footnote{Srovnání § 978 OZ.} a oproti jiným trust-like institutům i poměrně dopodrobna\footfullcite{schmidt_k_trust_2016}, zásadním faktorem konstrukce českého svěřenského fondu tak je úplná oddělenost a nezávislost vyčleněného majetku tvořící insolvenční efekt\footfullcite[Strana 995]{sro_httpwwwpraguebestcz_podrobneji_nodate-1}\textsuperscript{,}\footfullcite[Strana 36]{skuhravy_jan_trust_2010}.\\

% --------------------------------------------------------------------------
%Doplnit pojem svěřenský fond, inspirace z práce Lukáš Kolenský Obchodní korporace v.svěřenský fondz pohledu věřitele

%Věci ve svěřenském fondu náležejí svému účelu\footfullcite{svejkovsky_jaroslav_sverenske_2018}.\\

%Úprava institutu svěřenského fondu a obecná úprava správy cizího majetku vstoupila v účinnost 1. ledna 2014 jako součást Nového Občanského zákoníku. Svěřenský fond je speciálně upraven v rámci ustanovení paragrafů 1448 až 1474, které jsou zařazeny do oddílu 4, jež je součástí dílu 6, Správa cizího majetku, spadající do části třetí, hlavy druhé, Věcná práva. \\
% --------------------------------------------------------------------------

Úprava svěřenského fondu jest úpravou speciální, která doplňuje úpravu generální obsaženou v ustanoveních § 1400 až 1447 OZ, Správa cizího majetku. Český zákonodárce se inspiroval, jak již je vysvětleno ve výkladu historie\footnote{Kapitola 4 Správa cizího majetku - historie, 4.3 Recepce právní úpravy Svěřenského fondu z právního řádu provincie Quebec.}, občanským zákoníkem z Quebecu ("CCQ"), z kterého ve své podstatě transplantoval právní úpravu svěřenského fondu, která, přes určité výkladové změny\footnote{Například v pojmech majetek a jmění, srovnání, Vít Zvánovec, Založení svěřenského fondu se zvláštním důrazem na oddělení majetku \textit{in}: Svěřenský fond a trust - Jejich fungování v mezinárodním srovnání, TICHÝ Luboš, RONOVSKÁ Kateřina, KOCÍ Miloš.}, do značné míry zachovala původní znění úpravy trustu tak, jak jej známe z CCQ, jak dokládá Svejkovský, Marek a kol. v díle Správa cizího majetku v novém občanském zákoníku v části Srovnání právní úpravy ObčZ a CCQ se stručným popisem rozdílů\footnote{Viz Kapitola 4.3.}.\\ 

Samotná definice svěřenského fondu chybí, zákonodárce svěřenský fond pouze legislativně vymezuje v § 1448 OZ a pouze definuje jeho konstrukci jako\footfullcite[Komentář k § 1448]{svestka_j_obcansky_2014-1}:\\

\begin{itemize}
\item majetek vyčleněný zakladatelem svěřenského fondu,
\item jež tvoří samostatně určený majetek svěřený svému účelu (bez právní subjektivity),
\item svěřen k následné správě majetku a fondu svěřenským správcem,
\item při povinnosti svěřenského správce majetek držet a spravovat ho dle statutu, přičemž v kombinaci s § 1456 svěřenskému správci náleží plná správa.	
\end{itemize}

\subsubsection{Účel svěřenského fondu}

Svěřenský fond je zakládán, respektive musí být založen, za určitým účelem, ten může být buď veřejně prospěšný, nebo soukromý, přičemž lze založit pouze za jedním z těchto účelů, nelze je kombinovat\footnote{To ovšem neznamená, že například v rámci diskrečního soukromého svěřenského fondu nemůže svěřenský správce v souladu s přáním zakladatele určovat obmyšlené a jim plnit z filantropických pohnutek.}. Účel pak definuje ve statutu zakladatel.\\

\textbf{Soukromý svěřenský fond} lze založit jako entitu sloužící prospěchu určité osoby, nebo na její památku, soukromý svěřenský fond lze také založit za účelem dosahování zisku, tedy například i provozování obchodního závodu.\\

\textbf{Veřejně prospěšný svěřenský fond} pak nelze založit pouze za účelem dosahování zisku, nebo provozování obchodního závodu. Zde bych v rámci textace ustanovení § 1449 odst. 3 akcentoval slovo hlavní, lze tak dovodit, že svěřenský fond může dosahovat zisku i provozovat obchodní závod, nesmí to však být jeho hlavním účelem, obdobně se lze domnívat, že prostředky takto získané lze vynaložit opět a jen pouze za veřejně prospěšným účelem. Domnívám se, že samotnou definici veřejně prospěšného účelu lze analogicky vyvodit z ustanovení § 146 OZ\footnote{Měl by tak bez dalšího svou činností přispívat k obecnému blahu.} o veřejné prospěšnosti právnických osob.\\

Charakteristiku jednotlivých dalších dílčích částí zahrnuji do následující kapitoly/pod následující kapitolu, neb jsou důležité při zkoumání tendence zákonodárce přibližovat svěřenský fond právnické osobě a samotnému porovnání právnické osoby se svěřenským fondem.\\

\subsubsection{Označení svěřenského fondu}

Takto vzniknuvší svěřenský fond má povinnost mít vlastní označení\footfullcite[§ 1450]{noauthor_zakon_2012}, jedná se o povinnou náležitost zápisu do evidence svěřenských fondů. To musí obsahovat slovní spojení "svěřenský fond", zakladatel má rovněž povinnost zahrnout do označení svěřenského fondu i jeho účel. Je pak celkem logické, proč zákonodárce požaduje v označení svěřenského fondu slovní spojení "svěřenský fond". Logiku lze spatřovat v odlišení od jiných osob a zjednodušení rozpoznání svěřenského fondu\footfullcite[Strana 135]{lucie_joskova_sprava_2017}, obdobně obchodní firma obchodní korporace také musí obsahovat firemní dodatek, tedy například s.r.o. či a.s.\footfullcite[§ 132 odst. 2, § 243 odst. 2]{noauthor_zakon_nodate_ZOK}. Smysl zahrnutí účelu pak spatřuji v jednoduchém odlišení jednotlivých svěřenských fondů jeden od druhého\footfullcite[Komentář k § 1450]{spacil_j_a_kolektiv_obcansky_2013}.\\

Volba označení je jinak zcela na zakladateli, ustanovení OZ nijak dále nespecifikují jak by mělo označení vypadat, či v jakém pořadí by mělo obsahovat jednotlivé části. Na rozdíl od obchodních korporací nelze slovní spojení "svěřenský fond" ničím nahradit, lze se domnívat, že označení svěřenského fondu krom tohoto musí splňovat ještě obecné požadavky na něj kladené v § 424 OZ\footnote{Obdobně bude aplikovatelná ochrana jména podle § 77 OZ, srovnání ŠVESTKA J., DVOŘÁK J., FIALA J. a kolektiv. Občanský zákoník. Komentář. Svazek III. § 976 - 1474, komentář k § 1450.}.\\

\subsubsection{Zakladatelské právní jednání}

Existenci samotného svěřenského fondu musí rovněž předcházet zakladatelské právní jednání, tedy smlouva \textit{inter vivos}, nebo pořízení pro případ smrti \textit{mortis causa}. Do tohoto právního jednání je třeba zařadit i statut svěřenského fondu.\\

O samotném procesu vzniku svěřenského fondu je pojednáno v kapitole 6.3 Založení svěřenského fondu, v této kapitole se věnuji obligatorní náležitosti svěřenského fondu, a to statutu a požadavkům na něj.\\

Zákon pro statut obligatorně stanovuje formu veřejné listiny\footfullcite[§ 1452 odst. 3]{noauthor_zakon_2012}, tedy notářský zápis\footfullcite[§ 62]{noauthor_zakon_nodate_NR}. Samotný statut pak konstituuje zakladatel svěřenského fondu, Kocí k tomuto dodává, že svěřenský fond mohl vzniknout i při absenci statutu, neboť zákon stanovil okamžik vzniku k úmrtí zakladatele, nebo přijetí správy\footfullcite[Komentář k § 1452]{svestka_j_obcansky_2014-1}, osobně se však domnívám, že po novelizaci Občanského zákoníku zakládající povinnost registrace svěřenských fondů do evidence svěřenských fondů tomu takto již nemůže být, neboť podle § 66 RZ se do sbírky listin zakládá i statut\footfullcite[§ 66, odst. 1]{noauthor_zakon_vr}, svěřenský fond \textit{inter vivos} tak nemůže bez statutu platně vzniknout, neboť zápis do evidence je v tomto případě konstitutivní povahy. Stejně tak jako svěřenský fond \textit{mortis causa}, neboť v případě absence statutu nebude možné po zakladateli jeho doplnění požadovat\footnote{Srovnání komentář k § 1452, Wolters Kluwer.}, tedy i přes fakt, že zápis do evidence je v případě svěřenského fondu \textit{mortis causa} pouze deklaratorní povahy, absence statutu by tak měla vést k nemožnosti plnit účel svěřenského fondu a soud by ho tak měl dle § 1473 odst. 1 zrušit\footnote{Tamtéž.}, popřípadě by takovýto svěřenský fond vůbec nevznikl, neboť zakladatelské právní jednání by bylo stiženo absolutní neplatností\footfullcite[Komentář k § 1452]{spacil_j_a_kolektiv_obcansky_2013}.\\

Účel samotného statutu je vytyčení základního regulačního rámce v rámci svěřenského fondu, tedy účelu, správy a plnění ze svěřenského fondu. Statut je založen jednostranným jednáním zakladatele, přičemž je ho nutné jako statut označit a v rámci zakladatelského právního jednání ho řádně odlišit od samotného vyčlenění majetku. Nic však nebrání tomu, aby statut samotný byl součástí zakladatelského právního jednání, tedy smlouvy, nebo pořízení pro případ smrti, vždy je však nutné dodržet formu veřejné listiny, toho lze dosáhnout i obsažením statutu v dědické smlouvě\footnote{Tamtéž, komentář k § 1452.}. Mimo pořízení pro případ smrti nejsou na samotné zakladatelské právní jednání, mimo statutu, kladeny žádné zvláštní nároky, zakladatel má tedy obsah smlouvy plně ve svých rukou.\\

Jednotlivé náležitosti statutu jsou stanoveny v § 1452 odst. 2 následovně:\\

\begin{enumerate}[label=\alph*)]
\item označení svěřenského fondu,
\item označení majetku, který tvoří svěřenský fond při jeho vzniku,
\item vymezení účelu svěřenského fondu,
\item podmínky pro plnění ze svěřenského fondu,
\item údaj o době trvání svěřenského fondu; není-li uvedena, platí, že fond byl zřízen na dobu neurčitou,
\item má-li být ze svěřenského fondu plněno určité osobě jako obmyšlenému, určení této osoby, nebo způsobu, jak bude obmyšlený určen, a
\item počet svěřenských správců a způsob jejich jednání.	\\
\end{enumerate}

K jednotlivým bodům\footnote{Pihera pak v komentáři k § 1452 v komentáři od C.H.Becka zmiňuje, že se jedná o náležitosti nezbytné k tomu, aby byl statut považován za platný, lze je tedy pojmou jako \textit{essencialia negotii}, ve statutu je rovněž možné uvést další \textit{naturalia negotii}, těm se však dále nebudu věnovat.}:\\

\textbf{Označení} jest nezbytné pro samotnou existenci svěřenského fondu, tvoří povinnou součást statutu. Povinně se zapisuje do evidence svěřenských fondů\footfullcite[§ 65 písmeno d)]{noauthor_zakon_vr}.\\

\textbf{Označení majetku}, který slouží účelu svěřenského fondu. Svejkovský tuto náležitost popisu následovně: \textit{"část označující majetek, který tvoří svěřenský fond při jeho vzniku, musí být natolik určité a srozumitelné, aby i ze statutu bylo patrné, který majetek do svěřenského fondu patří, a který nikoliv. Tato identifikace, která je komentována i k § 1448, musí být dostatečně určitá i srozumitelná pro jednoznačnou interpretaci, co do svěřenského fondu je z vlastnictví zakladatele vyčleněno, a co do svěřenského fondu nepatří."}\footfullcite[Strana 429]{jaroslav_svejkovsky_sprava_2015}. Jak již je poznamenáno v kapitole 6.1 do svěřenského fondu lze vyčlenit majetek, respektive věc v mnoha formách, musí ale splňovat zákonné náležitosti týkající se označení dané věci. V případě nemovitost tak musí obsahovat povinné náležitosti nemovitost označující dle katastrálního zákona\footfullcite[§ 8]{noauthor_zakon_nodate_KZ}. Správné označení je pak zásadní i kvůli výše zmíněnému insolvenčnímu efektu, předejde se tak potenciálním sporům s věřiteli při úpadku osoby zúčastněné v rámci svěřenského fondu způsobeného špatným označením majetku\footfullcite[Strana 136]{lucie_joskova_sprava_2017}.\\

\textbf{Vymezení účelu} je pro existenci svěřenského fondu rovněž zásadní, záleží také na tom, zda je svěřenský fond založen za veřejným nebo soukromým účelem. Je nutné ho vymezit co nejpřesněji, konkrétně a specificky. Správné vymezení účelu je zásadní pro správnou správu svěřenského fondu v souladu s přáním zakladatele, rovněž je důležité pro plnění ze svěřenského fondu, o to víc je pak přesné vymezení účelu svěřenského fondu důležité při konstruování veřejně prospěšného svěřenského fondu.\\

\textbf{Podmínky pro plnění} jsou plně v diskreci zakladatele svěřenského fondu, spolu s vymezením účelu a určení obmyšlených jsou určující pro posouzení jednotlivých nároků za svěřenským fondem těch, kterým má být plněno. Svěřenský správce k nim musí vázat svoji povinnost vykonávat správu s péčí řádného hospodáře\footfullcite[Strana 430]{svejkovsky_jaroslav_sverenske_2018}. Plněno může být ze samotného kmenového majetku z fondu, nebo mohou být poskytovány plody a užitky ze svěřenského fondu (zejména zisk). Podmínky pro plnění pak mohou být jakékoliv povahy\footfullcite[Strana 138]{lucie_joskova_sprava_2017}, domnívám se však, že ke zjevně nemorální, nesplnitelné, obtěžující podmínce či podmínce odporující veřejnému pořádku se nebude přihlížet\footnote{Podobně například u vedlejších doložek v závěti.}. Podmínky také může v souladu se statutem u diskrečního svěřenského fondu určit svěřenský správce\footfullcite[§ 1457]{noauthor_zakon_2012}.\\

\textbf{Trvání}, lze dovodit, že pro vznik svěřenského fondu nemusí být uvedeno. V případě jeho opomenutí platí, že jse svěřenský fond zřízen na dobu neurčitou. Obecně ho lze stanovit časově, samozřejmě splněním účelu, či jakoukoliv jinou rozvazovací podmínkou. Veřejný svěřenský fond se dá zřídit na neomezenou dobu, v případě opomenutí určení trvání u soukromého svěřenského fondu je jeho trvání omezené v § 1460 OZ na 100 let, respektive právo na plnění vzniká maximálně dobou uplynutí 100 let od založení svěřenského fondu\footfullcite[Strana 278]{pihera_v_nejpodivnejsi_2012}.\\

\textbf{Určení obmyšlených}, k tomuto blíže Kocí dodává následující: \textit{"3. Zákon blíže nespecifikuje, kdo má způsobilost být obmyšleným, neboť se vychází z toho, že je zcela na uvážení zakladatele, koho bude považovat za způsobilého a koho nikoliv. Může se jednat o člověka i o právnickou osobu včetně osob nenarozených nebo právnických osob ještě nezaložených (§ 1464), kdy hovoříme o tzv. eventuálních obmyšlených (v Québecu bénnéficiaires éventuels)."}\footfullcite[Komentář k § 1457]{svestka_j_obcansky_2014-1}. Je tak zřejmé, že obmyšleným může být i sám zakladatel, počet obmyšlených rovněž není limitován. Výslovné právo k určení obmyšlených, či posloupnosti obmyšlených, má sám zakladatel, ten však může přenechat určení obmyšlených i na svěřenském správci. Zakladatel rovněž může stanovit okruh potenciálních obmyšlených, nebo způsob, jakým mají být vybráni\footnote{Srovnání § 1457.}. Zákon zakladateli rovněž ponechává možnost stanovit obmyšleného i po založení svěřenského fondu, pokud toto statut výslovně nevylučuje\footnote{Srovnání Kocí, komentář Wolters Kluwer, § 1457.}. Je však nutné upozornit na to, že v případě, kdy zakladatel ve statutu minimálně neurčí okruh potenciálních obmyšlených, nebo návod, podle kterého by je mohl svěřenský správce jmenovat, tak svěřenský správce nesmí žádné obmyšlené jmenovat, toto představuje u soukromého svěřenského fondu problém, a může vést až k jeho zrušení soudem podle ustanovení § 1469 odst. 1 ustanovení § 1471 OZ\footnote{Tamtéž.}. V případě veřejně prospěšného svěřenského fondu však tato podmínka zakladatelem být splněna nemusí, protože zákon kompetenci k určení obmyšlených presumuje u svěřenského správce\footfullcite[Komentář k § 1469]{spacil_j_a_kolektiv_obcansky_2013}. U diskrečního svěřenského fondu, tedy v případě, kdy zakladatel správci právo na určení obmyšlených výslovně určí, pak tato podmínka rovněž není aplikovatelná\footnote{Kocí, komentář Wolters Kluwer, § 1457.}.\\

\textbf{Svěřenský správce} musí rovněž být uveden ve statutu, je ponecháno opět zcela na zakladateli, koho určí. Zakladatel rovněž není omezen v počtu svěřenských správců, těch tak může určit libovolný počet. Ve statutu pak musí být uveden způsob jejich jednání. Výše zmíněné se následně zapíše do evidence svěřenských fondů.\\

% --------------------------------------------------------------------------
 %do nějž
 
 %\subsubsection{Správa svěřenského fondu}

%Úprava svěřenského fondu jest úpravou speciální, která doplňuje úpravu generální obsaženou v ustanoveních 1400 až 1447, Správa cizího majetku. Český zákonodárce se inspiropval, jak již je vysvětleno v kapitole Historie \footnote{Správa cizího majetku - historie; Recepce právní úpravy svěřenských fondů}, Občanským zákoníkem z Quebecu ("CCQ"), z kterého ve své podstatě transplantoval právní úpravu svěřenského fondu, která, přes určité výkladové změny \footnote{například v pojmech majetek a jmění, viz Svěřenský fond a trust - Jejich fungování v mezinárodním srovnání od Luboše Tichého}, do značné míry zachovala původní znění úpravy trustu tak, jak jej známe z CCQ, jak dokládá Svejkovský, Marek a kol. v díle Správa cizího majetku v novém občanském zákoníku v části Srovnání právní úpravy ObčZ a CCQ se stručným popisem rozdílů.\\ 

%\newpage
%\thispagestyle{smallertextinheader}

%KOMENT-------

%Důvod, pro který si Český zákonodárce vybral jako zdroj inspirace pro zavedení trustu do našeho právního řádu práve CCQ je ten, že Quebecký Občanský zákoník si zachoval výrazný charakter kontinentálního práva\footfullcite{Duvodovazprava}, ačkoli je obklopena státy s právním systémem Common law\footnote{Výjimku představují například DOPLNIT}.\\

%Výhodou je i vysoká flexibilita svěřenského fondu. Ve statutu lze do značné míry podle potřeby upravit vnitřní fungování svěřenského fondu podle představ zakladatele - např. počet svěřenských správců a obmyšlených, a způsob jejich jmenování, pravidla výplaty plnění obmyšleným, pravidla správy apod.\\

%Tento paragraf bude ještě potřeba přepsat, aby pasoval do této kapitoly.

%Svěřenský fond je tedy jakousi quazi právní entitou bez právní subjektivity, sestávající se z autonomního majetku bez vlastníka vyčleněného svému účelu\footnote{Jedná se o něco, co není v českém právu běžné, z toho důvodu existuje i mnoho výkladových rozporů.}, je spravován svěřenským správcem podle instrukcí zakladatele, které jsou součástí zakladatelského právního jednání, respektive statutu.\\

%Důležité je zmínit dobu, na kterou svěřenský fond může být zřízen. Pro svěřenské fondy založené za soukromým účelem je tato doba omezená. Pro svěřenské fondy založené za veřejným účelem nicméně není stanovena maximální doba. Při porovnání obou forem založení tak můžeme \textit{a contrario} dospět k tomu, že svěřenský fond založený za veřejným účelem není omezen maximální dobou trvání a může tak být založen jako takzvaná perpetuita.\\

%KOMENT---------

%Doplnit zdroj z OZ, tedy paragraf.

%PŘEPSAT
%Nabízí se otázka, co lze rozumět pod pojmem majetek. Občanský zákoník chápe majetek jako souhrn všeho, co osobě patří. Přitom musí být pro účely vyčlenění do svěřenského fondu předmětem právního obchodu. Tato definice nevylučuje ani pohledávky a neexistuje důvod, proč by nemohly být předmětem vyčleňovaného majetku. To se týká také pohledávek zakladatele vůči fondu. Zakladatel proto může vyčlenit majetek tak, že zřídí ve prospěch fondu určité právo, jímž se třeba sám zavazuje. Naproti tomu dluhy zakladatele rozhodně mezi majetek počítat nelze. Existuje však možnost, že bude vyčleněný majetek zatížen nějakými závazky. V tom případě se mohou společně s majetkem přesunout do svěřenského fondu také dluhy. Vyčleněným majetem tedy mohou být hmotné věci movité (pila) i nemovité (pozemek) i nehmotné (pohledávka). ZDROJ: https://books.google.cz/books?id=aeq4DwAAQBAJ&pg=PT164&lpg=PT164&dq=sv%C4%9B%C5%99ensk%C3%BD+fond+vy%C4%8Dlen%C4%9Bn%C3%AD+dluh%C5%AF&source=bl&ots=ah6SoELD-z&sig=ACfU3U1vKrdxRriyDfwSDEsHE9npgYfwpg&hl=en&sa=X&ved=2ahUKEwjul--0wvnpAhX2BGMBHeXQC74Q6AEwBHoECAoQAQ#v=onepage&q&f=false

%\subsubsection{Svěřenský fond a právní subjektivita}

% Dobrý zdroj je následující: https://advokatnidenik.cz/2019/05/23/sverenske-fondy-nejsou-obchodni-korporace-fikce-podle-zakona-o-skutecnych-majitelich/

%https://advokatnidenik.cz/2020/04/04/vlastnictvi-bez-vlastnika-v-ceskem-sverenskem-fondu-skica-na-margo-trustu-v-ceskem-pravu/

%Dále je ještě jeden zdroj v bakalářské práci kterou lze dohledat v repu, pak má také Michal Skuhrovec, ale od tamtud to nechci úplně celé opisovat.

%https://www.fulsoft.cz/33/pojem-a-vznik-sverenskeho-fondu-uniqueidmRRWSbk196FNf8-jVUh4EnhEI8yvqwHRaEpgGCfo3bkMO6VlftQMFg/

%Důležitou součástí definice svěřenského fondu je fakt, že svěřenský fond je koncipován jako oddělený majetek, který je svěřen svému účelu, jež nemá vlastníka a nemá ani právní subjektivitu. Přes jasnou definici, kterou obsahuje Občanský zákoník, důvodová zpráva a komentář k občanskému zákoníku a i samotný zdroj inspirace právní úpravy svěřenského fondu, tedy CCQ, není však na škodu podrobit svěřenský fond jisté komparaci s právnickou osobou. Není totiž tajemstvím, že určité právní normy na svěřenský fond pohlížejí jako na právnickou osobu\footfullcite[§17]{noauthor_zakon_zdp}\textsuperscript{,}\footfullcite[§4b]{noauthor_zakon_zdph} a právě tato tendence k subjektivizaci svěřenského fondu by se mohla projevit i v jiných aspektech svěřenského fondu. DOPLNIT\\

%Brilantní náhled na toto téma již vyslovil Michal Skuhrovec ve své diplomové práci, která z tohoto hlediska poskytuje vhodný náhled na tuto otázku a také patřičný úvod do této problematiky.\\

%Doplnit do začátku práce, že jsem se touto prací hodně inspiroval, ale myslím si, že neposkytuje odpovědi na všechny otázky a v určitých částech není aktuální, nebo odpovídající, doplnil bych jí tedy o svůj pohled, nicméně některé pohledy autora jsou výtečné a já se s nimi jen ztotožňuji, stejně tak jako se s nimi ztotožňují i pohledy právních expertů. Je nicméně potřeba jeho pohled rozšířit i o nové poznatky v souvislosti s institutem svěřenského fondu

%Na tomto místě je vhodné uvést dvě základní premisy související se svěřenským fondem, které budou v bližším detailu rozvedeny dále. Povaha svěřenského fondu předně vylučuje, aby měl svěřenský fond právní osobnost (svěřenský fond tedy není právnickou osobou) a stejně tak vylučuje, aby byl svěřenský fond věcí v právním slova smyslu. Uvedenému přirozeně neodporují závěry některých autorů, kteří nevylučují, aby měl majetek vložený do svěřenského fondu povahu věci hromadné. 2 HORN, K.: Podrobněji k svěřenskému fondu. Ad notam, č. 6, ročník 2014, s. 1.

%Lze nicméně uvést, že svěřenský fond představuje naprosto ojedinělou a zvláštní entitu, která se skládá z (bývalého) majetku zakladatele a v souvislosti s kterou je prostřednictvím činnosti správce svěřenského fondu naplňován zakladatelem určený cíl. Přitom základním cílem společným pro všechny svěřenské fondy je nepochybně ochrana majetku zejména před nežádoucími dispozicemi 260 260 AKONTinfo, Trust jako nejúčinnější nástroj ochrany majetku. [online]. 2010. On-line zdroj: http://www.akont.cz/cz/375.trust-jako-nejucinnejsi-nastroj-ochrany-majetku.

%PŘEPSAT
%Pihera v tomto konextu vytahuje na povrch diskuzi pocházející z germánského prostoru, kdy Jhering kritizoval koncepci samostatného jmění jako contradictio in adiecto, kdy svou kritiku zakládá na "subjektlose subjektive Rechte", neboli existenci subjektivního práva bez subjektu. Takové by mělo představovat nebezpečí popření ústředního postavení člověka v centru soukromého práva, jak dovozuje Becker. Nicméně, lze oběma autorům přisvědčit, že jeho uznání by při racionálním zhodnocení všech okolností nemělo bez dalšího vést k dotváření "Rechtlosen subjekten". PIHERA NEJPODIVNĚJŠÍ ZVÍŘE V LESE Komentář k OZ pak zároveň dodává, že určité samostatné majetky již v právním řádu 14 známe a vyrovnáváme se s nimi již tradičně, například ležící pozůstalost (hereditas iacens) v ABGB. Dalším právem bez subjektu jsou například práva nenarozeného dítěte.\\

%PŘEPSAT
%Rozdíl mezi úpravou CCQ a OZ pak tkví dále v samotné definici majetku a je otázkou, zda zde český zákonodárce udělal quebeckému trustu dobrou službu co do přípravy "podhoubí". CCQ totiž v § 915 stanoví: "Property belongs to persons or to the State or, in certain cases, is appropriated to a purpose." Majetek patří buď osobám, státu a nebo, v určitých případech, je vyčleněn účelu. Oproti tomu §495 ani §1011 nic takového nestanoví (viz výše). Je tak nutno vycházet z toho, že §1448 OZ je ustanovením speciálním a tvoří základnu teoretické koncepce svěřenského fondu. PIHERA NEJPODIVNĚJŠÍ ZVÍŘE V LESE Lze tak uzavřít, že svěřenský fond je entitou bez právní subjektivity, která je majetkem se samostatným účelovým určením. S ohledem na výše uvedené, podpořeno argumenty zaznívající v diskuzi odborné veřejnosti, se lze přiklonit na stranu zastánců přijaté koncepce v duchu "patrimoine d`affectation, neb zjevně skutečně představuje vhodný most mezi dvěma právními kulturami, podepřený zkušenostmi z Quebecu. Smith\\

%\subsection{Evidence svěřenských fondů}

%Přes mnohem striktnější úpravu svěřenského fondu, než kterou lze najít u jiných trust-like institutů nelze ji upřít, že zákonodárce omezil možnosti potenciálního zneužití svěřenského fondu ve smyslu vysloveného například ministerstvem financí\footnote{Viz níže.}. Zakotvení trustu v české právní úpravě ve formě svěřenského fondu s povinnými náležitostmi ve formě statutu, zápisu do evidence a dalších obdobných požadavků minimalizující zmenšující rizika jeho použití k nekalým finančním praktikám kvůli netransparentní konstrukci. I přes toto je mnohem snadnější užít k účelům transferu majetku či obchodním účelům svěřenský fond, než je zakládat jinou právnickou osobu.\\

%\subsubsection{Novelizace právní úpravy - zápis svěřenských fondů do rejstříku}

%Majetek ve svěřenském fondu tedy není ve vlastníctví nikoho, před novelizací Občanského zákoníku pak ve veřejných seznamech nebo evidencích mohla být dohledána pouze osoba svěřenského správce. Samotné svěřenské fondy se pak neměli povinnost se jakkoliv registrovat nebo evidovat, v souvislosti s tím pak nebyl veřejně dostupný ani statut svěřenského fondu, který přeci musel mít formu veřejné listiny, ale v souvislosti s absencí povinnosti evidence svěřenských fondů nebyl nikde zveřejněn. Absence jakékoliv evidence, vedla k anonimitě zakladatele, obmyšlených a přispívala k celkové netransparantnosti vlastnické struktury a majetku vyčleněného do svěřenského fondu.\\

%V souvislosti s tímto panoval oprávněný strach, že svěřenské fondy by mohli být využity k legalizaci výnosů z trestné činnosti, daňovým únikům nebo být jinak zneužity. V souvislosti s touto obavou hned tři instituce, kterými byli Ministerstvo financí, Ministerstvo vnitra a Vrchní státní zastupitelství volali po novelizace právní úpravy svěřenského fondu a povinné evidenci. Důvodová zpráva k novele pak stanovuje přímo následující: \textit{"V reakci na naléhavé připomínky Ministerstva financí, Ministerstva vnitra a Vrchního státního zastupitelství v Praze se navrhuje zavedení evidence svěřenských fondů z důvodu jejich netransparentní vlastnické struktury a vysokého potenciálu zneužití především k legalizaci výnosů z trestné činnosti."}\footfullcite{noauthor_4602016_nodate}.\\

%Zákon č. 33/2020 Sb. novela, se svěřenských fondů téměř vůbec nedotýká, pro tuto práci není důležitá, neboť se měni jen určité ustanovení ve vztahu k notářskému zápisu.

%Rozepsaná novelizace svěřenský fondů z roku 2018, zhodnotit, zda přispělo ke zmírnění možností využití svěřenských fondů jako prostředků k legalizaci výnosu z trestných činností.\\

%Poměrně zásaní změny v souvislosti se zakládáním svěřenských fondů tak v souvislosti s tímto přinesla novelizace č.z. 460/2006, která s sebou přinesla nutnost zápisu svěřenského fondu do evidence s tím, že tento zápis je konstitutivní. Svěřenský fond tedy vzniká tímto zápisem. Toto se nicméně dotklo pouze svěřenského fondu založeného mezi živými. Pro svěřenský fond založený pro případ smrti zákonodárce stanovil v § 1451 odstavci 3 následující:\\

%\textit{"Byl-li však svěřesnký fond zřízen pořízením pro případ smrti, vznikne smrtí zůstavitele. Do evidence svěřenských fondů se zapíše po svém vzniku."}\\

%Důvod, pro který zákonodárce dopřál toto, do jisté míry zvýhodňující, postavení způsobu založení svěřenského fondu \textit{mortis causa}, je v důvodové zprávě objasněno následovně:\\

%\textit{"Odlišný princip deklaratorního zápisu se má uplatnit u svěřenských fondů, které byly zřízeny pořízením pro případ smrti. Tento princip lépe vyhovuje fundamentu testovací volnosti a respektu k vůli zůstavitele."}
% --------------------------------------------------------------------------

\newpage

% --------------------------------------------------------------------------
%\thispagestyle{smallertextinheader}
% --------------------------------------------------------------------------

\subsection{Založení svěřenského fondu a jeho vznik}

S ohledem na téma práce, tedy mezigenerační převod majetku, je otázka založení svěřenského fondu, ať již mezi živými, tedy \textit{inter vivos}, či pořízením pro případ smrti \textit{mortis causa}, nebo velice specifickou možností založení svěřenského fondu \textit{inter vivos} s odkládací podmínkou smrti zakladatele a dalších náležitostí s tímto spojených, integrální součástí této práce a je nezbytné k dalšímu zkoumání jednotlivých možností, které nám s přijetím tohoto institutu český zákonodárce vložil do rukou a obecně i celé problematiky spojené se svěřenským fondem. \\

Pro tuto kapitolu, a následné návrhy na změny, je v první řadě podstatné vymezit,  v jakých fázích vzniká svěřenský fond, a tyto následně porovnat s průběhem pozůstalostního řízení upraveného v rámci dědického práva\footnote{Vymezení samotného vzniku tímto postupem lze dosáhnout vhodného startovního bodu pro zodpovězení hypotézy stanovené v úvodu práce, tedy zodpovězení otázky, zda mají dědicové v jednotlivých fázích možnosti bránit se proti jejich opomenutí ze strany zůstavitele a brát se o svá práva.} \\ 

Z počátku je vhodné předeslat, že považuji založení a vznik svěřenského fondu za dvě různé věci. Založením svěřenského fondu považuji všechny úkony před samotným vznikem svěřenského fondu, za vznik pak považuji okamžik dle § 1451 OZ.\\

Dle mého názoru, který potvrzuje Miloš Kocí v rámci komentáře k úpravě svěřenského fondu, lze hovořit o 3 způsobech vzniku svěřenského fondu. Těmito způsoby jsou, založení svěřenského fondu \textit{inter vivos}, \textit{mortis causa} a potenciální souběh těchto založení ve formě založení svěřenského fondu \textit{inter vivos} s odkládací podmínkou, dle § 584 až 549 OZ, spočívající ve smrti zůstavitele.\footfullcite[Komentář k § 1448]{svestka_j_obcansky_2014-1} \\

% --------------------------------------------------------------------------
%\newpage

%\thispagestyle{smallertextinheader}
% --------------------------------------------------------------------------

Vizuálně by se tedy průběh komparace napříč jednotlivými způsoby založení dal vyobrazit následovně, kde střed osy představuje čas úmrtí zůstavitele:

\begin{figure}[h]
\centering
\includegraphics[width=27cm,height=9cm,keepaspectratio]{Vznik_sverenskeho_fondu.png}
\caption{Vznik svěřenského fondu}
\label{fig:komparace}
\end{figure}

% --------------------------------------------------------------------------
%---Koment\\
%Doplnit další informace.
%Popsat dědické právo - co jde do dědictví, co jsou nepominutelní dědicové, jak funguje ědické řízení atp.
%Vysvětlit rozdíl mezi majetkem a jměním.\\
%---\\

%V případě založení svěřenského fondu mezi živými s odkládací podmínkou smrti zakladatele se ve své podstatě jedná o založení \textit{mortis causa} formou \textit{inter vivos}, s ohledem, že se nejedná o založení svěřenského fondu pořízením pro případ smrti je zřejmé, že v tomto případě nelze uvažovat ani nad ochranou dle §1492, pro ilustraci je možné porovnat tento způsob založení s další formou mezigeneračního předání majetku a to darováním pro případ smrti. Podklad pro toto tvrzení poskytují například autoři článku Darování pro případ smrti v Časopisu pro právní vědu a praxi Adam Talanda a Iveta Talandová. Ti přímo říkají: \textit{"Darování pro případ smrti je při splnění všech zákonných podmínek obligačním vztahem mezi dárcem a obdarovaným a neřídí se ustanoveními dědického práva. Zákon tak umožňuje užitím jednoho zákonného institutu obcházet soubor jiných zákonných institutů. Především se jedná o obcházení ustanovení omezujících zůstavitele v nakládání s jeho majetkem, která slouží k ochraně nepominutelných dědiců."}\footfullcite{noauthor_casopis_2015}.\\
% --------------------------------------------------------------------------






Ať již se budeme bavit o jakémkoliv způsobu založení svěřenského fondu, tak zákon klade na zakladatele v paragrafech 1448 - 1474 Občanského zákoníku následující čtyři povinnosti:

\begin{enumerate}
\item Označení a vyčlenění majetku k určitému účelu, tedy zakladatelské právní jednání,
\item statut\footnote{Ve formě veřejné listiny.},
\item přijetí správy svěřenského fondu ze strany svěřenského správce,
\item zápis svěřenského fondu do evidence.
\end{enumerate}

Samotné zakladatelské právní jednání pak může být učiněno ve formě smlouvy, nebo pořízením pro případ smrti, svěřenský fond tak nelze založit jednostranným právním jednáním s účinky vzniku ještě za života zakladatele\footfullcite[Komentář k § 1448]{svestka_j_obcansky_2014-1}.\\

Právní jednání rovněž nesmí být zdánlivé a musí být platné\footfullcite[Strana 9]{svejkovsky_jaroslav_sverenske_2018}.\\

Co se však bodu 4 týče, jde nutné rozlišovat mezi založením svěřenského fondu \textit{inter vivos} a \textit{mortis causa}, neboť v obou případech se bod číslo 4 liší. V případě založení \textit{inter vivos} je zápis svěřenského fondu konstitutivní, v případě založení \textit{mortis causa} je pouze deklaratorní, svěřenský fond totiž vznikne smrtí zůstavitele, o tomto však blíže u jednotlivých typů svěřenských fondů.\\

Obecné vysvětlení k pojmu a založení svěřenského, tedy k paragrafům 1448 až 1452 Občanského zákoníku, fondu je obsaženo již v důvodové zprávě, respektive v její části vtahující se ke svěřenským fondům.\\

\textit{"Svěřenský fond může být zřízen jednak bezúplatně rozhodnutím zakladatele věnovat část svého majetku určitému účelu, zřízení svěřenského fondu však může být sjednáno i za úplatu nebo jiné protiplnění (např. od osoby, která se má stát obmyšleným, ale i od osoby jiné). Podle toho se pak odvíjí i právní postavení obmyšleného, jak to reflektují následující ustanovení osnovy. Osnova rozlišuje svěřenské fondy podle účelového určení na veřejně prospěšné (zřízené k naplňování účelů kulturních, vzdělávacích, vědeckých, náboženských nebo podobných) a soukromé. Důvodem vzniku svěřenského fondu může být zákon, pravidelně však smlouva nebo ustanovení závěti. Protože svěřenskému fondu musí být ustaven svěřenský správce, vyžaduje § 1448 souhlas osoby označené za svěřenského správce se svým ustavením. I to je tedy právní podmínka vzniku svěřenského fondu (§ 1451). Svěřenský fond se spravuje statutem obsaženým ve smlouvě, závěti či vydaným samostatně."}\footfullcite{prof_dr_judr_karel_elias_a_kolektiv_novy_2012}\\

To hovoří i o vzniku svěřenského fondu ze zákonného důvodu, v občanském zákoníku však s touto možností zákonodárce nepočítá. Dále pak s touto částí důvodové zprávy přímo souvisí § 1448 OZ, ve kterém zákonodárce uvádí následující:\\

\noindent\textit{"(1) Svěřenský fond se vytváří vyčleněním majetku z vlastnictví zakladatele tak, že ten svěří správci majetek k určitému účelu smlouvou nebo pořízením pro případ smrti a svěřenský správce se zaváže tento majetek držet a spravovat.\\
(2) Vznikem svěřenského fondu vzniká oddělené a nezávislé vlastnictví vyčleněného majetku a svěřenský správce je povinen ujmout se tohoto majetku a jeho správy."}\\

A § 1451 OZ:\\

\noindent\textit{"(1) Svěřenský fond je zřízen, když svěřenský správce přijme pověření k jeho správě; je-li svěřenských správců více, postačí, pokud pověření přijme alespoň jeden z nich.\\
(2) Svěřenský fond vzniká dnem zápisu do evidence svěřenských fondů.\\
(3) Byl-li však svěřenský fond zřízen pořízením pro případ smrti, vznikne smrtí zůstavitele. Do evidence svěřenských fondů se zapíše po svém vzniku."}\\

Jak již jsem zmínil výše, v uvedeném ustanovení § 1448 OZ tedy zákonodárce nepočítá s tím, že důvodem vzniku svěřenského fondu by mohl být zákon, nýbrž umožňuje pouze explicitní založení svěřenského fondu dvěma způsoby a to \textit{inter vivos} a \textit{mortis causa}, to znamená, že svěřenský fond nemůže být založen rozhodnutím soudu tak, jako tomu je v jiných právních řádech. Jedním z těchto právních řádů je například právní řád používán v jednotlivých státech USA, který umožňuje i takzvané implicitní založení svěřenského fondu, respektive trustu, tedy založení na základě rozhodnutí soudu, nebo i ze zákona, pokud existuje zákonné zmocnění, například advokátní a notářské úschovy\footfullcite{claxton_john_b_studies_2005}. Další možností, která není explicitně uvedena v tomto paragrafu, ale je možné ji využít k založení svěřenského fondu, je založení svěřenského fondu mezi živými s odkládací podmínkou.\\

Co se samotného průběhu založení a vzniku svěřenského fondu týče, přikláním se stejně, jako komentářová literatura k multifázovému pojetí vzniku svěřenského fondu\footfullcite[Komentář k § 1448]{svestka_j_obcansky_2014-1}. Kocí v komentáři hovoří o možnost vzniku svěřenského fondu jednofázově a dvoufázově. Po zavedení povinné evidence svěřenských fondů bych však hovořil o zakládání minimálně dvoufázovém\footfullcite[Strana 151]{lucie_joskova_sprava_2017}\textsuperscript{,}\footnote{V případě založení svěřenského fondu mezi živými již nepostačí samotné přijetí pověření ke správě ze strany svěřenského správce, ale je třeba svěřenský fond pro vznik zapsat i do evidence. V případě svěřenského fondu zřízeného pro případ smrti je obdobně nutné právní jednání ve formě pořízení pro případ smrti a právní událost ve formě úmrtí zakladatele.}.\\

% --------------------------------------------------------------------------
%Svěřenský fond vzniká vyčleněním majetku z vlastnictví zakladatele. Obecným východiskem zde je, že v zásadě nemůže existovat svěřenský fond bez majetku nebo bez svěřenského správce. Za majetek se přitom považuje i pohledávka. Tedy asi když vzniká svěřenský fond mortis causa smrtí zůstavitele, tak do něho již musí být vložen majetek. Není tedy asi možné, aby byl majetek součástí pozůstalosti a svěřenskému fondu ho přiznal teprve až soud. To že však majetek není součástí pozůstalosti neznamená, že svěřenský fond není účastníkem dědického řízení. Dle zákona o zvláštních řízeních soudních jím totiž je, zastupuje ho zde svěřenský správce, který ale sám o sobě dědicem být nemůže, protože za dědice nebyl povolán. Svěřenský fond tedy v tomto případě má jekési quazi postavení právní subjektivity, tak jako tomu je například třeba i v případě daňového práva.

%Účelem se rozumí veřejný nebo soukromý.
% --------------------------------------------------------------------------

\subsubsection{Vznik svěřenského fondu \textit{inter vivos}}

S ohledem na historické konotace svěřenských fondů popsané výše a velkou flexibilitu tohoto institutu skutečně není divu, že český zákonodárce umožnil založení svěřenského fondu jak mezi živými, tak i pro případ smrti. Začnu s popisem založení svěřenského fondu mezi živými, neboť tento způsob založení je méně problematický než založení svěřenského fondu pro případ smrti a přináší obdobné možnosti mezigeneračního předání majetku.\\

V tomto rámci je podstatné zmínit, že se rozlišuje mezi vytvořením a vznikem svěřenského fondu, celkový proces konstituování svěřenského fondu, tedy od okamžiku samotného rozhodnutí vytvořit svěřenský fond, až k samotnému vzniku, má tedy více fází. Vytvořením se rozumí jednání jednoho, nebo i více zakladatelů\footfullcite{spacil_j_a_kolektiv_obcansky_2013}\textsuperscript{,}\footnote{S ohledem na fakt, že samotná právní úprava svěřenského fondu v §1452 odkazuje na založení nadace, je vhodné i jiné ustanovení, která se týkají nadací, využít analogicky ve vztahu ke svěřenským fondům. Jedním z těchto paragrafů je paragraf  309 odstavec 3, který stanoví, že pokud na straně zakladatelů stojí více osob, považují se obdobně za jednoho zakladatele.}, kteří vyčlení část svého majetku, popřípadě celý svůj majetek a tento poté převedou do jiného samostatně určeného jmění, tedy svěřenského fondu, který je zřízen za určitým účelem. Toto jednání zahrnuje v první řadě řádnou nabídku, tedy ofertu, vůči jiné osobě, která by se měla stát potenciálním správcem\footfullcite{jaroslav_svejkovsky_sprava_2015}.\\

% --------------------------------------------------------------------------
%Doplnit nějaké zdroje k založení svěřenského fondu ve více fázích
% --------------------------------------------------------------------------

Jak již je z podstaty svěřenského fondu založeného za života zakladatele zřejmé, jedná se o minimálně o dvoustranné právní jednání, přičemž samotná akceptace této oferty\footnote{Návrh na uzavření smlouvy je dále upraven v §1731 a dalších v Občanském zákoníku} vede dle §1448 a souvisejícího §1451 Občanského zákoníku jednáním vedoucím ke zřízení, tedy vytvoření, svěřenského fondu na základě takto uzavřené smlouvy, kterou je možné uzavřít jak s, tak i bez protiplnění. Tato oferta a akceptace zahrnující uzavření smlouvy o správě by se dala považovat za první fázi založení svěřenského fondu.\\ 

Občanský zákoník však již nestanoví, s kým může být daná smlouva uzavřena ačkoli předpokládá eventuální účast svěřenského správce v § 1448 odst. 1 OZ, ten je taky nezbytným článkem při funkci svěřenského fondu. Pihera k tomuto dodává: \textit{"Smlouva, jíž se vytváří svěřenský fond, bude zpravidla uzavírána mezi zakladatelem a svěřenským správcem. Zákon však nestanoví požadavek na to, s kým musí zakladatel smlouvu uzavřít. Protistranou nebude svěřenský správce především v případech, kdy zakladatel určí jiný způsob jeho jmenování. Stejně si lze představit, že bude svěřenský fond zřizován v době, kdy nebude budoucí svěřenský správce přítomen. Svěřenský fond může být též vytvářen na základě smlouvy mezi zakladatelem a obmyšleným, popřípadě jinou osobou (např. osobou pověřenou dohledem nad fondem podle § 1463)."}\footfullcite[Komentář k § 1448]{spacil_j_a_kolektiv_obcansky_2013}.\\

Kocí pak rovněž dodává: \textit{"Není tak vyloučeno, aby kromě zakladatele byli jejími stranami rovněž svěřenský správce či obmyšlený anebo oba zároveň. Lze si představit též takové uspořádání, kdy zakladatel smlouvu uzavře s osobou třetí a svěřenského správce ani obmyšleného do smluvního rámce nezahrne. Pak by se tedy mohlo jednat o smlouvu ve prospěch třetí osoby (pactum in favorem tertii) ve smyslu § 1767 až 1768, a pokud by se spolukontrahent např. zavázal, že se u zamýšleného svěřenského správce přimluví, aby se ujal správy vyčleněného majetku, tak by se jednalo o smlouvu o plnění třetí osoby (pactum in detrimentum tertii) ve smyslu § 1769."}\footfullcite[Komentář k § 1448]{svestka_j_obcansky_2014-1}.\\

V takovémto případě by tedy měla ze smlouvy jako signatář závazek třetí osoba, ta by se tak musel například přimluvit u jiné osoby, aby zakladatelem vyčleněný majetek, jež se zakladatel zavázal do svěřenského fondu vložit, začala v rámci svěřenského fondu spravovat, respektive, aby tedy přijala pověření ke správě majetku, bez toho, aby byla spolusignatářem smlouvy o vyčlenění majetku\footfullcite[Strana 41]{strelecek_tomas_sverenske_fondy}\textsuperscript{,}\footnote{Obdobně Pihera v komentáři C.H.Beck k § 1448: \textit{"V jiných případech bude přijetí závazku správy adresováno jiné osobě, která jmenuje správce na základě pravidel svěřenského fondu, nebo osobě, která z jiného titulu (např. osoby pověřené dohledem nad správou fondu) oznámí budoucímu správci, že byl jmenován správcem. Ve všech těchto případech nastanou účinky takového přijetí nejdříve v okamžiku, kdy relevantní projev souhlasu dojde takové osobě."}.}.\\

Kocí k tomuto dále dodává: \textit{"Jakmile zakladatel vyčlení majetek ve výše uvedeném smyslu smlouvou, je tedy v zásadě vznik svěřenského fondu již neodvratný. I kdyby totiž designovaný svěřenský správce odmítl ujmout se správy vyčleněného majetku, soud jmenuje jiného (§ 1455 odst. 2)."}\footnote{Komentář Wolters Kluwer, § 1448.}.\\

V případě, že by svěřenský správce tedy pověření ke správě nepřijal, aplikovalo by se ustanovení § 1455 odst. 2\footfullcite[Strana 152]{lucie_joskova_sprava_2017}, ve kterém se odráží zásada tradičního práva trustů, že svěřenský fond nemůže ztroskotat pro absenci svěřenského správce\footfullcite[Strana 145]{claxton_john_b_studies_2005}. Druhá fáze spočívá v samotném zápisu svěřenského fondu do evidence dle §1451 Občanského zákoníku a §65a Zákona o veřejných rejstřících právnických a fyzických osob, který je konstitutivní, svěřenský fond tedy formálně vzniká až tímto zápisem a zakončuje tak druhou fázi, návrh na zápis podává dle § 65b VR svěřenský správce.\\

Svěřenský správce může pověření ke správě přijmout i konkludentně\footnote{Kocí, komentář k § 1448}.\\

V případě, že by zakladatelské právní jednání bylo pouze dvoufázové a svěřenský správce tak přijal pověření ke správě již ve smlouvě, svěřenský fond by byl zřízen a následným zápisem do evidence svěřenských fondů by i vznikl\footnote{Tamtéž.}.\\

V případě postupu dle § 1731 až 1739 OZ nemůže v případě neakceptace oferty svěřenský fond vzniknout, neboť svěřenský fond zakládaný \textit{inter vivos} nemůže vzniknout jednostranným právním jednáním zakladatele, s ohledem na výše uvedené se tak rovněž domnívám, že v případě nepřijmutí oferty ze strany správce nemůže ani dojít k postupu dle § 1455 odst. 2 OZ, neboť ten bude aplikovatelný v případě uzavření smlouvy nikoli se svěřenským správcem, ale se třetí osobou.\\

Pro vznik svěřenského fondu \textit{inter vivos} jsou tak podstatné tři kroky\footnote{Nikoli však fáze.}: \\

\begin{itemize}
\item vyčlenění majetku z vlastnictví zakladatele,
\item přijetí pověření ke správě tohoto majetku svěřenským správcem (v tomto okamžiku je svěřenský fond zřízen) a
\item zápis svěřenského fondu do evidence svěřenských fondů (v tomto okamžiku svěřenský fond vzniká).\\	
\end{itemize}


Je zřejmé, že v případě založení svěřenského fondu \textit{inter vivos} je v zásadě bezpředmětné bavit se o možnosti ochrany nepominutelných dědiců. V rámci tohoto způsobu založení svěřenského fondu je nutné poznamenat, že ačkoli je možné, aby zakladatel vyčlenil svůj majetek do svěřenského fondu ještě během svého života s tím, že by tímto jednáním sledoval zájem poměrně ponížit o vyčleněný majetek hodnotu pozůstalosti a tím snížit hodnotu vyplacenou jednotlivým nepominutelným dědicům v rámci jejich povinného dědického podílu, není možné se tomuto jednání ze strany jednotlivých nepominutelných dědiců bránit ve smyslu zkrácení povinného podílu. Toto tvrzení lze doložit hned několika ustanoveními, předně Čl. 11 Listiny základních práv a svobod, která stanoví, že \textit{"Každý má právo vlastnit majetek...."}\footfullcite[čl. 11]{noauthor_zakon_lzps} a v návaznosti na toto poté Občanský zákoník v §996 \textit{"Poctiví držitel smí v mezích právního řádu věc držet a užívat ji, ba ji i zničit nebo s ní jinak nakládat..."} a v §1012 \textit{"Vlastník má právo se svým vlastnictvím v mezích právního řádu libovolně nakládat a jiné osoby z toho vyloučit...."}, z tohoto vyplývá, že vlastník může za svého života se svým majetkem jakkoliv nakládat, tedy má možnost ho zcizit. Ustanovení § 1475, které stanoví, že pozůstalost tvoří celé jmění zůstavitele už dále jen potvrzuje, že takto zcizené věci nejsou a nemohou být součástí pozůstalosti. S ohledem na fakt, že dle definice svěřenského fondu v § 1448 vzniká oddělené a nezávislé vlastnictví vyčleněného majetku a dále již tento vyčleněný majetek není ve vlastnictví zůstavitele a dále vzhledem k tomu, že toto právní jednání bylo učiněno ještě za života zakladatele, lze důvodně předpokládat, že tento majetek již nebude \textit{a priori} předmětem pozůstalosti a dědického řízení, tím pádem nemůže být ani připočten k jmění zakladatele v době smrti a tím pádem z tohoto majetku nemůže být počítán povinný díl nepominutelných dědiců.\\

% --------------------------------------------------------------------------
%Dopsat ještě ústavu, že je chráněno právo na to, aby si majitel se svým majetkem nakládal podle své vůle.

%Doplnit konstitutivní a deklaratorní
% --------------------------------------------------------------------------

Jako jednoduchá analogie k založení svěřenského \textit{inter vivos} a dopadu tohoto jednání na dědický podíl se může jevit jiná forma zcizení majetku, například darování, koupě, směna a další smluvní či mocenské formy zcizení věci, které působí na snížení, či zvýšení majetku zůstavitele. \\

Možnost bránit se proti převodům majetku, které byli učiněny ještě za života zůstavitele, mají dědici tak pouze v případě, že jednání učiněné zakladatelem není platné například v důsledku toho, že je zůstavitel fakticky nezpůsobilý činit jakékoliv majetkové dispozice, pokud však proběhlo vše podle zákona a zůstavitel měl vůli takto majetek převést, dědicové nemají možnost se proti takovýmto dispozicím jakkoliv bránit.\\

% --------------------------------------------------------------------------
%---Koment\\
%Doplnit další informace.\\
%---\\
%Je nicméně vhodné říci, že toto neznamená, že by dědici neměli možnost brát se o své právo jinou formou.\\
%---Koment\\
%Doplnit možnosti žaloby ve smyslu zneplatnění právního jednání pokud by bylo učiněno například pod tlakem, relativní a absolutní neúčinnost.\\
%---\\

%Claxton říká následující: \textit{"a trust should not fail for the want of a trustee."}\footfullcite[Strana 145]{claxton_john_b_studies_2005}.

%\newpage
%\thispagestyle{smallertextinheader}
% --------------------------------------------------------------------------

\subsubsection{Vznik svěřenského fondu \textit{mortis causa}}

Pokud se zakladatel rozhodne svěřenský fond zřídit pořízením pro případ smrti, tedy závětí, dědickou smlouvou, nebo dovětkem\footfullcite[§ 1491]{noauthor_zakon_2012}, jedná se o svěřenský fond \textit{mortis causa}, na rozdíl od svěřenského fondu založeného mezi živými, dochází ke vzniku svěřenského fondu založeného pro případ smrti již samotnou smrtí zakladatele\footnote{Tamtéž, § 1451 odst. 1.}, nikoli tedy zápisem do evidence. Tento fakt je důležitý k pozdějšímu zkoumání, zda tento majetek náleží do pozůstalosti, či nikoliv.\\

Majetek se tedy stává odděleným a nezávislým vlastnictvím již smrtí zůstavitele bez ohledu na to, zda zakladatelem pověřená osoba správu bude vykonávat. Zakladatel může pro tyto případy určit jiný způsob určení svěřenského správce, v ostatních případech jmenuje správce soud na návrh osoby, která na tom má právní zájem (typicky touto osobou bude obmyšlený)\footfullcite[Strana 153]{lucie_joskova_sprava_2017}.\\

Svěřenský fond vzniklý smrtí zůstavitele se nicméně také zapisuje do evidence svěřenských fondů, zapisuje se však až po svém vzniku a zápis je tak pouze deklaratorní povahy.\\

% --------------------------------------------------------------------------
%Další možnost pro převod majetku představuje založení svěřenského fondu \textit{morti causa}, která vedle dědického řízení slouží jako další prostředek, který zůstavitel může, krom pořízení pro případ smrti, užít ke kontrole způsobu nakládání se svým jměním po jeho smrti. 
% --------------------------------------------------------------------------

Občanský zákoník však blíže neupravuje vztah svěřenského fondu zřizovaného pořízením pro případ smrti a úpravou dědického práva. Není tedy zřejmé, zda zakladatel může zřízením svěřenského fondu rozhodovat o veškerém svém majetku bez zohlednění nepominutelných dědiců\footnote{Tamtéž.}, potenciální problém tedy nastává pokud by toto založení mělo směřovat k úmyslnému zmenšení jednotlivých povinných dílů náležejících nepominutelným dědicům. V aktuální právní úpravě totiž není jasně stanoveno jakým způsobem se majetek při smrti zůstavitele, který svěřenský fond zřídil pořízením pro případ smrti do tohoto fondu vyčlení, respektive, zda je tento majetek součástí pozůstalosti, nebo nikoliv. Nejasný výklad ustanovení svěřenského fondu ve vztahu k dědickému právu představuje právní nejistotu, nejen pro dědice, která není žádoucí. \\

S ohledem na popis dědického práva vyložený v předchozí kapitole je třeba s ohledem na založení svěřenského fondu \textit{mortis causa} nejdříve vyřešit jakým způsobem probíhá jeho založení pořízením pro případ smrti. S ohledem na § 1452, odstavec 1 je po výkladové stránce zřejmé, že zákonodárce zamýšlel způsob vyčlenění majetku do svěřenského fondu ve smyslu § 311, který upravuje vyčlenění majetku do nadace, a který zároveň stanovuje v odstavci 1 následující: \textit{"(1) Při založení nadace pořízením pro případ smrti se do nadace vnáší vklad povoláním nadace za dědice nebo nařízením odkazu. V takovém případě nabývá založení nadace účinnosti smrtí zůstavitele."}\footfullcite[§ 311, odst. 1]{noauthor_zakon_2012}. Z tohoto ustanovení, které je možné analogicky použít na svěřenský fond, je možné dovodit úmysl zákonodárce pohlížet na majetek, který má být vyčleněn do svěřenského fondu, jako na součást pozůstalosti. Toto tvrzení svým komentářem potvrzuje i Vlastimil Pihera, který ve svém komentáři k § 1448 dovozuje, že při zřízení svěřenského fondu pořízením pro případ smrti vyčlení zakladatel (rozuměj zůstavitel) majetek ve prospěch svěřenského fondu tak, že v jeho prospěch vyhradí dědické právo, nebo nařídí odkaz\footfullcite[komentář k § 1448]{spacil_j_a_kolektiv_obcansky_2013}. Problematika nepominutelného dědice je blíže rozebrána v Kapitole 7.4 Nejasný výklad ochrany nepominutelných dědiců.\\

% --------------------------------------------------------------------------
%Na svěřenský fond zřízený \textit{mortis causa} je v rámci občanského zákoníku kladený velký důraz, je nicméně vhodné poznamenat, ať již se jedná o svěřenský fond založený za veřejným a nebo soukromým účelem.

%V rámci pohledu na právní úpravy trustů ze zahraničí je možné se setkat s dalšími druhy trustů a mnohými jinými možnosti, za kterýchžto účelem může být svěřenský fond založený, tyto se však Český zákonodárce dále rozhodl v našem občanském zákoníku neupravovat.

%\newpage
%\thispagestyle{smallertextinheader}

%Tento majetek je dle Vlastimila Pihery a i dle mého výkladu a pochopení OZ tedy součástí pozůstalosti, tím pádem, dle § ...., který stanovuje, že pozůstalost tvoří veškeré jmění v držení zůstavitele ke dni jeho smrti, se z výše tohoto jmění spolu se započtením výše jmění, které není vyčleněno do svěřenského fondu počítají jednotlivé dědické podíly, které je správce pozůstalosti povinen převést jednotlivým nepominutelným dědicům. S ohledem na to, že další úprava upravující tetnto souběh v OZ chybí a zároveň k tomuto problému neexistuje žádná judikatura a není k němu ustálená soudní praxe.....

%a výklad komentáře.

%Popsat, že existují dva různé pohledy, já se ztotožňuji s druhým pohledem.

%Dopsat ty výhody svěřenského fondu mortis causa, na které narážím v tom druhém dokumentu s názvem Bakalářská práce poznámky.
% --------------------------------------------------------------------------

\subsubsection[Vznik svěřenského fondu \textit{inter vivos} s odkládací podmínkou]{Vznik svěřenského fondu \textit{inter vivos} s odkládací\\ podmínkou}

% --------------------------------------------------------------------------
%Dopsat co píše Tomáš Střeleček ve své práci, strana 46, 47 a 48.

%Viz zajímavý článek: https://dspace.cuni.cz/bitstream/handle/20.500.11956/117098/150047668.pdf?sequence=1&isAllowed=y o tomto rozhodně napsat

%Z dikce zákona je též možné dovodit možnost vyčlenění majetku smlouvou s odkládací podmínkou

%Zajímavé je dovození, že lze vyčlenit majetek svěřenského fondu smlouvou s odkládací podmínkou ve smyslu ustanovení § 548 až 549 OZ, jak činí Miloš Kocí232. Pokud by tou podmínkou byla smrt zakladatele, šlo by o svěřenský fond mortis causa, ale založený způsobem předpokládaným u svěřenských fondů inter vivos, tedy dvoufázově (případně jednofázově, pokud takové přijetí pověření ke správě je obsaženo ve smlouvě). Pokud nepřijme svěřenský správce pověření ke správě ke smrti zakladatele, takový svěřenský fond vůbec nevznikne.
% --------------------------------------------------------------------------

V souvislosti s ustanovením § 548, zejména odstavce (2), a 549 občanského zákoníku lze (obdobně) uvažovat o dalším způsobu vzniku svěřenského fondu, respektive vyčlení majetku. Takovéto vyčlenění by bylo vázáno na určitou odkládací podmínku. Zakladatel může stanovit mnoho různých druhů odkládacích podmínek, na jejichž splnění se váže vznik práv\footnote{Je samozřejmě nutné splnění dalších podmínek, které na založení svěřenského fondu klade zákon, tedy například přijetí svěřenským správce, srovnání viz výše.}. Při splnění pak vznikne svěřenský fond a obmyšlený bude mít bez dalšího, samozřejmě s ohledem na podmínky stanovené ve statutu, právo na vyplácení prostředků z tohoto svěřenského fondu.\\ 

% --------------------------------------------------------------------------
%v tomto případě daný student vysoké školy, bude mít bez dalšího právo na vyplácení prostředků z tohoto svěřenského fondu.\\

%může se tak jednat například o podmínku dokončení studia na vysoké škole, po jejímž splnění\footnote{Je samozřejmě nutné splnění dalších podmínek, které na založení svěřenského fondu klade zákon, tedy například přijetí svěřenským správcem, srovnání viz výše.} vznikne svěřenský fond a obmyšlený, v tomto případě daný student vysoké školy, bude mít bez dalšího právo na vyplácení prostředků z tohoto svěřenského fondu.\\
% --------------------------------------------------------------------------

Zajímavé je tak dovození, že svěřenský fond může být zřízen vyčleněním majetku smlouvou s odkládací podmínkou smrti\footfullcite[Komentář k § 1448]{svestka_j_obcansky_2014-1}, neboť pak by se jednalo ve své podstatě, samotnou právní skutečností, smrtí zůstavitele, o obdobu založení svěřenského fondu \textit{mortis causa}.\\

% --------------------------------------------------------------------------
%, přitom však věřím, že by se na toto založení pohlíželo jako na založení svěřenského fondu \textit{inter vivos}.\\ 
% --------------------------------------------------------------------------

Zakladatel by tak mohl ještě za svého života určil ve smlouvě majetek, který má být z jeho vůle svěřenským správcem držen a spravován za vymezeným účelem, v rámci první fáze zřizování, přičemž v rámci druhé fáze by toto pověření určený svěřenský správce přijal, okamžik zřízení by však nebyl vázán například na podpis smlouvy oběma stranami, nebo přijetí oferty, ale na budoucí okamžik smrti zakladatele, přičemž datum ke kterému by zakladatel zemřel by bylo datumem zřízení svěřenského fondu, neboť k tomu by byla splněna odkládací podmínka. Třetí fází by pak byl konstitutivní zápis svěřenského fondu do evidence a vznik.\\ 

Zřízení by také mohlo být jednofázové, pokud by přijetí takovéto správy bylo již obsažené jako součást smlouvy mezi zakladatelem a svěřenským správcem, respektive dvoufázové, neboť z principu (jedná se o založení \textit{inter vivos}) je nutné podat návrh na zápis svěřenského fondu do evidence, přičemž toto činí svěřenský správce. Přestože by samotným důvodem ke vzniku svěřenského fondu byla smrt zakladatele, nejednalo by se dle mého názoru o založení svěřenského fondu \textit{mortis causa}, ale o založení svěřenského fondu \textit{inter vivos} s tím, že by svěřenský fond dle novely č. z. 460/2016 Sb. vznikl až zápisem do evidence svěřenských fondů, tento zápis by tedy byl konstitutivní.\\

Kocí k tomuto dodává, že pokud by pověření ke správě svěřenský správce nepřijal, tak by svěřenský fond vůbec nevznikl a majetek by připadl do pozůstalostního řízení\footnote{Kocí tvrdí v komentáři Wolters Kluwer k § 1448 následující: \textit{"Jinak by svěřenský fond smrtí zakladatele platně nevznikl z toho důvodu, že zakladatelské právní jednání nebylo učiněno ve formě posledního pořízení (§ 1451 a contrario, viz níže), a vyčleňovaný majetek by proto byl součástí pozůstalosti."}}, toto by tedy představovalo velkou nevýhodu založení svěřenského fondu tímto způsobem oproti zřízení \textit{mortis causa}, kde svěřenský fond vznikne i v případě, že svěřenský správce pověření ke správě nepřijme. Domnívám se však, že takto tomu bude jen v případě nepřijetí oferty k uzavření smlouvy ze strany svěřenského správce, modality takovéto smlouvy jsou obdobné, jako u prostého zřízení svěřenského fondu \textit{inter vivos}, šlo by proto tuto smlouvu obdobně uzavřít i s osobou třetí, přičemž pokud by svěřenský správce pověření ke správě nepřijal, uplatnil by se postup dle § 1455 odst. 2 OZ, přičemž poslední fází zřízení svěřenského fondu \textit{inter vivos} s odkládací podmínkou smrti by byl zápis do evidence svěřenských fondů, na základě kterého by svěřenský fond vznikl.\\

V případě, že by svěřenský fond vznikl až po započetí pozůstalostního řízení se domnívám, že by se svěřenský správce mohl domáhat vydání majetku i za pomoci reivindikační žaloby vůči dědicům v případě, že by kvůli obsahu na vyčlenění majetku pohlíželo jako na stanovení odkazu\footfullcite[Strana 111]{noauthor_casopis_2015}. Svěřenský fond by v takovémto případě ani nebyl zcela bez majetku, neboť by zahrnoval pohledávky za dědici, viz výklad níže.\\

% --------------------------------------------------------------------------
 %kdy by se na vyčleněný majetek dalo pohlížet jako na odkaz. Svěřenský fond by v takovémto případě ani nebyl zcela bez majetku, neboť by zahrnoval pohledávky za dědici, viz výklad níže.\\

 %která může v tomto případě nastat. Ta souvisí, se samotným stavem zakladatele svěřenského fondu. Pokud by zakladatel zemřel a odkládací podmínka by nebyla vázána na jeho smrt ale na splnění jiné podmínky, hledělo by se na takto založený fond jako na \textit{inter vivos}.\\

%(Pokud je mrtev jedná se stále o založení svěřenského fondu \textit{inter vivos}? Já si myslím, že ano}.

%Dalším zajímavým příkladem je možnost stanovit smrt zakladatele jako odkládací podmínku, protože pak by se jednalo ve své podstatě, samotnou právní skutečností, smrtí zůstavitele, o obdobu založení svěřenského fondu \textit{mortis causa}, přitom však věřím, že by se na toto založení pohlíželo jako na založení svěřenského fodnu {inter vivos}. Kdy by tedy zakladatel ještě za svého života určil ve smlouvě majetek, který má být z jeho vůle svěřenským správcen držen a spravován za vymezeným účelem v rámci první fáze vzniku a následně by v rámci druhé fáze vzniku přijal toto pověření určený svěřenský správce. Založení by také mohlo být jednofázové, pokud by přijetí takovéto správy bylo již vedené jako součást smlouvy mezi zakladatelem a svěřenským správcem. Přestože by samotným důvodem ke vzniku svěřenského fondu byla smrt zakladatele, nejednalo by se dle mého názoru o založení svěřenského fondu \textit{mortis causa}, ale o založení svěřenského fondu \textit{inter vivos} s tím, že by svěřenský fond dle novely č. z. 460/2016 Sb. vznikl až zápisem do evidence svěřenských fondů, tento zápis by tedy byl konstitutivní.\\
% --------------------------------------------------------------------------

S ohledem na fakt, že takovéto založení svěřenského fondu stále spadá do kategorie \textit{inter vivos}, v rámci kterého takto vyčleněný majetek není součástí dědického řízení, respektive pozůstalosti po zesnulém\footnote{Viz výše.}. Je dle mého názoru možné využít tuto formu založení svěřenského fondu místo klasického založení \textit{mortis causa}, u kterého existuje právní nejistota, zda je takto vyčleněný majetek součástí pozůstalosti.\\

% --------------------------------------------------------------------------
%Tato forma založení svěřenského fondu ssebou ale nese i jisté nevýhody oproti formě založení pro případ smrti. Těmi jsou především.
% --------------------------------------------------------------------------

Pokud by tedy zakladatel chtěl docílit mezigeneračního převodu majetku avšak chtěl by \textit{ad libidum} fakticky vydědit jakéhokoliv svého dědice, nebo se vyhnout dědickému řízení, je cesta buď založení svěřenského fondu \textit{inter vivos}, nebo v případě, že by chtěl zakladatel se svým majetkem v mezích disponoval až do doby své smrti, pak je s ohledem na nejasnosti vztahující se k založení svěřenského fondu \textit{mortis causa} vhodné, vydat se cestou založení svěřenského fondu \textit{inter vivos} s odkládací podmínkou smrti zakladatele.\\

V případě založení svěřenského fondu mezi živými s odkládací podmínkou smrti zakladatele se ve své podstatě jedná o založení \textit{mortis causa}, jak již je zmíněno výše, avšak formou \textit{inter vivos}. S ohledem na to, že se nejedná o založení svěřenského fondu pořízením pro případ smrti je zřejmé, že v tomto případě nelze uvažovat ani nad ochranou dle § 1492 OZ, pro ilustraci je možné porovnat tento způsob založení s další formou mezigeneračního předání majetku a to darováním pro případ smrti.\\

Podklad pro toto tvrzení poskytují například autoři článku Darování pro případ smrti v Časopise pro právní vědu a praxi Adam Talanda a Iveta Talandová. Ti přímo říkají: \textit{"Darování pro případ smrti je při splnění všech zákonných podmínek obligačním vztahem mezi dárcem a obdarovaným a neřídí se ustanoveními dědického práva. Zákon tak umožňuje užitím jednoho zákonného institutu obcházet soubor jiných zákonných institutů. Především se jedná o obcházení ustanovení omezujících zůstavitele v nakládání s jeho majetkem, která slouží k ochraně nepominutelných dědiců."}\footfullcite{noauthor_casopis_2015}.\\

K tomuto zároveň dodávají:\\

\textit{"Za splnění těchto podmínek se darovací smlouva posuzuje jako darování pro případ smrti, na základě kterého nabývá obdarovaný vlastnické právo k předmětu daru už smrtí dárce (v případě nemovitosti má po smrti dárce právo podat návrh na vklad do katastru nemovitostí). V případě daru má obdarovaný nikoli obligační právo vůči dědicům, ale věcné právo. Nejsou-li splněny tyto náležitosti, posuzuje se smlouva jako stanovení odkazu, kterým se zřizuje obligační právo vůči dědicům."}\footnote{Tamtéž, strana 113.}\\

V tomto případě by podle Talandy povinnost z darování pro případ smrti měla být subsumována pod pojem dluhů a závad váznoucích na majetku již v době zůstavitelovy smrti, ergo odečtena od hodnoty pozůstalostních aktiv (§ 1655 odst. 1o.z.)\footfullcite[Strana 474]{talanda_a_darovani_2015}, domnívám se, že obdobně by tomu bylo u majetku vyčleňovaného do svěřenského fondu.\\

Taktéž:\\

\textit{"V případě, že je darování pro případ smrti platné, není předmětem řízení o pozůstalosti podle zákona o zvláštních řízeních soudních. Obdarovaný nabývá k předmětu daru vlastnické právo splněním odkládací podmínky, tedy smrtí dárce. V případě nemovitosti se vlastnické právo nabývá až zápisem do veřejného seznamu, ke kterému lze dát podnět po smrti dárce, přičemž druhou stranou řízení bude dědic nebo dědicové, kterým by jinak připadla nemovitost."}\footnote{Tamtéž, strana 117.}\\

Je dále důležité zdůraznit podmínky, za kterých je darovací smlouvou možné platně založit darování pro případ smrti. Talanda v tomto ohledu zmiňuje, že je zásadní vzdát se práva dar odvolat\footnote{Tamtéž, strana 113.}. Kocí v této souvislosti zmiňuje, že u bezplatného vyčlenění formou smlouvy se přiměřeně použijí ustanovení o darování, což znamená, že zakladatel může od smlouvy odstoupit a převod vyčleněného majetku odepřít, změní-li se po uzavření smlouvy zásadně okolnosti\footfullcite[Komentář k § 1448]{svestka_j_obcansky_2014-1}. Domnívám se proto, že zakladatel by se musel ve smlouvě rovněž vzdát práva dar odvolat, aby byla takováto smlouva o zřízení svěřenského fondu platná\footnote{Viz Příloha č. 1 § 1b Vyčlenění majetku.}. Dalším zásadním předpokladem pro zřízení svěřenského fondu \textit{inter vivos} s odkládací podmínkou bude nutná forma zakladatelského právního jednání jakožto veřejné listiny, neboť tuto formu vyžaduje statut\footfullcite[§ 1452 odst. 3]{noauthor_zakon_2012}.\\ 

Je tedy zřejmé, že v případě nepoužití pořízení pro případ smrti nelze aplikovat omezení vztahující se právě na pořizování. S ohledem na fakt, že obdarovaný nabývá vlastnické právo již smrtí dárce nemůže být daný předmět darování považován ani za součást pozůstalosti, obdobně nelze brát zřetel ani k započtení, neboť obdarovaný v tomto případě není dědicem\footfullcite{sro_httpwwwpraguebestcz_darovani_mortis_causa}. Rovněž v případech, že by vyčleněný majetek byl součástí pozůstalosti, považoval by se dle výkladu výše za obligační právo vůči dědicům.\\

% --------------------------------------------------------------------------
%Povinnost z darování pro případ smrti by podle jeho názoru měla být subsumována pod pojem dluhů a závad váznoucích na majetku jižv době zůstavitelovy smrti, ergo odečtena od hodnoty pozůstalostních aktiv (§ 1655 odst. 1o.z.)\footfullcite[Strana 474]{talanda_a_darovani_2015}.\\

%V takovém případě se domnívám, že by se svěřenský správce mohl domáhat vydání majetku i za pomoc reinvikační žaloby vůči dědicům, kdy by se na vyčleněný majetek dalo pohlížet jako na odkaz. Svěřenský fond by v takovémto případě ani nebyl zcela bez majetku, neboť by zahrnoval pohledávky za dědici.\\

%V nějaké zemi je možné žádat o to, aby bylo dědicům vyplaceno část dědictví i z toho, co dal zakladatel komukoliv 3 roky před jeho smrtí, najít kde a doplnit to sem a poznamenat, že u nás toto neplatí.

%Je vložení majetku do svěřenského fondu darováním? Pokud ano, muselo by v tomto případě ve smlouvě s odkládací podmínkou být stanoveno i nemožnost odvolání daru jak říká Talanda ve svém příspěvku.

%Možná bych měl smlouvu učinit neodvolatelnou.

%\newpage
%\thispagestyle{smallertextinheader}
% --------------------------------------------------------------------------

\subsubsection{Jakou formu vybrat pro mezigenerační převod majetku?}

Potenciálního zakladatele tak může napadnout otázka, jakou z těchto forem vybrat? Jak již jsem výše zmínit, svěřenský fond je velice flexibilní institut a šíře účelů je nezměrná. Jednoduchá odpověď tak neexistuje, výše jsem nastínil problém ve vztahu svěřenského fondu \textit{mortis causa} k nepominutelnému dědici, v situaci kdy by toto mohlo přestavovat problém je doporučení této formy problematičtější, takovéto doporučení však nelze generalizovat, neboť v určitých případech může být tato problematika spíše marginálního charakteru.\\

Pro úspěšný mezigenerační převod majetku je tak možné vybrat jakoukoliv z výše uvedených forem, zakladatel musí mít pouze na vědomí rozdíly mezi nimi.\\

% --------------------------------------------------------------------------
%Domnívám se, že i když bude nepominutelný dědic chráněn ve vztahu k založení svěřenského fondu \textit{mortis causa}, tedy založení touto formou bude znevýhodněno jak tvrdí Kocí, 
% --------------------------------------------------------------------------

Bez ohledu na otázku nepominutelných dědiců, jež ostatně řeším v kapitole 7.4, bude mít způsob založení \textit{mortis causa} stále benefity před formou svěřenského fondu \textit{inter vivos}. Mezi tyto benefity pak patří například to, že při založení \textit{inter vivos} v případě kdy zakladatel předpokládá, že může v brzké době zemřít, avšak nezemře, může fakticky zůstat bez prostředků po zbytek svého života, pokud je všechny převede do svěřenského fondu, komentářová literatura se pak rozchází v názoru na to, zda může být zrušení svěřenského fondu vázáno i na jiné krom zákonem taxativně vymezené podmínky, tedy zda je v případě těchto jinak kogentních ustanovení připuštěna jistá autonomie vůle stran a možnost upravit si důvody zániku ve statutu\footnote{Srovnání Pihera a Kocí, komentář k § 1471.}, otázkou pak zůstává, zda si může zakladatel například vyhradit právo svěřenský fond zrušit, jak tvrdí například Peter Kováčik\footnote{Článek, ve kterém je uvedeno toto tvrzení je dostupný zde: https://echo24.cz/a/wMCST/jak-rychle-se-da-agrofert-vyvest-ze-sverenskeho-fondu-klidne-okamzite, jednalo by se tedy o revocable trust, respektive zrušitelný svěřenský fond.}. Pokud by tak založil formu \textit{mortis causa}, mohl by s majetkem nakládat až do své smrti. Po své smrti by pak mohl stanovit, že zbývající majetek má být převeden do svěřenského fondu. Na druhou stranu svěřenský fond \textit{mortis causa} nechrání zůstavitelův majetek před věřiteli za jeho života, neboť se převádí až ve chvíli jeho smrti, tolik opěvovaný insolvenční efekt je tak značně posunut. A pak je tu samotná otázka nepominutelných dědiců.\\

% --------------------------------------------------------------------------
% a navíc by msuel být respektován nepominutelní díl nepominutelného dědice.
% --------------------------------------------------------------------------

Svěřenský fond zřízený \textit{inter vivos} s odkládací podmínkou smrti a svěřenský fond zřízený \textit{inter vivos} mezi sebou nemají z pohledu výhodnosti jednoho, nebo druhého zásadní rozdíly. Rozdíly tak spatřuji pouze v posunutí data vyčlenění majetku do svěřenského fondu, toto však může hrát pro zakladatele v jistých případech zásadní roli, problémem však zůstává absence insolvenčního efektu, neboť zakladatel je až do doby vyčlenění majetku jeho vlastníkem.\\

% --------------------------------------------------------------------------
%jen by mohl představovat, že majetek vyčleněný v dobe vypracování smlouvy už k době založení neexsituje (problém pro správce jak to bude v případě jednání zůstavitele protivícího se smlouvě vymáhat) obdobně je takovýto majetek ohrožen zůstavitelovými věřiteli i v případě zajíštění věcným právem.\\
% --------------------------------------------------------------------------

Mohlo by se tak zdát, že svěřenský fond zřízený \textit{inter vivos} s odkládací podmínkou smrti má oproti zřízení svěřenského fondu \textit{mortis causa} jen samé výhody\footnote{Věřitelé v tomto případě nejsou bráni v potaz, neboť stejně tak, či tak, se budou z majetku moci uspokojovat, neboť majetek bude do svěřenského fondu v obou případech vyčleněn až smrtí zůstavitele. Zatímco však toto je velká nevýhoda svěřenského fondu zřízeného \textit{mortis causa} a \textit{inter vivos} s odkládací podmínkou smrti, neboť při těchto způsobech zřízení je doba, po kterou se mohou věřitelé z tohoto pro budoucno vyčleněného majetku uspokojovat, delší, takže v případě ochrany majetku před věřiteli vychází lépe svěřenský fond zřízený \textit{inter vivos} se zřízením ke dni účinnosti smlouvy.}, neboť zakladatel bude moci nakládat se svým majetkem, nebude muset řešit povinný díl nepominutelných dědiců a věřitelé se v obou případech budou moci uspokojovat po stejně dlouhou dobu. Je tedy otázkou proč vůbec mít svěřenský fond mortis causa? Na tuto otázku bych odpověděl velice jednoduše: "protože proč ne?". Ve chvíli, kdy zůstavitel platně vydědí nepominutelné dědice, nebo nebude mít nepominutelných dědiců, může si vybrat oba způsoby, neboť oba pro něj budou mít stejné důsledky. Tzn. možnost zřízení svěřenského fondu \textit{inter vivos} s odkládací podmínkou smrti jasně nemusí vylučovat možnost zřízení svěřenského fondu \textit{mortis causa}. Navíc zakladatel může mít své vlastní pohnutky proč by nemusel chtít, aby byl svěřenský fond zřízen smlouvou za jeho života s odkládací podmínkou, ale aby bylo zakladatelské právní jednání vtěleno do pořízení pro případ smrti (takovouto možnost zřízení \textit{inter vivos} pro případ smrti vylučuje), může například mít zájem na tom, aby zakladatelské právní jednání bylo jednostranné, což u zřizování smlouvou není možné, s tímto souvisí i fakt, že svěřenský fond zřízení pořízením pro případ smrti vzniká bez ohledu na svěřenského správce, zatímco v případě zřízení svěřenského fondu smlouvou je dané přijetí vyžadováno, v případě, že daného přijetí není, svěřenský fond pak nemusí vzniknout. V tomto ohledu pak poskytuje zakladateli větší jistotu, že s jeho majetkem bude naloženo podle jeho vůle, zřízení svěřenského fondu pořízením pro případ smrti.\\

\subsection{Evidence svěřenských fondů}

Přes mnohem striktnější úpravu svěřenského fondu, než kterou lze najít u jiných trust-like institutů jí nelze upřít, že zákonodárce omezil možnosti potenciálního zneužití svěřenského fondu ve smyslu vysloveného například ministerstvem financí\footnote{Viz níže.}. Zakotvení trustu v české právní úpravě ve formě svěřenského fondu s povinnými náležitostmi ve formě statutu, zápisu do evidence a dalších obdobných požadavků minimalizuje rizika jeho použití k nekalým finančním praktikám kvůli netransparentní konstrukci. I přes toto je mnohem snadnější užít k účelům transferu majetku či obchodním účelům svěřenský fond, než je zakládat jinou právnickou osobu, tedy obchodní korporaci, nebo za jistých okolností fundaci\footnote{O ústavu neuvažuji.}.\\

\subsubsection{Novelizace právní úpravy - zápis svěřenských fondů do rejstříku}

Majetek ve svěřenském fondu tedy není ve vlastnictví nikoho, před novelizací Občanského zákoníku pak ve veřejných seznamech nebo evidencích mohla být dohledána pouze osoba svěřenského správce. Pro samotné svěřenské fondy pak nebyla stanovena povinnost se jakkoliv registrovat nebo evidovat, v souvislosti s tím pak nebyl veřejně dostupný ani statut svěřenského fondu, který přeci musel mít formu veřejné listiny, ale v souvislosti s absencí povinnosti evidence svěřenských fondů nebyl nikde zveřejněn. Absence jakékoliv evidence, vedla k anonymitě zakladatele, obmyšlených a přispívala k celkové netransparentnosti vlastnické struktury a majetku vyčleněného do svěřenského fondu.\\

V souvislosti panoval oprávněný strach, že svěřenské fondy by mohli být využity k legalizaci výnosů z trestné činnosti, daňovým únikům nebo být jinak zneužity. V souvislosti s touto obavou hned tři instituce, kterými byli Ministerstvo financí, Ministerstvo vnitra a Vrchní státní zastupitelství volali po novelizaci právní úpravy svěřenského fondu a povinné evidenci. Důvodová zpráva k novele pak stanovuje přímo následující: \textit{"V reakci na naléhavé připomínky Ministerstva financí, Ministerstva vnitra a Vrchního státního zastupitelství v Praze se navrhuje zavedení evidence svěřenských fondů z důvodu jejich netransparentní vlastnické struktury a vysokého potenciálu zneužití především k legalizaci výnosů z trestné činnosti."}\footfullcite{noauthor_4602016_nodate}.\\

% --------------------------------------------------------------------------
%Zákon č. 33/2020 Sb. novela, se svěřenských fondů téměř vůbec nedotýká, pro tuto práci není důležitá, neboť se měni jen určité ustanovení ve vztahu k notářskému zápisu.

%Rozepsaná novelizace svěřenský fondů z roku 2018, zhodnotit, zda přispělo ke zmírnění možností využití svěřenských fondů jako prostředků k legalizaci výnosu z trestných činností.\\
% --------------------------------------------------------------------------

I přes původní prohlášení v důvodové zprávě k novému občanskému zákoníku, ve které byla jasně stanovena "absence veřejnoprávního dohledu u svěřenských fondů"\footfullcite[Komentář k § 1448 až 1452]{Duvodovazprava} tak zákonodárce, kvůli výše zmíněným obavám, zavedl po účinnosti zákona 460/2006, kterým novelizoval občanský zákoník, povinnou evidenci svěřenských fondů.\\

Absence veřejnoprávního dohledu ale rovněž znamenala, mimo jiné\footnote{Neevidování svěřenských fondů představoval problém i pro věřitele zakladatele, těm nemusel být znám vznik, což umožňovalo potenciální antedataci vzniku a zhoršení dobytnosti pohledávek.}, maření úmyslu zákonodárce sledovaných v dalších zákonech namířených proti praní špinavých peněz, či netransparentnosti vlastníka u akcií a tím pádem i netransparentní vlastnické struktuře akciových společností\footnote{Zejméně tedy zákony č. 134/2013 Sb. a č. 253/2008 Sb.}. Praní špinavých peněz zákonodárce přímo zmiňuje a zároveň dodává:\\

\textit{"Zřízení svěřenského fondu ve stávající podobě umožňuje zakladateli a beneficientům fondu zůstat v anonymitě. Nic na tom nemění ani skutečnost, že zákon předepisuje svěřenskému fondu mít statut vydaný zakladatelem ve formě veřejné listiny (§ 1452 odst. 1 a 3). Statut totiž není předmětem žádné evidence."}\footfullcite[Komentář k § 1448]{noauthor_4602016_nodate}. \\

%Poměrně zásaní změny v souvislosti se zakládáním svěřenských fondů tak v souvislosti s tímto přinesla novelizace č.z. 460/2006, která s sebou 

Novela tedy s sebou přinesla nutnost zápisu svěřenského fondu do evidence s tím, že tento zápis je konstitutivní. Svěřenský fond tedy vzniká tímto zápisem. Toto se nicméně dotklo pouze svěřenského fondu založeného mezi živými. Pro svěřenský fond založený pro případ smrti zákonodárce stanovil v § 1451 odstavci 3 následující:\\

\textit{"Byl-li však svěřenský fond zřízen pořízením pro případ smrti, vznikne smrtí zůstavitele. Do evidence svěřenských fondů se zapíše po svém vzniku."}\footfullcite[§ 1451 odst. 3]{noauthor_zakon_2012}\\

Důvod, pro který zákonodárce dopřál toto, do jisté míry zvýhodňující, postavení způsobu založení svěřenského fondu \textit{mortis causa}, je v důvodové zprávě objasněn následovně:\\

\textit{"Odlišný princip deklaratorního zápisu se má uplatnit u svěřenských fondů, které byly zřízeny pořízením pro případ smrti. Tento princip lépe vyhovuje fundamentu testovací volnosti a respektu k vůli zůstavitele."}\footfullcite[Komentář k § 1448]{noauthor_4602016_nodate}\\

Rovněž je do evidence zapisován svěřenský správce a způsob jeho jednání, toto zaručuje právní jistotu pro třetí osoby se svěřenským správcem jednajícími\footfullcite[§ 65d písmeno e) a písmeno f)]{noauthor_zakon_vr}.\\

Je zřejmé, že evidence je jednou z nejsilnějších zbraní zabraňujících netransparentnosti a její zneužití. Zákonodárce bezesporu zavedením povinné evidence snížil riziko zneužití svěřenských fondů.\\

Je pravdou, že tato nově zavedená evidence byla mnohými považována za značný zásah do soukromí a představovala značné narušení flexibility svěřenského fondu, přičemž považovali tuto novelizaci za kontraproduktivní, jedním z nich je i jeden z otců nového občanského zákoníku Karel Eliáš\footfullcite[Strana 23]{elias_bulletin_novelizacni_tlaky}. Je si však nutné uvědomit, že ve finální podobě evidence svěřenských fondů do soukromý tolik nezasahuje, neboť ve veřejně přístupné části jsou pouze obecné údaje týkající se svěřenského fondu a svěřenského správce. Ostatní údaje, zejména o dalších osobách zapisovaných do evidence nejsou veřejně přístupny a soud je poskytne na žádost pouze osobě, která prokáže právní zájem\footfullcite[§ 65e odst. 1 a odst. 3]{noauthor_zakon_vr}.\\

Domnívám se tak, že došlo k vyvážení snahy o anonymizaci svěřenského fondu, ale zároveň i ochrany veřejného zájmu ve smyslu majetkové trestné činnosti a věřitelů. Rovněž se domnívám, že povinná evidence značně ztížila použití svěřenských fondů za nekalými účely. Ve vztahu k samotné flexibilitě a jednoduchosti svěřenských fondů, kterou toliko považuji za velkou výhodu ve světle jeho použití za mezigeneračním převodem majetku, oproti jiným právnickým osobám taktéž s povinností zápisu, nepovažuji nutnost evidence za nikterak ztěžující podmínku, o to více, že v případě testamentárního svěřenského fondu tento vznikne i bez nutnosti jeho zápisu. Je ovšem pravdou, že společně s daňovými předpisy zákonodárce povinnou evidencí přiblížil svěřenský fond zase o krok blíže k právnickým osobám, je pak otázkou, zda je vhodné se tímto směrem v budoucím vývoji svěřenského fondu vydat.\\

\subsection{Správa svěřenského fondu}

Dle ustanovení § 1453 může správu svěřenského fondu vykonávat jakákoliv svéprávná fyzická osoba. Správu může rovněž vykonávat osoba právnická, nicméně pouze v zákonem stanovených případech\footnote{Tyto případy jsou upravené v zákoně č. 240/2013 Sb. zákon o investičních společnostech a investičních fondech.}. V případě veřejného svěřenského fondu pak jde spolu s ustanovením § 146 OZ dovodit, že správce musí být bezúhonný. Správci náleží dle ustanovení § 1456 plná správa majetku ve svěřenském fondu, dle § 1409 a 1410 OZ tedy správce dbá o udržení a rozmnožení majetku v zájmu beneficienta svěřenského fondu, přičemž správce může činit se spravovaným majetkem vše, co je nutné a užitečné.\\

Zásadní podmínkou pro vznik správy je její akceptace ze strany osoby správou pověřené, bez tohoto nemůže správa svěřenského fondu započít, neboť nikdo se nemůže stát svěřenským správcem proti své vůli.\\

Na samotný výkon správy svěřenského správce se subsidiárně aplikuje obecná úprava správy cizího majetku, tedy § 1440 až 1447 OZ. Zakladatel může určit správce jednoho, či několik správců a určit způsob jejich jednání. V některých případech pak zákon ukládá povinnost, aby svěřenský fond měl minimálně dva správce\footfullcite[§ 1454]{noauthor_zakon_2012}, a to v případech, kdy by došlo ke kumulaci pozic svěřenského správce a zakladatele, nebo svěřenského správce a obmyšleného, respektive kdyby v obou těchto funkcích stála totožná osoba, tato povinnost by měla zajistit nestranné rozhodování, v reálu je ale těžké domyslet všechny konsekvence a není absolutně vyloučeno, že druhý správce bude ve výkonu správy ovlivnitelný. Větší množství svěřenských správců pak může přinést jak výhody, tak nevýhody. Výhodou může být nezaujaté rozhodování v případě řádně nastaveného způsobu jednání a odbornost správců specializujících se na jednotlivé části správy majetku a tím i důkladnější a lepší správa svěřenského fondu, problematickou stránku věci pak ale například může představovat složitější kontrola svěřenských správců a beze sporu větší, jednoduše řečeno, papírování. Do evidence se pak budou zapisovat všichni správci\footnote{Srovnání §65d písmeno e) a písmeno f) zákona č. 304/2013 Sb., § 11 a) zákona č. 256/2013 Sb.}.\\

\subsection{Dohled nad svěřenským fondem}

Dohled nad správou svěřenského fondu, tedy dohled na výkonem povinností správce, vykonává zakladatel nebo obmyšlený nebo jiná osoba\footfullcite[Strana 36]{svejkovsky_jaroslav_sverenske_2018}.\\

Zakladatel jako nástroj pro tento dohled může v podstatě použít cokoliv v souladu se statutem a zákonem, co mu umožní kontrolovat to, zda správce postupuje v souladu se statutem a se zákonem, tedy například zda vykonává správu s péčí řádného hospodáře\footnote{Tamtéž, strana 37.}. V případě, že by mu správce tuto kontrolu neumožnil, má možnost postupovat dle § 1466 OZ\footnote{Zákaz určitého jednání správci, nebo uložení povinnosti správci.}\\

Svěřenský správce má povinnost vést záznamy o spravovaném majetku a tyto předkládat kontrolující osobě.\\

Svěřenský správce rovněž může v rámci vykonávání dohledu nad správou svěřenského fondu svěřenského správce odvolat a to i bez udání důvodu.\\

V rámci správy je však v odborné literatuře vyloučena možnost, aby zakladatel udílel správci závazné pokyny. Správce má plnou správu majetku, má správce v smyslu § 1409 OZ povinnost dbát o rozmnožení majetku a jeho uplatnění v zájmu beneficienta, může tak činit vše, co je nutné a užitečné ke splnění tohoto účelu, a to s péčí řádného hospodáře. Jediným pokynem, respektive mezí pro výkon správy je zakladatelova vůle vyjádřena ve statutu. Jenom soud má tak právo posuzovat právní jednání správce, soud tak může správci udělit na návrh osoby, která na tom má právní zájem, určité jednání v souladu se statutem, a nebo něco zakázat. Osoby mající právní zájem se rovněž mohou obrátit na soud s návrhem na odvolání správce a jmenování nového správce. Oprávněné osoby se také mohou dovolat neplatnosti právního jednání správce bez nutnosti naplnění § 545 a následujících OZ, neboť podle § 1466 OZ odst. 1 OZ se mohou jmenované osoby dovolat neplatnosti právního jednání, kterým svěřenský správce poškozuje svěřenský fond nebo práva obmyšleného\footnote{Tamtéž, strana 37.}.\\

Podle § 1467 OZ \textit{"zakladatel by si rovněž mohl přivodit odpovědnost za případné poškození práv zakladatelova věřitele nebo poškození svěřenského fondu, pokud by se se svěřenským správcem (ale i obmyšleným) dopustil činů sledujících úmyslné poškození takových práv."}\footnote{Tamtéž.}. V takovémto případě by byli dané osoby zavázány společně a nerozdílně.\\

Přestože zakladatel nejspíše nemůže udělovat svěřenskému správci pokyny, je možné větší kontroly dosáhnout pomocí vyhrazení oprávnění zakladatele ve statutu schvalovat určité činnosti správce. Jak tedy poznamenává Lucie Josková, takováto kontrola je zřejmě možná i preventivně, je tedy možné vymezit určitý okruh právních jednání, s nimiž musí vyslovit zakladatel souhlas\footfullcite[Postavení svěřenského správce, strana 180]{trust_2014}. \\

Rovněž se domnívám, že je možné ve statutu toto právo vyhradit i obmyšlenému, nebo třetí osobě\footfullcite[Strana 61]{strelecek_tomas_sverenske_fondy}.\\

% --------------------------------------------------------------------------
% Kro samotné kontroly správce je ještě vhodné uvažovat o oprávnění zakladatele ve vztahu k činnosti správce. Jak poznamenává Lucie Josková, takováto kontrola je zřejmě možná i preventivně, je tedy možné vymezit určitý okruh právních jednání, s nimiž musí vyslovit zakladatel souhlas\footfullcite[Postavení svěřenského správce, strana 180]{trust_2014}. \\
% --------------------------------------------------------------------------

Další osobou, která vykonává \textit{ex lege} dohled je obmyšlený. Obmyšlený má totožné nástroje k výkonu dohledu jako svěřenský správce, nemá však možnost svěřenského správce odvolat, pokud mu toto právo ovšem nevyhradí zkladatel svěřenského fondu ve statutu.\\

Krom zakladatele a obmyšleného může být dohledem podle statutu pověřena rovněž třetí osoba. V případě, že v době vzniku svěřenského fondu ještě není určen obmyšlený, nebo jej v době vzniku fondu nelze určit, jmenuje zakladatel osobu oprávněnou dohlížet na svěřenský fond v zájmu obmyšleného\footfullcite[Strana 38]{svejkovsky_jaroslav_sverenske_2018}.

\subsection{Změny svěřenského fondu}

Zákon vymezuje dva způsoby změny svěřenského fondu\footnote{§ 1468 a § 1469 OZ.}:\\

\begin{itemize}
\item zvýšení majetku ve svěřenském fondu smlouvou, nebo pořízením pro případ smrti\footnote{Avšak ten, kdo zvýší majetek svěřenského fondu není považován za jeho zakladatele.},
\item změnou statutu soudem.	
\end{itemize}

Další potenciální způsob je změna statutu samotným zakladatelem. Tuto možnost zákon explicitně nepřipouští, nicméně v teoretické rovině tuto možnost autoři dovozují\footfullcite[Strana 143]{lucie_joskova_sprava_2017}. Lze se ovšem setkat i s názory, které tuto možnost zcela vylučují\footfullcite[Strana 75]{svejkovsky_jaroslav_sverenske_2018}, osobně se pak ztotožňuji spíše s druhým pohledem, a to nejen kvůli kogentnosti těchto ustanovení, ale i kvůli samotnému jazykovému výkladu, z toho v žádné části nelze dovodit, že by zákon měl umožňovat změny statutu zakladateli svěřenského fondu, nebo komukoliv jinému krom soudu, i samotný soud by při potenciálním provádění změn měl brát dle Kocího\footfullcite[Komentář k § 1469]{svestka_j_obcansky_2014-1} zřetel k vůli zakladatele, která je vyjádřena ve statutu, domnívám se, že zásadní změny statutu jsou tak vyloučeny i ze strany soudu.\\

Pokud se za změny svěřenského fondu dá považovat zvýšení majetku ve svěřenském fondu, domnívám se, že se dá za změny považovat obdobně i plnění ze svěřenského fondu, respektive snížení majetku ve svěřenském fondu, přesto, že i tuto možnost, obdobně jako změnu statutu zakladatele, zákon explicitně nepřipouští\footfullcite[Strana 79]{strelecek_tomas_sverenske_fondy}.\\

\subsection{Zánik svěřenského fondu}

Je potřeba rozlišit ukončení správy svěřenského fondu a samotný zánik svěřenského fondu. Je tak zřejmé, že samotné trvání výkonu správy konkrétním správcem nemůže mít vliv na zánik svěřenského fondu, neboť ten může být odvolán, či výkon jeho funkce může být skončen jinými právem předvídanými způsoby. Svěřenský fond tak může po omezenou dobu fungovat i bez existence správce, dokud nebude nový správce určen\footnote{Obdobně dle Claxtona: \textit{"a trust should not fail for the want of a trustee."}, srovnání kapitola 6.2 Vznik svěřenského fondu.}.\\

Samotné jednotlivé důvody vedoucí k zániku jsou vymezeny v § 1471 OZ.\\

Zanikne tedy z jednak z důvodu \textbf{uplynutí doby}, na kterou byl zřízen, jednak \textbf{dosažením účelu}, pro který byl zřízen\footfullcite[Strana 77]{svejkovsky_jaroslav_sverenske_2018}. K dosažení účelu může dojít u obou druhů svěřenského fondu, tedy založeného za veřejným i z soukromým účelem, obdobně svěřenský fond zaniká, pokud jeho účelem bylo plnění obmyšleným a žádného obmyšleného již není.\\

Pod pojem \textbf{dosažení účelu} lze také podřadit \textbf{rozvazovací podmínku}, při jejímž splnění svěřenský fond zaniká a majetek do něj vyčleněný je převedený na určenou osobu. Takovouto rozvazovací podmínkou může být například dokončení studia, dosažení určitého věku, narození dítěte apod. Samozřejmě tato podmínka může být koncipována i jen k samotnému plnění ze svěřenského fondu. V případě, že obmyšlený danou podmínku splní s tím, že v návaznosti na dané splnění může být předán majetek ze svěřenského fondu obmyšlenému, svěřenský fond může zaniknout, bude-li tak upraveno ve statutu\footnote{Tamtéž.}.\\

Dále může svěřenský fond zaniknout \textbf{rozhodnutím soudu} podle § 1469 OZ je-li dosažení svěřenského fondu nemožné, nebo obtížně dosažitelné po návrhu osoby, jež má právní zájem\footnote{Tamtéž.}.\\

Další důvod, pro který může svěřenský fond zaniknou je, pokud se všichni obmyšlení \textbf{vzdají práva na plnění} u fondu založeného pro soukromý účel\footnote{Tamtéž.}.\\

% --------------------------------------------------------------------------
%Pokud svěřenský správce se zakladatelm sledují účel poškození věřitelů, budou zavázáni společně a nerozdílně.

%Doplnit i zdroje od Zími.
% --------------------------------------------------------------------------

Jak již jsem přiblížil výše, je potřeba od sebe odlišit zánik správy a zánik svěřenského fondu. Totiž v případě, kdy nastane jeden z výše uvedených důvodů, skončí primárně správa svěřenského fondu.\\

Následný proces je v teoretické rovině diskutován, lze se setkat s názorem, že by svěřenský správce měl primárně vypořádat závazky svěřenského fondu, následně by měl vydat zbylý majetek ve svěřenském fondu osobě určené ve statutu, zákon stanovuje vyvratitelnou domněnku, že primárně se vyplácí obmyšlenému, poté zakladateli, není-li jich, tak státu\footfullcite[§ 1472]{noauthor_zakon_2012}, avšak pouze v případě, že ve statutu není určena osoba, které má být po zániku svěřenského fondu majetek vydán.\\

Způsob naložení s majetkem při zániku svěřenského fondu založeného pro veřejný účel je jasně popsán v § 1473 OZ.\\

Po následném vydání zbylého majetku svěřenský fond zaniká. Proces zániku svěřenského fondu se tedy sestává ze 4 fází:\\

\begin{enumerate}
\item skončení správy svěřenského fondu, pro jeden z důvodů uvedených výše,
\item vypořádání závazků\footnote{Srovnání kapitola 7.7 Možnosti použití svěřenského fondu k obcházení věřitelů při zániku správy.},
\item zánik svěřenského fondu, naložení s majetkem dle statutu, nebo dle zákona,
\item výmaz svěřenského fondu z evidence svěřenských fondů, faktický zánik svěřenského fondu.	
\end{enumerate}


\subsection{Svěřenský fond a právní subjektivita}

% --------------------------------------------------------------------------
% Dobrý zdroj je následující: https://advokatnidenik.cz/2019/05/23/sverenske-fondy-nejsou-obchodni-korporace-fikce-podle-zakona-o-skutecnych-majitelich/

%https://advokatnidenik.cz/2020/04/04/vlastnictvi-bez-vlastnika-v-ceskem-sverenskem-fondu-skica-na-margo-trustu-v-ceskem-pravu/

%Dále je ještě jeden zdroj v bakalářské práci kterou lze dohledat v repu, pak má také Michal Skuhrovec, ale od tamtud to nechci úplně celé opisovat.

%https://www.fulsoft.cz/33/pojem-a-vznik-sverenskeho-fondu-uniqueidmRRWSbk196FNf8-jVUh4EnhEI8yvqwHRaEpgGCfo3bkMO6VlftQMFg/
% --------------------------------------------------------------------------

Důležitou součástí definice svěřenského fondu je fakt, že svěřenský fond je koncipován jako oddělený majetek, který je svěřen svému účelu, jež nemá vlastníka a nemá ani právní subjektivitu. Přes jasnou definici, kterou obsahuje Občanský zákoník, důvodová zpráva a komentář k občanskému zákoníku a i samotný zdroj inspirace právní úpravy svěřenského fondu, tedy CCQ, není však na škodu podrobit svěřenský fond jisté komparaci s právnickou osobou. Není totiž tajemstvím, že určité právní normy na svěřenský fond pohlížejí jako na právnickou osobu\footfullcite[§17]{noauthor_zakon_zdp}\textsuperscript{,}\footfullcite[§4b]{noauthor_zakon_zdph} a právě tato tendence k subjektivizaci svěřenského fondu by se mohla projevit i v jiných aspektech svěřenského fondu, tedy například ve vztahu k dědickému právu a pozůstalostnímu řízení. \\

Jak uvádí Bohumil Havel: \textit{česká jurisdikce je běžnou civil law jurisdickí, vystavěnou na recepci římského práva do středověkého práva městského a později do občanského zákoníku Habsburské monarchie z roku 1811 (ABGB)"}\footfullcite[Právní osobnost (právní subjektivita) svěřenského fondu ve světle českého občanského zákoníku, HAVEL Bohumil.]{trust_2014}, je tak logické, že nelze zcela aplikovat duální vlastnictví vlastní common law trustům avšak aplikace bezsubjektového vlastnictví rovněž přinesla do českého právního řádu problémy, i proto na svěřenský fond zákonodárce mnohdy pohlíží jako na právnickou osobu. \\

Pro porovnání svěřenského fondu s obecnou definicí právnický osob je nejprve nutné danou definici vymezit. Zákonodárce tak činí v § 20 odst. 1 OZ: \\

\textit{"(1) Právnická osoba je organizovaný útvar, o kterém zákon stanoví, že má právní osobnost, nebo jehož právní osobnost zákon uzná. Právnická osoba může bez zřetele na předmět své činnosti mít práva a povinnosti, které se slučují s její právní povahou."}\footfullcite[§ 20 odst. 1]{noauthor_zakon_2012}\\

Právnickou osobou je tedy: (A) o které to zákon stanoví, (B) nebo jejíž právní osobnost zákon uzná.\\

Podle (A) tedy nelze hovořit o tom, že by svěřenský fond měl právní subjektivitu, neboť to o něm zákon nestanoví. Dle (B) už se však o právní subjektivitě uvažovat dá.\\

Havel ve svém výkladu právní subjektivitu svěřenského fondu dovozuje\footnote{HAVEL B. \textit{in}: Trust a srovnatelné instituty v Evropě, strana 205.}, nicméně jak zmiňuje Tichý\footfullcite{noauthor_vlastnictvi_2020}, ve své druhé publikaci však právní subjektivitu svěřenského fondu znovu vyvrací a tvrdí že právní subjektivitu svěřenského fondu nelze dovodit, neboť jí zákonodárce vědomě a záměrně nestanovil\footnote{Tamtéž.}.\\

V rámci sousloví \textit{"zákon uzná"} je ve vztahu k definici právnické osoby ve smyslu teorie fikce a organické teorie\footnote{Právnická osoba je organizovaný útvar.} možné právní subjektivitu u svěřenského fondu dovodit, neboť: \textit{"lze konstatovat, že svěřenský fond splňuje veškeré formální předpoklady (znaky) právnické osoby. Je totiž nepochybně organizovaným útvarem, jehož vznik, změna a zánik jsou kogentním způsobem zákonem stanoveny (§ 1448, 1451, § 1468 a násl., resp. § 1471 a násl.), který je zřízen za určitým účelem (§ 1449), a musí mít vlastní označení a strukturu (viz § 1450 a 1452). Jako vlastník (...) má jak majetkovou samostatnost, tak i smluvní a deliktní odpovědnost, kterou v jeho zastoupení vykonává „zákonný komisionář“ – svěřenský správce (§ 1448, 1456). Svěřenský fond tak má majetkovou samostatnost."}\footnote{Tamtéž.}. U svěřenského fondu lze spatřovat s ohledem na výše zmíněné jasné prvky přibližující ho k právnické osobě, stejně jako právnické osoby musí mít stanovy, společenskou smlouvu, nebo zakladatelskou listinu, svěřenský fond musí mít statut, je rovněž evidován a má označení. Jak jsem dále výše zmínil, jisté zákony ho dokonce jako právnickou osobu fingují.\\

V případě pojetí svěřenského fondu jako právnické osoby by tak nebyl v postavení "zákonného komisionáře", ale za svěřenský fond by jednal. Je pak zřejmé, že v případě takovéhoto pojetí svěřenského fondu by byla celá zkoumaná problematika v kapitole 7.4 zodpovězena a svěřenský fond by tak nepochybně byl ve vztahu k pozůstalostnímu řízení dědicem.\\

Jak tedy tvrdí Tichý, lze soudit, že svěřenský fond má právní subjektivitu, která splývá se způsobilostí k právnímu jednání, tedy se svéprávností, dle (B) by tedy zákon měl svěřenský fond \textbf{uznat} jako právnickou osobu. Proti tomuto však stojí důvodová správa, která jasně stanovuje, že svěřenský fond právní subjektivitu nemá\footfullcite[Komentář k § 1448 až 1452]{Duvodovazprava}, rovněž § 1448 odst. 3 konstituuje oddělené vlastnictví nikoli právnickou osobu, přímo však nevylučuje vlastnictví samotným svěřenským fondem.\\

Rovněž se domnívám, že oproti právnickým osobám svěřenský fond stále postrádá mnoho dalších pojmových znaků například ve formě deliktní a procesní způsobilosti, která je nahrazena osobou svěřenského správce\footfullcite[Právní osobnost (právní subjektivita) svěřenského fondu ve světle českého občanského zákoníku, HAVEL Bohumil.]{trust_2014}, který je aktivně i pasivně legitimován ve všech věcech svěřenského fondu\footnote{Srovnání § 79 odst. 1 a § 267b odst. 2 zákona č. 99/1963 Sb.}. Samotnou odpovědnost svěřenského fondu za veřejnoprávní delikty nelze dovodit ani ve vztahu k Zákonu o odpovědnosti za přestupky a řízení o nich, či ze zákona o trestní odpovědnosti právnických osob a řízení proti nim\footnote{Srovnání § 1 odst. 1 zákona č. 418/2011 Sb. a § 20 odst. 1 a odst. 5 zákona č. 250/2016 Sb.}, nesouhlasím tedy zcela s Lubošem Tichým, že například deliktní odpovědnost lze spatřovat přímo u svěřenského fondu.\\

I přes pojmové znaky právnické osoby, jejichž naplnění lze sledovat u svěřenského fondu nelze o svěřenském fondu s ohledem na současnou právní úpravu prohlásit, že by se o právnickou osobu jednalo, jak již jsem popsal v kapitole 6.1, jedná se o bezsubjektové vlastnictví, které nemá právní subjektivitu, ačkoli tedy naplňuje určité znaky právnické osoby. Odpovědi na problematiku plynoucí z pojetí svěřenského fondu jako bezsubjektového vlastnictví, například tedy problematiku nepominutelného dědice (viz kapitola 7.4), nelze nalézt u právní subjektivity svěřenského fondu.\\

Je však pravdou, že mnozí autoři spatřují ve svěřenském fondu a i v jeho předloze ve formě quebecké fiducie právní osobnost a z dogmatické stránky mu ji i přisuzují\footnote{Srovnání TICHÝ Luboš, Vlastnictví bez vlastníka v českém svěřenském fondu (skica na margo trustu v českém právu)}\textsuperscript{,}\footfullcite[POPOVICI A. Trust québeckého a českého práva:autonomní vlastnictví?]{lubos_tichy_sverensky_2017}. Domnívám se, že zde existuje právní podhloubí připravené na to, aby se na něm vystavěl svěřenský fond jako právnická osobnost, ostatně zákonodárce může rozvíjet právo i po samotné dogmatické stránce. Otázkou však je, zda by změnou pohledu na svěřenský fond udělal zákonodárce dobrou službu jeho předloze a nezničil by flexibilitu a jednoduchost svěřenského fondu. Pokud by tak učinil, mohlo by se ztížit uplatnění svěřenského fondu i jako způsob mezigeneračního převodu majetku. V současné chvíli se tedy o právnickou osobu nejedná, na problémy spojené se svěřenským fondem je tak třeba nahlížet s vědomím, že v současné právní úpravě se o právnickou osobu zkrátka nejedná.\\

% --------------------------------------------------------------------------
%Brilantní náhled na toto téma již vyslovil Michal Skuhrovec ve své diplomové práci, která z tohoto hlediska poskytuje vhodný náhled na tuto otázku a také patřičný úvod do této problematiky.\\

%Doplnit do začátku práce, že jsem se touto prací hodně inspiroval, ale myslím si, že neposkytuje odpovědi na všechny otázky a v určitých částech není aktuální, nebo odpovídající, doplnil bych jí tedy o svůj pohled, nicméně některé pohledy autora jsou výtečné a já se s nimi jen ztotožňuji, stejně tak jako se s nimi ztotožňují i pohledy právních expertů. Je nicméně potřeba jeho pohled rozšířit i o nové poznatky v souvislosti s institutem svěřenského fondu

%Na tomto místě je vhodné uvést dvě základní premisy související se svěřenským fondem, které budou v bližším detailu rozvedeny dále. Povaha svěřenského fondu předně vylučuje, aby měl svěřenský fond právní osobnost (svěřenský fond tedy není právnickou osobou) a stejně tak vylučuje, aby byl svěřenský fond věcí v právním slova smyslu. Uvedenému přirozeně neodporují závěry některých autorů, kteří nevylučují, aby měl majetek vložený do svěřenského fondu povahu věci hromadné. 2 HORN, K.: Podrobněji k svěřenskému fondu. Ad notam, č. 6, ročník 2014, s. 1.

%Lze nicméně uvést, že svěřenský fond představuje naprosto ojedinělou a zvláštní entitu, která se skládá z (bývalého) majetku zakladatele a v souvislosti s kterou je prostřednictvím činnosti správce svěřenského fondu naplňován zakladatelem určený cíl. Přitom základním cílem společným pro všechny svěřenské fondy je nepochybně ochrana majetku zejména před nežádoucími dispozicemi 260 260 AKONTinfo, Trust jako nejúčinnější nástroj ochrany majetku. [online]. 2010. On-line zdroj: http://www.akont.cz/cz/375.trust-jako-nejucinnejsi-nastroj-ochrany-majetku.

%PŘEPSAT
%Pihera v tomto konextu vytahuje na povrch diskuzi pocházející z germánského prostoru, kdy Jhering kritizoval koncepci samostatného jmění jako contradictio in adiecto, kdy svou kritiku zakládá na "subjektlose subjektive Rechte", neboli existenci subjektivního práva bez subjektu. Takové by mělo představovat nebezpečí popření ústředního postavení člověka v centru soukromého práva, jak dovozuje Becker. Nicméně, lze oběma autorům přisvědčit, že jeho uznání by při racionálním zhodnocení všech okolností nemělo bez dalšího vést k dotváření "Rechtlosen subjekten". PIHERA NEJPODIVNĚJŠÍ ZVÍŘE V LESE Komentář k OZ pak zároveň dodává, že určité samostatné majetky již v právním řádu 14 známe a vyrovnáváme se s nimi již tradičně, například ležící pozůstalost (hereditas iacens) v ABGB. Dalším právem bez subjektu jsou například práva nenarozeného dítěte.\\

%PŘEPSAT
%Rozdíl mezi úpravou CCQ a OZ pak tkví dále v samotné definici majetku a je otázkou, zda zde český zákonodárce udělal quebeckému trustu dobrou službu co do přípravy "podhoubí". CCQ totiž v § 915 stanoví: "Property belongs to persons or to the State or, in certain cases, is appropriated to a purpose." Majetek patří buď osobám, státu a nebo, v určitých případech, je vyčleněn účelu. Oproti tomu §495 ani §1011 nic takového nestanoví (viz výše). Je tak nutno vycházet z toho, že §1448 OZ je ustanovením speciálním a tvoří základnu teoretické koncepce svěřenského fondu. PIHERA NEJPODIVNĚJŠÍ ZVÍŘE V LESE Lze tak uzavřít, že svěřenský fond je entitou bez právní subjektivity, která je majetkem se samostatným účelovým určením. S ohledem na výše uvedené, podpořeno argumenty zaznívající v diskuzi odborné veřejnosti, se lze přiklonit na stranu zastánců přijaté koncepce v duchu "patrimoine d`affectation, neb zjevně skutečně představuje vhodný most mezi dvěma právními kulturami, podepřený zkušenostmi z Quebecu. Smith\\
% --------------------------------------------------------------------------

\newpage

% --------------------------------------------------------------------------
%\thispagestyle{smallertextinheader}
% --------------------------------------------------------------------------

\subsection{Typy svěřenských fondů}

% --------------------------------------------------------------------------
%http://www.bulletin-advokacie.cz/funkcni-kategorizace-sverenskych-fondu#ftn10

%Podívat se ještě na to, zda existují v českém právu odvolatelné a neodvolatelné svěřenské fondy, a toto také doplnit do těla textu, respektive do této kapitoly

%https://www.obczan.cz/clanky/institut-sverenskeho-fondu-v-noz?do=detail-export
% --------------------------------------------------------------------------

Trust je založen na velké, řekl bych téměř neomezené, variabilitě účelů. Z toho důvodu je možné očekávat, že v rámci české obdoby trustu, jeho využití k velké škále předem nedefinovaných účelů. Z tohoto důvodu není, krom zákonného vymezení, možné poskytnout celkový a komplexní list zahrnující seznam všech potenciálních účelů svěřenského fondu, nebo oblastí jejich využití. V rámci publikací zaměřujících se na toto téma se však lze setkat s určitou kategorizací\footnote{Srovnání jednotlivých zdrojů použitých ke kompilaci formou jejich syntézy, Švestka J., Dvořák J., and Fiala J. a kolektiv. Občanský zákoník. Komentář. Svazek III. § 976 - 1474 . Praha: Wolters Kluwer, 2014. isbn : 978-80-7478-546, Spáčil J. a kolektiv.Občanský zákoník III. Věcná práva (§ 976–1474) . 1. vydání. Praha: C.H. Beck, 2013. isbn : 978-80-7400-499-5, KOCÍ, Miloš. Institut svěřenského fondu v NOZ. Bulletin advokacie, 1-2/2014.} v rámci jejich účelu a využití, s přihlédnutím k liteře zákona a záměru zákonodárce a zahraničního využití svěřenských fondů. S ohledem na toto tedy odborná a komentářová literatura nabízí následující kategorizaci oblastí využití svěřenských fondů:

\begin{enumerate}
\item obchodní svěřenský fond,
\item rodinný svěřenský fond,
\item diskreční svěřenský fond,
\item testamentární svěřenský fond,
\item bezúplatný svěřenský fond mezi živými,
\item zajišťovací svěřenský fond,
\item ochranný svěřenský fond,
\item přidružený fond nadace.
\end{enumerate}

Pro samotný účel ochrany a mezigeneračního převodu majetku je pak při takovéto kategorizaci možné použít svěřenský fond rodinný, testamentární a do značné míry i ochranný. Detailní popis jednotlivých typů svěřenských fondů není vzhledem k rozsahu práce důležitý, níže tedy shrnuji pouze základní vymezení a charakteristiku\footfullcite{koci_m_institut_nodate}\textsuperscript{,}\footfullcite{vitoul_v_funkcni_nodate}. Výjimku tvoří pouze tři výše zmíněné typy a svěřenský fond diskreční, které jsou popsány o něco více dopodrobna.\\

\textbf{Obchodní svěřenský fond} je zakládán k investičním účelům, jeho účelem je vytvářet zisk, může provozovat například obchodní závod.\\

\textbf{Rodinný svěřenský fond} řeší problémy spojené se správou rozsáhlého rodinného majetku a jeho mezigeneračním přechodem, umožňuje širší následná kontrola a intervence zakladatele než jiné nástroje umožňující převod majetku\footfullcite{koci_m_institut_nodate}.\\

Rodinný svěřenský fond jako takový je zásadně zřizován \textit{inter vivos}. Samotné správě rodinného majetku, nebo transferu rodinného majetku to však není na škodu, neboť případná smrt zakladatele nikterak nemusí limitovat možnost naplnění jeho účelu.\\

Pokud by úmysl zakladatele spočíval v založení svěřenského fondu za správou rodinného majetku, nebo mezigeneračního transferu, je založení rodinného svěřenského fondu společně s testamentárním svěřenským fondem jednou z nejlepších možností.\\

\textbf{Diskreční svěřenský fond} poskytuje svěřenskému správci širokou diskreci co do správy fondu, tedy například ve smyslu plnění obmyšleným, či jejich určení. Jde připodobnit k blind trustu\footnote{Jedná s o trust, ve kterém má správce plnou diskreci, přičemž sami obmyšlení nemusí mít přehled o majetkové podstatě trustu a nemohou zasahovat do správy majetku, může být použitý jako řešení možného střetu zájmů.}.\\

V rámci správy a transferu rodinného majetku pak připadá do úvahy jeho užití například v případech, kdy chce zakladatel ponechat na správci určení obmyšlených a obdobně i výše či pravidelnost plnění na základě toho, zda daný okruh osob, ze kterých svěřenský správce obmyšlené může jmenovat, zde se může jednat například o zakladatelovu rodinu, plní určitá zakladatelem stanovená pravidla a kvůli tomu mají nárok na určitou odměnu ve formě plnění ze svěřenského fondu.\\

\textbf{Testamentární svěřenský fond} je založen pořízením pro případ smrti, účel může být opět rozličný, od jednoduchého fondu založeného jen pro posunutí okamžiku předání majetku obmyšlenému (jinak například dědici), tedy o jednoduchou formu (alternativu) bare\footnote{Úkolem správce je majetek pouze udržet, správce ve své podstatě nemá žádné povinnosti krom udržení majetku a pouze vyčkává, až přijde chvíle, kdy má majetek převést dál, touto chvílí může být například dosáhnutí určitého věku obmyšleným. Srovnání De Mond Jr. v. The Queen, 1999 CanLII 466 (TCC), dostupné z: https://www.canlii.org/en/ca/tcc/doc/1999/1999canlii466/1999canlii466.html.} trustu, až po široké ovlivnění majetku po smrti zůstavitele ve smyslu stanovení různých podmínek v rámci statutu.\\

\textbf{Bezúplatný svěřenský fond mezi živými} se dá považovat za specifickou formu darování, nicméně umožňuje určitým osobám čerpat benefit z vyčleněného majetku bez toho, aby se sami stali vlastníky.\\

\textbf{Zajišťovací svěřenský fond} je fondem zřízeným za účelem zajistit splnění dluhu, nebo vyčlenění majetku určeného ke splnění nějakého dluhu.\\

\textbf{Ochranný svěřenský fond} je fondem založeným pro ochranu majetku, stojí na insolvenčním efektu a daný majetek je tak chráněn před věřiteli zakladatele, takovýto převod majetku však nesmí naplnit předpoklady, při jejichž splnění by dané vyčlenění mohlo být věřiteli zakladatele úspěšně odporováno.\\

Ve smyslu správy rodinného majetku a jeho potenciálnímu transferu může sloužit například ve chvíli, kdy se zakladatel rozhodne například vstoupit do rizikové oblasti podnikání, která spolu nese zvýšené riziko neúspěchu a s tím souvisejících potenciálních dluhů.\\

Ochranný svěřenský fond tak může zajistit zakladatelovu rodinu i v případě obchodních a finančních neúspěchů. Je však důležité poznamenat, že neochrání po všech směrech zakladatele, neboť jeho věřitelé se nemohou uspokojovat přímo z majetku ve svěřenském fondu, ale mohou se uspokojovat z plnění, které by zakladatel mohl obdržet jako obmyšlený.\\

\textbf{Přidružený fond nadace} je vnímán jako velmi specifický druh svěřenského fondu. Řídí se úpravou nadace, a to specifiky § 349 až § 352 OZ, tato úprava je ve vztahu k obecné úpravě svěřenského fondu úpravou speciální.\\

\subsection{Ostatní formy správy majetku, rozhodující faktory pro vhodný výběr a z toho vycházející problémy, doporučení a výhody}

% --------------------------------------------------------------------------
%Nadace a nadační fondy, rozdíly oproti svěřenskému fondu. Ty ale mají právní subjektivitu.

%Ona tam hovoří o tom, že je možné založit fond i podle cizího práva, to je možné na základě mezinárodního práva soukromého, nikoli na základě Haagské úmluvy o uznávání trustů. Napíšu tedy, že pro potenciálního zakladatele by mohlo být dobrou volbou i založení právě zahraničního trustu. O tomto pak více například zde: https://dspace.cuni.cz/bitstream/handle/20.500.11956/79456/150028353.pdf?sequence=1&isAllowed=y
% --------------------------------------------------------------------------

Pro zhodnocení, zda je svěřenský fond výhodnější formou mezigeneračního převodu majetku je nejprve nutné stanovit si faktory, dle kterých bude možné obě formy převodu majetku porovnat.\\

Níže vymezené faktory jsou aplikovány specificky na svěřenský fond a pozůstalostní řízení.\\

Vhodný výčet rozhodujících faktorů, který je přímo aplikovatelný na výběr, pro zakladatele, vhodné formy správy majetku dobře shrnuje ve svém příspěvku do Bulletinu advokacie Kateřina Ronovská\footfullcite[Vydání 7-8 2016]{ronovska_bulletin_vyber_formy_spravy_majetku}.\\

S těmito faktory se osobně ztotožňuji, mé potenciální doporučení a upozornění plynou zejména z aplikování těchto faktorů na právní úpravu svěřenských fondů.\\

Docentka Ronovská uvádí následující faktory, které aplikuje na větší množství forem správy majetku:

\begin{itemize}
	\item Faktor 1: stabilita právního prostředí, důvěra v právní řád jako celek, předvídatelnost soudního rozhodování,
	\item Faktor 2: vhodně a proporcionálně nastavený právní rámec; v oblasti práva soukromého zejména ponechání dostatečného prostoru pro realizaci představ vlastníka majetku (při vymezení účelu, způsobu správy, zakotvení práv osob beneficientů), jakož i efekt (ochrana majetku) vůči třetím osobám,
	\item Faktor 3: míra diskrétnosti, resp. ochrany soukromí,
	\item Faktor 4: způsob a adekvátnost dohledu nad správou (důvěra v osobu správce majetku a eliminace možnosti jeho selhání),
	\item Faktor 5: nákladnost zřízení určitého instrumentu, jakož i následné "provozní" náklady spojené se správou,
	\item Faktor 6: "příznivý" nebo alespoň "neutrální", resp. "nedemotivující" daňový režim.
\end{itemize}

% --------------------------------------------------------------------------
%Do takto vymezených faktorů lze zařadit všechny mnou nalezené nedostatky.

% Možná přidám nákladnost/nenákladnost řízení do poslední části práce, možná taky přidám zavedení evidence svěřenských fondů.

%\newpage
%\thispagestyle{smallertextinheader}
%\begin{tabularx}{\textwidth}{X|X|X|}
%\hline
%šmuí & cipí & kuí \\
%\hline
%\end{tabularx}
% --------------------------------------------------------------------------

Je tak nutné porovnat svěřenský fond s více možnostmi převodu majetku, při porovnání s převodem v pozůstalostním řízení má svěřenský fond následující výhody a nevýhody zanesené v tabulce níže.\\

Je však nutné vzít na vědomí, že určité problematické body závisí a odvíjí se od individuální zakladatelovi situace, problém nepominutelných dědiců tak nemusí být problematický v případě, že nepominutelných dědiců nebude, či se svého práva zřeknou\footnote{Viz kapitola 5.1 Stručný úvod do právní úpravy dědictví, vyloučení dědice z dědického práva.}.\\

V úvahu je také nutné vzít, zda se na výhody svěřenského fondu díváme z pozice zakladatele nebo obmyšleného, protože co je pro jednoho výhoda, tedy možnost pořízení a autonomie vůle ze strany zakladatele, může být pro dědice oproti převedení majetku za pomocí instrumentů dědického práva nevýhodou, neboť například pozbyde možností s majetkem plně disponovat. \\

\newpage

\textbf{K jednotlivým faktorům:}\\

\textbf{Faktor 1}: stabilita právního řádu, čili nepozměňování nastavených pravidel až do chvíle, kdy je dle empirických poznatků patrné, že není jiná možnost, je pro správu majetku a mezigenerační převod majetku zásadním faktorem, neboť je jasné, že v mnohých případech se jedná o dlouhodobou správu, či proces předání majetku, který může být podmíněn mnoha skutečnostmi, jež mohou nastat i dlouhou dobu po zřízení institutu správy majetku a mohou dávat vzniknout mnoha budoucím uspořádání majetkových poměrů. Pro potenciální zakladatele je tak stabilita právního řádu a s tím související právní jistota nanejvýš důležitá, pokud právní řád nebude stabilní, atraktivita nových institutů typu svěřenských fondů se značně snižuje.\\

Kateřina Ronovská k tomuto dodává: \textit{"Často se nabízejí různé interpretační varianty, které nelze na první pohled prohlásit za správné či nesprávné. Praxe tak stojí před nelehkou volbou, jaký z výkladů je „ten správný“ či alespoň „udržitelný“? Jaké následky bude mít, pokud se soudní praxe vydá odlišnou cestou?"}\footfullcite[Vydání 7-8 2016]{ronovska_bulletin_vyber_formy_spravy_majetku}.\\

Domnívám se, že u svěřenských fondů je toto zásadním faktorem zejména ve vztahu k nepominutelnému dědici, neboť kolizi dědického práva a práva svěřenského fondu zákon nikterak neupravuje, přičemž soudní praxe se v tomto ohledu žádnou cestou nevydala. Ve vztahu ke svěřenskému fondu tak toto může představovat zásadní právní nejistotu.\\

\textbf{Faktor 2}: ponechání dostatečného prostoru pro projevení a realizaci vůle vlastníka je pro institut, který by chtěl vlastník majetku použít za mezigeneračním převodem majetku, zásadním faktorem. S ohledem na principy občanského práva a relativně široké možnosti dědického práva je stanovení vůle vlastník odlišně od zákona zásadní pro atraktivitu daného nástroje.\\

% --------------------------------------------------------------------------
%co se týče ponechání dostatečného prostoru pro projevení a realizaci vůle vlastníka, poskytují svěřenské fondy téměř neomezenou variabilitu účelů a je tak možné za jejich pomoci dosáhnout téměř jakéhokoliv majetkového uspořádání. Ronovská hovoří\\
% --------------------------------------------------------------------------

Dalším důležitým bodem je jasné vymezení práv beneficientů zákonem nebo zakladatelským právním jednáním, zda je dané právo vybavené právním nárokem a jak je u beneficientů řešen převod samotného majetku v rámci dané struktury při jejím zániku. Společně s tímto je třeba mít na paměti nepominutelné dědice.\\

% --------------------------------------------------------------------------
%Co se týče jasně vymezené pozice beneficientů
% --------------------------------------------------------------------------

Dalším subfaktorem je pak dle Ronovské ochrana majetku vůči nárokům třetích osob.\\

\textbf{Faktor 3}: určitá míra diskrétnosti a anonymity může rovněž hrát důležitou roli. Ve veřejném zájmu je však důležitá i publicita, nicméně pro nástroj převodu majetku by tak obě složky měly být vhodně vybalancovány. V případě právnických osob je nastavena poměrně vysoká míra transparentnosti, což pro převod nebo správu rodinného majetku nemusí být ideální, ale přímo kontraproduktivní. Záleží však hodně na požadavcích vlastníka.\\

\textbf{Faktor 4}: je zásadní pro nástroje, které majetek oddělí od majetku původního vlastníka. Je zřejmé, že takovýto vlastník bude chtít mít přehled a kontrolu nad majetkem, který vyčlenil, aby byla zajištěna dobrá správa a dosahování zakladatelem stanoveného účelu. Funkce takového dohledu spočívá tedy nejen v kontrole správce daného nástroje, ale i ochraně majetku. Úzce souvisí s faktorem 3, neboť kontrola fungování nástroje je do značné míry i ve veřejném zájmu, a to zejména u veřejně prospěšných institutů správy majetku.\\

\textbf{Faktor 5}: důležitou roli pro výběr správné formy správy a převodu rodinného majetku hrají i transakční náklady. Je zřejmé, že nástroj, který by při translaci majetku spolknul jeho velkou část jen na svůj provoz nebude vhodným nástrojem převodů a správy rodinných majetků, jež jsou často svým rozsahem spíše menšího charakteru. Jedná se tak o náklady spojené s ustavením daného nástroje, jeho správou a jeho zánikem.\\

\textbf{Faktor 6}: posledním a neméně důležitým faktorem, který Kateřina Ronovská zmiňuje je daňová zátěž. Systém zdanění takovýchto nástrojů by neměl být příliš znevýhodňující, ale naopak jasný a proporcionální k jejich funkci.\\

Za pomocí těchto faktorů, mnohdy působících protikladně, by bylo možné porovnat svěřenský fond s mnoha způsoby správy a transferu majetku. I přes výhody svěřenského fondu, které lze za podpory empirických dat jasně dovodit není porovnání jednoduché, vždy záleží na nastavení svěřenského fondu a na jeho účelu a správě, jeho porovnání s jinými nástroji převodu majetku tak může býti značně subjektivní. Já tak činím v následující tabulce ve vztahu k pozůstalostnímu řízení.\\

% --------------------------------------------------------------------------
%Neméně zajímavý fakt, který Kateřina Ronovská zmiňuje napříč svým příspěvkem je využití zahraničních struktur pro správu majetku.
% --------------------------------------------------------------------------

\newpage

% --------------------------------------------------------------------------
%zda na svěřenský fond pohlížím optikou zakladatele, nebo obmyšleného, je zřejmé, že

%\begin{table}[]
% --------------------------------------------------------------------------

\noindent\begin{tabularx}{\textwidth}{l|X|X|X|X|}
\cline{2-5}
                               & \multicolumn{2}{c|}{Svěřenský fond} & \multicolumn{2}{c|}{Dědictví} \\ \hline
\multicolumn{1}{|l|}{\textbf{Faktor}}   & \textbf{Výhody}           & \textbf{Problémy}          & \textbf{Výhody}        & \textbf{Problémy}       \\ \hline
\multicolumn{1}{|l|}{Faktor 1} &         

% --------------------------------------------------------------------------
%\begin{itemize}
%\item Nejasná regulace a chybějící judikatura 
%\item Rozdílné pojetí vlastnictví
%\item Zákonné nejasnosti změn svěřenských fondů před a za jejich trvání
% --------------------------------------------------------------------------

Bohužel, k tomuto datu žádná. Útěchou však může být, že institut projde nanejvýše novelizací. Domnívám se, že zákonodárce nemá v úmyslu zrušit možnost zakládání svěřenských fondů.

% --------------------------------------------------------------------------
%\end{itemize}          
% --------------------------------------------------------------------------

&      Strohá právní úprava, neexistence soudní praxe v této oblasti.
Problematika ve vztahu k započtení na povinný díl a k nepominutelným dědicům\footnote{Tamtéž.}. Není zcela iluzorní očekávat, že se právní úprava svěřenského fondu může změnit, neboť se stále jedná o poměrně nový institut.  

% --------------------------------------------------------------------------
%Ochrana dědictví\footnote{Kapitola 6.7.1, Ochrana dědictví před rozmařilostí, nezkušeností a nezodpovědností ze strany dědiců, aneb bod 1.}         
% --------------------------------------------------------------------------

&      Dlouhodobá, stálá a rigidní úprava, poskytuje veskrze všechny potenciální aspekty.          &       Není vyřešena kolize se svěřenským fondem, tento problém je ale možné rovněž přičítat svěřenskému fondu.       \\ \hline
\multicolumn{1}{|l|}{Faktor 2} &         Zde se dostáváme k dominanci svěřenského fondu, možnosti účelů a nastavení statutu svěřenského fondu, stejně tak jako vymezení postavení obmyšleného, jsou opravdu velké, to vše hraje společně s insolvenčním efektem důležitou roli pro určité specifické převody majetku, například v případě obchodního závodu, popřípadě třeba i za filantropickým účelům\footnote{Tamtéž.}.          &        Žádné ve smyslu autonomie vůle zakladatele. Avšak negativum lze spatřovat v tom, že po vzniku svěřenského fondu zaniká právo zakladatela a nevzniká ani právo obmyšleného jakožto vlastnické právo k vyčleněnému majetku, možnost s ním tak disponovat neexistuje. Naopak v případě konečného převodu pomocí instrumentů dědického práva získá dědic daný majetek do výlučného vlastnictví a vlastnická triáda je tak plně uplatnitelná. Rovněž vztahy uvnitř svěřenského fondu jsou pouze obligační, což může sehrát negativní roli například v případě, že by svěřenský správce s majetkem neoprávněně naložil. Jako nevýhodu lze rovněž považovat nemožnost změny statutu\footnote{Tamtéž.}.      &       Také velké možnosti autonomie ze strany zůstavitele, například ve formě dovětku, rovněž placení veškerých pohledávek u předluženého dědictví je možné se vyhout za pomoci výhrady soupisu. Převedení majetku je finální, nový vlastník má jak věcná práva, tak práva vlastnická. Může tak například majetek použít jako zástavu.       &      Autonomie vůle není tak vysoká jako u svěřenského fondu. Možnost nakládání s převedeným majetek je však mnohonásobně vyšší, u svěřenského fondu lze přeci čistě v teoretické rovině hovořit o tom, že obmyšlený je konečný vlastník majetku, toto ale nemá jakoukoliv oporu v zákoně, obmyšlený tak nemá k majetku jakékoliv právo, má pouze právo na plnění a na vydání při zániku svěřenského fondu, určí-li tak statut, a na kontrole činnosti svěřenského správce\footnote{\textit{Stricto sensum} lze tak práva obmyšleného na samotný majetek za doby trvání svěřenského fondu připodobnit k právům člověka k bytové jednotce jeho souseda, nemá tedy žádná věcné nebo vlastnická práva, jedná se o zcela oddělený majetek. Rozdíl je samozřejmě v plnění, finálním transferu majetku a oprávněním obmyšleného/zakladatele vůči správci a tedy v jistých ohledech správě svěřenského fondu.}.        \\ \hline
\end{tabularx}

\newpage

\noindent\begin{tabularx}{\textwidth}{l|X|X|X|X|}
\cline{2-5}
                               & \multicolumn{2}{c|}{Svěřenský fond} & \multicolumn{2}{c|}{Dědictví} \\ \hline
\multicolumn{1}{|l|}{Faktor 3} &         Svěřenský fond poskytuje i přes povinnou evidenci značnou míru anonymity, neboť jsou zveřejněny pouze základní informace o svěřenském fondu a svěřenském správci, listiny obsahující neveřejné údaje se rovněž nezveřejňují\footfullcite[§ 65e odst. 2]{noauthor_zakon_vr}. Princip materiální publicity ve vztahu ke zveřejněným údajům pak rovněž uspokojuje veřejnou důvěru.          &    Samotnou existenci evidence lze však považovat za zásah do soukromí uživatelů svěřenského fondu, osobně se však domnívám, že se o značný zásah do soukromí nejedná.             &        Převáděný majetek nemusí být předmětem žádné evidence, v jistých případech tak může být anonymita větší.        &     V evidencích je zapsaný samotný vlastník, například v katastru nemovitostí je tak možné přímo dohledat jméno majitele, v případě svěřenského fondu je možné dohledat pouze osobu svěřenského správce. V určitých případech tak může být pro majitele vhodnější, aby nebyl zapsán a jeho jméno tak nebylo veřejně dohledatelné.         \\ \hline
\multicolumn{1}{|l|}{Faktor 4} &         Je možné stanovit poměrně široký způsob dohledu a to dokonce aktivního.          &        Zakladatel ani obmyšlený nemají právo udělovat svěřenskému správci přímo pokyny, jejich přímý vliv na majetek je tak tímto omezen. Z povahy svěřenského fondu je to logické, neboť osamostatnění majetku a zvýšená anonymita společně s insolvenčním efektem z logiky věci musí omezit možnosti dispozic s majetkem původním vlastníkem. Ronovská rovněž zmiňuje nedostatečnou regulaci postavení a pravomocí svěřenských správců\footfullcite[Vydání 7-8 2016]{ronovska_bulletin_vyber_formy_spravy_majetku}, osobně se však domnívám, že s ohledem na možnost stanovit ve statutu poměrně extenzivní možnosti dohledu toto nepředstavuje zásadní problém. Další nevýhodu pak spatřuji v nemožnosti ustanovení správcem právnickou osobu specializující se na srpávu cizího majetku tak, jako je možné učinit v rámci jiných zahraničních právních řádů.        &        Majetek nabyde přímo další vlastník, kontrola nad majetkem je tak absolutní.        &    Kontrola nad specifickým majetkem může být velice složitá, například v případě obchodního závodu, v případě svěřenského fondu je možné docílit toho, že majetek bude spravovat profesionální správce lépe, než dědic nebo obmyšlený.          \\ \hline
\end{tabularx}

\newpage

\noindent\begin{tabularx}{\textwidth}{l|X|X|X|X|}
\cline{2-5}
                               & \multicolumn{2}{c|}{Svěřenský fond} & \multicolumn{2}{c|}{Dědictví} \\ \hline
\multicolumn{1}{|l|}{Faktor 5} &         Současná občanskoprávní úprava svěřenského fondu jej činí dostupný pro každého, kdo by o ně měl zájem a vlastní majetek, který by do něj mohl vyčlenit. U svěřenského fondu není před potenciálního zakladatele kladená žádná překážka, a to ať již institucionální, mocenská a dokonce ani finanční, neboť k jejich zřízení a vzniku nejsou kladeny, oproti jiným právnickým osobám, žádné nároky na minimální vložený majetek. Jediné co tak potenciální zakladatel musí mít na paměti jsou základní správní poplatky, které je nutné uhradit, ty však nejsou nikterak závratné výše, a náklady spojené se správou svěřenského fondu, ty budou zahrnovat zejména odměnu správce. Naopak při správném nastavení svěřenského fondu a výběru správce se majetek může kvůli vhodné správě zvyšovat.         & Svěřenský fond je levným institutem, který se hodí i pro správu menších majetků, žádné zásadní výdaje nelze v souvislosti s jeho správou očekávat, právě naopak transakční náklady jsou například oproti právnickým osobám nižší\footfullcite[Strana 312]{hayton_d_j_extending_2002}.                &        Finanční náklady v souvislosti s projednáním dědictví nejsou rovněž nikterak vysoké, bude se jednat zejména o odměnu notáři, jakožto soudnímu komisaři. Naopak následná správa majetku může být levnější než u svěřenského fondu, neboť jí vykonává sám vlastník.        &       Autor si není vědom zásadnějších nevýhod spojených s finanční nákladností projednání dědictví.       \\ \hline
\multicolumn{1}{|l|}{Faktor 6} &         Obecně řečeno jsou svěřenské fondy při plnění z majetku postaveny na roveň dědění, v případě osoby se vztahem k zakladateli\footnote{Viz kapitola 6.11.1, bod 4.}, nebo při zřízení \textit{mortis causa} jsou totiž od daně osvobozeny.          &  Avšak v případě plnění ze zisku je již uplatňována srážková daň společně se zdaněním zisku svěřenského fondu. Může tak docházek k dvojímu zdanění. Jedná se tak v podstatě o stejný režim, který se uplatňuje u jiných právnických osob\footnote{Za předpokladu, že se nejedná o veřejně prospěšného poplatníka.}.              &        Dědictví je rovněž osvobozeno od daně jakožto bezúplatné nabytí\footfullcite[§ 4a]{noauthor_zakon_zdp}, dědicové tak platí pouze odměnu notáři. Navíc je zamezeno dvojímu zdanění v případech, kdy k němu dochází u svěřenských fondů.       &       Ve smyslu zdanění dědictví, nebo zdanění po nabytí majetku žádné nespatřuji.       \\ \hline
\end{tabularx}

\newpage

% --------------------------------------------------------------------------
%\end{table}

%\noindent\begin{tabularx}{\textwidth}{|X|X|X|}
%\hline
%\textbf{Faktor} & \textbf{Problém} & \textbf{Výhoda} \\
%\hline
%Faktor 1 & 
%\begin{itemize}
%\item Nejasná regulace a chybějící judikatura 
%\item Rozdílné pojetí vlastnictví
%\item Zákonné nejasnosti změn svěřenských fondů před a za jejich trvání
%\end{itemize}
%& \\
%\hline
%Faktor 2 & 
%\begin{itemize} 
%\item Systémové nedostatky vztahující se k dědickému právu,
%\item Započtení na dědický podíl,
%\item Nejasný výklad ochrany nepominutelných dědiců,
%\item Nejasný výklad ve vztahu k vyživovací povinnosti
%\end{itemize} 
%& 
%\begin{itemize} 
%\item Systémové nedostatky vztahující se k dědickému právu,
%\end{itemize} 
%\\
%\hline
%Faktor 3 & 
%\begin{itemize}
%\item Zavedení evidence svěřenských fondů (snížení anonimity)
%\end{itemize}
%& 
%\begin{itemize}
%\item I přes zavedení evidence, větší anonimita než v případě převodu a vlastnictví majetku
%\end{itemize}
%\\
%\hline
%Faktor 4 & 
%\begin{itemize}
%\item Neúmyslné chybné rozhodování správce, úmyslné poškození svěřenského fondu a zájmů obmyšlených
%\end{itemize}
%& 
%\begin{itemize}
%\item Možnost stanovení poměrně značné kontroly správce, správa profesionálem, omezení důsledku chybného rozhodování dědiců
%\end{itemize}
%\\
%\hline
%Faktor 5 & 
%\begin{itemize}
%\item Nenákladnost zřízení ale potenciálně vyšší provozní náklady spojené se správou
%\end{itemize}
%& 
%\begin{itemize}
%\item Větší nákladnost řízení, ale potenciálně nižší pozdější náklady
%\end{itemize}
%\\
%\hline
%Faktor 6 & 
%\begin{itemize}
%\item Nejasnosti v rámci daňových aspektů svěřenského fondu, možnost dvojího zdanění
%\end{itemize}
%& 
%\begin{itemize}
%\item Nejasnosti v rámci daňových aspektů svěřenského fondu, možnost dvojího zdanění
%\end{itemize}
%\\
%\hline
%\end{tabularx}
% --------------------------------------------------------------------------

\newpage

\subsubsection{Využití zahraničních trust-like institutů}

Neméně zajímavý fakt, který Kateřina Ronovská zmiňuje napříč svým příspěvkem je využití zahraničních struktur pro správu majetku.\\

ZMPS se svěřenským fondům v části čtvrté, hlavě \rom{7} věnuje v § 73, jež nese název \textit{Svěřenský fond nebo obdobné zařízení}. Fond zřízený podle zahraničního práva se bude řídit daným zahraničním právem, za logické podmínky, že daný řád zvolený trust-like institut upravuje. Pro určení rozhodného práva se uplatní odst. 2 zmíněného paragrafu s přihlédnutím k místu správy, místu, kde je převážně umístěn majetek jej vytvářející, místa sídla nebo obvyklého pobytu svěřenského správce, nebo účelům, za kterým by daný fond vytvořen a místům, kde má být daných účelů dosaženo.\\

Odstavec 4 pak uznává možnost založení zahraničního trustu. Ten se podle něj uzná v případě, že vykazuje zásadní prvky požadované pro něj českým právem\footfullcite[§ 73 odst. 4]{noauthor_zakon_nodate_ZMPS}.\\

 Je samozřejmě otázkou, co zákonodárce považuje za tyto zásadní prvky. S ohledem na článek \rom{2} Haagské úmluvy o uznávání trustů může poskytnout určitou představu, jako takový však není pro samotnou definici v rámci ZMPS určující. S ohledem na fakt, že bezsubjektové vlastnictví je specifikem úpravy CCQ a občanského zákoníku se domnívám, že považovat bezsubjektové vlastnictví jako zásadní prvek trustu není možné, neboť pak by ve své podstatě ZMPS neumožňovalo vytvoření ničeho jiného, než českého svěřenského fondu\footnote{Což z logiky věci nemusí, neboť to umožňuje samotný Občanský zákoník.} nebo quebecké fiducie.\\
 
 Dle mého názoru se tedy budou zásadní prvky trustu shodovat s definicí původního čistě \textit{common law} trustu. Při pohledu zpět na Haytonovu definici bych za esenciální prvek označil do značné míry osobu zakladatele\footnote{Osoba zakladatele nemusí být nutně esenciálním prvkem v případě implied trustů, ty ale dle logiky § 73 odst. 1 ZMPS není možné v České republice konstruovat ani za použití zahraničního práva.}, osobu správce a i osobu obmyšleného, přičemž v tomto rámci je vykonávána správa určitého majetku správcem za nějakým účelem.\\
 
 Zakladatel by tak mohl uvažovat i o zřízení trustů podle cizího práva. Nabízí se řada důvodů proč by tak mohl učinit, jedním z nich pak může být například větší anonymita\footfullcite[SCHRUFF F., strana 159]{trust_2014}.\\ 
 
 S ohledem na fakt, že při zřízení zahraničního trustu \textit{mortis causa} hraje rovněž roli právo dědické by bylo vhodné na závěr dodat, že kolizní úpravu v členských státech EU upravuje nařízení Evropského parlamentu a Rady (EU) č. 650/2012, neboli takzvané \textit{"nařízení o dědictví"}. To však dle Článku 1 vylučuje z působnosti otázky související se zakládáním, správou a rušením trustů. Článek 13 úvodních ustanovení k tomuto však stanoví, že: \textit{"Toto ustanovení by však nemělo být vykládáno tak, že trusty jsou vyloučeny obecně. Byl-li trust založen na základě závěti nebo podle právního předpisu v souvislosti s děděním ze zákona, mělo by se ve vztahu k přechodu majetku a k určení oprávněných osob uplatnit právo rozhodné pro dědění podle tohoto nařízení."}\footfullcite[Článek 13 úvodních ustanovení]{noauthor_narizeni_dedictvi}.\\
 
 Článek 21 jakožto obecné pravidlo pak stanoví následující:\\
 
\textit{"1. Nestanoví-li toto nařízení jinak, je rozhodným právem pro dědění jako celek právo státu, ve kterém měl zůstavitel svůj obvyklý pobyt v době smrti."}\\
 
\textit{"2. Je-li ze všech okolností věci zřejmé, že v době smrti měl
zůstavitel zjevně užší spojení s jiným státem, než je stát, jehož
právo by se použilo na základě odstavce 1, bude právem
rozhodným pro dědění výjimečně právo tohoto jiného státu."}\\

Je tak zřejmé, že až na výjimečné případy, se bude vyčlenění majetku do trust formou \textit{mortis causa} řídit podle českého dědického práva.\\

\newpage

\subsection{Potenciální možnosti použití v rámci převodu majetku a jeho přínos oproti klasickým způsobům pořízení pro případ smrti a jiným způsobům převodu práv}

% --------------------------------------------------------------------------
%Je prokázáno, že pokud je rodina hodně bohatá, tak její našetřené peníze v průměru utratí již 6 generace, nicméně asi to nebude výhoda ve vztahu ke svěřenskému fondu s holedem na maximální vyplácení po 100 letech a kvůli tomu, že svěřenský fond nemůže zakládat další svěřenský fond. Ale raději zjistím, mohl bych totiž potenciálně přednést jako další výhodu a benefit.

%Možná toto pojmu s ohledem na jakékoliv jiné formy převodu majetku, ne jenom dědění, částečně to tak již mam, ale mohl bych to i rozšířit. Musel bych tedy tuto část rozšířit i o porovnání s těmito způsoby převodu majetku.

%Měl bych zde vyjít i z Barbory bednaříkové, dále doplnit zdroje od Lucie Joskové ze Správy Cizího Majetku.
% --------------------------------------------------------------------------

Nový občanský zákoník, který vstoupil v účinnost prvního ledna roku 2014 okruh nástrojů sloužících pro převody a správu majetku nadále rozšířil zavedením obecné úpravy správy cizího majetku a možnosti zřizování svěřenských fondů. S ohledem na téma práce je důležité zmínit, že tato nová právní úprava občanského práva přispěla k rozšíření možností, které lze využít při mezigeneračním převodu majetku, tedy především v případech, které byli dříve řešeny výhradně v rámci dědického řízení popřípadě pomocí darování.\\

Tyto mezigenerační převody majetku spočívají nejčastěji v situacích, kdy takovéto nakládání s majetkem je činěno ve vztahu k mladší generaci v jedné rodině, tedy především k potomkům a prapotomkům. 

% --------------------------------------------------------------------------
%Lingvistický výklad slova mezigenerační v sobě nicméně nezahrnuje pouze pokrevně či právně zpřízněny generace, ale i různé demografické ročníky. Je tak klidně možné tento převod učinit 
% --------------------------------------------------------------------------
Tento převod majetku samozřejmě nemusí být jenom mezigenerační, jedná se i o případy, kdy byl majetek převáděn manželovi, partnerovi, rodinnému příslušníkovi, či kterémukoliv dalšímu příbuznému. Není nicméně vyloučeno, aby takovéto jednání spočívalo v převedení majetku jiným blízkým osobám, jako jsou například přátelé, nebo klidně i osobám zůstaviteli naprosto cizím.\\

% --------------------------------------------------------------------------
%či k jiným blízkým osobám - přátelé, není nicméně vyloučeno, aby takovéto jednání spočívalo ve snaze převést majetek na osobu, která je zůstavitelovi zcela cizí. \\

%Toto budu asi muset ještě upravit a lépe navázat na další odstavec.
% --------------------------------------------------------------------------

Převod majetku z účelem jeho mezigeneračního předání lze samozřejmě dosáhnout i jinými způsoby. Tyto způsoby zahrnují klasické způsoby převodu ve smyslu dědění ale i jiné způsoby nabytí vlastnických práv. Způsobů nabytí vlastnických práv upravuje Občanský zákoník v části třetí, Absolutní majetková práva, hlavě II., věcná práva, dílu 3, oddílu 2 hned několik. V rámci mezigeneračního převodu je však použitelný výhradně pododdíl 6, převod vlastnického práva, na kterém stojí i všechny odlišné způsoby mezigeneračního převodu majetku.

% --------------------------------------------------------------------------
%Doplnit všechny odlišené způsoby převodu majetku a okomentovat je v souvislosti s výhodami mezigeneračního převodu majetku za použití svěřenského fondu
% --------------------------------------------------------------------------

S ohledem na to, že institut svěřenského fondu představuje v české právní úpravě stále poměrně nový přírůstek, tak nejsou zásadněji představeny, natož používány, možnosti, kterých lze prostřednictvím založení svěřenského fondu dosáhnout v rámci mezigeneračního převodu majetku. S ohledem na výše zmíněné a na, i díky událostem poslední doby, větší zájem široké veřejnosti o toto téma, se tato kapitola bude soustředit na možnosti a výhody svěřenského fondu právě v oblasti převodu majetku mezi generacemi a zároveň s ohledem na tyto body, které do značné míry vychází z pojetí svěřenského fondu z pohledu Kocího, načne tedy i potenciální problémy, především s ohledem na kogentní úpravu dědického práva a absenci právní úpravy svěřenských fondů, která by upravovala kolizi mezi kogentní úpravou dědického práva a úpravou svěřenských fondů, které mohou vzniknout při zvolení této cesty společně s nevýhodami, které může tato forma převodu majetku pro všechny zúčastněné strany mít. Část textu se bude věnovat i čistě jurisprudenčnímu pohledu na recepci právní úpravy trustu, spočívající především v pohledu na jednotlivé otázky vyvstávající s přijetím institutu běžnému právní úpravě common law do systému civil law, například ve smyslu rozdílného pojetí vlastnického práva, právní osobnosti.

% --------------------------------------------------------------------------
%Poslední větu bude asi nejspíše třeba upravit. V rámci právní osobnosti myslím právnickou osobu. V případě potřeby upravit tento odstavec.
%Hlavně bude tuto poslední část potřeba přesunout do vhodné sekce v rámci mé bakalářské práce.
% --------------------------------------------------------------------------

\subsubsection{Výhody svěřenského fondu jako alternativního nástroje dědické sukcese}

% --------------------------------------------------------------------------
%Doplnit porovnání i s ostatními prostředky převodu majetku mimo dědictví, to je v další kapitole, doplnit tedy tak například, darování, založení nadace, či jiné právnické osoby a tak dále.

%S ohledem na hlavní téma této práce, tedy svěřenský fond jako mezigenerační nástroj převodu majetku, představují způsoby založení svěřenského fondu a jejich specifika vhodný startovní bod pro tuto část práce a obecně pro celou zkoumanou problematiku. \\

%Základním faktem týkajícím se založení svěřenského fondu je, i s ohledem na popsanou flexibilut instutut trustu, ze kterého český zákonodárce vycházel, je nepřekvapivě možnost založení jak mezi živými (\textit{inter vivos}), tak i pro případ smrti (\textit{mortis causa}).
% --------------------------------------------------------------------------

Kocí pro celý institut svěřenského fondu zmiňuje například následující efekty a výhody\footfullcite[Komentář k § 1448]{svestka_j_obcansky_2014-1}:\\

\begin{enumerate}[label=\alph*)]
\item napomáhá realizovat zejména následující právně-politické cíle a účely,
\item zabraňuje relokaci majetku do zemí, které umožňují flexibilní a efektivní správu majetku formou klasického trustu;
\item obohacuje právní řád o moderní zajišťovací instrument usnadňující zapojení do přeshraničních finančních projektů a globálních bankovních obchodů včetně participace na islámském bankovnictví;
\item umožňuje lidem, které z jakéhokoliv důvodu tíží břemeno správy jejich majetku, aby tuto správu delegovali na jinou osobu;
\item slouží lidem, kteří chtějí mít po své smrti větší vliv na nakládání s jejich majetkem, než by jim umožňovalo dědické právo;
\item usnadňuje mezigenerační přechod majetku;
\item umožňuje lidem sebe nebo své blízké v době blahobytu zabezpečit proti finan­čním těžkostem vzniklým v budoucnosti (k typickým účelům svěřenského fondu viz dále komentář k § 1449),
\item extenze individuální sféry vlastnické svobody a zatraktivnění investičního prostředí.
\end{enumerate}

Na tyto výhody lze pohlédnout jako na zastřešující, níže jsou vyloženy výhody svěřenského fondu ve vztahu k mezigeneračnímu převodu majetku, tedy především ve vztahu k bodům: d), e), f), g) a h). \\ 

Pokud se podíváme na internet a položíme jednomu z mnoha internetových vyhledávačů otázku "Jaké výhody má založení svěřenského fondu", tak nás ohromí nepřeberné množství odkazů na jednotlivé stránky, které jen překypují popsanými benefity. Při mé první exkurzi do světa svěřenských fondů po novelizaci Občanského zákoníku mě toto nepřeberné množství výhod opravdu překvapilo. Od té doby jsem samozřejmě vystřízlivěl a na jednotlivé body již nahlížím mnohem kritičtějším pohledem. Jedna věc se ovšem musí svěřenským fondům nechat, možnosti, které jejich založení přináší zakladateli ve smyslu možností, jak ovlivňovat nakládání se svým majetkem i po své smrti jsou opravdu praktické.\\

Přes toto vše není založení svěřenského fondu pro všechny, neboť má i jisté nevýhody a vyplácí se jen v určitých situacích. Dále nepovažuji zakotvení této právní úpravy v Občanském zákoníku za konečné a troufám si tvrdit, že s ohledem na různé výkladové nejasnosti, je velice pravděpodobné, že dojde ke změnám právní úpravy. Tyto změny, jejichž příklady uvádím v poslední kapitole, jsou nutné k tomu, aby tyto alternativní nástroje mohli zcela plnit svůj účel jako alternativy, nebo suplementy ke klasickým způsobům převodu majetku zejména v oblasti dědického práva. Z tohoto hlediska je možné v tomto směru pochybovat o právní jistotě, neboť se dle mého názoru dočkáme jak hmotněprávních, tak i procesněprávních změn.\\

 Problematické části svěřenských fondů budou nicméně zhodnoceny níže. Tato část je věnována, jak již název napovídá, výhodám tohoto alternativního způsobu předání majetku. 
 
 Pokud se budeme bavit o mezigeneračním převodu majetku, tak se nám jistě primárně vybaví rodinný svěřenský fond. Jak již jsem výše poznamenal, obmyšlený může být samozřejmě kdokoliv, i rodině naprosto cizí člověk. Věřím však, že svěřenský fond jako nástroj mezigeneračního převodu majetku bude sloužit především převodu v rámci rodiny, a z tohoto důvodu se budu dále soustředit výhradně na tento typ svěřenského fondu a na ukázku příkladů k jednotlivým situacím, které je výhodné řešit právě pomocí založení svěřenského fondu.\\
 
% --------------------------------------------------------------------------
 %Popsat i ostatní typy svěřenských fondů
% --------------------------------------------------------------------------
 
 Rodina dozajista představuje největší radost v životě a dá se považovat i za jistý cíl naší existence. Je proto logické, že pro starší generaci představuje rodina motivaci, kvůli které se snaží finančně zajistit jak sebe, tak i svoji rodinu a společně pak pro všechny členy najít bezpečí po všech směrech a ve smyslu všech lidských potřeb. Finanční zajištění saturuje mnoho těchto potřeb a pro mnoho rodičů tedy představuje zachování a předání jistého odkazu ve formě dědictví pro další generace téměř životní poslání. Předání tohoto odkazu může ale mít i potenciální negativní konsekvence, jak říká i známe české přísloví \textit{"Rychle nabyl, rychle pozbyl"}, velkou neznámou tedy může představovat fakt, zda by takto rychle nabyté dědictví mladší generace zvládla. Řešení nejen tohoto problému může představovat právě založení svěřenského fondu, ať už za života zakladatele, či pro případ jeh smrti.\\
 
 Věřím, že nejlepší pojetí této kapitoly představuje kompilace mnoha různých výhod svěřenského fondu a k nim uvedeným příkladům, proto jsem se rozhodl zvolit obdobnou formu jako Barbora Bednaříková ve svém díle Institut pro uchování a převody rodinného majetku. Začal bych tedy, z mého pohledu, největší výhodou svěřenského fondu. Zde se bez jakýchkoliv výhrad ztotožňuji s pohledem Lucie Joskové a Lukáše Pěsny, kteří tvrdí, že základní výhodou svěřenského fondu je samostatné vyčlenění části majetku\footfullcite{lucie_joskova_sprava_2017}. Tato výhoda tedy spočívá ve faktu, že daný majetek je v podstatě majetkem \textit{sui generis}, který nepatří nikomu, ale pouze je spravován svěřenským správcem. Je tedy chráněn před počínáním obmyšlených a jedná se mimo jiné o jistou formu ochrany dědictví. Jedná se zde tedy nejen o základní výhodu samotného svěřenského fondu, která svojí použitelností přesahuje i do mnoho dalších právních odvětví, ale specificky ve vztahu k dědickému právu z této obecné výhody vyplývá i, dle mého názoru, nejzásadnější výhoda a důvod, pro které je vhodné využít svěřenský fond k mezigeneračnímu předání majetku. Touto výhodou je ochrana dědictví jako takového.\\

% -------------------------------------------------------------------------- 
 %Tato ochrana je dle m=ho názoru základní výhodou svěřenského fondu, pokud ho speciálně vztáhneme na tuto funkic. 
 
 %mezigeneračního převodu majetku za pomoci svěřenského fondu, kterou jsem již nakousl výše, tedy jako jistou ochranu dědictví.
 
 %Další výhoda, možnost dát si podmínku, aby se o zakladatele obmyšlení starali v rámci statutu svěřenského fondu, samozřejmě pro svěřenské fondy inter vivos, porovnat s dovětkem
 
 %Ochrana nejen před věřiteli, ale i před třetími osobami, které by mohli chtít obmyšlené, nebo zakladatele využít
 
 %Hodně těchto benefitů lze dosáhnou i klasickými smlouvami, ale všechny jednotlivé benefity svěřenských fondů působí společně s těmi dvoumi nebo třemi hlavními benefity a těmi jsou oddělené vlastnictví, ochrana před věřiteli a anonimizace vlastnictví, pokud by byla tedy zvolena cesta pouze smluvní, tyto benefity by vedle těchto smluvně ujednaných náležitostí, které poskytují tyto benefity dále nepůsobily.\\
% -------------------------------------------------------------------------- 
 
 \begin{enumerate}
 {\Large\item Ochrana dědictví před rozmařilostí, nezkušeností a nezodpovědností ze strany dědiců}
 \item[1.] Mezigenerační transfer a uchování majetku
 \end{enumerate}

% -------------------------------------------------------------------------- 
 %Ochranu lze hodnotit jako jeden ze základních atributů svěřenského fondu. AKONTinfo, Trust jako nejúčinnější nástroj ochrany majetku. [online]. 2010 [vid. 2016-05-29].Dostupné z WWW http://www.akont.cz/cz/375.trust-jako-nejucinnejsi-nastroj-ochrany-majetku. 
% -------------------------------------------------------------------------- 
 
 Přes zákonné nejasnosti\footnote{ V první řadě je důležité zmínit, že následující výklad nemusí být kvůli nedostatečnostem jak lingvistickým, tak i věcným v právní úpravě svěřenských fondů absolutně přesný a je možné, že díky potenciálním problémům vytyčeným níže se může zákonná úprava změnit a s tím společně i tento bod. S ohledem na pohled právních expertů na problematiku svěřenských fondů a mého názoru k datu publikace této práce je vhodné tento bod zařadit v podobě, která se s ohledem na názorovou disputaci a pohledem na procesněprávní náležitosti svěřenského fondu kloní spíše k názoru Miloše Kocího. Tedy, že majetek vyčleněný do svěřenského fondu není součástí pozůstalosti, další informace o tomto problému jsou popsány v další kapitole.} v úpravě svěřenských fondů, se možnost založení svěřenského fondu za účelem ochrany dědictví jeví jako vhodná v případech, kdy se jedná o větší majetek.\\
 
 Nestává se zřídka, že rodiče se strachují, zda jejich potomci bezpečně zvládnou po nich nabyté dědictví. Založení svěřenského fondu tak v tomto případě může pomoci nejen s ohledem na fakt, že potomci nedostanou jednorázově velký finanční obnos, ale i v mnoha dalších ohledech.\\
 
 Nezvládnutí takto nově nabytého dědictví se může projevovat mnoha způsoby, které zahrnují; neuvážlivé a rozmařilé utrácení peněz, ztrátu motivace k vlastnímu výdělku, zanedbávání povinností, špatná finanční rozhodnutí a nerozvíjení finanční gramotnosti, nepoznání hodnoty peněz, práce a času. Tyto rizika lze mitigovat právě založením svěřenského fondu, díky kterému nejen, že dědicové nedostanou jednorázově všechno jmění\footnote{Srovnání další bod, ochrana před věřiteli}, ale v rámci statutu svěřenského fondu jim lze stanovit podmínky, které musí pro plnění ze svěřenského fondu splňovat a plnit\footfullcite[§1452, odst. 2, bod d)]{noauthor_zakon_2012}. Těmito mohou být například požadavek studia, podmínka dosáhnutí určitého věku k vyplácení vyšší části a podobně.\\
 
 Dále lze jako součást procesu založení svěřenského fondu definovat i jisté rodinné poslání, které může sloužit jako opora při stanovení společných hodnot a cílů. Založení svěřenského fondu tak ve svém důsledku nemusí chránit jednotlivce před majetkem a majetek před jedinci, ale může sloužit jako ochrana pro celou rodinu, přičemž tuto ochranu lze hodnotit jako jeden ze základních atributů svěřenského fondu\footfullcite{noauthor_trust_2020-1}. Díky svěřenskému fondu lze sledovat společné cíle všech členů rodiny a vytvořit tak díky tomuto společný mezigenerační plán. Správa rodinného majetku je tak svěřena všem jejím členům, kteří se na ní podílejí a sdílí mezi sebou odpovědnost za rozhodnutí učiněná při péči o vyčleněn\-ý majetek. Tento proces pozitivně formuje potomky a v konečném důsledku je připraví na zodpovědnou správu rodinného jmění.\\
 
 Zde lze opět pochválit zákonodárce, za přijetí svěřenského fondu, neboť absence úpravy ochrany, před marnotratnou osobou před přijetím nového Občanského zákoníku, byla opravdu velkou chybou, neboť tato absence může vést k pauperizaci jednotlivých státem nechráněných osob.\\
 
 Je zde samozřejmě na místě diskuze, do jaké míry by měl stát chránit občany a dělat jim faktického opatrovníka a do jaké by měl nechat jejich počínání čistě na jejich vůli. Nicméně právní stát založený na vzájemné úctě a solidaritě by se měl snažit své občany před chudnutím chránit, a to zejména v případech, kdy by tato marnotratnost mohla dopadnout i na jiné osoby, což by v případě dědického práva do značné míry určitě mohla. Již v antickém Římě byl majetek před tímto počínáním chráněn ve formě opatrovnictví pro marnotratné\footfullcite[Strana 54]{michaela_zidlicka_personae_2015}, kteréžto se svým obsahem nejblíže rovná spendthrift trustu. Zavedení takovéto právní úpravy tedy vede k facilitaci ochrany majetku před marnotratností a posílení ochrany zůstavitelovi vůle vůči tomuto majetku i po jeho smrti a ochrany dalších dědiců, zejména marnotratníkových, nebo obmyšlených.\\
 
% -------------------------------------------------------------------------- 
 %cura furiosi, cura debilium a cura prodigy
% -------------------------------------------------------------------------- 
 
 Je třeba zmínit, že tento bod představuje opravdový benefit pouze v případě, faktické možnosti vydědění pomocí svěřenského fondu, pokud bychom na problematiku nepominutelných dědiců pohlíželi jinou optikou, a to takovou, že nepominutelný dědic má právo i z majetku vyčleněného do svěřenského fondu, tento benefit by pak nebyl tak zřetelný, i v takovém případě by však mohl chránit majetek více, než pokud by zůstavitel ponechal převod majetku na dědickém právu bez pořízení pro případ smrti, tedy za pomoci intestátní posloupnosti.\\
 
 Věřím, že tento účel může představovat jeden z hlavních důvodů, pro který by mohli rodiče zakládat rodinný svěřenský fond \textit{inter vivos} i \textit{mortis causa}, neboť v českém právním řádu nebyli před přijetím NOZ umožněny podobné konstrukce, respektive právní uspořádání, které by takto majetek chránili. Velkou oblibu si tento typ fondu získal například ve Spojených státech amerických pod pojmem "spendthrift" trust\footfullcite[Strana 165]{bednarikova_barbora_sverenske_2014}.\\
 
 Jistou alternativu spatřuji v použití dovětku, ve kterém zůstavitel může stanovit určitou podmínku, lze ale jasně říci, že použití dovětku nelze docílit takové univerzálnosti a všestrannosti, kterou poskytuje svěřenský fond\footnote{Srovnání § 1498 OZ, rovněž přichází v úvahu použití vedlejších doložek v závěti, srovnání § 1551 a § 1552 OZ. Další srovnání: \detokenize{https://www.nkcr.cz/sluzby/dedicke-pravo/dovetek}.}. Obdobně lze spatřit možnost využití paragrafů 724 a 1647 Občanského zákoníku, opět se ale jedná o použití ve specifických situacích, přičemž u druhého z případů musí dojít k vydědění, kdy zůstavitel takovéto marnotratné počínání ze strany svého dědice musí spatřovat již za svého života a nemůže ho použít pouze jako prevenci potenciálního marnotratného počínání.\\
 
 \begin{enumerate}
 \item[2.] Ochrana dědictví před jinými dědici
 \end{enumerate}
 
 Přes obdobné zákonné nejasnosti\footnote{Při tomto bodu lze také namítat nejasnost právní úpravy svěřenských fondů, srovnání kapitola 7.4 Nejasný výklado chrany nepominutelných dědiců.} lze jako na potenciální benefit použití svěřen\-ského fondu k předání majetku nahlížet i na možnost zvolení beneficienta dle uvážení zakladatele. Přes kogentní ustanovení dědického práva je tedy možné stanovit obmyšlené, kteří nespadají do kategorie nepominutelných dědiců, či dokonce obmyšlených, kteří nejsou v žádném příbuzenském vztahu se zakladatelem, a tímto způsobem tak omezit právo těchto dědiců na podíl z pozůstalosti, bez jejich faktického vydědění\footnote{Srovnání s kapitolou Dědické právo v České republice, důvody k vydědění.}\textsuperscript{,}\footnote{Toto lze samozřejmě považovat i za nevýhodu, pokud se na tuto možnost pohlížíme z pohledu dědice.}. Tento způsob se jeví při současném výkladu právní úpravy svěřenského fondu, s porovnáním s právní úpravou dědického práva, jako vhodnější oproti vydědění. Dědické právo nenabízí mnoho důvodů vydědění a není výjimkou, že i když takovéto důvody existují, je složité je prokázat. Tato možnost je vhodná v případech, kdy chce zakladatel zamezit dědickému nároku osob, u kterých neshledává, že by měli mít na část pozůstalosti nárok. Na majetek vložený do svěřenského fondu by tedy tyto dědicové neměli již nárok, neboť takto vyčleněný majetek přestává být ve vlastnictví zakladatele a stává se výlučným vlastnictvím svěřenského fondu a plody a užitky z tohoto majetku náleží výhradně obmyšlenému. I v případě, že by byli nepominutelní dědicové před opomenutím ze strany zůstavitele chráněni i ve vztahu k svěřenskému fondu, má takto zakladatel možnost vyjmout z dědictví velkou část svého majetku.
 
% -------------------------------------------------------------------------- 
 %Hodnota peněz a práce a času
% -------------------------------------------------------------------------- 
 
 \newpage
 
% -------------------------------------------------------------------------- 
 %\thispagestyle{smallertextinheader}
% -------------------------------------------------------------------------- 
  \begin{enumerate}
 {\Large\item[2.] Ochrana před věřiteli}
 \end{enumerate}
 
 Další důležitou výhodu uvádí ve svém díle Správa cizího majetku (dále jen Správa cizího majetku) Lucie Joková a Lukáš Pěšna. Tento bod úzce souvisí s výhodou, kterou jsem již nastínil v předchozím bodu, tedy s převedením majetku na jiný quazi subjekt a tím subsekventní oddělení daného majetku od majetku zakladatele, svěřenského správce i obmyšleného. Věřitelé kohokoliv z výše vymezených osob, v rámci tohoto bodu je jako tato osoba přednostně myšlena osoba zakladatele svěřenského fondu, se tedy nemohou v rámci svých pohledávek za jakoukoliv z těchto osob uspokojovat z majetku, který byl vyčleněn do svěřenského fondu\footnote{Obdobně Barbora Bednaříková, Svěřenské fondy, Institut pro mezigenerační uchování a převody rodinného majetku, strana 168.}, toto platí i v případě úpadku těchto osob. Je nicméně důležité stanovit, že plnění z tohoto fondu ve prospěch obmyšlených může sloužit k uspokojení věřitelů\footfullcite{lucie_joskova_sprava_2017} dle exekučního řádu\footfullcite[§59]{noauthor_zakon_2001}.\\
 
 Výborný příklad je uveden v tomtéž díle: \textit{"Zakladatel se zamýšlí pustit do rizikového podnikání. Vyčlení proto část svého majetku do svěřenského fondu a jako obmyšlené určí sebe a svou rodinu. Tím majetek vyjme z dosahu svých potenciálních budoucích věřitelů. V případě podnikatelského neúspěchu se věřitelé budou moci uspokojit z majetku zakladatele, nikoli však z majetku vyčleněného do svěřenského fondu. To nevylučuje právo věřitelů domáhat se uspokojení z jednotlivých nároků, které zakladateli ze svěřenskému fondu plynou jako obmyšlenému."}\\
 
 Jak dále uvádí autoři, je nicméně nutné podotknout, že si pod založením svěřenského fondu nelze představovat jakýsi nástroj ochrany zakladatelova majetku pro každou situaci. Obecně tak například nelze do svěřenského fondu vyčlenit majetek po vzniku pohledávky, neboť takto vyčleněný majetek by mohl zhoršit dobyvatelnost pohledávky a toto vyčlenění by tak bylo zatíženo relativní neúčinností dle § 589 až 599 obč. zák., popřípadě by mohl být věřiteli odporován dle paragrafů insolvenčního práva, §235 až §243 insolvenčního zákona.\\
 
 Tomuto přisvědčuje nejen Kocí v Komentáři k Občanskému zákoníku jak uvádí správně autoři, ale i Svejkovský a Marek ve svém díle Správa cizího majetku v novém občanském zákoníku v komentáři k paragrafu 1467\footfullcite{jaroslav_svejkovsky_sprava_2015}.\\
 
% -------------------------------------------------------------------------- 
 %\newpage
 %\thispagestyle{smallertextinheader}
% -------------------------------------------------------------------------- 
 
  \begin{enumerate}
 {\Large\item[3.] Anonymizace či diskrétnost}
 \end{enumerate}
 
 V případech snahy o zachování jisté diskrétnosti, pro jednoduchost uveďme, převodce a nabyvatele určitých práv, poskytuje svěřenský fond v tomto ohledu značnou výhodu oproti klasickým způsobům převodu, kdy je ve veřejně přístupných seznamech možné vidět obě osoby ať již přímo, či ve sbírce listin. Novela Občanského zákoníku a Zákona o veřejných rejstřících právnických a fyzických osob z roku 2018 tento bod poměrně zkomplikovala povinným zápisem svěřenských fondů do evidence. Měla ovšem i pozitivní přínos v jiných ohledech, zejména pokud se jednalo o oprávněné obavy, že by svěřenský fond mohl být použit k legalizaci výnosu z trestné činnosti, či dalších nelegálních aktivit. V mnohém tedy tato novela zbourala představu absolutní anonymizace zakladatelé a obmyšlených, kteří jsou teď povinně zapisováni do evidence. Není tomu ovšem tak, je sice pravda, že osoba zakladatele a jednotliví obmyšlení jsou zapsání do evidence, ale data z této evidence jsou poskytnuty jiné osobě pouze pokud prokáže právní zájem\footfullcite[§65e odst. 3]{noauthor_zakon_vr}. Na první pohled je ve všech relevantních seznamech uvedena pouze osoba svěřenského správce.\\
 
 Tato dodatečná bariéra, mezi daty a subjekty, které by dané informace chtěli získat slouží jako další ochrana před věřiteli, či jinými osobami, které by tyto informace mohli využít ke svému prospěchu a poškození zakladatele či obmyšlených, slouží tedy jako firewall, který zajišťuje zakladateli a obmyšleným větší bezpečí a pohodu. Kvůli nečitelnosti vlastnictví je tedy ztíženo zjistit informace týkající se převodu majetku, distribuci výnosů z majetku ať již vně rodinného kruhu, či mimo rodinný kruh\footfullcite[Strana 168]{bednarikova_barbora_sverenske_2014}.\\
 
% -------------------------------------------------------------------------- 
 % \begin{enumerate}
 %{\Large\item[4.] Mezigenerační transfer a uchování majetku}
 %\end{enumerate}
 
 %Vypuštěno, neboť si myslím, že jsem tento bod zahrnul pod prvním bodem
% -------------------------------------------------------------------------- 
 
 \newpage
 
% -------------------------------------------------------------------------- 
 %\thispagestyle{smallertextinheader}
% -------------------------------------------------------------------------- 
 
  \begin{enumerate}
 {\Large\item[4.] Udržení celistvosti a ochrana majetku}
 \end{enumerate}
 
 Založení rodinného svěřenského fondu má také pozitivní vliv na udržení celistvosti a tedy i na ochranu majetku. Toto je důležité zejména v případech, kdy existuje více dědiců a pozůstalost by mělo tvořit jmění, u kterého si lze jen stěží představit možnost jeho rozdělení mezi dědice, popřípadě by takové rozdělení představovalo například značnou konkurenční nevýhodu\footfullcite[Strana 163]{bednarikova_barbora_sverenske_2014}, popřípadě přenesení odpovědnosti na všechny dědice, aby se zabránilo nezodpovědnému počínání jen některých z nich, ovšem bez nutnosti ponižovat podíly náležející těmto obmyšleným. Takovéto jmění může představovat například rodinný závod, nemovitý majetek či různé investiční fondy a podobně. Je samozřejmě možné tento potenciální problém vyřešit dědickou smlouvou, závětí, či darováním. S ohledem na bod 1 si však myslím, že vložení tohoto majetku do svěřenského fondu mu zajistí ochranu a v souvislosti s bodem 1 může obmyšlené i v uvozovkách vzdělat v uvědomělé správě majetku. Vyčlenění majetku do svěřenského fondu také může sloužit k zajištění řádné správy tohoto majetku v případě, kdy by větší problém, než potenciální rozdrobení majetku mohla představovat například potenciální budoucí nemožnost nebo další neochota ze strany zakladatele se o majetek dále starat nebo odůvodněné obavy, že právní nástupci nejsou, ať už z důvodu nedostatku zkušeností nebo nízkého věku, připraveni na povinnosti spojené se zodpovědnou správou majetku.\\
 
S tímto souvisí i v podstatě téměř neomezené možnosti na straně zůstavitele spočívající v obecné výhodě svěřenského fondu a to té, že zakladatel může nastavit správu svěřenského fondu podle svého uvážení, samotná škála variant, které může zvolit je opravdu nepřeberná a záleží pouze na zakladatelově vůli jakou z nich, při uvážení svých potřeb, ochrany majetku a předpokládaných schopností dědice, či obmyšleného, zvolí, respektive nastaví ve formě výplaty ze svěřenského fondu, převádění majetku, zániku či volnosti uvážení ohledně těchto činností ze strany správce.\\

% --------------------------------------------------------------------------
%Diskrece.
 
 %Podívat se jak probíha dědění rodinného závodu a jak rodinný závod funguje ve svěřenském fondu.
% -------------------------------------------------------------------------- 
 
 Příklad uvedený v díle Správa cizího majetku: \textit{"Zakladatel může  do svěřenského fondu vyčlenit závod na výrobu automobilů a potomkům přiznat právo podílet se rovným dílem na zisku dosaženým tímto závodem."} Tento závod může býti přitom spravován profesionálem, potřeba jisté odbornosti dědiců ve směru řízení podniku tak nemusí být nikterak vysoká, je jim pouze přiznáno právo na plnění, aby byly zajištěny jejich potřeby. A dále uvádím pro druhý dílčí přínos zvolení předání majetku ve formě svěřenského fondu tento příklad:\textit{"Otec vyčlení část majetku - akcií a určí jako obmyšlené své nezletilé děti. Správou majetku pověří profesionála zběhlého ve správě cenných papírů a děti budou až do dosažení vysokoškolského vzdělání oprávněni pouze k výnosu z akcií."} V tomto příkladu bude zabráněno rozdrobení majetku vlivem rozdílný názorů a ekonomických zájmů jednotlivých dědiců, či jak již je výše zmíněno nepřipravenosti\footfullcite{lucie_joskova_sprava_2017}.\\
 
% -------------------------------------------------------------------------- 
 %S tímto samozřejmě do značné míry souvisí i obecný výhoda svěřenského fondu spočívající v tom, že zakladatel může nastavit správu svěřenského fondu pdole svého uvážení, samotná škála variant, které může zvolit je opravdu nepřeberná a záleží pouze na zakladatelově vůli jakou z nich zvolí, respektive nastaví. 
 
 %\newpage
 %\thispagestyle{smallertextinheader}
% -------------------------------------------------------------------------- 
 
 Svěřenský fond také dále může být praktickým nástrojem pro správu velkého majetku, který lze rozdělit a toto rozdělení by potenciálně ani nepředstavovalo vážnou konkurenční nevýhodu, ale je praktičtější majetek držet pohromadě než ho drobit mezi dědice. Dobrým příkladem tak může být syndikovaný úvěr\footnote{Syndikovaný úvěr je takový úvěr, kdy je úvěr poskytnut syndikátem, tedy nějakou skupinou věřitelů, tyto úvěry jsou poskytovány v případech, kdy je třeba vysokého objemu zapůjčených finančních prostředků, například při fúzi, akvizici nebo expanzi korporace, refinancování stávajících závazků, nebo realizaci projektového financování jak správně uvádí na svých stránkách například Komerční banka,https://www.kb.cz/cs/firmy-a-instituce/produkty/uvery-a-financovani/strukturovane-financovani/syndikovany-uver}, kdy majetek slouží ku prospěchu velkého počtu osob, jež jsou spojeny společným zájmem, v těchto případech je se správou práv a povinností spojena velká administrativní zátěž, která může být vytvořením svěřenského fondu zmírněna jak pro věřitele, tak i pro dlužníky\footfullcite{lucie_joskova_sprava_2017}. Další příklad uvádí Barbora Bednaříková u rodiny Rockefellerů. Jedná se o velice rozsáhlou rodinu, která svým počtem převyšuje 100 členů. I přes fakt, že mnoho majetku bylo rozdrobeno do menších fondů, je obecně majetek rodiny držen v celku místo osobního vlastnictví jednotlivých členů. Společně s takto značným majetkem, který je centralizovaně spravován, a správnou investiční strategií, dosahují výnosy z takovýchto fondů nemalé částky, z nichž jednotliví členové mohou pohodlně žít. I přes fakt, že by tento majetek tedy potenciálně mohl být rozdělen mezi členy rodiny, tak pokud je soustředěn v celku, výnosy z něho budou mnohem větší a mohou sloužit rodinným členům k jejich potřebám lépe, než kdyby se majetek rozdělil\footfullcite{bednarikova_barbora_sverenske_2014}.\\
 
% -------------------------------------------------------------------------- 
 %v případě, že zakladatel již nechce nebo nemůže řádně majetek spravovat, stejně tak pokud by měl pochybnosti o tom, zdaby jeho právní nástupci například pro nedostatek zkušeností, či nízký věk tuto zprávu zvládli, m, 
% -------------------------------------------------------------------------- 
 
 \newpage
 
 V souvislosti s výše definovanými faktory dle Kateřiny Ronovské je však důležité uvést potenciální negativní daňové dopady vyčlenění majetku do svěřenského fondu oproti dědickému právu. Pro informaci uvádím následující tabulku\footnote{Tabulka převzata z: https://www.svsp.cz/magazin/clanky/vyplati-se-sverensky-fond-z-pohledu-dani, překontrolováno vůči ZDP ke květnu 2020.}:\\
 
 \noindent\begin{tabularx}{\textwidth}{|X|X|X|X|}
\hline
\textbf{Vstup majetku do fondu} & \textbf{Typ obmyšleného} & \textbf{Zdanění transakce} & \textbf{Vysvětlení}\\
\hline
Zisk vytvořený aktivitou SF & bez omezení & 31\% & 19\% daň ze zisku SF; 15\% daň z příjmu obmyšleného \\
\hline
Vyčlenění/zvýšení za života & právnická osoba & 19\% & 19\% daň z příjmu právnické osoby obmyšleného \\
\hline
Vyčlenění/zvýšení za života & fyzická osoba * bez vztahu * & 15\% & 	15\% daň z příjmu fyzické osoby obmyšleného \\
\hline
Vyčlenění/zvýšení za života & fyzická osoba * se vztahem * & 0\% & osvobození výplaty obmyšleného * se vztahem *\\
\hline
Vyčlenění/zvýšení po smrti & bez omezení & 0\% & 	osvobození výplaty obmyšleného z majetku vyčleněného po smrti \\
\hline
\end{tabularx}\\

Obecně lze tedy příjmy související se svěřenským fondem dělit na:\\

\begin{enumerate}
\item Příjmy svěřenského fondu samotného,
\begin{itemize}
\item příjmy z vyčlenění majetku do fondu,
\item příjmy ze zisku fondu,	
\end{itemize}
\item příjmy obmyšleného ze svěřenského fondu,
\begin{itemize}
\item příjmy z nabytí majetku z fondu,
\item příjmy ze zisku fondu,	
\end{itemize}
\item příjmy správce svěřenského fondu.	
\end{enumerate}

Je logické že příjmy fondu, a tedy i jejich zdanění předchází příjmu obmyšleného. Je také důležité mít na paměti, že při plnění ze svěřenského fondu se nejdříve plní ze zisku svěřenského fondu a teprve potom z jeho majetku\footfullcite[§ 21c odst. 1]{noauthor_zakon_ZOZP}. Na vyčlenění majetku do svěřenského fondu se hledí jako na vklad do obchodní korporace, nejedná se tedy o příjem zdanitelný daní z příjmu\footnote{Tamtéž, § 21f.}.\\

Je zřejmé, že z hlediska daňového zatížení není svěřenský fond nikterak výhodný, neboť jeho zdanění je v podstatě obdobné jako zdanění u obchodních korporací, jeho výhody spočívají jinde. Daňovou výhodnost tedy u svěřenského fondu shledat nelze, co ale daňovou nevýhodnost?\\

Zůstavitel by si měl uvědomit, že v případě vyčlenění majetku do svěřenského fondu mohou na majetek a na následné plnění ze zisku z takto vyčleněného majetku dopadnout nepříznivá daňová ustanovení vycházející právě z toho, že zákonodárce na svěřenský fond z pohledu ZDP nahlíží jako na právnickou osobu\footnote{Tamtéž, § 17 odst. 1.}.\\

Lze si tak představit například situaci, kdy by zakladatel do svěřenského fondu vyčlenil 10 bytů, které pronajímá v centru Prahy. Z pohledu samotného majetku, který má hodnotu \textit{X}, neplynou pro obmyšlené žádné negativní daňové konsekvence při plnění z majetku. V případě, že by tak zakladatel stanovil, že majetek má přejít na obmyšlené ve chvíli, kdy dosáhnout 25 roku života, je daňové zatížení postavené na roveň dědictví, pokud by zakladatel měl k obmyšlenému vztah, například by byl jeho otec, bude obmyšlenému vyplacen majetek v hodnotě \textit{X} bez jakéhokoli daňového zatížení\footnote{Tamtéž, § 4a písmeno b).} a je tak potřeba uvažovat pouze nad náklady spojené se správou svěřenského fondu (například odměna správce, náklady by ale vznikly i v rámci pozůstalostního řízení).\\

Jiná situace však nastává v případě, kdy by z pronájmů daných bytů plynul svěřenskému fondu zisk \textit{Y}. Ten by měl tak podat daňové přiznání z příjmu právnických osob a uhradit daň ve výši 19\% z \textit{Y}. V případě, že by zakladatel zároveň stanovil, že obmyšlenému má být vyplácen podíl na zisku, bude tato částka podléhat také srážkové dani ve výši 15\%\footnote{Tamtéž, § 36 odst. 2 písmeno a).}. Obmyšlenému by tak náležela částka ve výši $(Y * 0,81) * 0,85$, při $Y = 100000$ by se tak částka vyplacená obmyšlenému rovnala 68 850 Kč.\\

Pokud by tak svěřenský fond měl existovat po delší dobu, obmyšlený by přišel o poměrně značnou část peněz na dani z příjmu, neboť pokud by dané nemovitosti pronajímal sám jako fyzická osoba, hradil by pouze daň z příjmu fyzických osob ve výši 15\%, měsíčně by tak měl příjem ve výši 85 000 Kč, ušetřil by tak 16 150 Kč měsíčně, což představuje v delším časovém horizontu poměrně značný finanční obnos.\\

Obdobná situace nastává při vyčlenění obchodního závodu do svěřenského fondu. Nejenže svěřenský fond má povinnost hradit daň z příjmu právnických osob ve výši 19\% ze zisku, ale v případě podílu na zisku bude plnění obmyšlenému také podléhat srážkové dani.\\

O kolik peněz může obmyšlený takto přijít lze spočítat následovně:

\begin{center}
$$ M = \left( \frac{Y}{\frac{100}{85}} - \frac{(Y * 0,81)}{\frac{100}{85}} \right)  $$ 	
\end{center}

Kde \textit{M} je množství peněz, o které obmyšlený v případě dvojího zdanění přijde a \textit{Y} je zisk fondu před zdaněním.\\

Je tak na podnikateli a jeho dědici, aby se domluvili, zda jim výhody ve smyslu větší nedotknutelnosti majetku ve svěřenském fondu, vyšší anonymitě, či rozsáhlejší možnosti projevení vůle ze strany podnikatele stojí za to, že faktický zisk bude pro dědice jako obmyšleného menší.\\

Možnost postupu dle § 13 živnostenského zákona tak může být z finančního hlediska pro dědice lákavějším, a to i v případě, že nesplňuje podmínky podle § 6, § 7 a § 8 živnostenského zákona pokud ustanoví odpovědného zástupce\footfullcite[§ 6 až 8, § 13]{noauthor_zakon_ZOZP}.\\

Problém však nepředstavuje vyplácení dividend a jiných podílů na zisku obchodních korporací přes svěřenský fond, neboť v tomto případě je při splnění zákonem daných podmínek zamezeno dvojímu zdanění\footnote{Tamtéž, § 19 bod ze), zi), viz pokyn GFŘ D-22 z 6.2.2015 k § 19 ZDP, dostupné z \detokenize{https://www.financnisprava.cz/assets/cs/prilohy/d-zakony/Pokyn_GFR_D-22.pdf}.}.\\

% --------------------------------------------------------------------------
%povinnost daný příjem \textit{Y} zdanit jako
% --------------------------------------------------------------------------
 
 \newpage
 
% -------------------------------------------------------------------------- 
 %\thispagestyle{smallertextinheader}
% --------------------------------------------------------------------------
 
  \begin{enumerate}
 {\Large\item[5.] Péče o členy rodiny}
 \end{enumerate}
 
 Další výhodu svěřenského fondu spatřuji v možnosti jeho nastavení tak, aby se dal využít pro péči o jednotlivé členy rodiny, kteří danou péči potřebují, ať již díky fyzickému nebo mentálnímu postižení nebo kvůli pokročilému, nebo naopak nízkému věku, a tedy zajištění jejich hmotného zabezpečení\footnote{Na tento bod je rovné možno nahlížet z širšího hlediska, tedy například čistě snaha o hmotné zabezpečení, nebo finanční pomoc obmyšlenému, viz příklad od Barbory Bednaříkové, babička se zřízením svěřenského fondu snaží pomoci vnukovi Adamovi v dokončení jeho studia, viz BEDNAŘÍKOVÁ B., Svěřenský fond, strana 154. Takovýto případ pak do značné míry souvisí se samotným mezigeneračním předáním majetku, srovnání bod 1.}. Těmito mohou býti jak zakladatelé, tak jiní členové rodiny, či dokonce osoby mimo rodinný kruh. Tato výhoda je velice posílena flexibilitou svěřenského fondu, kterou vložil zákonodárce do rukou jeho zakladateli a ten prostřednictvím statutu do rukou svěřenského správce, který poté může vykonávat profesionální správu svěřenského fondu bez toho, aby se jakýkoliv člen rodiny musel angažovat. Díky flexibilitě svěřenského fondu, je možné přesně popsat a vymezit situace a podmínky, za kterých mají jednotliví členové rodiny právo na plnění ze svěřenského fondu. Vhodným nastavením je možné docílit hmotného zabezpečení a péči o osoby, které to nejvíce potřebují po celý jejich život.\\
 
 Toto využití je tak například vhodné v případech, kdy je zakladatel stižen vážnou chorobou, která v budoucnu zapříčiní znemožnění vlastního racionálního úsudku. Založením svěřenského fondu tak může zakladatel docílit toho, že potomci budou mít motivaci se o něj starat a prostředky z majetku takto vyčleněného budou distribuovány výhradně osobám, které o zakladatele projevují zájem a starají se o něj v jakékoliv nelehké situaci a podíl na majetku nebudou mít osoby, které si takovéto plnění nezaslouží\footnote{Tamtéž strana 113}\textsuperscript{,}\footfullcite[Kocí M. Strana 1208]{svestka_j_obcansky_2014-1}. Vedle tohoto benefitu také samozřejmě stojí i výhody popsané v bodech 1, 4 a 6. Toto myslím tak, že například s ohledem na výhody uvedené v bodě 1 se může zakladatel takto rozhodnout, aby uchránil majetek před dědici, či jinými osobami, které by chtěli využít jeho ztížené situace, kdy by dle platných předpisů nebyla ještě učiněna žádná podpůrná opatření při narušení schopnosti jednat\footnote{§38-65 Občanského zákoníku} a byl by tak plně způsobilý jednat po právní stránce, nicméně po stránce mentální by tohoto již nemusel být schopen.\\
 
 Obdobně je možné svěřenského fondu využít jako zajištění pro postižené dítě po smrti rodičů společně s přiznáním odměny za tuto péči jednotlivým osobám.
 
% -------------------------------------------------------------------------- 
 %Je praktické vyčlenit určitou část do svěřenského fondu a jako obmyšlené určit potomky. Tento způsob\\
 
 %Dědic je v rámci postavení k ostatním dědicům v postavení věřitele, je možné se bránit jako dědic jednání zůstavitele za jeho života, v podstatě to nejde, jen pokud by takovéto jednání bylo učiněno pod nátlakem, nebo pokud bychom se bavili o kolaci, tzn. započtení na dědický podíl díl?
 
 %Pokud by byli dva dědicové a zůstavitel by založil svěřenský fond jen pro jednoho dědice, mohl by druhý dědic požadovat, aby se tento majetek započetl na dědický podíl?
 
 %V první řadě je důležité zmínit, že následující výklad nemusí být kvůli nedostatečnostem jak lingvistickým, tak i věcným v právní úpravě svěřenských fondů absolutně přesný a je možné, že díky potenciálním problémům vytyčeným níže se může zákonná úprava změnit a s tím společně i tento bod. S ohledem na pohled právních expertů na problematiku svěřenských fondů a mého názoru k datu publikace této práce je vhodné tento bod zařadit v podobě, která se s ohledem na názorovou disputaci a pohledem na procesněprávní náležitosti svěřenského fondu kloní spíše k názoru Miloše Kocího. Tedy, že majetek vyčleněný do svěřenského fondu není součástí pozůstalosti, další informace o tomto problému jsou popsány v další kapitole.\\ 
 %DOPLNIT

%Ukázat příklady, ke kterým se hodí svěřenský fond, v souvislosti s tématem práce by zde bylo vhodné popsat institut svěřenského fondu jako nástroj k mezigeneračnímu převodu majetku.

%S ohledem na hlavní téma této práce, tedy svěřenský fond jako mezigenerační nástroj převodu majetku, představují způsoby založení svěřenského fondu a jejich specifika vhodný startovní bod pro tuto část práce a obecně pro celou zkoumanou problematiku. \\

%Základním faktem týkajícím se založení svěřenského fondu je, i s ohledem na popsanou flexibilut instutut trustu, ze kterého český zákonodárce vycházel, je nepřekvapivě možnost založení jak mezi živými (\textit{inter vivos}), tak i pro případ smrti (\textit{mortis causa}).
% --------------------------------------------------------------------------

  \begin{enumerate}
 {\Large\item[6.] Přiznání různých práv oprávněným osobám}
 \end{enumerate}
 
 S ohledem na široké možnosti účelu svěřenského fondu a podmínky plnění, které je možné stanovit ve statutu, dále svěřenský fond umožňuje efektivní využívání majetku. Toto efektivní využívání majetku spočívá ve své podstatě v zahrnutí možností poskytnutých dědickým právem v §1507-1524 Občanského zákoníku, tedy možností svěřenského náhradnictví a svěřenského nástupnictví. V rámci statutu je možné, aby zakladatel přiznal jednotlivým obmyšleným různá práva jednorázově i postupně\footfullcite{lucie_joskova_sprava_2017}. Toto se dá využít například tak, že zakladatel přizná určité právo na věci vložené do svěřenského fondu jedné osobě po dobu jejího života a po její smrti či naplnění jiného zakladatelem stanoveného kritéria přejde toto právo na osobu jinou. Pokud by se tak například zůstavitel nemohl spolehnout na svého potomka, kvůli jakémukoliv důvodu, že se nebude schopen o dědictví vhodně postarat, ale chtěl by aby dané dědictví následně přešlo na potomka tohoto potomka, přičemž by ale chtěl zabezpečit i svého potomka, tak použití institutu svěřenského nástupnictví není vhodné ani v případech, kdy by podle §1521 nebo §1522 Občanského zákoníku nebylo svěřeno dědici s dědictvím volně nakládat. V takovéto obdobné situaci je nicméně zřízení svěřenského fondu ve prospěch potomka se stanovením podmínky vyplacení celé podstaty svěřenského fondu po jeho smrti jeho potomkovi\footnote{Toto je plně v souladu s §1460 Občanského zákoníku} velice praktické, neboť majetek vložený do svěřenského fondu může spravovat profesionální svěřenský správce a v souladu s bodem 1 - Ochrana majetku, bude majetek zachován pro další generace. V souladu s důrazem na pořizovací volnost tedy takto bude zajištěno, že zakladatel svěřenského fondu bude mít nepřímý vliv na svůj majetek i po své smrti a o tento majetek bude postaráno dle jeho vůle.\\
 
 Vhodný příklad opět uvádí autoři v díle Správa cizího majetku, pro ilustraci výše napsaného si však dovolím tuto ukázku trochu pozměnit: \textit{"Po smrti zakladatele má vyčleněný majetek (bytovou jednotku) využívat jeho syn s jeho vnukem. Po smrti jeho syna má majetek získat vnuk zakladatele."}
 
   \begin{enumerate}
 {\Large\item[6.] Možnost převedení majetku na jinak nezpůsobilé dědice}
 \end{enumerate}
 
Tato možnost využití svěřenského fondu, respektive účel, za kterým by mohl být založen, souvisí významně s bodem 5. Nejprve je asi vhodné poznamenat, že tuto výhodu spatřuji zejména ve vztahu založení svěřenského fondu \textit{inter vivos} u založení svěřenského fondu \textit{mortis causa} ji považuji spíše za okrajovou.\\

Toto činím především z důvodů, které spatřuji jako potenciální důvod pro založení svěřenského fondu za tímto účelem pro případ smrti. Zde lze považovat za ony důvody především určitou snahu zůstavitele "morálně" vzdělat dědice po své smrti v případech, kdy se oni dědicové za zůstavitelova života k němu zachovali nepatřičně, nebo přímo učinili něco, co by je z dědění po zůstaviteli, bez jeho výslovného odpuštění, přímo vylučovalo, pro tento bod připomeňme kapitolu Dědické právo v České republice a k tomu specifikum dědické nezpůsobilosti.\\

Pokud bychom tedy uvažovali, že dědic byl zákonně vyloučen z dědického práva a zůstavitel mu výslovně toto jednání neprominul, mohl by zůstavitel založit svěřenský fond jako jakousi formu vzdělání dědice, tedy obmyšleného, po své smrti a nastavit možnost výplaty tak, aby správce mohl, při splnění určitých podmínek, takovémuto dědici vyplácet plnění ze svěřenského fondu. Nastavení těchto podmínek je pak plně na vůli zakladatele a od toho se odvíjí i diskrece správce zhodnotit plnění takto nastavených podmínek a přiznávat výplatu ze svěřenského fondu.\\

Tímto způsobem a pro tento účel pak lze obejít dědickou nezpůsobilost bez výslovného prominutí činu vedoucího \textit{a priori} k dědické nezpůsobilosti ze strany zůstavitele, převést na něj majetek a čistě teoreticky dědice i po tomto směru, například tedy po morální stránce, vzdělat.\\

Větší výhody tohoto bodu však spatřuji při založení svěřenského fondu \textit{inter vivos}, zde tedy narážím na podobnost s bodem 5. Příkladem budiž případ, kdy je zakladatel stížen vážnou chorobou a uvědomuje si, že se o sebe nebude schopný v budoucnu postarat. V takovém případě může použít svěřenský fond jako motivaci pro své dědice, aby se o něj postarali. Za oplátku pak při splnění jím stanovených podmínek nechá na správci posouzení, zda byly podmínky opravdu splněny a pokud ano, tak bude dědicům, tedy obmyšleným, vypláceno plnění ze svěřenského fondu i ve chvíli, kdy by podle dědického práva jinak na dědictví neměli nárok. Tento svěřenský fond tedy může sloužit jako jakási šance pro dědice ukázat se v lepším světle a odčinit skutky, které napáchali vůči zůstaviteli a za oplátku získat něco z původního zůstavitelova majetku, na který by jinak v případě pozůstalostního řízení neměli nárok.\\

Při porovnání se samotným prominutím tohoto jinak vylučujícího jednání z dědického práva tak docházíme k závěru, že při použití svěřenského fondu lze vhodně působit na osobnost daného dědice a lze nad ním uvažovat jako nad určitou motivací k polepšení a je z tohoto hlediska i společensky prospěšný. Předpokladem je, že nejspíše nebude mnoho fondů založených za tímto účelem a mnoho potenciálních zakladatelů ani nad podobným účelem, pro který založit svěřenský fond, nebude přemýšlet, tento bod je proto pojmut pouze jako jistá inspirace pro lidi čelící této situaci.\\

% --------------------------------------------------------------------------
%\clearpage
%}
% --------------------------------------------------------------------------
\newpage

% --------------------------------------------------------------------------
%\thispagestyle{smallertextinheader}
%Data.
%\newpage
% --------------------------------------------------------------------------

\section{Systémové nedostatky a návrhy na řešení}

% --------------------------------------------------------------------------
%Zneužití svěřenského fondu viz Babiš, nebo k obcházení věřitelů, dále rozdílné pojetí vlastnictví, započtení, obcházení nepominutelného dědice

%Právní jednání za svěřenský fond, úvěr, závazky, vložení závoud do svěřenského fondu, změna statutu za trvání svěřenského fondu

%Je možné převést do svěřenských fondů například majetek, na kterém vázne zástavní právo? Je tedy možné převést jmění? Tedy nejanom aktiva ale i pasiva

%Pokud je zakladatel i správce, tak je neovlivnitelnost druhého správce, kterého musí povinně mít velkou otázkou.

%Daně nejsou nevýhodou, v rámci příbuzenských vztahů je plnění bezúplatné, problém by byl při plnění ze zisku, nebo pokud by obmyšlenou byla cizí osoba. Při zániku svěřenského fondu a převedení majetku na obmyšleného by se žádná daň platit neměla, tohle ale musím ještě ověřit. Dle článku(portál pohoda), který jsem našel by stejné principy, které platí na plnění ze svěřenského fondu měli platit i na zánik a tedy převedení majetku ze svěřenského fondu obmyšlenému, nebo zakladateli. Jedná se o bezúplatný příjem na základě Zákona o dani z příjmu.

%Co se daní týče tak zde je to hezky shrnuto v oblasti daní z pozemků a nemovitostí

%Daňové aspekty

%Daňová problematika týkající se svěřenských fondů je rozsáhlá a vzhledem k absenci relevantní judikatury či ustálené praxe ne vždy jednoznačná. S ohledem na zaměření tohoto článku se proto omezíme na daňové povinnosti svěřenského fondu u daně z pozemků a ze staveb a daně z nabytí nemovitých věcí.

%Svěřenský fond je poplatníkem daně z nabytí nemovitosti dle ust. § 1 odst. 1 a 2 zákonného opatření Senátu č. 340/2013 Sb. o dani z nabytí nemovitých věcí („zákon o dani z nabytí nemovitosti“) a poplatníkem daně z pozemků a staveb dle ust. § 3 odst. 2 písm. b) a § 8 odst. 2 písm. b) zákona č. 338/1992 Sb., o dani z nemovitých věcí.

%Co se týče daně z pozemků a staveb je svěřenský fond standardním poplatníkem a platí pro něj srovnatelná pravidla jako pro ostatní daňové poplatníky.

%Často řešenou otázkou je však vklad nemovitosti do svěřenského fondu z pohledu daně z nabytí nemovité věci. Jak je uvedeno výše, svěřenský fond je poplatníkem této daně. Předmětem daně je úplatné nabytí vlastnického práva k nemovitosti.

%Vzhledem k tomu, že se při vyčlenění nemovitosti z majetku zakladatele či zvýšení majetku svěřenského fondu o nemovitost dle ust. § 1468 občanského zákoníku nejedná o úplatné nabytí vlastnického práva, není tento převod ani zatížen daní z nemovitých věcí. Nelze zde analogicky uplatňovat úpravu pro nepeněžité vklady do obchodních korporací, neboť zakladatel/vkladatel nezískává podíl na svěřenském fondu ani obdobné protiplnění.[18] K tomuto závěru dochází i odborná literatura.[19]

%Závěrem

%Využití svěřenského fondu pro řešení rozličných situací již nevyvolává údiv (ani úlek), nýbrž je akceptovaným a efektivním východiskem mnohých situací. Výše nastíněný scénář má za cíl ukázat, že i z praktického hlediska se nejedná o nic složitého.

%[19] SVEJKOVSKÝ, Jaroslav a Ivan KOVÁŘ. Svěřenské fondy: příležitosti a rizika. V Praze: C. H. Beck, 2018. Právní praxe. ISBN 978-80-7400-726-2, str. 56.

%Další nejasnost, jak naložit s majetkem ve chvíli, kdy všechni obmyšlení zemřeli? Vrátí se následně do dědictví a dědicům bude náležet povinný díl? Popsaáno v práci Ivany Vladyková, ochrana nepominutelných dědiců a jejich vydědění - Dle https://www.altaxo.cz/poradna/sverenecke-fondy/zanik-a-zruseni-sverenskeho-fondu připadne státu, možná se ještě podívat, že pokud není možné dosahovat účelu svěřenského fondu, tak soud může pozměnit statut a tím pádem i účel nebo majetek převést nějakému subjektu, který má podobný účel jako měl ten svěřenský fond, tak jestli se nestane toto v případě že nebude ani obmyšleného a ani zůstavitele místo toho aby takovýto majetek připadl státu. Jinak by to ale smysl dávalo.

%---Koment\\
%Zhodnocení faktu, že by taková právní úprava měla existovat, poznámka, že tento problém zatím u nás není řešen soudní praxí. Pokusit se navrhnout legislativní změny, vedoucí ke zlepšení postavení nepominutelných dědiců v rámci dědického řízení za předpokladu, že zároveň vznikne i svěřenský fond mortis causa, do kterého se vyčlení majetek, který by jinak byl součástí dědického řízení. Zkusit navrhnout řešení na základě inspirace jinými právními systémy, ve kterých existuje jak institut trustu, tak i ins  titut nepominutelného dědice, mohlo by se jednat například o provincii Quebec, nebo stát Louisiana, pokusit se tyto změny navrhnout na základě studia zákonů a dále se pokusit najít soudní rozhodnutí, v Angličtině se institut nepominutelného dědice nazývá forced heir. \\
%---\\
% --------------------------------------------------------------------------

Je nepochybné, že s přijetím institutu svěřenského fondu vzniklo mnoho otázek, na které nám zákonodárce neposkytl patřičné odpovědi. Zajímavé ovšem je, že zákonodárci v rámci procesu přejímání právní úpravy svěřenského fondu, a její implementací do českého právního řádu, muselo, u mnoha z níže popsaných problémů, být jasné, že takovéto vzniknou. Otázka tedy je, proč byl svěřenský fond takto upraven a proč jeho přijetí nevedlo k odpovídajícím změnám v rámci českého práva, které by takovéto problémy řešily. Věřím, že pomocí analýzy těchto problémů a zodpovězením otázek z nich plynoucích, budu schopen v závěrečné kapitole navrhnou patřičná opatření a změny, které by mohly vést k vyjasnění, či přímému odstranění těchto problémů české právní úpravy.\\

Následující část se zabývá a analyzuje primárně problematiku spojenou s vlastnictvím vyčleněného majetku do svěřenského fondu, které není ani vlastnictvím zakladatele, ani vlastnictvím správce a ani vlastnictvím obmyšleného, se vší jistotou však nepostihují veškeré potenciální problémy, které mohou v souvislosti s použitím svěřenského fondu vzniknout a ani si to nekladou za cíl. Účelem následujících kapitol je tak pouze představit zásadní problémy spojené se svěřenským fondem především ve vztahu k jeho použití za účelem mezigeneračního převodu majetku. Poskytují tak potenciálnímu zakladateli jistý přehled o množství možných typů problémů spojených se svěřenským fondem nad rámec samotného zkoumání problému vytyčeného v hypotéze této práce.\\

% --------------------------------------------------------------------------
%V dílčím závěru této části si také kladu za úkol navrhnout možná řešení této problematiky, tyto jsou směřována vůči zákonodárci.\\

%\newpage

%Ověřit si, že jsem nezapomněl na žádné nedostatky, nebo nejasnosti
% --------------------------------------------------------------------------

\subsection{Nejasná regulace a chybějící judikatura}

Implementace trustu byla bezesporu potřebná a obohatila český právní prostor o mnohé možnosti, které zkrátka nebyli před účinností NOZ použitelné. Je však pravdou, že s ohledem na spíše strohou právní úpravu a dosud chybějící judikaturu\footnote{O tomto se dá přesvědčit při procházení jednotlivých soudních rozhodnutí na portálu justice, dostupné zde: https://www.justice.cz/web/msp/rozhodnuti-soudu-judikatura-.} však vznikají nejasnosti při výkladu některých zákonných ustanovení.\\

Tyto se týkají například pozice svěřenského fondu v rámci projednávání pozůstalosti, problematikou typizace vlastnictví a jednání svěřenského správce či zřízení svěřenského fondu pořízením pro případ smrti.\\

Ani novela občanského zákoníku z roku 2016 nevedla v tomto směru k žádnému zásadnímu zlepšení a naopak zavedla pouze povinnou registraci svěřenských fondů, kterou mnoho právních expertů kritizovalo jako vzdalování se původní angloamerické formě trustu\footfullcite{sro_httpwwwpraguebestcz_prakticke_nodate-1}.\\

Ani další novely občanského zákoníku ze začátku roku 2020\footnote{Novela 163/2020 Sb. a novela 33/2020 Sb.} v tomto ohledu nepřinesly toliko potřebné změny. Žádné zásadní změny institutu svěřenského fondu, nebo jiných ustanovení, která jsou s právní úpravou svěřenského fondu v kolizi a absence konstantní soudní praxe znamená, že výklad zákonných ustanovení svěřenského fondu je zatím stále vytvářen jednotlivými právními experty, tyto názory jsou však velice často v rozporu a nelze se tak na ně vždy stoprocentně spolehnout, výkladem níže bych tedy chtěl vnést trochu světla do temnoty týkající se výkladu ustanovení upravujících svěřenský fond.\\

Snahu o zodpovězení tohoto problému neulehčuje ani fakt, že standardně používané komentáře si při pohledu na tuto problematiku vzájemně odporují.

% --------------------------------------------------------------------------
%výhradně v rukou jednotlivých občanů
% --------------------------------------------------------------------------

\subsection{Systémové nedostatky vztahující se k dědickému právu}

% --------------------------------------------------------------------------
%Zde je můj pohled obecně následující: V případě svěřenského fondu založeného inter vivos se není o čem bavit, zde je to jasné. V případě založení svěřenského fondu inter vivos s odkládací podmínkou smrti zakladatele, respektive zůstavitele, je to složitější, zde musím zjistit, zda je majetek takto vložený do svěřenského fondu součástí pozůstalosti.

%V případě založení svěřenského fondu mortis causa nemá dle mého názoru dle výkladu občanského zákoníku nemá dědic právo na povinný díl a zakladatel se takto efektivně vyhne institutu nepominutelného dědice. Jak je to však s nezletilým dědicem? Majetek je vložen do svěřenského fondu okamžikem smrti zůstavitele, protože svěřenský fond vzniká okamžikem smrti zůstavitele a nemůže existovat bez vloženého majetku. Dědic by se mohl dle mého názoru bránit pouze v souvislosti s neplatností právního jednání vedoucího k založení svěřenského fondu vztahujícímu se k například donucení zůstavitele takovýto fond založit. Dědic by se možná mohl domáhat na správci vydání svého povinného dílu v souvislosti s ustanovením paragrafu 1492 o zákazu zkrácení nepominutelného dědice, pokud by však dostal částku minimálně rovnající se jeho dílu, který by mu jinak dle zákona minimálně připadnul, myslím si, že by se nemohl bránit. Obdobně si myslím, že plnění ze svěřenského fondu se ani nezapočítává na dědický podíl, toto však musím ještě zjistit. Paragraf §1452 přece stanoví, že se použije ustanovení paragrafu 311, ale podle mě v souvislosti s paragrafem 311 proběhne jenom jakýsi quazi stejný proces jako u dědického řízení, tzn. majetek je vložen v rámci dědického řízení, ale není součástí pozůstalosti a fond, protože fond nemá právní subjektivitu, zastupuje při tomto jednání svěřenský správce. Takto vyčleněný majetek ale dle mého názoru stále není součástí pozůstalosti a svěřenský správce plní v řízení jenom jakousi informační povinnosti a o tomto majetku se ani nevede dědické řízení, takže je součástí jenom jakoby, ale nejedná se o něj.

%Důležitý je taky okamžik smrti, poněvadž se majetek vylčení do svěřenského fondu ve stejnou cvhíli jako vznikne pozůstalost. Dobrá je argumentace Michala Skuhrovce, ale zapomněl na paragrafy 311 a 1492 a 1452.

%Tady bych chtěl ještě doplnit toto: https://cs.wikipedia.org/wiki/V%C3%BDklad_pr%C3%A1va.
% --------------------------------------------------------------------------

Věřím, že na problematiku spojenou s kolizí norem dědického práva a svěřenského fondu lze nahlížet dvěma způsoby, podle kterých záleží, na jakou stranu názorového spektra se v souvislosti se zodpovězením mé hypotézy přikloním.\\

První pohled bych nazval "výkladový", pomocí jeho optiky nahlížím na kolizní problematiku z pohledu formálního výkladu jednotlivých norem a jejich vzájemných vztahů s ohledem na základní principy svěřenského fondu, práva a logiky.\\

% --------------------------------------------------------------------------
%Pohled jak je to napsané v tom samotném předpisu a jak tomu lze po výkladové stránce v souvislosti s porovnáním kolizních norem porozumnět. Dle tohoto pohledu je například chybný paragraf 311, protože není v souladu s tím, že svěřenský fond není jako takový účasten dědického řízení a vyčlenění majetku do něj nepotvrzuje soud, majetek takto vyčleněný tedy ani není součástí pozůstalosti.

%Pokud bychom tak porovnávali dva paragrafy, kdy jeden by nedával smysl, tak bychom ho při vyložení hypotézy nebrali v potaz.
% --------------------------------------------------------------------------

Druhý pohled bych nazval "úmyslový", který je vykládán jako mé pochopení potenciálního úmyslu zákonodárce a jeho směřování při implementaci tohoto nového právního institutu, krom samotného výkladu smyslu jednotlivých právních ustanovení, se zde zaměřuji i na logicky vyvozené korelace mezi jednotlivými zákonnými ustanoveními a z toho plynoucí kauzalitu, kterou spatřuji v několika zákonných ustanoveních, jejichž výklad nelze ztotožňovat s výkladovým pohledem, neboť s ohledem na transplantaci právní úpravy svěřenského fondu je vyloučený. V jeho rámci se zaměřím spíše na teoretickou rozpravu s ohledem na základní principy dědického práva a svěřenského fondu v souvislosti s materiálními prameny práva.\\

% --------------------------------------------------------------------------
%Pohled jak to myslel zákonodárce, zmínil paragraf 311, proč by to dělal, pokud by nechtěl, aby byl aplikovatelný. S ohledem na to že dle zákona o zvláštních řízeních soudních je například svěřenský správce účastníkem dědického řízení a plní tedy takovou quazi informační a informativní funkci, je možné, že by mohl zastupovat samotný svěřenský fond a ten by potom majetek nabýval potvrzením soudu.

%Tohle budu muset možná trochu ještě přepsat.

%a smysl respektive základní principy svěřenského fondu, pokud budeme připodobňovat svěřenský fond k právnické osobě kdykoliv budeme chtít, popíráme tím základní principy svěřenského fondu a lepaullova vlastnický na jehož základě svěřenský fond stojí a na jehož základě jako účelově nezávislé jmění jsme ho přijali. Pokud bychom mu přiřazovali quazi právní subjektivitu kdykoliv, kdy by se nám to hodilo, tak proč bychom ho tedy vůbec měli, mohli bycho mu buď přiřadit vlastní právní subjektivitu, která by vyřešila mnoho problémů, ale proč bychom to potom nazývali svěřenským fondem, ale místo toho jsme jen neuvolnili podmínky, které je nutné splnit při zakládání nadace. A proč bychom tedy svěřenský fond takto přijímali pokud bychom jeho výklad směřovali k tomu, že by měl právní osobnost.

%Soud by tak mohl rozhodovat i s ohledem na neplatné paragrafy kvůli veřejnému pořádku.

%a budu prezentovat svůj názor
% --------------------------------------------------------------------------

\subsection{Započtení na povinný díl a dědický podíl}

Obdobně jako problematiku ochrany nepominutelného dědice je nutné minimálně stručně nahlédnout i na problematiku započtení na povinný díl a dědický podíl.\\

Problematika se dle mého názoru ilustruje lépe na příkladu. Představme si tak situaci, kdy po zůstaviteli zbydou dva dědicové a čistý majetek ve výši 3 000 000 korun. Zůstavitel zároveň v pořízení pro případ smrti přikáže zřídit svěřenský fond a vložit do něj 1 000 000 korun. Na dva zletilé nepominutelné dědice, syny zůstavitele, tak zbude 2 000 000 korun. V tomto případě je s ohledem na § 1643 OZ (2) splněn zákonný požadavek, aby se zletilému dědici dostalo minimálně jedné čtvrtiny jeho zákonného podílu. Problém však může nastat, pokud by měl jen jeden z těchto dědiců právo na plnění ze svěřenského fondu. Zákona jasně nestanovuje, zda je takovýto majetek, nebo plnění z něj poskytnuté možno započíst na povinný díl.\\

Pro analýzu tohoto potenciálního problému je nejprve nutné rozdělit takto vyčleněný majetek na majetek vložený do svěřenského fondu a na následné plnění ze svěřenského fondu.\\

Při pohledu na majetek vložený do svěřenského fondu by započtení na povinný díl nemělo být dle mého názoru možné, a to kvůli samotné definici toho, co lze započíst na dědický podíl. Důležitá část § 1660 OZ je pak následující: \textit{"Na povinný díl se započte vše, co nepominutelný dědic z pozůstalosti skutečně nabyl..."}, o nepominutelném dědici, nebo potomkovi se hovoří i v ostatních ustanoveních započtení na povinný díl. Důležité je si však uvědomit, že v případě pohledu na svěřenský fond jako na dědice, má postavení dědice tento svěřenský fond, či svěřenský správce, nikoli druhý nepominutelný dědic. Toto píše ve svém příspěvku i Horn: \textit{"Proto autor nesouhlasí s názorem, podle kterého by bylo nutné započtení plnění poskytnutých obmyšlenému ze svěřenského fondu na jeho souběžný dědický podíl podle § 1662– 1664 OZ, neboť postavení dědice sui generis zde má svěřenský správce a nikoliv obmyšlený."}\footfullcite{sro_httpwwwpraguebestcz_prakticke_nodate-1}. Na dědický podíl druhého dědice tak není možné tento majetek započítat, neboť byl odkázán jinému dědici, což je plně v souladu s ustanoveními o nepominutelných dědicích a testovací volností zůstavitele.\\

Obdobného pohledu na věc bychom docílili i za předpokladu, že bychom majetek vyčleňovaný do svěřenského fondu nepovažovali za součást pozůstalosti a svěřenský fond, nebo svěřenského správce za dědice. I v takovémto případě by totiž nebyl "obdarován" druhý nepominutelný dědic, ale jiná entita. Stejně lze tak proto nahlížet i na svěřenský fond \textit{inter vivos}, ani zde se nedá, i přes existenci § 1660 (2), započíst takovýto majetek na dědický podíl neboť byl obdobně převeden do vlastnictví jiné entity nežli daného nepominutelného dědice na jehož podíl by bylo žádáno o započtení.\\

Je tedy možné započíst alespoň plnění ze svěřenského fondu obmyšlenému, který by zároveň byl jedním z nepominutelných dědiců? Dle mého názoru rovněž nikoli, neboť ani v případě plnění ze svěřenského fondu založeného \textit{inter vivos} se nejedná o věc, kterou dědic obdržel od zůstavitele jak stanoví § 1660 OZ, obdržel ji totiž od svěřenského fondu, nikoli od zůstavitele. V případě svěřenského fondu zřízeného \textit{mortis causa} pak rovněž nebude možné uvažovat o započtení, neboť o plnění ze svěřenského fondu se nedá uvažovat jako o plnění z titulu odkazu, nebo jiného zůstavitelova opatření, navíc v časové ose bude následovat až po započetí pozůstalostního řízení.\\

Není však zcela bezpředmětné uvažovat nad tím, zda by toto plnění přeci jenom nemohlo být započteno na dědický podíl, jistý prostor pro diskuzi totiž zákon připouští. S jistou dávkou imaginace je proto možné považovat plnění ze svěřenského fondu, zejména zřízeného \textit{inter vivos} za "obdarování" nepominutelného dědice, neboť svěřenský správce koná a vyplácí obmyšleným dle statutu svěřenského fondu, který konec konců utváří zakladatel. Problematika kolačního práva tedy může představovat potenciální problém, který by mohl narušit konstruování svěřenského fondu, zákon totiž jasně nestanovuje, zda je možné v případě konstruování svěřenského fondu vyčleňovaný majetek, nebo plnění započítávat na povinný díl nebo dědický podíl.\\

K tomuto například Kocí uvádí: \textit{"Pokud by s majetkem ve svěřenském fondu mělo být zacházeno jako se součástí pozůstalosti, bylo by případně nutné provést zejména započtení plnění poskytnutého obmyšleným na dědický podíl (§ 1662 až 1664), ačkoliv tito nemají k plnění ze svěřenského fondu vlastnické právo (§ 1448 odst. 3) a u diskrečních fondů navíc ani nemusí být jisté, zda jim právo na plnění vůbec vznikne, popřípadě v jaké výši (viz níže sub II). To by v rámci dědického řízení vedlo ke značným komplikacím a průtahům."}\footfullcite[Komentář k § 1457]{svestka_j_obcansky_2014-1}.

% --------------------------------------------------------------------------
%Specificky zmiňuji pouze započtení na povinný díl, neboť započtení na dědický podíl se provede jen přikázal-li to zůstavitel projevem vůle učiněným ve formě předepsané pro pořízení závěti\footfullcite[§ 1663]{noauthor_zakon_2012}, je tedy třeba aktivity ze strany zůstavitele, vyjma § 1664, je pak v zásadě bezpředmětné se těmto ustanovením věnovat v případě opomenutí nebo omezení neopominutelného dědice, ačkoli je možné dovodit problematické situace i zde, s ohledem na téma práce se budu věnovat pouze započtení na povinný díl.

%Způsobilost být obmyšleným

%3. Zákon blíže nespecifikuje, kdo má způsobilost být obmyšleným, neboť se vychází z toho, že je zcela na uvážení zakladatele, koho bude považovat za způsobilého a koho nikoliv. Může se jednat o člověka i o právnickou osobu včetně osob nenarozených nebo právnických osob ještě nezaložených (§ 1464), kdy hovoříme o tzv. eventuálních obmyšlených (v Québecu bénnéficiaires éventuels).

%4. V této souvislosti je třeba uvést, že v zájmu spravedlivého uspořádání vztahů vzniklých založením svěřenského fondu důvodová zpráva proponuje, že „obmyšleným svěřenského fondu zřízeného bezúplatně může být jen osoba, která je způsobilá po zakladateli v den vzniku svěřenského fondu dědit“. Zde důvodová zpráva evidentně vychází z reality CCQ, který ve svém čl. 1279 stanoví, že „beneficientem bezúplatného svěřenského fondu může být pouze osoba, která má v době vzniku svého práva vlastnosti potřebné pro obdržení daru nebo dědictví“. Toto ustanovení však v účinném znění občanského zákoníku obsaženo není, a proto by s ohledem na zásadu „vše je dovoleno, co není zákonem výslovně zakázáno“ mělo být zakladateli umožněno, aby učinil obmyšleným z testamentárního svěřenského fondu i osobu, která není způsobilá být dědicem. Obmyšlený z testamentárního svěřenského fondu totiž nutně dědicem být nemusí, a proto není důvodu, aby byl podroben stejným standardům jako dědic jen proto, že bylo ke zřízení svěřenského fondu mortis causa využito závěti, tj. institutu dědického práva (viz § 1448 odst. 1). Tento závěr by měl platit tím spíše, že majetek vyčleněný do svěřenského fondu je vyloučen z dědického řízení, neboť již samotnou smrtí zůstavitele je vyčleněn do svěřenského fondu, a není tedy součástí pozůstalosti.

%5. Pokud by s majetkem ve svěřenském fondu mělo být zacházeno jako se součástí pozůstalosti, bylo by případně nutné provést zejména započtení plnění poskytnutého obmyšleným na dědický podíl (§ 1662 až 1664), ačkoliv tito nemají k plnění ze svěřenského fondu vlastnické právo (§ 1448 odst. 3) a u diskrečních fondů navíc ani nemusí být jisté, zda jim právo na plnění vůbec vznikne, popřípadě v jaké výši (viz níže sub II). To by v rámci dědického řízení vedlo ke značným komplikacím a průtahům. Zakladatel by způsobilost obmyšleného k plnění ze svěřenského fondu teoreticky mohl vázat např. i na mnohem přísnější podmínky, než jsou podmínky předvídané zákonodárcem pro dědickou nezpůsobilost (§ 1481 až 1483). Pokud tak učinit z nějakého důvodu nechce, měla by být jeho vůle respektována, stejně jako jeho rozhodnutí nepodrobit obmyšleného žádným podmínkám. Konkrétně se může jednat např. o situaci, kdy chce zakladatel obmyšlenému kompenzovat zřeknutí se dědického práva (§ 1484), které je důvodem nezpůsobilosti k dědění v širším slova smyslu a které má v konkrétní situaci např. umožnit hladší mezigenerační přechod majetku zakladatele v duchu rodinných tradic, anebo o situaci, kdy se zakladatel obává, že by jeho potenciální dědic mohl být z dědického práva vyloučen z důvodu spočívajícího v dědické nezpůsobilosti v užším slova smyslu (§ 1481 až 1483), o němž zatím nemá vědomost, a chce tak tvrdému právnímu následku předejít tím, že pro takový případ ze svého dědice učiní zároveň i obmyšleného ze svěřenského fondu. Je třeba zmínit i skutečnost, že občanský zákoník nikde nepočítá s tím, že by svěřenský správce nebo obmyšlený měli být v jakékoliv roli účastníky dědického řízení. Zřízení svěřenského fondu přitom v žádném případě nelze považovat za svěřenské nástupnictví nebo jeho druh (§ 1512).

%6. Konečně dodejme, že u svěřenských fondů založených za soukromým účelem představuje určitou časovou limitaci způsobilosti obmyšleného ustanovení § 1460 (viz níže).
% --------------------------------------------------------------------------

\subsection{Nejasný výklad ochrany nepominutelných dědiců}

% --------------------------------------------------------------------------
%Doplnit ještě zdroje z práce od Tomáše Střelečka

% První bod popsat ty moje dva rozdílné pohledy, druhý bod pak posat ty jednotlivé pohledy autorů, 4 pohledy, pak napsat, že si nemyslím, že je takový problém znevýhodnění svěřenského fondu mortis causa oproti svěřenskému fondu inter vivos. Viz poznámky, které jsem si poslal na mail a napsal na druhém počítači. Pak také popsat to, že již nebude neplatné celé pořízení pro případ smrti, ale pouze vznikne opominutému dědici pohledávka za jinými dědici (tady je další problém, je svěřenský fond dědicem?) tohle je v souvislosti s paragrafem upravujícím pořízení pro případ smrti, a který zároveň stanovuje nutnost respektovat povinný díl nepominutelného dědice.
% --------------------------------------------------------------------------

V České republice, obdobně jako ve většině zemí používajících kontinentální systém práva, právní řád chrání nepominutelné dědice před opomenutím ze strany zůstavitele a přiznává jim nárok minimálně na jejich povinný díl v zákonné výši. Ovšem jak již je zmíněno v kapitole výše, v rámci české právní úpravy bohužel chybí explicitní úprava funkce svěřenského fondu, pokud se nachází ve střetu s institutem nepominutelného dědice. Pro úplnost připomínám kapitolu 5.2 Nepominutelný dědic.\\

Inkorporace svěřenského fondu, dle mnohých, pak přinesla možnost zřídit obecný nelimitovaný fiduciární institut. Tato možnost společně na straně jedné s vůlí zákonodárce chránit nepominutelného dědice na straně druhé logicky nejsou slučitelné\footfullcite[Strana 333]{jaroslav_svejkovsky_sprava_2015}. Diskuze ohledem aplikace právních předpisů upravujících povinný díl a jejich souvislost se založením svěřenského fondu \textit{mortis causa} je tedy na místě.\\

Jak již jsem nastínil výše, dle mého názoru je na tuto problematiku možné nahlédnout dvojí optikou. V prvé řadě bych však chtěl poznamenat, že se v této části budu věnovat čistě svěřenskému fondu založenému pro případ smrti, neboť jak již jsem zmínil v kapitolách 6.3.1 a 6.3.3 omezení založená právní úpravou dědického práva se na svěřenský fond \textit{inter vivos} nevztahují.\\

Dobrý postoj k tomuto problému zaujal například Michal Skuhrovec ve své diplomové práci, která pojednává o podobném tématu jako má práce\footfullcite[Strana 54]{skuhrovec_michal_sverensky_2018}. Nemyslím si však, že názor jím vyslovený je kompletní a zahrnuje širší problematiku kolize právní úpravy nepominutelného dědice a svěřenského fondu. Z tohoto důvodu bych chtěl názor vyslovený v dané práci konfrontovat a rozšířit o můj pohled na problematiku, na níž jsem pohlédl z o něco širší perspektivy.\\

Skuhrovec na začátku hovoří o svázanosti evropských právních úprav dědickým právem, toto je do značné míry pravdivé a je faktem, že je třeba nalézt jistou harmonii mezi common law institutem trustu a kontinentálním dědickým právem. Dále poukazuje na to, že quebecký Občanský zákoník neobsahuje žádnou ochranu nepominutelného dědice, tím pádem ani neobsahuje ustanovení řešící případy touto kolizí vznikající, tato úprava tedy nemohla být součástí transplantace, nebo sloužit jako inspirační zdroj, či zdroj pro komparaci a obdobně ani jiné common law systémy tuto úpravu neobsahují a nemohou tak sloužit jako inspirační zdroj pro českého zákonodárce\footnote{Tamtéž.}. Ačkoli s první částí o svázanosti evropských právních systémů dědickým právem se dá autorovi přisvědčit, druhá část o neexistenci ustanovení sloužících jako pojistky řešení kolize trustu a jiných nároků připodobitelných k nárokům nepominutelných dědiců se již dle mého pohledu nezakládá zcela na pravdě. V této souvislosti bych chtěl upozornit na kapitolu 7.4.5, kde popisuji ochranu nároků, které lze s trochou představivosti připodobnit k institutu nepominutelného dědice v Občanském zákoníku provincie Quebec.\\

Zcela souhlasím s dalším tvrzením, že český zákonodárce si této potenciální kolize musel být vědom a přesto neučinil žádná opatření ve formě ustanovení, která by tyto kolize řešila. Ta tedy bohužel v celé části Občanského zákoníku týkající se správy cizího majetku z, pro mě nepochopitelných důvodů, absentují a zapříčiňují značnou právní nejistotu při zakládání testamentárního svěřenského fondu.\\

Není nicméně vhodné hned zpočátku tvrdit, že ze současné právní úpravy nelze v žádném případě vyčíst řešení této situace, jen že nemusí být, a také není, na první pohled zřejmé. Na následujících řádcích je tedy vhodné podrobit tuto premisu hlubší analýze za pomocí dvou mnou výše vytyčených optik.\\

   \begin{enumerate}
 {\Large\item[1.] Ochrana nepominutelných dědiců s ohledem na výklad jednotlivých ustanovení}
 \end{enumerate}
 
 Než přejdu k samotnému pohledu na věc za pomoci této optiky, je vhodné nejprve představit jednotlivé názorové proudy.\\
 
 Vhodný vhled do problematiky svěřenského fondu a nepominutelného dědice poskytují čtyři autoři, z jejichž názorů také budu čerpat a postavím na nich svojí analýzu. Těmito autory jsou Miloš Kocí, Vlastimil Pihera, Kryštof Horn a Ivan Kolář.\\
 
 Z počátku pojďme podrobit analýze pohledy Miloše Kocího z komentáře k Občanskému zákoníku od Wolters Kluwer a Vlastimila Pihery z komentáře k Občanskému zákoníku od vydavatele C.H.Beck.\\
 
 Kocí se profiluje jakožto podporovatel názorového proudu, který nepovažuje majetek vyčleněný do svěřenského fondu za součást pozůstalosti a tím pádem zavrhuje možnost jeho ponížení o povinný díl nepominutelného dědice, přičemž svojí argumentaci staví na dvou základních bodech.\\
 
 Prvním bodem pak jest konstatování, že majetek vložený do svěřenského fondu nemůže být součástí pozůstalosti, neboť přestává být k okamžiku smrti zůstavitele v jeho vlastnictví: \textit{"16. Lze pochybovat o tom, že svěřenský správce může být považován za dědice nebo za osobu, která by v dědickém řízení vykonávala jeho práva, a to již proto, že takový majetek žádného dědice nemá. Jestliže tedy svěřenský fond zřizovaný pořízením pro případ smrti vzniká k okamžiku smrti zůstavitele bez ohledu na to, zda svěřenský správce přijal pověření k jeho správě či nikoliv, pak není pojmově možné, aby byl tento oddělený a nezávislý majetek předmětem řízení, jehož účelem je potvrzení vlastníka. V opačném případě by de facto potvrzoval vznik svěřenského fondu soud (§ 1670) a správu svěřenského fondu by do té doby vykonával správce pozůstalosti nebo vykonavatel závěti, pokud by byli zůstavitelem ke správě pozůstalosti povoláni (§ 1677). Ani s jednou z těchto variant však právní úprava svěřenského fondu nepočítá. Na svěřenský fond by tak neměly přecházet ani dluhy zůstavitele, což je v souladu s odst. 1 komentovaného ustanovení, podle něhož lze vyčlenit pouze „majetek“, tedy aktiva ve vlastnictví zakladatele (§ 495). V tomto ohledu není důvodu, proč by se na pořízení pro případ smrti mělo nahlížet jinak než na zřízení svěřenského fondu inter vivos, kde je zakladateli rovněž umožněno, aby vyčlenil pouze svoje aktiva bez ohledu na dluhy. Takové pořízení pro případ smrti, které by věřitele zůstavitele poškozovalo na jejich právech, by však mohlo být napadeno jako neplatné pro rozpor s dobrými mravy (§ 580 odst. 1), resp. by mohlo být neplatné částečně (§ 576), přičemž za určitých okolností by k neplatnosti takového právního jednání soud přihlédl i bez návrhu (§ 588)."}\footfullcite[Komentář k § 1448]{svestka_j_obcansky_2014-1}\\
 
 Tento argument s ohledem na ustanovení paragrafu 1475 Občanského zákoníku, pohlížející na pozůstalost jako na celé jmění zůstavitele ke dni jeho smrti, tak Kocího pohledu přisvědčuje\footnote{Viz kapitola 5.4 Pozůstalost.}. Dále Kocí svůj pohled vykládá v souvislost s důrazem Občanského zákoníku na respektování vůle člověka, respektive zůstavitele, která byla posílena i mnoha jinými nově zavedenými nástroji dědického práva v NOZ\footnote{Například odkazy či příkazy.}. Zároveň však zcela opomíjí existenci § 1452 OZ stanovujícího v bodě 1 následující: \textit{"Každý svěřenský fond musí mít statut. Statut svěřenského fondu vydává zakladatel. Zřizuje-li se svěřenský fond pořízením pro případ smrti, použije se přiměřeně § 311."}, při pohledu na paragraf § 311 OZ pak zákonodárce stanovuje následující: \textit{"Při založení nadace pořízením pro případ smrti se do nadace vnáší vklad povoláním nadace za dědice nebo nařízením odkazu. V takovém případě nabývá založení nadace účinnosti smrtí zůstavitele."}. Toto ustanovení by tedy analogicky mělo být použito při vyčleňování majetku do svěřenského fondu. Zde se však v rámci pohledu na svěřenský fond první optikou dostáváme do částečného rozporu s jednou z základních premis, na které je vystavěn svěřenský fond v České republice, touto premisou je samotná forma svěřenského fondu, v rámci které rozhodně nelze na svěřenský fond pohlížet jako na osobu mající právní subjektivitu, jak by tedy mohl svěřenský fond být dědicem, účastnit se dědického řízení a tím pádem majetek do něj vyčleněný spadat do pozůstalosti? \\
 
 Odpovědí by mohla být osoba svěřenského správce, který by svěřenský fond v pozůstalostním řízení zastupoval a ustanovení § 311 by tak ve své podstatě mohlo být pouze jakýmsi pojmovým obratem. Problémem však zůstává, že ani svěřenskému správci vlastnické právo k majetku vyčleněnému do svěřenského fondu nesvědčí\footfullcite[POPOVICI A. Trust québeckého a českého práva:autonomní vlastnictví?]{lubos_tichy_sverensky_2017}. Dále s ohledem na úpravu řízení o pozůstalosti v ZŘS, oddílu 4, účastníci řízení je možné dovodit, že samotný svěřenský správce by neměl být ani účastníkem řízení, tím by ve své podstatě měl být svěřenský fond, jehož by měl svěřenský správce zastupovat. Kvůli nedostatku právní osobnosti svěřenského fondu však zase dospíváme k závěru, že samotný svěřenský fond nemůže být dědicem a tedy se ani řízení účastnit.\\
 
 S ohledem na výše zmíněné je třeba rozlišit, zda je svěřenský správce sám dědicem \textit{sui generis} kvůli nedostatku právní osobnosti svěřenského fondu, nebo zda je svěřenský fond pro tuto příležitost nadán právní subjektivitou, respektive tedy v uvozovkách považován za právnickou osobou, tak jako tomu zákonodárce činí  v daňových předpisech\footfullcite[§ 4b (2)]{noauthor_zakon_zdph} a tedy i možností účastnit se tohoto řízení a svěřenský správce ho pouze zastupuje (tak jak by zastupoval statutární orgán jinou právnickou osobu v případě, že by byla určena dědicem).\\
 
 OZ ani ZŘS toto nikde explicitně neupravují, naopak sám OZ v § 1448 hovoří odděleném a nezávislém vlastnictví, tedy o entitě bez právní subjektivity, znovu tedy vzniká problém spočívající v neexistenci právní subjektivity svěřenského fondu a tedy nemožnosti být účastníkem jakéhokoliv řízení.\\
 
 Nejenom k tomuto problému se vyjádřil Kryštof Horn: \textit{"Dle názoru autora je potřeba vycházet z toho, že svěřenský fond mortis causa se zřizuje pořízením pro případ smrti. Půjde tedy o závěť či dědickou smlouvu (popř. dovětek), které v sobě budou zřejmě muset obsahovat náležitosti statutu a současně budou muset splňovat formální požadavky na ně kladené. Pokud tedy majetek součástí pozůstalosti je, lze mít za to, že zřizovací instrument (tj. pořízení pro případ smrti) musí respektovat ustanovení o povinném dílu, a to zřejmě i v případě, kdy by zůstavitel zamýšlel zřídit svěřenský fond ohledně veškerého svého majetku ve prospěch nepominutelného dědice. Předně proto, že povinný díl může být zůstaven jen v podobě dědického podílu nebo odkazu a především musí být zcela nezatížen, ale také proto, že samotným vytvořením fondu se nepominutelnému dědici ničeho nezůstavuje, a tak se dědic dle názoru autora nemůže při použití § 1644 odst. 3 OZ ani rozhodnout pro přijetí takto omezeného „dědictví“ namísto povinného dílu. Proto autor nesouhlasí s názorem, podle kterého by bylo nutné započtení plnění poskytnutých obmyšlenému ze svěřenského fondu na jeho souběžný dědický podíl podle § 1662– 1664 OZ, neboť postavení dědice sui generis zde má svěřenský správce a nikoliv obmyšlený."}\footfullcite{sro_httpwwwpraguebestcz_prakticke_nodate-1}.\\
 
 Dle Horna má tedy fakticky dědit svěřenský správce ve prospěch svěřenského fondu, což je do značné míry vhodnější výklad, než opět pohlédnout na svěřenský fond jako na právnickou osobu, neboť konstantní subjektivizace svěřenského fondu nesvědčí jeho účelu jakožto common law institutu. Z pohledu civil law však tato subjektivizace může být s ohledem na konformitu s jeho principy možným řešením.\\
 
 Je zajímavé zmínit, že v komentáři k Občanskému zákoníku od C.H.Becka autoři zmiňují následující: \textit{"Uvedený případ zřízení svěřenského nástupnictví nelze zaměňovat s tím, kdy zůstavitel pořízením pro případ smrti zřizuje nadaci (§ 309 odst. 1), nebo nadační fond (§ 395), nebo ústav (§ 405) anebo svěřenský fond (§ 1452 odst. 1), neboť tyto subjekty jsou (za splnění zákonem stanovených podmínek) způsobilými přímými dědici."}\footfullcite[Komentář k § 1513]{fiala_r_drapal_l_a_kol_obcansky_2015}. Autoři tedy hovoří o svěřenském fondu jako o přímém dědici, jsou tedy ve shodě s Piherou, který totéž zmiňuje ve svazku třetím.\\
 
 S Hornem souhlasí i Pihera, který se řadí do opačného názorová spektra něž Kocí, kdy tvrdí následující: \textit{"Poněkud odlišnou povahu má vyčlenění majetku do svěřenského fondu zřizovaného pořízením pro případ smrti (mortis causa). V těchto případech vyčlení zakladatel majetek ve prospěch svěřenského fondu tak, že v jeho prospěch vyhradí dědické právo nebo nařídí odkaz (§ 1452 odst. 1 ve spojení s § 311 odst. 1). Správce svěřenského fondu mortis causa pak vykonává práva dědice majetku vyčleněného do svěřenského fondu nebo práva odkazovníka. Majetek vyčleněný zakladatelem (zůstavitelem) do svěřenského fondu mortis causa je součástí pozůstalosti a nabytí tohoto majetku do svěřenského fondu buď potvrzuje soud, nebo je převeden ve prospěch svěřenského fondu dědicem obtíženým odkazem. Z hlediska dědického práva je na postavení svěřenského fondu nutné pohlížet „jako na dědice“ nebo „jako na odkazovníka“, a to se vším všudy, co se s tímto postavením spojuje (např. odpovědnost za dluhy zůstavitele)."}\footfullcite[Komentář k § 1448]{spacil_j_a_kolektiv_obcansky_2013}\\
 
 Samotná odpovědnost za dluhy je však taky problematická, neboť zákon hovoří o vyčleňování majetku do svěřenského fondu, jedná se tedy zásadně o aktiva. Jak jsem popsal v kapitole 5.4, české dědické právo stojí na principu univerzální sukcese s možností výhrady soupisu. V případě, že by byl svěřenský fond dědicem, tento princip by se musel logicky aplikovat. Do svěřenského fondu by se tak vyčleňovalo jmění, nikoli majetek. V rámci výkladu svěřenského fondu je zřejmé, že do svěřenského fondu jdou vyčlenit dluhy, půjde o případy podle § 1107 OZ, nejčastěji v případech, kdy bude vyčleňován obchodní závod, jehož složkou jsou i dluhy, neboť se jedná o soubor jmění\footfullcite[Komentář k § 1488]{spacil_j_a_kolektiv_obcansky_2013}. Mimo tyto případy však lze vyčlenit zásadně majetek, což principu univerzální sukcese odporuje. Tento fakt tak rovněž přisvědčuje Kocího pohledu, který nepovažuje vyčleněný majetek za součást pozůstalosti.\\
 
 Při užití výkladového pohledu a favorizaci logiky je však nutné podrobit tyto názory na problematiku hlubší analýze. Při pohledu touto optikou na výše uvedené tři argumenty se lze ztotožnit s vysvětlením Michaela Skuhrovce, který vysvětluje, proč se zdá Kocího pohled ve vztahu k základní konstrukci svěřenského fondu lépe aplikovatelnější: \textit{"Vyjdeme-li však zásadně z předpokladu, že ve smyslu § 1448 OZ se svěřenský fond vytváří vyčleněním majetku, v případě testamentárního svěřenského fondu pak tedy pořízením pro případ smrti, a vzniká smrtí zůstavitele, opora pro Hornovu argumentaci se nachází pouze s těžkostmi. Především je totiž nutné důsledně aplikovat esenciální obsah konstrukce Lepaullova vlastnictví. Jak Horn sám správně uvádí, ve smyslu §1475 odst. 2 OZ pozůstalost tvoří celé jmění zůstavitele. Součástí tohoto jmění je pak bezpochyby i majetek, který je vyčleněn k určitému účelu a tvoří tak svěřenský fond (v případě testamentu pak jakýsi „spící” svěřenský fond, neb vznikne až smrtí zůstavitele, která představuje dies certus an, incertus quando, a vzhledem k možnostem dispozice s majetkem je nejasný i jaký majetek bude součástí svěřenského fondu). Pokud je zároveň postaveno najisto, že svěřenský fond vznikne smrtí zůstavitele (na rozdíl od CCQ, které sice neobsahuje institut nepominutelného dědice, ale stále váže vznik svěřenského fondu na přijetí správce, byť s retroaktivním účinkem), nelze považovat tento majetek dále za vlastnictví zůstavitele, nýbrž za entitu bez právní subjektivity a se samostatným účelovým určením. Ve spojení s § 1479, který ohledně dědického nápadu stanoví „Dědické právo vzniká smrtí zůstavitele.”, pak dostáváme určitou konkurenci norem OZ sui generis. To ovšem pokud na ni tak budeme chtít nahlížet. Bezpochyby nezle tvrdit, že v tomto jinak právně uceleném okamžiku (smrti) by šlo předřadit vznik jednoho práva před druhým. Rozumnější a konformnější by bylo nahlížet na vznik obou konformně a ponechat je koexistovat, kde se v ten samý okamžik majetek vyčleněný do svěřenského fondu stane majetkem samostatným a „zbytek” (je-li takového) pak zůstává zůstaviteli, ergo jen ten se může stát předmětem dědického řízení. Výkladem zákona také nelze rozumně dojít k závěru, že by mohl být obmyšlený či správce účastníkem dědického řízení, neb ani jednomu vlastnické právo nesvědčí, ke svěřenskému fondu (trustu dle CCQ) totiž nemá majetková práva nikdo."}\footfullcite[Strana 56]{skuhrovec_michal_sverensky_2018}\\
 
 Skuhrovec tedy také favorizuje Kocího pohled na problematiku, tato citace také nepřímo reaguje i na komentář Pihery, jež pro jím popsanou funkci svěřenského fondu v rámci pozůstalostního řízení předpokládá, že svěřenský fond bude dědicem, jehož práva bude vykonávat svěřenský správce, to ale není s ohledem na konstrukci svěřenského fondu bez dalšího zřejmé, jak jsem již ostatně vyložil výše a jak poznamenal i Skuhrovec.\\
 
Z výše uvedené citace tedy vychází, že při porovnání těchto zcela opačných názorových proudů na věc nepominutelného dědice a svěřenského fondu za pomocí výkladové optiky, do níž smyslově zcela spadá názor Skuhrovce, se lze ztotožnit spíše s názorem Kocího a nejpravděpodobnějším výsledkem tedy je, že majetek vložený do svěřenského fondu (nejspíše) nebude součástí pozůstalosti, stejně tak jako svěřenský fond nebude i přes velmi nešťastné ustanovení § 311 dědicem. Mohlo by se tedy zdát, že svěřenský fond \textit{mortis causa} může být úspěšně použitý k faktickému vydědění nepominutelného dědice a má hypotéza by tedy byla zodpovězena. I v případě neúspěchu ochrany dědice ve vztahu k pozůstalostnímu řízení je však stále třeba zaměřit se i na samotný instrument použitý ke zřízení svěřenského fondu a omezení na něj kladená.\\

% --------------------------------------------------------------------------
%lze vyvodit, že z pozice udržení konstrukce svěřenského fondu a logiky nasvědčující tomu, že správným výkladem problematiky jest Kocího pohled

%je tedy ?zřejmé?, že majetek vložený do svěřenského fondu nejspíše nebude součástí pozůstalosti, stejně tak jako svěřenský fond nebude i přes velmi nešťastně ustanovení § 311 dědic. Mohlo by se tedy zdát, že svěřenský fond \textit{mortis causa} může být úspěšně použit k faktickému vydědění nepominutelného dědice a má hypotéze je tedy zodpovězena. I v případě neúspěchu ochrany dědice ve vztahu k pozůstalostnímu řízení je však stále třeba zaměřit se i na samotný instrument použití k založení svěřenského fondu a omezení na něj kladená.
% --------------------------------------------------------------------------
 
 Skuhrovec je ve svém výkladu velmi konzistentní a svojí argumentaci zakládá na pevných základech, já s ním po všech směrech souhlasím. Opomněl však zohlednit následující část Hornova argumentu: \textit{"lze mít za to, že zřizovací instrument (tj. pořízení pro případ smrti) musí respektovat ustanovení o povinném dílu"}, Horn k tomu v jedné větě ovšem současně dodává, že je tak třeba učinit v případě, že vyčleněný majetek je součástí pozůstalosti.\\ 
 
 V tomto ohledu neobstojí ani argument: \textit{"Bezpochyby nezle tvrdit, že v tomto jinak právně uceleném okamžiku (smrti) by šlo předřadit vznik jednoho práva před druhým."}. Neboť úkon pořízení pro případ smrti lze jednoznačně předřadit jakýmkoliv jiným právním událostem nebo úkonům závislých právě na daném pořízení.\\
 
% -------------------------------------------------------------------------- 
 %Obstál by jen argument, který by tvrdil, že takováto pohledávka je nevymahatelná, neboť nikdy nemůže vzniknout pohledávka vůči někomu kdo fakticky nedědil, pokud by tedy svěřenský fond nebyl dědicem, nejspíše ani zde by se nedala hledat možnost ochrany nepominutelného dědice. Tomu by tak zbývalo pouze žalování na neplatnost pořízení pro případ smrti, které by popřel co do pravosti, nebo co do platnosti.
 
 %Tady to asi může popřít co do jiných věcí, musím zřejmě opravit.
% -------------------------------------------------------------------------- 
 
 Dle výše uvedeného tedy i v případě, kdy bychom uplatnili jednotlivá ustanovení v zákonem postupu, jak navrhují autoři v díle Svěřenské fondy, Příležitosti a rizika\footnote{Viz následující citace.}, v časové linii je samotné zůstavitelovo pořízení na prvním místě, bez ohledu na následný průběh pozůstalostí řízení by tak mělo splňovat nároky na něj kladené bez jakýchkoliv výjimek, pokud je ovšem zákonodárce v jiné části nestanoví, v tomto případě tak neučinil. Lze tedy dovodit, že podmínky omezují zůstavitelovu dispozici v § 1492 OZ o nemožnost zkrácení povinného dílu nepominutelných dědiců musí být dodržena a aplikována. K tomuto poskytuje dobré stanovisko Kolář: \textit{"Při kolizi svěřenského fondu zřizovaného pro případ smrti s nároky nepominutelných dědiců je třeba příslušená ustanovení, to je § 1451 odst. 3 ObčZ (svěřenský fond vzniká smrtí zůstavitele) a § 1479 ObčZ (dědické právo vzniká smrtí zůstavitele) vykládat tak, jak je upraveno v zákoně v logickém postupu. Postavení nepominutelných dědiců je upraveno v § 1642 a násl. ObčZ s tím, že nepominutelnému dědici náleží z pozůstalosti povinný díl. Nepominutelný dědic proto bude mít nárok proti správci svěřenského fondu a bude uspokojen ze svěřenského fondu právě na tento povinný díl, nikoli však na věci vyčleněné do svěřenského fondu zřizovaného pro případ smrti, který vznikne smrtí zůstavitele, čímž se naplní § 1448 odst. 2, resp. 3 ObčZ, podle kterého majetek ve svěřenském fondu není ani vlastnictvím správce, ani vlastnictvím osoby, ani vlastnictvím obmyšleného."}\footfullcite[Strana 82]{svejkovsky_jaroslav_sverenske_2018}.\\
 
 Ochrana nepominutelných dědiců by tak měla být zajištěna před samotným pozůstalostním řízením faktickým omezením možných dispozic s majetkem pořízením pro případ smrti, respektive omezením možnosti zůstavitele pořídit o veškerém svém majetku bez zohledněním nepominutelného dědice, respektive jeho omezením nebo opomenutím. Samotné opomenutí nepominutelného dědice však nezakládá neplatnost pořízení pro případ smrti tak jako tomu bylo za účinnosti OZ z roku 1964\footfullcite{sro_httpwwwpraguebestcz_postaveni_nodate}, i v případě opomenutí nepominutelného dědice není založen důvod, pro který by pořízení pro případ smrti nemělo být platné a z toho titulu by neměl vzniknout svěřenský fond. Je zřejmé, že v případě neplatnosti pořízení pro případ smrti by svěřenský fond nevznikl, jak ostatně i sám Kocí uznává \footfullcite[Komentář k § 1448]{svestka_j_obcansky_2014-1}. Toto dovození ale není v obdobném případě uplatnitelné.\\
 
 V NOZ vzniká dědici při jeho opomenutí v pořízení pro případ smrti ze strany zůstavitele pouze pohledávka za ostatními dědici\footfullcite[Komentář k § 311]{lavicky_p_a_kol_obcansky_2014}, ten se tak stává jejich věřitelem\footfullcite[§ 1696]{Duvodovazprava}. V případě toho, že by tedy svěřenský fond a ani svěřenský správce nebyli účastni pozůstalostního řízení jako dědicové a nebyli by za ně tedy ani považováni, je k diskuzi, zda by takováto pohledávka, jdoucí oproti ostatním dědicům, obstála. Kovář tvrdí ve svém příspěvku, že ano, opět ale narážíme na to, že explicitní úprava této kolize chybí, stejně tak jako soudní rozhodnutí tento problém řešící, tento názor je tedy jednoduše zpochybnitelný.\\
 
 Obstát by tedy mohl argument, který by tvrdil, že takováto pohledávka nikdy nemůže vzniknout vůči někomu, kdo fakticky nedědil, pokud by tedy svěřenský fond nebyl dědicem, jejímž dle § 1475 OZ dosti možná, i přes existenci § 311 OZ není, neboť dědicem může být pouze člověk, nebo právnická osoba, svěřenský fond však není ani jedním a zákonodárce nikde nenaznačuje, že by se na něj pro účely pozůstalostního řízení mělo jako na právnickou osobu hledět, svěřenský správce pak vykonává jen vlastnická práva a jak již jsem popsal výše, dědické právu mu rovněž nesvědčí\footfullcite[O tomto zmínku neobsahuje]{svestka_j_dvorak_j_fiala_j_a_kol_obcansky_2014}, nejspíše ani zde by se nedala hledat možnost ochrany nepominutelného dědice. Tomu by tak zbývalo pouze žalování na neplatnost pořízení pro případ smrti, které by popřel co do pravosti, nebo co do platnosti.\\
  
  Je samozřejmě možné uvažovat nad tím, že s ohledem na veřejný pořádek, jehož zjevnost lze doložit patrně s ohledem na zájem jasného určení subjektu, za nímž by tato pohledávka měla jít, by soud v případném soudním sporu mohl tuto pohledávku za svěřenským fondem postavit na roveň jiným pohledávkám za dědici, o nichž nelze mít pochybnosti, že dědici jsou, a z toho důvodu by opominutý dědic mohl mít právo dožadovat se svého povinného dílu po svěřenském fondu. Je to ale také obdobně diskutabilní.\\
  
  Přes toto všechno se mnoho autorů domnívá, že opomenutým dědicům vznikne z titulu opomenutí dle paragrafu 1492 OZ pohledávka za svěřenským fondem. Tento paragraf tedy poskytuje jistou možnost bránit se opomenutí.\\
  
  Další možnost bránit se zkrácení svého dílu, nebo opomenutí ze strany zůstavitele má nepominutelný dědic také vyvoláním sporu o dědické právo podle § 168 - 170 ZŘS jak již jsem letmo zmínil výše. Napadnout lze platnost a pravost pořízení pro případ smrti. Námitka nepravosti spočívá v tom, že dědic bude tvrdit, že pořízení pro případ smrti nepochází od zůstavitele, vůle v něm učiněná tak není vůlí zůstavitele, v takovém případě ho soud odkáže k podání žaloby dle § 170 ZŘS. V případě úspěšného prokázání, že pořízení pro případ smrti nepochází od zůstavitele, se k takovémuto pořízení pro případ smrti nebude přihlížet a v případě, že tímto neobživne jiné pořízení pro případ smrti, nastane dědění podle zákonné posloupnosti a svěřenský fond tak vůbec nevznikne. Dále může opominutý dědic namítat platnost pořízení pro případ smrti, při úspěchu nastává obdobných účinků jako u popření co do pravosti\footfullcite{sro_httpwwwpraguebestcz_postaveni_nodate}. Prokázat něco takového však může být nesmírně náročné, v případě nemožnosti bránit se ani podle § 1492 se ale patrně jedná o jedinou možnost, pomocí které by se dědic mohl domoci svého práva.\\
  
  Neplatnost je možné namítat i proti svěřenskému fondu založeného \textit{inter vivos} podle § 574 - 588 OZ, neplatnost právních jednání, například v případě uvedeném v §587 OZ.\\

% --------------------------------------------------------------------------  
  %ještě podáním žaloby dle § 170 ZŘS. Napadnout lze platnost a pravost pořízení pro případ smrti. Námitka nepravosti spočívá v tom, že dědic bude tvrdit, že pořízení pro případ smrti nepochází od zůstavitele, vůle v něm učiněná tak není vůlí zůstavitele.
 
 
 %Napsat o tom, že při takovémto opomenutí může být majetek pozůstalosti (pokud by nezahrnoval majetek vložený do svěřenského fondu absolutně nedostatečný k samotnému otevření pozůstalostního řízení, opravné prostředky použitelné v jeho průběhu by tak nebylo možné uplatnit, krom ustanoveních, o kterých zákonodárce říká, že jdou uplatnit v každém případě, § 153 ZŘS.
 
 
  %tak nespočívá na základě pozůstalostního řízení (není z titulu pozůstalostního řízení tzn. dědic proti dědici), ale z titulu samotné ochrany nepominutelnému dědici poskytované už v rámci možných dispozic s majetem v rámci pořízení pro případ smrti ze strany zůstavitele.
  
  %Možnost tak při tomto pohledu, kdy neuvažujeme, že samotná pohledávka za svěřenským fondem by obstála představuje pouze pouze vyvolaný spor o dědické právo, kdy by dědic popřel pořízení pro případ smrti.
 
 %Dále napsat co myslí Kovář tím, že pohledávka půjde za svěřenským fondem, ne za věc vyčleňovanou v rámci pozůstalostního řízení (nebo i mimo ně), toto souvisí s tím, že se již nevyhlašuje celé pořízení pro případ smrti neplatné, ale pouze opomenutému dědici vzniká pohledávka za ostatními dědici. Pořízení pro případ smrti je § 1492.
 
 %Rozdíl mezi opomenutím a omezením, opomenutí neuvedu dědice v závěti a 9 milionů z 10 odkážu svěřenskému fondu, omezení - omezím povinný díl nepominutelného dědice tak, že mu například odkážu pouze jeden milion a zbytek odkážu svěřenskému fondu (§ 1644)
 
 %Pokud by dědic napadl platnost závěti mohl by soud závěť i zrušit v tom případě by jak říká Kocí svěřenský fond nevznikl, dále k tomuto napříkl https://is.muni.cz/th/c82g3/Veronika_Kubatova_TEXT_RIGOROZNI_PRACE.pdf.
 
 %V případě výkladu těchto ustanovení tak, že svěřenský fond nemůže být dědicem a obdobně svěřenský správce nemůže být dědicem \textit{sui generis} dědící ve prospěch svěřenského fondu, omezila by se účast svěřenského správce podle § 164 ZŘS pouze na informační a informativní.\\
 
 %V jiných řízeních ve věcech svěřenského fondu je pak svěřenský správce plnohodnotným účastníkem, ve věcech pozůstalostního řízení je však jeho účast diskutabilní. Je pak pouze logické, že v případě, že by svěřenský fond nebyl účastníkem řízení o pozůstalosti, pak by takový vyčleněný majetek nebyl ani součástí pozůstalosti, nabytí jehož části ve formě dědictví by následně nepotvrzoval soud a z toho důvodu by nepominutelní dědici neměli nárok na svůj povinný díl ani v případě, že byl do svěřenského fondu vyčleněn veškerý zůstavitelův majetek.\\
 
 %na kterou rozhodně nělze pohlížet jao na subjekt
 
 %Pokud tedy navážu na Kocím vyslovený první bod, druhým budem jeho argumentace je pak jakási nutná rovnost mezi oboumi formami založení svěřenského fondu: \textit{"17. Pokud by měl být majetek vyčleněný do svěřenského fondu součástí pozů­stalosti, mohl by být teoreticky snížen o povinný díl nepominutelného dědice (§ 1642). To by znamenalo, že by zůstavitel pořízením pro případ smrti při existenci nepominutelného dědice nemohl do svěřenského fondu vyčlenit všechen svůj majetek, což by testamentární svěřenský fond značně znevýhodňovalo ve srovnání s ostatními typy svěřenských fondů (viz komentář k § 1449), a to bez rozumného důvodu a jakékoliv opory v ustanoveních o jeho vzniku, která na obě formy zakladatelského aktu, tj. na smlouvu a pořízení pro případ smrti, nahlížejí jako na rovnocenné."}\footnote{Tamtéž.}\\
 
 %Přes fakt, že při pohledu na úpravu svěřenského fondu a její konkurenci s právní úpravou dědického práva touto optikou docházím k obdobnému závěru jako Michal Skuhrovec ve své práci, nemohu se v žádném případě ztotožnit s tímto bodem. V případě ochrany nepominutelného dědice i ve chvíli vyčlenění majetku do svěřenského fondu \textit{mortis causa} ať už dovozené ze současné litery zákona, nebo potenciálně v budoucnu implementované formou novely si nemyslím, že dochází k zásadnímu znevýhodnění formy založení \textit{inter vivos} oproti \textit{mortis causa}. Neboť obě dvě formy, přestože se jedná o stejný isntitut, jsou založeny jiným způsobem a mohou být založeny často za značně jiným účelem (například ochrana majetku před věřiteli nebo za účelem businessu vs mezigenerační převod majetku). Obě tyto formy lze paralelně připodobnit k možnostem převedení majetku za života zůstavitele, jako je například darování, směna, prodej (neuvažuji daňové zatížení, pouze ve vztahu k institutu nepominutelného dědice), nebo obdobným, podmínkamy dědického práva nezatíženým, formám převodu, jakými jsou například darování pro případ smrti\footnote{Viz kapitola 6.3.3 Vznik svěřenského fondu \textit{inter vivos} s odkládací podmínkou} a převedení majetku za pomoci dědického práva, tedy v rámci pozůstalostního řízení, po smrti zůstavitele (kde pokud exituje nepominutelný dědic a není platně vyděděn, má nárok na svůj nepominutelný díl).\\
 
 %Přestože se tedy jedná o stejný institut, lze předpokládat, že účelem založení svěřenského fondu pořízením pro případ smrti bude s největší pravděpodobností mezigenerační převod majetku, zatímco u založení svěřenského fondu smlouvou mezi živými může být množina účelů širší. Toto samozřejmě nevylučuje možnost založení svěřenského fondu pořízením pro případ smrti za stejnými účely jako u \textit{inter vivos}, tyto jsou však mnohem výhodnější za zůstavitelova života, už jen zkrátka proto, že je zakladatel naživu. Stejně tak i pokud bychom se bavili o jedné výhodě, o kterou by tak testamentární svěřenský fond při ochraně nepominutelných dědiců přišel, tedy o faktickou možnost jeho pomocí vydědit, stále má tato forma mnoho výhod oproti svěřenskému fondu založeného \textit{inter vivos}, o tomto ostatně hovořím v kapitole 6.3 u jednotlivých forem založení svěřenského fondu.\\
 
 %Nárážím zde na to, že účely, které může zakladatel sledovat založením svěřenského fondu \textit{inter vivos}
 
 %Pihera se pak řadí do opačného názorová spektra, kdy tvrdí následující: \textit{"Poněkud odlišnou povahu má vyčlenění majetku do svěřenského fondu zřizovaného pořízením pro případ smrti (mortis causa). V těchto případech vyčlení zakladatel majetek ve prospěch svěřenského fondu tak, že v jeho prospěch vyhradí dědické právo nebo nařídí odkaz (§ 1452 odst. 1 ve spojení s § 311 odst. 1). Správce svěřenského fondu mortis causa pak vykonává práva dědice majetku vyčleněného do svěřenského fondu nebo práva odkazovníka. Majetek vyčleněný zakladatelem (zůstavitelem) do svěřenského fondu mortis causa je součástí pozůstalosti a nabytí tohoto majetku do svěřenského fondu buď potvrzuje soud, nebo je převeden ve prospěch svěřenského fondu dědicem obtíženým odkazem. Z hlediska dědického práva je na postavení svěřenského fondu nutné pohlížet „jako na dědice“ nebo „jako na odkazovníka“, a to se vším všudy, co se s tímto postavením spojuje (např. odpovědnost za dluhy zůstavitele)."}\footfullcite[Komentář k § 1448]{spacil_j_a_kolektiv_obcansky_2013}\\
 
 %Základním předpokladem, který musí být splněn je tedy nutnost považovat svěřenský fond za účastníka pozůstalostího řízení, přičemž za dědice \textit{sui generis} bude považován svěřenský správce, k tomu ostatně § 311 OZ přímo vybízí.
 
 %Je třeba rozlišit, zda je svěřenský správce sám dědicem sui generis kvůli nedostatku právní osobnosti svěřenského fondu, nebo zda je svěřenský fond pro tuto příložitost nadán právní subjektivitou, respektive tedy v uvozovkách považován za právnickou osobou, tak jako tomu činí zákonodárce v daňovýc předpisech (Uvést ZDROJ) a tedy i možností účastnit se tohoto řízení a svěřenský správce ho pouze zastupuje (tak jak by zastupoval statutární orgán jinou právnickou osobu v případě, že by byla určena dědicem.
 
 %\textit{"Svěřenský fond může být zřízen též pořízením pro případ smrti, tedy závětí nebo dědickou smlouvou; jeho podmínky mohou být moderovány též dovětkem (§ 1492). Závěť nebo dědická smlouva v takovém případě musí obsahovat náležitosti statutu svěřenského fondu (§ 1452 odst. 1, 2) a musí splňovat formální požadavky kladené na ně dědickým právem (§ 1533–1593). To platí též o vedlejších doložkách v závěti. Správce takto zřízeného svěřenského fondu pak vykonává práva dědice nebo odkazovníka, aniž by však sám byl osobním dědicem či odkazovníkem (dále srov. komentář k § 1451 a § 1452 odst. 1)."}\footnote{Tamtéž.}\\
% -------------------------------------------------------------------------- 
 
    \begin{enumerate}
 {\Large\item[2.] Ochrana nepominutelných dědiců s ohledem na úmysl specifických ustanovení}
 \end{enumerate}

% -------------------------------------------------------------------------- 
 %I přes 
% -------------------------------------------------------------------------- 
 
 S ohledem na výše uvedenou analýzu problematiky optikou prvního pohledu, z níž vychází s ohledem na výklad těchto ustanovení to, že svěřenský fond není dědicem a obdobně svěřenský správce nemůže být dědicem \textit{sui generis} dědící ve prospěch svěřenského fondu, omezila by se účast svěřenského správce podle § 164 ZŘS pouze na informační a informativní.\\
 
 V jiných řízeních ve věcech svěřenského fondu je pak svěřenský správce plnohodnotným účastníkem, ve věcech pozůstalostního řízení je však jeho účast diskutabilní. Je pak pouze logické, že v případě, že by svěřenský fond nebyl účastníkem řízení o pozůstalosti, pak by takový vyčleněný majetek nebyl ani součástí pozůstalosti, nabytí jehož části ve formě dědictví by následně nepotvrzoval soud a z toho důvodu by nepominutelní dědici neměli nárok na svůj povinný díl ani v případě, že byl do svěřenského fondu vyčleněn veškerý zůstavitelův majetek.\\
 
 I přes toto se však nelze ubránit pocitu, že daná ustanovení, která by dle výkladu výše neměla být jaksi uplatnitelná, tedy zejména § 1452 a s ním souvisejícího § 311 a § 1492 OZ, by měla mít nějaký smysl, ostatně proč by zákonodárce vůbec odkazoval na § 311, pokud by nechtěl, aby na svěřenský fond bylo pohlíženo jako na dědice.\\
 
 Kvůli této intenci zákonodárce nelze tak zcela zavrhovat ani názory Pihery nebo Horna, i ty se také zakládají na solidních základech a nelze stoprocentně vyloučit, že právě těmito autory vyslovený pohled je ten správný. Samotné nejasnosti ohledně testamentárního svěřenského fondu však pravděpodobně trvale vyřeší až konstantní praxe s tímto institutem spojená nebo legislativní činnost zákonodárce směřující k toliko potřebné novelizaci.\\
 
 Přes chybějící explicitní ustanovení přiznávající svěřenskému fondu právní subjektivitu ve vztahu k pozůstalostnímu řízení a dalších ustanovení upravující přímo kolize mezi právní úpravou svěřenského fondu a dědického práva, je i dle mnohých autorů, postavení dědice u svěřenského fondu, nebo svěřenského správce dovodit a je tak možné hovořit o tom, že majetek vyčleňovaný do svěřenského fondu je součástí pozůstalosti a může tak teoreticky být ponížen o povinný díl nepominutelného dědice, neboť svěřenský fond je dědicem.\\
 
 Nicméně samotný fakt, že takováto explicitní úprava chybí a je možné jí dovodit jen s notnou dávkou imaginace a i přes to si člověk nemůže být svým výkladem tohoto problému jistý jasně dokládá potřebu, že je tyto kolize nezbytné explicitně upravit. Já tento samotný návrh činím v kapitole 7.12.4.\\
 
 Na začátku této kapitoly jsem poznamenal, že Kocí staví svojí argumentaci na dvou bodech, pokud tedy navážu na Kocím vyslovený první bod, druhým bodem jeho argumentace je pak jakási nutná rovnost mezi oběmi formami založení svěřenského fondu: \textit{"17. Pokud by měl být majetek vyčleněný do svěřenského fondu součástí pozů­stalosti, mohl by být teoreticky snížen o povinný díl nepominutelného dědice (§ 1642). To by znamenalo, že by zůstavitel pořízením pro případ smrti při existenci nepominutelného dědice nemohl do svěřenského fondu vyčlenit všechen svůj majetek, což by testamentární svěřenský fond značně znevýhodňovalo ve srovnání s ostatními typy svěřenských fondů (viz komentář k § 1449), a to bez rozumného důvodu a jakékoliv opory v ustanoveních o jeho vzniku, která na obě formy zakladatelského aktu, tj. na smlouvu a pořízení pro případ smrti, nahlížejí jako na rovnocenné."}\footfullcite[Komentář k § 1448]{svestka_j_obcansky_2014-1}\\
 
 Přes fakt, že při pohledu na úpravu svěřenského fondu a její konkurenci s právní úpravou dědického práva první vymezenou optikou docházím k obdobnému závěru jako Michal Skuhrovec ve své práci, nemohu se v žádném případě ztotožnit s tímto bodem. V případě ochrany nepominutelného dědice i ve chvíli vyčlenění majetku do svěřenského fondu \textit{mortis causa} ať už dovozené ze současné litery zákona, nebo potenciálně v budoucnu implementované formou novely si nemyslím, že dochází k zásadnímu znevýhodnění formy založení \textit{inter vivos} oproti \textit{mortis causa}. Neboť obě dvě formy, přestože se jedná o stejný institut, jsou založeny jiným způsobem a mohou být založeny často za značně jiným účelem (například ochrana majetku před věřiteli nebo za účelem businessu versus mezigenerační převod majetku). Obě tyto formy lze paralelně připodobnit k možnostem převedení majetku za života zůstavitele, jako je například darování, směna, prodej (neuvažuji daňové zatížení, pouze ve vztahu k institutu nepominutelného dědice), nebo obdobným, podmínkami dědického práva nezatíženým, formám převodu, jakými jsou například darování pro případ smrti\footnote{Viz kapitola 6.3.3 Vznik svěřenského fondu \textit{inter vivos} s odkládací podmínkou} a převedení majetku za pomoci dědického práva, tedy v rámci pozůstalostního řízení, po smrti zůstavitele (kde pokud existuje nepominutelný dědic a není platně vyděděn, má nárok na svůj nepominutelný díl).\\
 
 Přestože se tedy jedná o stejný institut, lze předpokládat, že účelem založení svěřenského fondu pořízením pro případ smrti bude s největší pravděpodobností mezigenerační převod majetku, zatímco u založení svěřenského fondu smlouvou mezi živými může být množina účelů širší. Toto samozřejmě nevylučuje možnost založení svěřenského fondu pořízením pro případ smrti za stejnými účely jako u \textit{inter vivos}, tyto jsou však mnohem výhodnější za zůstavitelova života, už jen zkrátka proto, že je zakladatel naživu. Stejně tak i pokud bychom se bavili o jedné výhodě, o kterou by tak testamentární svěřenský fond při ochraně nepominutelných dědiců přišel, tedy o faktickou možnost jeho pomocí vydědit, stále má tato forma mnoho výhod oproti svěřenskému fondu založeného \textit{inter vivos}, o tomto ostatně hovořím v kapitole 6.3 u jednotlivých forem založení svěřenského fondu.\\
 
 Při pohledu druhou optikou, ostatně u tohoto bodu zastávám stejný postoj i při pohledu optikou první, si nemyslím, že potenciální znevýhodnění svěřenského fondu \textit{mortis causa} je dostatečným důvodem proč nechránit nepominutelné dědice, nebo jejich již existující ochranu vyvracet.\\
 
 V rámci pohledu touto optikou pak neobstojí ani Kocího první bod, připomenu přitom opět výše citovaný názor Pihery, specificky jeho část: \textit{"Majetek vyčleněný zakladatelem (zůstavitelem) do svěřenského fondu mortis causa je součástí pozůstalosti a nabytí tohoto majetku do svěřenského fondu buď potvrzuje soud, nebo je převeden ve prospěch svěřenského fondu dědicem obtíženým odkazem. Z hlediska dědického práva je na postavení svěřenského fondu nutné pohlížet „jako na dědice“ nebo „jako na odkazovníka“, a to se vším všudy, co se s tímto postavením spojuje (např. odpovědnost za dluhy zůstavitele)."}\footfullcite[Komentář k § 1448]{spacil_j_a_kolektiv_obcansky_2013}\\

% -------------------------------------------------------------------------- 
 %Pokud bych měl znova citovat Piheru se pak řadí do opačného názorová spektra, kdy tvrdí následující: \textit{"Poněkud odlišnou povahu má vyčlenění majetku do svěřenského fondu zřizovaného pořízením pro případ smrti (mortis causa). V těchto případech vyčlení zakladatel majetek ve prospěch svěřenského fondu tak, že v jeho prospěch vyhradí dědické právo nebo nařídí odkaz (§ 1452 odst. 1 ve spojení s § 311 odst. 1). Správce svěřenského fondu mortis causa pak vykonává práva dědice majetku vyčleněného do svěřenského fondu nebo práva odkazovníka. Majetek vyčleněný zakladatelem (zůstavitelem) do svěřenského fondu mortis causa je součástí pozůstalosti a nabytí tohoto majetku do svěřenského fondu buď potvrzuje soud, nebo je převeden ve prospěch svěřenského fondu dědicem obtíženým odkazem. Z hlediska dědického práva je na postavení svěřenského fondu nutné pohlížet „jako na dědice“ nebo „jako na odkazovníka“, a to se vším všudy, co se s tímto postavením spojuje (např. odpovědnost za dluhy zůstavitele)."}\footfullcite[Komentář k § 1448]{spacil_j_a_kolektiv_obcansky_2013}\\
% -------------------------------------------------------------------------- 
 
 Základním předpokladem, který musí být pro platnost Piherova názoru splněn je tedy nutnost považovat svěřenský fond za účastníka pozůstalostního řízení, přičemž za dědice \textit{sui generis} bude považován svěřenský správce, k tomu ostatně § 311 OZ přímo vybízí.\\
 
 \textit{"Svěřenský fond může být zřízen též pořízením pro případ smrti, tedy závětí nebo dědickou smlouvou; jeho podmínky mohou být moderovány též dovětkem (§ 1492). Závěť nebo dědická smlouva v takovém případě musí obsahovat náležitosti statutu svěřenského fondu (§ 1452 odst. 1, 2) a musí splňovat formální požadavky kladené na ně dědickým právem (§ 1533–1593). To platí též o vedlejších doložkách v závěti. Správce takto zřízeného svěřenského fondu pak vykonává práva dědice nebo odkazovníka, aniž by však sám byl osobním dědicem či odkazovníkem (dále srov. komentář k § 1451 a § 1452 odst. 1)."}\footnote{Tamtéž.}\\
 
 Pihera tedy krom samotné zákonné ochrany nepominutelného dědice akcentuje i výše zmíněnou ochranu již při konstruování pořízení pro případ smrti, pokud by byl svěřenský fond považován za dědice je tato ochrana bez dalšího uplatnitelná a v případě opomenutí by nepominutelným dědicům vznikla pohledávka za svěřenským fondem, jak již jsem ostatně popsal výše.\\
 
% -------------------------------------------------------------------------- 
 %Je tam vidět intence zákonodárce, nelze tak zcela zavrhovat ani názory Pihery nebo Horna i ty se také zakládají na solidních základech a nelze stoprocentně vyloučit, že právě těmito autory vyslovený pohled je ten správný. Samotné nejasnosti ohledně testamentárního svěřenského fondu však pravděpodobně trvale vyřeší až konstantní praxe s tímto institutem spojená nebo legislativní činnost zákonodárce.
 
 %Zmínit tento článek: https://advokatnidenik.cz/2020/04/04/vlastnictvi-bez-vlastnika-v-ceskem-sverenskem-fondu-skica-na-margo-trustu-v-ceskem-pravu/ , ten řiká, že je v podstatě svěřenský fond právnickou osobu, pak by byla ochrana nepominutelného dědice zajištěna.

%je proto logické, že v současném Občanském zákoníku tyto ustanovení ochybý, tedy nemohla být součástíV této souvislosti bych však chtěl upozornit na kapitolu 7.5.1, kde popisuji ochranu nároků, které lze připodobnit k institutu nepominutelného dědice v Občanském zákoníku provincie Quebec. Tím pádem nelze zcela souhlasit s druhou částí

%kdy bych chtěl jeho názor konfrontovat a rozšířit o můj názor, nicméně z jednoho pohledu s ním však souhlasím.

%Neboť s ním minimálně v rámci jednoho z mých pohledů souhlasím.

%Snahu o zodpovězení tohoto problému neulehčuje ani fakt, že standardně používáné komentáře si při pohledu na tuto problematiku vzájemně odporují.

%upravuje právní řádve svém právním řádu 
% --------------------------------------------------------------------------

Je tedy zřejmé, že pokud bychom důsledně aplikovali konstrukci svěřenského fondu, při pohledu první optikou by tak dědic neměl téměř žádnou možnost dožadovat se svého povinného dílu a jediná diskutabilní obrana by mohla být v souvislosti s § 1492 OZ. Přes toto všechno je však možné si povšimnout jisté tendence zákonodárce za pomoci všech výše zmíněných ustanovení nepominutelného dědice chránit i v takovém případě, to že tato ochrana není ve vztahu ke konstrukci svěřenského fondu moc dobře aplikovatelná je věc jiná. Dle mého pohledu je úmysl zákonodárce nějak dědice chránit i v tomto případě z těchto ustanovení jasný a tedy i při úmyslovém pohledu na věc si myslím, že dědice chránit chtěl. Je tedy obdobně možné, že by soud rozhodl ve prospěch takovéhoto opomenutého dědice právě na základě výše zmíněných ustanovení i přes samotný ohled na konstrukci svěřenského fondu.\\

Pokud by mu tedy přiznal jakousi právní subjektivitu, tak jak tomu jednou udělal v daňových zákonech, bylo by zřejmé, že tato ochrana by dědici byla poskytnuta. S ohledem, že tak zákonodárce jednou i učinil, není absolutně vyloučeno, aby tak udělal i ve vztahu k dědickému právu, ačkoli by to bylo skrytě a mlčky.\\

% --------------------------------------------------------------------------
% ZÁSADNÍ PROBLÉM: Za všechny tyto problémy může samotné pojetí vlastnictví u svěřenského fondu tak, jak jsme jej přejali z právní úpravy Quebecu. Řešením je tedy změna pojetí, návrhy poskytuje například Luboš Tichý v: https://advokatnidenik.cz/2020/04/04/vlastnictvi-bez-vlastnika-v-ceskem-sverenskem-fondu-skica-na-margo-trustu-v-ceskem-pravu/








%Obdobné stanovisko zaujímá například i Kryštof Horn ve svém příspěvku do Časopisu Ad Notam. Ve svém příspěvku říká následující: \textit{"Příkladem takové nejasnosti může být vztah fondu k institutu nepominutelného dědice, který je neznámý jak v zemích common law, kde má své kořeny trust, tak bohužel i v Québecu, odkud zákonodárce přejal úpravu svěřenských fondů. Výslovná ustanovení o vztahu svěřenského fondu mortis causa a dědického práva totiž chybí."}. \footfullcite{sro_httpwwwpraguebestcz_prakticke_nodate-1} \\

%Tento příspěvek se při pohledu na úpravu svěřenského fondu zakládá na pravdě a autor v něm potvrzuje fakt, že takováto právní úprava v českém právu chybí a bylo by vhodné jí prostřednictvím novelizace doplnit. S výrokem nicméně nemohu souhlasit jako s celkem, protože přesto, že je pravda, že institut nepominutelného dědice jako takový v Quebeckém právu neexistuje a Quebecká právní úprava se dále vyznačuje téměr úplnou testamentální volností, tak Quebec Civil Code připouští určitou ochranu dědiců před opomenutím v závěti. Abych byl přesný, tak Quebecký právní řád upravuje tři případy ochrany dědiců, které by se daly přirovnat k institutu nepominutelného dědice v našem právním řádu. Je tedy možné říci, že, přestože tak zákonodárce neučinil, mohl se u Quebecká právní úpravy inspirovat i v rámci ochrany nepominutelných dědiců, protože Quebecké právní úprava, narozdíl od té České, upravuje i případy, ve kterých by tyto nároky nemohly být splněny kvůli právnímu jednání učiněnému zůstavitelem, mezi toto právní jednání je právě možné zařadit i založení trustu.\footfullcite{leclercq_inheritance_2014} \\

%Obdobně je také možné hledat inspiraci pro vyřešení této nejastnosti ve vztahu k dědickému právu i v případech z Itálie, či v právní úpravě institutu obdobného svěřenskému fondu v Lichtenštejnsku, tento fakt a jeho možné aplikace na českou právní úpravu jsou zhodnoceny v kapitolách. \\

%Z výše uvedeného tedy vychází, že při kolizi práva nepominutelného dědice a vytvoření svěřenského fondu \textit{mortis causa} není zřejmé, zda je právo na povinný díl chráněno, nebo nikoli. Při pohledu mojí první definovanou optikou se naprosto ztotožňuji s výkladem Skuhrovce a to do té míry, že nepominutelný dědic není nikterak chráněn. Není ale možné přehlížet ustanovení, která by přes svojí zjevnou nekompatibilitu s konstrukcí svěřenského fondu přesto měla být aplikována, z těchto vyvozuji potenciální úmysl zákonodárce chránit nepominutelné dědice i v případě založení svěřenského fondu \textit{mortis causa} pořízením pro případ smrti a při užítí druhé definované optiky je tak možné dojít i protichůdného názoru a tvrdit, že nepominutelný dědic chráněn je. Zůstává tedy otázkou, zda by soud v případě sporu prioritizoval naprostou testovací volnost zakladatele, nebo důkladně aplikova ustanovení o nepominutelném dědici i přes jejich zjevnou nedotaženost a nekompatibilitu ve světle úpravy svěřenského fondu. Samotná absence norem, které by toto explicitně upravovali je alarmující, výklad ve prospěch nepominutelného dědice tedy při důsledné aplikaci konstrukce svěřenského fondu nemůže obstát, bylo by ale nutné naprosto ignorovat jiná ustanovení omezující konání zůstavitele jako je § 1492 OZ, či ustanovení uznačující svěřenský fond jako dědice, tedy § 311 OZ. Minimálně ustanovení § 1492 OZ by však mělo být bez dalšího zůstavitelem respektováno v případě, že nepominutelného dědice platně nevydědil, lze tedy obdobně tvridit, že toto ustanovení samo o sobě ochranu nepominutelnému dědici poskytuje a je do jisté míry kompatibilní i s absencí právní osobnosti svěřenského fondu, v takovém případě je však nutné pohlížet na svěřenský fond jako na dědice (entitu) \textit{sui generis}, vůči které může uplatnit opomenutý dědic svoji pohledávku a být tak jeho věřitelem, při takovémto pohledu tak nemůže váhu argumentů ustát ani první tvrzení. Nezbývá však než nechat řešení této otázky na zákonodárci, či na výkladu soudu, ani jeden výklad totiž není bez dalšího superiorní a nelze považovat za jediné správné řešení.\\
% --------------------------------------------------------------------------

\subsubsection{Dílčí závěr, komparace právní úpravy svěřenského fondu s právní úpravou dědictví}

Z výše uvedeného tedy vychází, že při kolizi práva nepominutelného dědice a vytvoření svěřenského fondu \textit{mortis causa} není zřejmé, zda je právo na povinný díl chráněno, nebo nikoli. Při pohledu mojí první definovanou optikou se naprosto ztotožňuji s výkladem Skuhrovce a to do té míry, že nepominutelný dědic není nikterak chráněn. Není ale možné přehlížet ustanovení, která by přes svojí zjevnou nekompatibilitu s konstrukcí svěřenského fondu přesto měla být aplikována, z těchto vyvozuji potenciální úmysl zákonodárce chránit nepominutelné dědice i v případě založení svěřenského fondu \textit{mortis causa} pořízením pro případ smrti a při užití druhé definované optiky je tak možné dojít i protichůdného názoru a tvrdit, že nepominutelný dědic chráněn je. Zůstává tedy otázkou, zda by soud v případě sporu prioritizoval naprostou testovací volnost zakladatele, nebo důkladně aplikoval ustanovení o nepominutelném dědici i přes jejich zjevnou nedotaženost a nekompatibilitu ve světle úpravy svěřenského fondu. Samotná absence norem, které by toto explicitně upravovali je alarmující, výklad ve prospěch nepominutelného dědice tedy při důsledné aplikaci konstrukce svěřenského fondu nemůže obstát, bylo by ale nutné naprosto ignorovat jiná ustanovení omezující konání zůstavitele jako je § 1492 OZ, či ustanovení označující svěřenský fond jako dědice, tedy § 311 OZ. Minimálně ustanovení § 1492 OZ by však mělo být bez dalšího zůstavitelem respektováno v případě, že nepominutelného dědice platně nevydědil, lze tedy obdobně tvridit, že toto ustanovení samo o sobě ochranu nepominutelnému dědici poskytuje a je do jisté míry kompatibilní i s absencí právní osobnosti svěřenského fondu, v takovém případě je však nutné pohlížet na svěřenský fond jako na dědice (entitu) \textit{sui generis}, vůči které může uplatnit opomenutý dědic svoji pohledávku a být tak jeho věřitelem, při takovémto pohledu tak nemůže váhu argumentů ustát ani první tvrzení. Nezbývá však než nechat řešení této otázky na zákonodárci, či na výkladu soudu, ani jeden výklad totiž není bez dalšího superiorní a nelze považovat za jediné správné řešení.\\

V procesu nalézání úmyslu zákonodárce je samozřejmě možné říci, že zákonodárce chtěl pro svěřenský fond zřizovaný \textit{mortis causa} pouze konstruovat způsob, jakým tak může občan učinit a zvolil k tomu poněkud nešťastně způsob pořízení pro případ smrti, čímž společně s § 1452 mohl vytvořit neúmyslnou falešnou představu, že svěřenský fond je dědicem a na právní úkon učiněný za jeho zřízením, tedy pořízení pro případ smrti, je nutné vztahovat i omezení tohoto pořízení se týkající. Je možné na problematiku tedy pohlédnout tak, že zákonodárce pouze pojal zřízení svěřenského fondu pomocí pořízení pro případ smrti jen jako způsob, který ve skutečnosti nechtěl omezit instituty dědického práva, nebo omezení kladenými na pořízení pro případ smrti. Je tedy k diskuzi, zda neměl zákonodárce raději zvolit možnost zřízení svěřenského fondu \textit{mortis causa} pouze smlouvou pro případ smrti místo pořízení, zejména i proto, že tento způsob již občanský zákoník upravuje ve formě darování pro případ smrti. Takováto by tak totiž potom nebyla svázána instituty dědického práva, z tohoto pohledu je třeba zmínit, že není vyloučena možnost zřízení svěřenského fondu \textit{inter vivos} s odkládací podmínkou smrti, která by tak tento způsob založení naplnila (jednalo by se v podstatě o založení svěřenského fondu \textit{mortis causa}, avšak bez omezení na něj kladených), avšak pokud by bylo umožněno zřizovat svěřenský fond pouze smlouvou, nebylo by možné svěřenský fond zřídit jednostranně, pokud by tak chtěl zákonodárce tuto možnost zachovat, musel by provést patřičné změny.\\

Pak by se beze všech pochybností dalo nahlížet na svěřenský fond jako na možnost vyloučení majetku zakladatele zcela z dědického řízení s podstatnými následky např. i pro nepominutelné dědice, jak rozvádí Kocí.\\

S ohledem na vše výše uvedené bych doporučil potenciálnímu zakladateli založení svěřenského fondu \textit{inter vivos}, popřípadě s odkládací podmínkou smrti. V případě absence nepominutelných dědiců, nebo určení nepominutelných dědiců jako obmyšlených svěřenského fondu či platného vydědění nebude tato kolizní problematika hrát takovou roli, v tomto případě tak není třeba odrazovat od založení svěřenského fondu \textit{mortis causa}.\\

Závěrem bych chtěl dodat, že přes všechno výše vyřčené se osobně přikláním k druhému pohledu. Tedy přikláním se k existenci ochrany nepominutelného dědice vzhledem k principu koherence právního řádu, který stanovuje, že je nepřípustné, aby jediný institut zcela popřel ucelený komplex institutů jiných, v tomto případě dědického práva, které lze navíc vykládat jako institut požívající ústavní ochranu\footnote{Srovnání Listina základních práv a svobod čl. 11 odst. 1, Dědění se zaručuje.}, aniž by pro to byla výslovná a jednoznačná zákonná opora, ta však dle mého výkladu nelze hledat v samotné konstrukci svěřenského fondu, ale měla by být v zákoně jasně stanovena, tak tomu však není a § 1452 OZ, § 1492 OZ a § 311 OZ dokonce stanovují přímý opak, tedy jistou ochranu nepominutelných dědiců. Do doby novelizace Občanského zákoníku, či rozhodnutí soudu měnící momentální paradigma by se měla právní úprava v pochybnostech dle mého vykládat ve prospěch nepominutelných dědiců.\\

% --------------------------------------------------------------------------
%Možnost jak se vyhnout takovéto nejasnosti pak představuje založení svěřenského fondu \textit{inter vivos}, při zobrazení časové osy vzniku svěřenského fondu je toto dobře doložitelné, samotná vizualizace pomůže i k následnému návrhu

%Tady je potřeba to trochu upravit, divně to navazuje, ale jinak je to fajn.
% --------------------------------------------------------------------------

Jak tedy tvrdí Kryštof Horn ve svém příspěvku do Časopisu Ad Notam: \textit{"...nejasnosti může být vztah fondu k institutu nepominutelného dědice, který je neznámý jak v zemích common law, kde má své kořeny trust, tak bohužel i v Québecu, odkud zákonodárce přejal úpravu svěřenských fondů. Výslovná ustanovení o vztahu svěřenského fondu mortis causa a dědického práva totiž chybí."}. \footfullcite{sro_httpwwwpraguebestcz_prakticke_nodate-1} \\

Tento příspěvek se při pohledu na úpravu svěřenského fondu a na výše provedenou analýzu zakládá na pravdě a autor v něm potvrzuje fakt, že takováto jasná právní úprava v českém právu chybí a bylo by vhodné jí prostřednictvím novelizace doplnit.\\

S ohledem na další část však s výrokem nemohu souhlasit jako s celkem, jak již jsem výše letmo zmínil, protože přesto, že je pravda, že institut nepominutelného dědice jako takový v Quebeckém právu neexistuje a Quebecká právní úprava se dále vyznačuje téměř úplnou testamentární volností, tak Quebec Civil Code připouští určitou ochranu dědiců před opomenutím v závěti, kterou ve stručnosti přiblížím, jak již jsem slíbil v předchozí kapitole. Abych byl přesný, tak Quebecký právní řád upravuje tři případy ochrany dědiců, které by se daly přirovnat k institutu nepominutelného dědice v našem právním řádu. Je tedy možné říci, že, přestože tak zákonodárce neučinil, mohl se u Quebecká právní úpravy inspirovat i v rámci ochrany nepominutelných dědiců, protože Quebecké právní úprava, na rozdíl od té České, upravuje i případy, ve kterých by tyto nároky nemohly být splněny kvůli právnímu jednání učiněnému zůstavitelem, mezi toto právní jednání je právě možné zařadit i založení trustu.\footfullcite{leclercq_inheritance_2014} \\

Obdobně je také možné hledat inspiraci pro vyřešení této nejasnosti ve vztahu k dědickému právu i v případech z Itálie, či v právní úpravě institutu obdobného svěřenskému fondu v Lichtenštejnsku, tento fakt a jeho možné aplikace na českou právní úpravu jsou zhodnoceny v následujících kapitolách. \\

% --------------------------------------------------------------------------
%Domnívám se, že problém s nepominutelným dědicem by potenciálně mohl zakladatel obejít, pokud by stanovil odkaz.

%Tuto kapitolu možná přejmenuji jako na dílčí závěr problematiky nepominutelného dědice zkoumané v této části.

%---Koment\\
%Zde proběhne komparace, musím zde napsat, že neexistuje právní úpravy, která by upravovala situaci, ve které bude založen svěřenský fond mortis causa a zároveň bude probíhat dědické řízení, ve kterém budou figurovat nepominutelní dědici, na které v souvislosti se založením svěřenského fondu nevyjde jejich zákonný podíl. Konstatovat, že těmto problémům se dá vyhnout tím, že se svěřenský fond vyčlení už za života, dokázat toto tvrzení na základě jiných institutů převodu majetku za zůstavitelova života a poukázat na to, že tento majetek následně není předmětem dědického řízení a tudíž ani nemůže připadnout jakýmkoliv dědicům, neboť pokud budeme mít na časové ose bod x, ve kterém zůstavitel zemře, a svěřenský fond vznikne v bodě x-y, tak v době zůstavitelova úmrtí již není součástí jmění a tudíž nemůže být součástí dědického řízení. Pokud by nicméně zůstavitel zemřel v bodě x a svěřenský fond by také vznikl v bodě x, tak lze polemizovat o tom, zda by dědici nemohli mít právo na snahu obsáhnout i tento vyčleněný majetek v soupisu dědického řízení.\\
%---\\

%---Koment\\
%Dále se pokusit obhájit můj názor, že pokud budou obmyšlenými svěřenského fondu i nepominutelní dědicové, tak plněním ze svěřenského fondu právě tímto dědicům dojde k neformálnímu doplnění jejich povinného podílu ze strany zůstavitele - pokusit se najít nějakou právní zásadu, která by něco takového říkala. Pokud by však z fondu mělo být plněno nikomu jinému, tak se pokusit obhájit fakt, že tato situace není zákonně obhájena a není na ní ani žádná judikatura či ustálená soudní praxe, a potenciálně by se tedy dědicové mohli bránit soudní cestou. Pokusit se stanovit, jak by takovýto spor mohl v rámci našich existujících právních norem dopadnout, pokud to půjde tak se zkusit podívat, zda toto neni řešeno na mezinárodní úrovni například v rámci mezinárodního práva soukromého. \\
%---\\

%Pro tuto kapitolu, a následné návrhy na změny, je v první řadě podstatné vymezit,  v jakých fázích vzniká svěřenský fond, a tyto následně porovnat s úpravou dědického práva. Tímto postupem lze dosáhnout vhodného a dostatečného porovnání možností, které mají dědici v jednotlivých fázích, jak jsou chráněni a jak se potenciálně mohou brát o svá práva. \\ 

%Dle mého názoru, který potvrzuje Miloš Kocí v rámci komentáře k úpravě svěřenského fondu, lze hovořit o 3 způsobech vzniku svěřenského fondu. Těmito způsoby jsou, založení svěřenského fondu \textit{inter vivos}, \textit{mortis causa} a potenciální souběh těchto založení ve formě založení svěřenského fondu \textit{inter vivos} s odkládací podmínkou, dle § 584 až 549 OZ, spočívající ve smrti zůstavitele.\footfullcite{svestka_j_obcansky_2014-1} \\

%\newpage

%\thispagestyle{smallertextinheader}

%Vizuálně by se tedy průběh komparace napříč jednotlivými způsoby založení dal vyobrazit následovně, kde střed osy představuje čas úmrtí zůstavitele:

%\begin{figure}[h]
%\centering
%\includegraphics[width=18cm,height=6cm,keepaspectratio]{Vznik_sverenskeho_fondu.png}
%\caption{Vznik svěřenského fondu}
%\label{fig:komparace}
%\end{figure}

%---Koment\\
%Doplnit další informace.
%Popsat dědické právo - co jde do dědictví, co jsou nepominutelní dědicové, jak funguje ědické řízení atp.
%Vysvětlit rozdíl mezi majetkem a jměním.\\
%---\\

%V případě založení svěřenského fondu mezi živými s odkládací podmínkou smrti zakladatele se ve své podstatě jedná o založení \textit{mortis causa} formou \textit{inter vivos}, s ohledem, že se nejedná o založení svěřenského fondu pořízením pro případ smrti je zřejmé, že v tomto případě nelze uvažovat ani nad ochranou dle §1492, pro ilustraci je možné porovnat tento způsob založení s další formou mezigeneračního předání majetku a to darováním pro případ smrti. Podklad pro toto tvrzení poskytují například autoři článku Darování pro případ smrti v Časopisu pro právní vědu a praxi Adam Talanda a Iveta Talandová. Ti přímo říkají: \textit{"Darování pro případ smrti je při splnění všech zákonných podmínek obligačním vztahem mezi dárcem a obdarovaným a neřídí se ustanoveními dědického práva. Zákon tak umožňuje užitím jednoho zákonného institutu obcházet soubor jiných zákonných institutů. Především se jedná o obcházení ustanovení omezujících zůstavitele v nakládání s jeho majetkem, která slouží k ochraně nepominutelných dědiců."}\footfullcite{noauthor_casopis_2015}.

%Dědic by se s ohledem na paragraf 1492 dal brát jako věřitel zůstavitele. Musím na to nahlížet z hlediska veřejného pořádku a základních morálních principů.

%Je darování vždy oboustraný vztah, co když mu daruji například peníze jednostranně, bude se pak také jednan o obligatorní právní vztah založený na smlouvě? Nebo poté bude předmět takovéhoto darování součástí pozůstalosti a bude se na něj například nahlížet jako na odkaz? Nikoliv je to dvoustranné právní jednání, jedná se tedy o obligatorní vztah.
% --------------------------------------------------------------------------

\subsubsection{Ochrana zájmů nepominutelných dědiců zakladatele trustu v Itálii}

Itálie se vydala oproti jiným právním řádům jinou cestou, když přijala Haagskou úmluvu o uznávání trustů. Účelem této úmluvy pak není pozitivní povinnost upravit trust v jednotlivých právních řádech signatářských zemí, ale pouze na určení rozhodného práva a uznání trustů\footfullcite{noauthor_hcch_nodate}.\\

V Itálii užívaný trust \textit{interno}, na něž lze úspěšně aplikovat jak právo trustů, v čemž se italské soudy ukázaly jako velice efektivní\footfullcite{lupoi_m_country_2012}, tak i kolizní ustanovení italského práva, například o nepominutelném dědici, ačkoli opět řešené u soudu\footfullcite[Riziko zneužití trustu, poučení z italské zkušenosti, strana 197]{lubos_tichy_sverensky_2017}, poskytují v Itálii možnosti, které samotné italské právo mnohdy nenabízí.\\

Z principu možnosti použití ustanovení práva určeného kolizními normami právního státu místa soudu, jako jsou například kogentní ustanovení dědického práva tak dochází v Itálii nejenom možnosti založení trust, ale zároveň s tím i ochrany nepominutelných dědiců. Založení trustu i pokud by vedlo k "obcházení" práv nepominutelných dědiců však nevede k neplatnosti jejich založení jak v několika soudních řízeních stanovil italský soud v Udine či Turíně \footnote{Tamtéž, strana 198}, stejně tak, jako opomenutí dědice nezakládá neplatnost pořízení pro případ smrti ani v českém právu. Dle rozhodnutí italských soudů tedy nelze žalobou požadovat zrušení založeného trustu, ale zároveň je možné dožadovat se svého dědického nároku žalobou na vydání, ta totiž může směřovat proti nepřímým darům, přičemž na převody majetku do trustového zřízení je v Itálii nahlíženo jako na nepřímé darování.\\

Zajímavé je pak zmínit, že žalobě na vydání nebylo vyhověno rozhodnutím soudu v Urbine, i přes opomenutí nepominutelných dědiců. Soud totiž shledal, že zřízení trustu bylo učiněno za účelem plánování mezigeneračního převodu majetku a žalobce tak nemá právo žalovat na vydání, neboť věc mu bude plněna jako obmyšlenému v rámci trustového zřízení\footnote{Tribunale di Urbino, 31.ledna 2015, z Luboš Tichý, Svěřenský fond a trust - jejich fungování v mezinárodním srovnání.}. Pokud bychom tedy takovéto rozhodnutí aplikovali v mezích českého práva, bylo by možné bavit se nejenom o ochraně dědiců a jejich možnosti žalovat na vydání, ale také o započtení plnění ze svěřenského fondu na dědický podíl.\\

Problémem však zůstává to, že Česká republika není signatářem výše zmíněné úmluvy, jednotlivé její články tak u nás nejsou akceptovatelné a aplikovatelné. Na svěřenský fond je tedy nutné pohlížet optikou ustanovení občanského zákoníku, která nepostihují všechny kolize. Je však možné, že se podobného rozhodnutí jako v Itálii dočkáme i od některého z českých soudů, který buď potvrdí, nebo vyvrátí ochranu nepominutelného dědice například za pomocí fingování svěřenského fondu jako subjektu práva.\\

% --------------------------------------------------------------------------
%Napsat ještě o možnosti založení cizího trustu na našem území pomocí aplikací jednotlivých ustanovení zákona o mezinárodním právu soukromém. Asi ne, napíšu o tom výše v rozhodujících faktorech pro výběr formy mezigeneračního převodu majetku.

%V této kolizi italské soudy již také rozhodly.

%Napsat o situaci v Itálii.\footfullcite{lubos_tichy_sverensky_2017}
% --------------------------------------------------------------------------

\subsubsection{Právní úprava v Lichtenštejnsku}

Zásadním rozdílem mezi českým svěřenským fondem a Lichtenštejnským \textit{die Treuhänderschaften} je postavení správce. Čl. 919 odst. 3 PGR\footnote{Personen- und Gesellschaftsrecht, zákon o osobách a společnostech.} pak stanoví následující: \textit{"Der Treuhänder ist unter Vorbehalt seiner Verpflichtungen aus der Treuhandurkunde berechtigt, über das Treuhandgut gleich einem selbständigen Träger von Rechten und Pflichten, wie namentlich ein Eigentümer, Gläubiger, Mitglied oder Organ einer Verbandsperson oder Gesellschaft oder dergleichen zu verfügen, für das Treugut vor allen Behörden und in allen Verfahren im eigenen Namen als Partei, Beteiligter, Beigeladener, Intervenient und dergleichen aufzutreten, die zu ihm gehörigen Rechte gegen alle Dritte gemäss der Treuhandurkunde zu verwalten und auszuüben und, soweit nötig, zu versilbern und neu anzulegen, wenn es sich aus dem Treuhandzweck nicht anders ergibt."}\footfullcite[Článek 919]{noauthor_1926004_nodate}.\\

Správce je tedy s výhradou povinností, které plynou ze zakladatelské listiny oprávněn se majetkem trustu disponovat v rovném postavení jako samotný nositel práv a povinností, zákon pak přímo stanovuje jako vlastník, věřitel, či člen nebo orgán společnosti a korporace a je oprávněn za majetek trustu vystupovat v každém řízení před všemi orgány veřejné moci v každém procesním postavení, jedná se tedy o vlastnictví správce \textit{sui generis} oproti autonomnímu vlastnictví v českém Občanském zákoníku. Toto v souvislosti s § 762 ABGB\footnote{Allgemeines bürgerliches Gesetzbuch, občanský zákoník.}, které zakládá pro děti zůstavitele obdobné postavení, jako má nepominutelný dědic v českém právu, a § 785 ABGB, které podle odstavce 1 stanoví, že převody učiněné zůstavitelem v období 2 roky před smrtí se započítávají na dědický podíl jednotlivých nepominutelných dědiců. Ti tak vůči správci mohou uplatnit své právo na zvýšení jejich dědického podílu\footfullcite{noauthor_abgb_nodate}.\\

Lichtenštejnská úprava se tedy na první pohled zdá být dobrým inspiračním zdrojem pro českého zákonodárce k vyřešení problémů spojené s nepominutelným dědicem. Občanský zákoník žádnou takovouto výslovnou ochranu neposkytuje, jistou analogii lze spatřovat v § 1588 a § 1492 jakožto ustanovení poskytující ochranu při konstruování pořízení pro případ smrti, jak je ale zmíněno výše, toto naráží na problém potenciální neaplikovatelnosti ustanovení dědického práva na svěřenský fond. Tento problém vychází z problematiky pojetí vlastnictví ve svěřenském fondu, čerpání z Lichtenštejnského ABGB by tak bylo značně problematické, kvůli neaplikování Lepaullova pojetí vlastnictví, je nutné si uvědomit, že konstrukce vlastnictví je zcela jiná\footfullcite[Fiduciární smlouvy v lucemburském právu, strana 109]{trust_2014}.\\

% --------------------------------------------------------------------------
%PŘEPSAT
%Zajímavé je srování s Lichtenštějsnkou právní konstrukcí, která obsahuje (byť s odchylkami, viz kapitola 3.2.3.) oba dva instituty a zároveň řeší jejich vzájemný vztah. Tzv. „Noterben” (nepominutelní dědici) mohou vůči správci trustu uplatnit své dědické právo, aby zvýšili svůj dědický podíl. Ve smyslu § 785 (3) ABGB takové právo mají pouze stran převodů majetku ze zakladatele na správce do 2 let před smrtí zůstavitele, resp. zakladatele. Promlčecí lhůta je pak podle § 1487 ABGB objektivní ve lhůtě 3 roky.
%123 POPOVICI, Alexandra in TICHÝ, Luboš, ed. Svěřenský fond a trust - jejich fungování v mezinárodním srovnání. Praha: Centrum právní komparatistiky Právnické fakulty Univerzity Karlovy, 2016, s. 36.
 %56
%Jakkoliv je taková konstrukce inspirativní, vůči OZ je zde hned několik odlišností. První z nich je fakt, že OZ ekvivalentní ustanovení neobsahuje. Obsahově se mu podobá § 1588 odst. 2, který poskytuje určitou ochranu smluvnímu dědici. Činí tak však pouze pro smluvního dědice a navíc pouze co do jednání pro případ smrti (pořízení, smlouva), a contrario žádná podobná ochrana jiným dědicům poskytována není. Druhou je pak zásadní rozdíl v lichtenštejnské Treuhänderschaft a svěřenském fondu. Lichtenštejnská úprava počítá s převodem na správce a jeho vlastnickým právem sui generis. V českém právním prostředí by úprava opět narážela na problém autonomního vlastnictví (majetkové právo nikomu nesvědčí). Obdobně stavěná je i ochrana nepominutelného dědice v Lucembursku.124

%Toto bych chtěl asi předělat na subsubsection, myslím ty nároky.
% --------------------------------------------------------------------------

\subsubsection{Jiné obdobné nároky v Quebecu připodobitelné k nároku na povinný díl}

K ilustrování možností inspirace pro řešení problémů spojených s nedostatečnou právní úpravou vztahující se na kolizi právní úpravy svěřenského fondu a právní úpravy nepominutelného dědice je možné nahlédnout do toho samého právního řádu, z něhož byla transplantována právní úprava trustu.\\

% --------------------------------------------------------------------------
%------ Tady hovořím o výživném, nejspíše nemá cenu používat

%Okrajově není na škodu zmínit ani vyživovací povinnost. Problematika s ní spojená se týká zřejmě jen svěřenského fondu založeného \textit{inter vivos}, v případě svěřenského fondu \textit{mortis causa} se dostáváme do dědické sféry, vyživovací povinnost totiž smrtí povnného končí, jejíž problematiku jsem popsal výše.\\

%Účelem této části není nalézání problémů týkajících se svěřenského fondu a jeho dopadu na vyživovací povinnost a \textit{vice versa} ani vytváření arbitrárních problémů s tímto spojených. V tomto ohledu také ani neshledávám zásadní nedostatky současné právní úpravy.\\

%S ohledem na veřejný zájem je však vhodné uvažovat o problému, který může vzniknout po určení výživného ze strany soudu, kdy by povinný následně založil svěřenský fond a převedl do něj veškerý svůj majetek. V případě, že by byl i obmyšlený, mohla by případná pohledávka dopadnout na plnění ze svěřenského fondu. V případě, že by ale obmyšlený nebyl, mohla by tato skutečnost zapřičinit zhoršení dobytnosti pohledávky výživného. Stejná situace by však mohla nastat i při jakýchkoliv jiných dispozicích s majetkem ze strany povinného, toto je tedy obecný problém a nelze jej vztahovat pouze ke svěřenským fondům.\\

%Samotný svěřenský fond nicméně může být dobře použitelný k plnění vyživovací povinnosti, viz článek: https://www.svsp.cz/magazin/clanky/sverensky-fond-jako-nastroj-ke-sprave-vyzivneho-urceneho-k-tvorbe-uspor.

%https://www.svsp.cz/magazin/clanky/sverensky-fond-jako-nastroj-ke-sprave-vyzivneho-urceneho-k-tvorbe-uspor

%Doplnit doktrína

%------
% --------------------------------------------------------------------------

\subsubsection{Typy ochrany dědických nároků v rámci Quebec Civil Code}

V Quebeckém právní řádu existují dva instituty, které by se svým účelem a formou daly přirovnat k českému institutu nepominutelného dědice. Ve své podstatě se jedná o instituty, které mají chránit osobu, které je zůstavitel povinen poskytnou nějaké plnění před opomenutím této osoby ze strany zůstavitele v rámci závěti, nebo nevzetí v potaz této osoby pokud se dědí na základě intestátní posloupnosti. Zákonodárce zároveň upravil vztah těchto institutů k právnímu jednání zůstavitele, které mění výši zůstavitelova majetku již v období před jeho smrtí a s účinností k datu jeho smrti\footfullcite[Strana 373]{palmer_mixed_2012-1}. \\

\vspace{5 mm}

Obě ochrany, které jsou upravené v Quebeckém právním řádu:

\begin{enumerate}
\item Survival of the Obligation to Provide Support\footnote{Jistou obdobu pak představuje v OZ právo některých osob na zaopatření, srovnání § 1665 OZ.}
\item Family Patrimony	
\end{enumerate}

% --------------------------------------------------------------------------
%\item Compensatory Allowance
% --------------------------------------------------------------------------

\vspace{5 mm}

\underline{\textbf{Survival of The Obligation to Provide Support}}\\

Prvním způsobem ochrany jistých dědických nároků vůči majetku zachované\-mu zůstavitelem spočívá v institutu zvaném \textit{The Survival of the Obligation to Provide Support}, který zaručuje příjemci určité finanční podpory vyplácené danému příjemci za zůstavitelova života tuto podporu i po smrti zůstavitele. \\

Důležitou roli pro dokázání argumentu, že Quebecká právní úprava obsahuje v jisté formě ochranu dědiců a upravuje i případy zmenšení zůstavitelova majetku v období před smrtí popřípadě účinné ke dni smrti zůstavitele je článek 689, který stanovuje následující: \textit{"Where the assets of the succession are insufficient to make full payment of the contributions due to the spouse or to a descendant, as a result of liberalities made by acts inter vivos during the three years preceding the death or having the death as a term, the court may order the liberalities reduced.}\footfullcite[Článek 689, odstavec 1]{noauthor__nodate}\\ 

Tento odstavec ve zkratce zaručuje osobám s nárokem vůči zůstavitelovi ochranu tohoto nároku v případě nedostatečného jmění v rámci dědictví, které bylo jistým právním jednáním zůstavitele sníženo na tuto nedostatečnou hranici. Soud může v tomto případě omezit právní jednání učiněné zůstavitelem dokonce i za jeho života, co je ale důležité pro založení trustu mortis cause je fakt, že může omezit i právní jednání s účinností ke dni zůstavitelovi smrti. Tímto způsobem tedy může omezit i vyčlenění majetku do svěřenského fondu, respektive trustu. \\

Dále je potřebné a důležité poznamenat i to, že článek 689 ve své druhé části stanovuje, že v případě, že z těchto právních jednání bude mít prospěch právě osoba, která se domáhá svého podílu na zůstavitelově majetku, tak tento prospěch bude započten v rámci potenciálního omezení právního jednání zůstavitele\footfullcite[Článek 689, odstavec 2]{noauthor__nodate}. Z této části je tedy možné dovodit, že uspokojit tento právní nárok je možné i založením trustu, který bude daný prospěch oprávněným osobám poskytovat místo likvidátora dědictví\footnote{V Quebeckém právu je likvidátor osobou, která vypořádává dědictví, její úlohu můžou plnit i jednotliví dědicové, viz https://www.justice.gouv.qc.ca/en/your-money-and-your-possessions/successions/what-needs-to-be-done-if-a-close-relative-or-friend-dies/appointing-a-liquidator/.}. Dá se tak tedy fakticky hovořit o obdobě započtení na dědický podíl v českém právu.\\

Dle článku 693 pak takováto reduction of liberalities směřuje vůči jednotlivým stranám, které od zůstavitele něco obdrželi ve formě "gratuitous dispositions"\footfullcite[Strana 48]{braun_passing_2016}, jedná se tedy i o založený trust\footfullcite[Strana 13]{lionel_d_smith_re-imagining_2012}. Trustu by pak nemělo nic bránit v tom stát se fakticky dlužníkem\footfullcite{noauthor_vlastnictvi_2020}, jež by měl následně povinnost splnit správce\footfullcite[Strana 342]{jaroslav_svejkovsky_sprava_2015}.\\

Novelizaci, která by se tímto inspirovala a vložila do rukou soudům možnost omezit dispozice s majetkem minimálně k datu smrti zůstavitele by tak vedla k vyřešení problémů kolize jednotlivých ustanovení. V případě uznání ochrany nepominutelného dědice podle § 1492 by pak takováto úprava byla nadbytečná, pokud by tedy zákonodárce nezamýšlel omezení testovací volnosti zůstavitele v určitém období před jeho smrtí, jak tomu částečně činí, avšak bez stanovení lhůty, v § 1588 OZ.\\

\underline{\textbf{Family Patrinomy}}\\

Pro účely family patrimony je důležité ustanovení 416: \textit{"In the event of separation from bed and board, or the dissolution or nullity of a marriage, the value of the family patrimony of the spouses, after deducting the debts contracted for the acquisition, improvement, maintenance or preservation of the property composing it, is equally divided between the spouses or between the surviving spouse and the heirs, as the case may be."}\footfullcite[Článek 416, odstavec 1]{noauthor__nodate}.\\

Článek 421 pak stanoví následující: \textit{"Where property included in the family patrimony was alienated or misappropriated in the year preceding the death of one of the spouses or the institution of proceedings for separation from bed and board, divorce or annulment of marriage and was not replaced, the court may order that a compensatory payment be made to the spouse who would have benefited from the inclusion of that property in the family patrimony."}\footfullcite[Článek 421, odstavec 1]{noauthor__nodate}.\\

Je tedy opět možné dovodit, že by správce trustu, do něhož vyčlenil manžel část "společného jmění manželů" měl poskytnou jednomu z manželů určité kompenzační plnění. Principiálně se tedy jedná o stejnou ochranu jako v předchozím případě.\\

Pokud by tak do trustu byla vložena například nemovitost, měl by manžel nárok na určitou část této hodnoty došlo-li by k jednomu z důvodů zakládající nutnost vyrovnání v rámci family patrimony. Toto ostatně stanovil nejvyšší soud ve svém rozhodnutí 2019 SCC 62, kdy rozhodl, že hodnota rodinného obydlí v trustu je do trustu zahrnuta v případě práva jí užívat\footfullcite{noauthor_yared_2019}.\\

% --------------------------------------------------------------------------
%\underline{\textbf{Compensatory Allowance}}

%\subsubsection{Ochrana v Anglii}
% --------------------------------------------------------------------------

\subsection{Další systémové nedostatky}

Mimo výše uvedené nedostatky české právní úpravy, respektive potenciálními problémy svěřenských fondů spojené s úpravou svěřenských fondů, je minimálně vhodné zmínit i množinu dalších problémů, kterým mohou zúčastněné osoby v souvislosti se založením svěřenského fondu čelit. Hlavní nosnou částí je pak popis problému samotného rozdílného pojetí vlastnictví svěřenského fondu, respektive použití Lepaullova vlastnictví, které považuji za společného jmenovatele ve své podstatě veškerých, v této části popsaných, problémů. Popis tohoto problému je nadmíru důležitý s ohledem na návrh řešení, neboť ty spočívají především ve změně pojetí autonomního vlastnictví, nebo subjektivizace svěřenského fondu.\\

% --------------------------------------------------------------------------
%Slouží jako další výběr z problémů spojených se svěřenským fondem, které považuji v souvislosti s účelem mezigeneračního převodu majetku za důležité zmínit.
% --------------------------------------------------------------------------

\subsubsection{Nejasnosti v rámci daňových aspektů svěřenského fondu}

% --------------------------------------------------------------------------
%Samotná daňová úprava je relativně OK, rpoblém nastává při tom ukončení (jak je tam zdaněný převod), jak jsou chráněni věřitelé když neprobíhá likvidace. Daňové zatížení respektive proces zdanění při ukončování svěřenského fondu by měl být stejný jako při vyplácení ze svěřenského fondu za jeho existence, našel jsem ale pouze jeden názor toto tvrdící někde na nějaké jedné internetové stránce.
% --------------------------------------------------------------------------

Jisté nejasnosti lze pozorovat i v právní úpravě svěřenského fondu ve vztahu k daním. V souvislosti s přijetím nového Občanského zákoníku a jeho účinností bylo potřeba novelizovat i další právní předpisy. Jednou z těchto novelizací prošly právní normy upravující zdaňování, především tedy zákon o daních z příjmu. I přes tyto novelizace však shledávám v daňové úpravě vztahující se ke svěřenským fondům jisté nedostatky. Nemyslím si, že s ohledem na téma práce je vhodné hlouběji zabředávat do tématu zdanění svěřenského fondu a jeho výhodnosti z pohledu daňového zatížení, neboť žádná takováto výhodnost neexistuje, avšak v tomto rámci se vyskytuje i jeden problém týkající se zániku svěřenského fondu a následného převedení zbylého majetku obmyšlenému, nebo zakladateli, který je dle mého názoru vhodný uvést kvůli potenciálním nejasnostem vztahujícím se k zániku svěřenského fondu.\\

% --------------------------------------------------------------------------
%Daňová úprava svěřenských fondů, by se dala nejspíše shrnout následovně.
% --------------------------------------------------------------------------

Problém dle mého názoru nastává při zániku svěřenského fondu, kdy není zákonně zcela jasně upraven postup. Respektive daňové zákony neobsahují žádná zvláštní ustanovení, která by se jeho zániku týkala. Ustanovení § 24 odst. 6 Daňového řádu říká následující: \textit{"Ustanovení týkající se právnických osob se použijí obdobně i na organizační složky státu nebo pobočky nebo jiné organizační složky obchodního závodu zahraniční osoby a na jiné jednotky, kterým zákon svěřuje výkon práv a povinností osob zúčastněných na správě daní."}\footfullcite[§ 24, odst. 6]{noauthor_zakon_dan_rad}. V případě zániku svěřenského fondu by se tak obdobně měla použít ustanovení týkající se zániku právnických osob.\\

Takovýto zánik ve vztahu k daňové povinnosti upravují paragrafy 240 až 240d. Problémem však je, že v případě zániku svěřenského fondu nelze hovořit o zániku s likvidací, ani o zániku bez likvidace s právním nástupcem. Lze tedy dovodit, že tyto ustanovení jsou na zánik svěřenského fondu neaplikovatelná. Dle § 20 odst. 3 DŘ, má svěřenský správce stejné povinnosti jako daňový subjekt\footfullcite[§ 20, odst. 3]{noauthor_zakon_dan_rad}. Vše tedy nasvědčuje, že daňové tvrzení by měl zpracovat a podat svěřenský správce i po samotném zániku za období do zániku svěřenského fondu\footnote{Obdobně v případě, že by byl svěřenský fond plátcem daně z přidané hodnoty, dojde jeho zrušením také ke zrušení jeho registrace jako plátce.}. V případě vzniku daňové povinnosti po dni zániku, byl by správce povinen zpracovat a podat dodatečné daňové tvrzení. Na takovouto povinnost by se zřejmě pohlíželo jako na daňovou povinnost vzniklou za jeho trvání.\\

Obdobně problém vzniká v případě, kdy za existence svěřenského fondu bylo plněno s daní, či bylo plnění od daně osvobozeno. V případě zániku fondu není na první pohled zřejmé, zda plnění v souvislosti se zánikem svěřenského fondu podléhá stejným pravidlům, jako plnění v průběhu existence svěřenského fondu.\\

Domnívám se, že pravidla, která se uplatňují na plnění ze svěřenského fondu v průběhu jeho existence se stejně uplatní i pro plnění v souvislosti se zánikem svěřenského fondu a vypořádáním jeho majetku\footnote{Ke stejnému závěru například dospěl i autor diplomové práce na téma zdanění svěřenských fondů, dostupné zde: \detokenize{https://is.muni.cz/th/caitw/Diplomova_prace-_pravni_a_danove_aspekty_sverenskeho_fondu.pdf}.}.\\

Dle mého názoru by však i tyto aspekty měly být v zákoně jasně vydefinovány a upraveny, v rámci rozhodovacího procesu zakladatele i toto může hrát značnou roli, jasná definice mi tedy v zákoně schází. Pomocí \textit{argumentum per analogiam} přece můžeme dovodit, že při zániku svěřenského fondu by se měla uplatňovat stejná pravidla, zákon by však měl být jasný i po samotné gramatické stránce, to zde však není splněno.\\


% --------------------------------------------------------------------------
 %a trápí mě neexistující úprava stanovující jakým způsobem má být zdaněna finální výplata ze svěřenského fondu, má to být stejně jako za existence svěřenského fondu, já se domnívám, že pokud aplikujeme § 24 daňového řádu, tak ano(protože v jistých případech výplata ze svěřenského fondu zdaněna není, v jiných případech zdaněna je), 
 
 %Dle mého názoru by bylo vhodné toto jasně a srozumitelně v zákoně definovat, neboť i tento nedostatek může přispět k nerozhodnosti, zda je vhodné svěřenský fond založit, či nikoli, ze strany zakladatele. S ohledem na to, že svěřenský fond je představován i jako nová forma umožňující mezigenerační převod majetku, tak tady to nelze jasně aplikovat, protože ve chvíli kdy by tato vůle směřovala například k tomu, že se svěřenský fond má po nějaké době zrušti, že je tedy založen jen na určitou dobu, může zakladatel při založení takovéhoto svěřenského fondu jasně předpokládat jak na následné převedení dopadne daňové zatížení? Já si myslím, že ne a to tedy ve svém konečném důsledku omezuje možnosti zakladatele plánovat jak má být s jeho majetkem po jeho smrti naloženo. Právě proto by tedy bylo dobré tyto nedostatky a nejasnosti napravit a objasnit.

%Mohlo by se namítat, že problémy se zdaňením převáděného majetku ze svěřenského fondu obmyšleným nebo zůstavitelům již nejsou problémem zakladatele, neboť ten již bude při založen ísvěřenského fondu mortis causa nebo inter vivos s odkládací podmínkou smrti po smrti, ale je vhodné si uvědomit, že tento problém dopadá i na svěřenské fondy založené inter vivos, kdy zakladatel může být stále na živu a majetek tak může připadnou po zániku svěřenského fondu jemu. Je tedy v zájmu zakladatele o tomto problému přemýšlet a je vhodné, aby toto bylo tedy zákonnem patřičně upraveno.\\

%Obdobně může zakladatel o tomto problému přemýšlet i ve vztahu k založení mortis causa či s odkládací podmínkou smrti ve vztahu k obmyšleným svěřenského fondu (tedy těm v jejichž prospěch se mezigenerační majetek převádí a zda daňové zatížení v takovémto případě nebude horší než v případě převedení majetku pomocí klasických instrumentů dědického práva, kdy se pokud se majetek převádí osobám v kategorii, které jsou osvobozeni od placení daně a paltí tak pouze notáři).

%Samotné daňové zatížení v rámci funkčního svěřenského fondu není horší než například pokud by se majetek předával pomocí založení právnické osoby, jak například stanovuje tady v té práci, ale při zániku mohou vznikat nejistoty a potenciální problémy a vyšší zdanění, než které by vzniklo například při předání majetku v klasickém dědickém řízení.

%Nakonec je problém i v likvidaci svěřenského fondu a tedy ochranou věřitelů (toto je obecný zájem, ne zájem zakladatele, obmyšlených, nebo svěřenského správce), protože likvidaci, ačkoli by mohla nějakým způsobem být vložena do statutu, jak říká například následující práce, tak standardně u svěřenských fondů neprobíhá.

%Tady asi ta ochrana věřitelů, nemožnost likvidace, a otázky spojené se zánikem svěřenského fondu.

%\subsubsection{Ochrana v Anglii}

%file:///Users/dominikbalint/Downloads/RPTX_2011_1__0_370222_0_120694.pdf

%\subsection{Další systémové nedostatky}

%Mimo výše uvedené nedostatky české právní úpravy, respektive potenciálními problémy svěřenských fondů spojené s úpravou svěřenských fondů, je minimálně vhodné zmínit i množinu dalších problémů, kterým můžou zůčastněné osoby v souvislosti se založením svěřenského fondu čelit. Hlavní nosnou částí je pak popis problému samotného rozdílného pojetí vlastnictví svěřenského fondu, respektive použití Lepaullova vlastnictví, které považuji za společného jmenovatele ve své podstatě veškerých, v této části popsaných, problémů. Popis tohoto problému je nadmíru důležitý s ohledem na návrh řešení, neboť ty spočívají především ve změně pojetí autonomního vlastnictví, nebo subjektivizace svěřenského fondu.\\

%Napsat že subjektivizace asi není to pravé ořechové, ve vztahu k implementaci tohoto institutu do právních systému civil law to ale je jedno z možných řešení, otázka pak samozřejmě je, jak budou takto omezované instituty (zde můžu vyjít z práce Michala Skuhrovce) konkurenceschopné na celosvětové scéně.
% --------------------------------------------------------------------------

\subsubsection{Možnosti použití svěřenského fondu k obcházení věřitelů při zániku správy}

Obavu nejenom o potenciální možnosti obcházení věřitelů zmiňovalo při přijetí trustu do české právní mnoho autorů\footnote{Srovnání Miloš Kocí, Institut svěřenského fondu v NOZ. Bulletin advokacie. 2014, 1-2.}. I toto tak vedlo zákonodárce k novelizaci právní úpravy svěřenských fondů a zavedení jejich registrace do rejstříku svěřenských fondů.\\

Obsahem této části není rozprava o potenciální možnosti založení svěřenského fondu pro znemožnění věřitelům uspokojit se na majetku zakladatele, přes negativní ohlasy v tomto směru považuji právní ochranu s ohledem na možnost dovolání se relativní neúčinnosti za dostačující\footnote{Viz paragrafy 589 OZ, 1467 OZ, či úprava v insolvenčním zákoně, § 235 InsZ a § 242 InsZ, či v trestněprávní rovině dokonce trestný čin poškození věřitele.}.\\

Obsahem je tedy pojednání o možných rizicích a namítání právě této neúčinnosti v případech zániku správy a svěřenského fondu jednotlivými věřiteli majetku fondu, neboť ten může být bezpochyby dlužníkem\footnote{Samotný proces zániku je pak popsán v kapitole 6.4 Zánik svěřenského fondu}.\\

% --------------------------------------------------------------------------
%Tady to přepíšu, že může bezpochyby dlužit.
% --------------------------------------------------------------------------

V případě založení svěřenského fondu za účelem transferu majetku jeho pomocí na zůstavitelovi dědice, nebo jiné osoby, kterým chce zůstavitel odkázat určitý majetek, toto dle mého názoru může představovat větší problém, o kterém by potenciální zakladatel měl rozhodně vědět.\\

V případě, že by měla zaniknout správa svěřenského fondu a následně i svěřenský fond samotný\footnote{Viz kapitola 6.4 Zánik svěřenského fondu.}, se domnívám, že obdobně, jako při zrušení fondu kvůli zániku majetku v něm vyčleněného, zánik takovéhoto účelově určeného majetku musí pak být úplný a nezhojitelný, tedy k majetkové podstatě například nesmí náležet ani korespondující pohledávka na úhradu újmy, či výplatu pojistného plnění\footfullcite{sro_httpwwwpraguebestcz_podrobneji_nodate-1}. V rámci veřejného zájmu je pak vhodné zajistit vypořádání všech pohledávek vůči svěřenskému fondu před jeho samotným zánikem.\\

Jak již bývá u svěřenského fondu zvykem, i zde lze nalézt mnoho protichůdných názorů týkajících se procesu vypořádání závazků při zániku svěřenského fondu. Jak již bylo popsáno u daňové problematiky, u svěřenského fondu se při zániku neprovádí likvidace, ani obdobný proces, respektive zákon ji nepředpokládá, i přes to však převládá právní názor, že správce musí před vydáním majetku dle § 1472 OZ provést vypořádání stále trvajících dluhů, které vznikly při správě svěřenského fondu\footfullcite[Komentář k § 1472]{spacil_j_a_kolektiv_obcansky_2013}.\\

Osobně dovozuji, že svěřenský správce by mohl provést likvidaci, respektive proces odpovídající likvidaci, v případě, že by byl zakotven ve statutu svěřenského fondu. Ten by dle mého mohl stanovit, jakým způsobem má být vypořádáno zbylé jmění ve svěřenském fondu, mohl by tedy stanovit povinnost provést likvidaci, vypořádání dluhů a následné převedení likvidačního zůstatku na obmyšlené. Určení kdy by se s takovouto likvidací mělo začít však může být velice obtížné. Při zániku svěřenského fondu uplynutím doby, na kterou byl založen by tak mohl správce jistou formu likvidace provést, v případě zániku svěřenského fondu vzdáním se práva na plnění ze strany všech obmyšlených, nebo naplněním účelu by však automaticky došlo k zániku správy, svěřenský správce by tak pro likvidaci prostor neměl. Možnost řešení takovéhoto případu by opět mohl představovat statut určení dodatečného účelu svěřenského fondu, kterým by byla jeho likvidace, kumulativním splněním obou účelů by tak mohlo dojít k zániku správy a konsekventně tedy i k zániku svěřenského fondu.\\

% --------------------------------------------------------------------------
%Doplnit tuto věc do smlouvy, tedy nějaké přání zakladatele, že se mají vypořádat před skončením správy dluhy svěřenského fondu.
% --------------------------------------------------------------------------

Avšak bez ohledu na výše zmíněné, tedy i v případech, kdyby statut s likvidací nepočítal, se mnoho autorů\footnote{Srovnání Pihera v Komentáři k Občanskému zákoníku, svazek III, či Lucie Josková, Správa cizího majetku, 4 fáze zániku svěřenského fondu.} přiklání k teorii, že vypořádání dluhů zkrátka musí proběhnout. S ohledem na § 1472, který stanoví, že se vydává majetek, nikoli jmění, je tento názor zcela logický. Samotná likvidace\footnote{Účelem likvidace jest vypořádat majetek zrušené společnosti, vyrovnat dluhy věřitelům a naložit s tzv. likvidačním zůstatkem v souladu se zákonem.} se pak principiálně jeví jako nejvhodnější způsob vypořádání dluhů svěřenského fondu. Až po skončení likvidace by majetek má být zbývající majetek, respektive likvidační zůstatek, dovoleno bezformálního a bezúplatného převodu na obmyšlené(či jiné osoby statutem určené).\\

Právní teorie však již opomíjí prostředky obrany věřitele svěřenského fondu pro případy, že by svěřenský správce proces obdobný likvidaci neprovedl. Toto představuje problém nejenom z hlediska veřejného pořádku, ale i z pohledu zakladatele, obmyšlených a jiných statutem určených osob.\\

Tato obrana je o to důležitější v souvislosti s tím, že osoba, na kterou přechází majetek ze svěřenského fondu neodpovídá a neručí za dluhy náležející k majetkové podstatě svěřenského fondu jak dovozuje Pihera: \textit{"Osoba, jíž se podle komentovaného ustanovení vydává majetek svěřenského fondu, nepřebírá odpovědnost za dluhy, které vznikly při správě svěřenského fondu."}\footfullcite[Komentář k § 1472]{spacil_j_a_kolektiv_obcansky_2013}. Otázkou však je, zda i přes autory dovozovanou nutnost vypořádání a nemožnosti převzetí dluhů náležející k majetku, by v určitých případech nemohlo k tomuto převodu dluhů skutečně dojít\footnote{Například při převodu jednotky - § 1186 odst. 2 OZ, nebo družstevního podílu - § 736 odst. 2 ZOK, v případě nevypořádání těchto pohledávek by pak bylo možné uvažovat o odpovědnosti správce za škodu způsobenou jejich nevypořádáním.}. Tuto otázku však vyřeší až souvislá praxe se svěřenskými fondy spojená.\\

Ochrany se tak věřitel bude moci dovolat zřejmě opět ve vztahu k § 1466 OZ, tzn. pomocí namítnutí relativní neplatnosti konání správce v případě\footnote{Pohledem Pihery nikoli relativní neúčinnosti, v komentáři k § 1466 OZ. Relativní neplatnost tedy v tomto případě zakládá neplatnost jednání svěřenského správce vůči všem, oproti relativní neúčinnosti, kdy působí pouze vůči tomu, kdo se jí dovolává. Porušení pravidel upravující vydání majetku oprávněným osobám zakládá protiprávnost a v tomto smyslu tedy neplatnost.}, že by dané jednání převedení majetku svěřenského fondu oprávněné osobě věřitele poškozovalo na jeho právech. V případě úspěchu, by tak majetek majetkovou prostatu svěřenského fondu nikdy neopustil. Věřitel by se tak ní mohl nadále uspokojovat.\\

Z pohledu nedostatečné úpravy tak nezbývá než jen doporučit, aby zakladatel vymezil ve statutu nějakou formu likvidace, aby nemohlo dojít k potenciálnímu, ačkoli nepravděpodobnému, převzetí dluhů osobou oprávněnou k vydání majetku ze svěřenského fondu, či problematického procesu zániku svěřenského fondu, kdy by vydání takovéhoto majetku mohlo být značně ztíženo jednáním věřitelů svěřenského fondu.\\

Samotná absence výslovného ustanovení upravujícího vypořádání dluhů svěřenského fondu před zánikem správy a následným zánikem svěřenského fondu a možností obrany před situací, kdy by svěřenský správce vypořádání neprovedl, je dle mého názoru také znepokojující.\\

% --------------------------------------------------------------------------
%Existuje však i názor, že

% takovémto případě by mohl předsta


%Autoři se však shodují na tom, že 
% --------------------------------------------------------------------------

\subsection{Rozdílné pojetí vlastnictví}

Další otázku, kterou si mnozí právní experti položili je samotné pojetí vlastnictví v systému civil law, které je dle mnohých jen těžko kompatibilní s děleným vlastnictvím, které je inherentně spojeno se systémy common law.\\

% --------------------------------------------------------------------------
%které je vlastní právním systémům common law.
% -------------------------------------------------------------------------- 

Jedná se o velice dobrou otázku, neboť na svěřenský fond by se dalo nahlížet jako na dělené vlastnictví. Tento fakt byl dlouhou dobu považován za překážku vstupu trustu do kontinentálních právních řádů, neboť kontinentální právní řády jsou založené na konceptu nedělitelnosti vlastnictví\footfullcite{lucie_joskova_sprava_2017}, zejména od doby, kdy se na Evropském kontinentu opustil koncept feudálních vztahů\footfullcite{michael_milo_trusts_2000}.\\

Pojem dělené vlastnictví by se dal jednoduše popsat na příkladu reálných vlastnických vztahů několika rodin, bydlící ve vícegeneračním domu. V takovémto případě by tyto rodiny fungovaly na principu děleného vlastnictví, dalo by se říci tedy i na principu rodinného fideikomisu, kdy o nemovitosti, jako o větší majetkové hodnotě, nerozhoduje pouze formální vlastník, který je zapsán v katastru nemovitostí, ale rodina jako celek\footfullcite{ondrej_horak_vd_2010}.\\

Samotná Francouzská úprava, z níž z historického hlediska vychází i právo Kanadské provincie Quebec, nepoužívá dělené vlastnictví ale nahlíží na něj jako na obdobu českého jmění. České právo tento koncept obdobně nevyužívá a samotný občanský zákoník hovoří jen o majetku, jmění \textit{§ 495 Souhrn všeho, co osobě patří, tvoří její majetek. Jmění osoby tvoří souhrn jejího majetku a jejích dluhů.} a vlastnictví \textit{§ 1011 Vše, co někomu patří, všechny jeho věci hmotné i nehmotné, je jeho vlastnictvím.}, kde zásadnější rozdíl představuje pouze majetek a jmění, přičemž majetek a vlastnictví nemají mezi sebou téměř žádný rozdíl.\\

Není však absolutní pravdou, že by bylo dělené vlastnictví kontinentálním řádům cizí úplně, například Rakouský občanský zákoník ABGB obsahoval koncept děleného vlastnictví\footnote{Na rozdíl však třeba od Francouzského Code civil, které naopak instituty založené na děleném vlastnictví až do Napoleonovi prohry rušilo, po konci Napoleonovi vlády se však i ve Francii přistoupilo o obnovení těchto institutů. Ve větším množství se například fideikomisy začaly rušit až po první světové válce například ve Výmarské ústavě, nebo v Československu po vyhlášení samostatného státu}, kdy k jednomu majetku existovalo právo vrchní a právo užitkové, při hlubším zamyšlení lze dojít k tomu, že právě takovýmto způsobem je postaven trust, kdy právo vrchní, tedy formální, zapsané, má svěřenský správce, ale právo užitkové má obmyšlený\footnote{Tamtéž.}. V CCQ byl tento problém vyřešen poměrně netradičně, a to takovým způsobem, že vyčleněný majetek není nikým vlastněn, ale že se jedná o autonomní vlastnictví, tedy respektive majetek svěřen vlastnímu účelu\footfullcite[POPOVICI A. Trust québeckého a českého práva:autonomní vlastnictví?]{lubos_tichy_sverensky_2017}. Tento koncept je založen na \textit{patrimoine d'affectation}, v CCQ patrimony by appropriation, v občanském zákoníku účelově určený majetek\footnote{Tento definuje Miloš Kocí v komentáři k Občanskému zákoníku od Wolters Kluwer jako "majetek bez vlastníka vyčleněný zakladatelem k naplňování určitého účelu a spravovaný svěřenským správcem podle instrukcí vtělených zakladatelem do zakladatelského právního jednání, resp. do zvláštního statutu"}, od Francouzského právníka Pierra Lepaullea\footfullcite{skuhravy_jan_trust_2010}.\\ 

V rámci tohoto pojetí vlastnictví spočívá Québecký trust na obligačních vztazích mezi správcem a beneficientem, nikoli na principu odděleného vlastnictví. Podle CCQ, které se v tomto přímo inspirovalo Lepaullem, nenáleží ani zakladateli, ani správci a ani obmyšlenému\footfullcite[Článek 1261]{noauthor__nodate}, ale existuje samo o sobě společně se závazky, které k němu vznikly v souvislosti s plněním účelu trustu\footfullcite{skuhravy_jan_trust_2010}.\\

Trust je tedy v CCQ vnímán jako samostatná a oddělená samosprávná quasi entita, která obsahuje jak majetek, tak závazky s ním spojené\footfullcite[Tomáš Richter: Mezi smlouvou, vlastnictvím a korporací: právní úprava trustu v návrhu nového občanského zákoníku]{stenglova_ivana_pocta_2006}, tedy obsahuje vlastní jmění vyhrazené svému vlastnímu účelu. Tam, kde majetek samozřejmě nemá vlastníka je vhodné zamyšlení nad jeho samotným odlišením od právnické osoby\footfullcite[Kenneth G.C. Reid: Patrimony not Equity: the Trust in Scotland]{michael_milo_trusts_2000}\textsuperscript{,}\footnote{Srovnání s kapitolou Svěřenský fond a právní subjektivita}, přes fakt, že zákonodárce nezavedl pro trust klasifikaci právnické osoby, tak trust tímto způsobem funguje\footfullcite[Madeleine Cantin-Cumyn: The Quebec Trust]{michael_milo_trusts_2000}, neboť je sám vlastníkem majetku a závazků v něm obsažených. Z formálního hlediska je však správce zapsán v registrech a vykonává správu nad majetkem ve svěřenském fondu\footfullcite[Článek 1260]{noauthor__nodate}\textsuperscript{,}\footfullcite{skuhravy_jan_trust_2010}.\\

Takto byl do jisté míry vyřešen problém děleného vlastnictví a mnoho tedy již, i přes existující teorie, které problém děleného vlastnictví vyvrací\footfullcite{lucie_joskova_sprava_2017}, nebránilo tomu, aby mohl být tento institut přijat do našeho kontinentálního právního řádu s obdobným pohledem na majetek takto do něj vložený.\\

Takovéto řešení však není bezchybné a není na podivu, že při transplantaci quebeckého trustu, který doktrinálně vychází z anglického trustu mohou vzniknout jisté výkladové problémy. Tyto zmiňuji především z důvodů, že v rámci rozhodovacího procesu o založení svěřenského fondu pro jakýkoliv účel by si měl být zakladatel těchto problémů, ačkoliv jen potenciálních, být vědom.\\

Dle Bolgár je zásadní problém princip \textit{numerus clausus} v kontinentálním právu\footfullcite{vera_bolgar_why_1953}. Tento princip vymezuje rozdíly mezi věcnými (\textit{iura in rem})právy a závazkovými právy (\textit{iura in personam}), z práva závazkového tak vyčleňuje okruh věcných práv. Podstata tohoto principu tedy spočívá v taxativním vymezení věcných práv a omezením smluvní dispozice v rámci závazkového práva. Toto faktické omezení možnosti vytváření nových věcných práv z práv závazkových, respektive ze smluvních ujednání. Toto omezení je však neslučitelné s institutem trustu, který umožňuje v zásadě neomezené smluvní nakládání s majetkem, které se vytváří na základě právních vztahů konstituentů utvářených výhradně dle jejich vůle\footfullcite{noauthor_sverensky_nodate-2}.\\

Z této logiky se tak přejímání institutu trustu, historicky vytvářeného v rámci rozdílného právního systému, a jeho transponování do národních řádů zemí používajících systém práva civil law jeví jako problematické. Mnohé národní státy tak přistoupily k vytváření analogických konstrukcí stojících na principech trustu, které neodporují principům systému civil law. Toto ostatně učinil zákonodárce v Quebecu, kdy přijal do svého právního řádu trust ale koncepci vlastnictví majetku v trustu vystavěl na Lepaullově vlastnictví. Takovéto pojetí se pak s transplantací Quebecké právní úpravy trustu dostalo i do českého Občanského zákoníku\footnote{Tamtéž.}.\\

Samotný vývoj trustu je pak spjat s dualitou vlastnictví založené na common law a equity. Kdy termín dualita vlastnictví má pod sebou podřazené právní schéma rozdělení vlastnických práv k majetku, kdy jednomu plynou práva z common law a druhému z equity. Zásadní dopad duality vlastnictví pak na trust dopadá v tom smyslu, že svěřenský správce má vlastnické právo z titulu common law a obmyšlený z titulu equity. Tato dualita vlastnictví ale není kompatibilní s principem autonomie a nedělitelnosti vlastnictví common law. \\

Filip Čeladník tak ve svém článku například vyslovil otázku: \textit{"Určitou analogii trustu lze sice přísně vzato nalézt ve feudální době v institutu fideikomis. Zde se však nabízí otázka, zda se zákonodárce ve snaze pokračovat v tradici kontinentálního práva neměl inspirovat spíše tímto institutem majícím návaznost na římské právo, než tak specifickým produktem anglické jurisprudence. Zmíněná civil law tradice nepřipouštějící fragmentaci vlastnických práv má pak ve vztahu k vlastníku věci vyjádření v tom, že ten je absolutním a výhradním vlastníkem celé věci, popř. podílu na věci v případě podílového spoluvlastnictví."}\footnote{Tamtéž.}. Já osobně věřím, že i přes problémy spojené s přijímáním institutu doktrinálně vycházejícího z common law trustu český zákonodárce učinil dobře, viz kapitola 4.2 Historie svěřenského fondu u nás a ve světě, kde jsem vyložil argumentaci proč přijmout trust a neinspirovat se místo toho například v německém právním řádu. To však nezbavuje svěřenský fond problémů s tímto přejetím spojených, zejména tedy ve spojením s dualitou vlastnictví, která bezesporu představovala problém přejímání trustu do civilistických jurisdikcí\footfullcite[Strana 328]{jaroslav_svejkovsky_sprava_2015} ačkoli v uvozovkách vyřešených aplikací Lepaullova vlastnictví.\\

Zatímco v rámci anglického pojetí trustu je majetek jak ve vlastnictví správce fondu tak i obmyšleného, v českém a quebeckém právu není ve vlastnictví nikoho\footnote{Rozhodně se však nedá hovořit o \textit{res nullius} dle paragrafu 1045 OZ. Podle ustanovení paragrafu 1456 bude totiž vlastnická práva k majetku vykonávat svěřenský správce, který bude majetek spravovat a ten i bude zapsán v příslušné evidenci.}. Zatímco však v Quebecu zákonodárce stanovil, že k danému majetku nemá nikdo ani vlastnická, ani věcná práva, český však stanoví pouze, že k takovémuto vyčleněnému majetku nemá nikdo vlastnická práva, o věcných právech však mlčí. Na český svěřenský fond by tedy také šlo pohlížet tak, že věcná práva k takto vyčleněnému majetku existují a za použití vyvratitelné právní domněnky uvedené v paragrafu 1400 odstavce 2, by je bylo možné nalézt u svěřenského správce, ten by tak mohl mít neomezené právo dispozice s majetkem \textit{erga omnes}\footnote{Tamtéž strana 328}. Zde by tedy mohla být nalezena věcnost práv správce projevená například v možnosti podal reivindikační žalobu.\\

Oním problémem ve vztahu přijetí institutu založeného na dualitě vlastnictví je tedy spočívá v samotném konání správce a možností bránit se proti takovémuto konání. Nemohu než jen důrazně poznamenat, že správný výběr svěřenského správce je nad míru důležitý, špatný výběr totiž může vést nejen k neúmyslnému ohrožování majetkové hodnoty vložené do fondu, ale také k úmyslnému jednání svěřenský fond, nebo práva obmyšlených poškozující.\\

Příkladem budiž situace kdy svěřenský správce z fondu založeného za soukromým účelem převede majetek na třetí osobu a to v rozporu povinnostmi a právy uložené mu s ohledem na správu svěřenského fondu, tedy například s majetkem nakládat či jakkoliv ho zcizovat. V takovémto případě svědčí obmyšlené osobě v anglickém právu postih vůči správci plynoucího z titulu závazkového, tedy například žaloba proti správci na náhradu škody tak i z titulu věcného práva právě kvůli dualitě vlastnictví. Z tohoto titulu má obmyšlený možnost žalovat i takovouto třetí osobu, například žalobou na vydání pokud osoba nebyla v dobré víře\footfullcite{noauthor_sverensky_nodate-2}.\\

Tato konstrukce anglosaského trustu poskytuje obmyšlenému větší ochranu než, kterou má s ohledem na českou právní úpravu. V rámci českého práva totiž takto poškozenému svědčí pouze deliktní postih vůči správci svěřenského fondu v rámci závazkového práva. Nebude se však možno domáhat vydání majetku po třetí osobě, které byl majetek zcizením převeden. Toto nebude možné patrně ani v případě, že třetí osoba nejednala v dobré víře, neboť v české konstrukci nesvědčí obmyšlenému, či zakladateli žádná věcná, ani vlastnická, práva k majetku vyčleněnému do svěřenského fondu. Ta totiž poskytuje pouze deliktní ochranu.\\

Jak autor článku také poznamenává: \textit{"Posledně zmíněné odráží princip, že jednou vyčleněný majetek si nelze po vzniku svěřenského fondu přisvojit zpět ani zrušit svěřenský fond, jestliže taková možnost nebyla ve statutu stanovena; naproti tomu v závazkovém právním vztahu lze po uzavření smlouvy její ustanovení měnit i smlouvu vypovědět popř. zrušit. (...) Pro úplnost je třeba dodat, že osoba oprávněná ze svěřeneckého i zakladatel mohou podle § 1466 NOZ rovněž vznést námitku neplatnosti právního jednání správce, kterým poškozuje svěřenecký fond; tento nárok je však limitován absencí dobré víry třetí osoby, což je ovšem krajně nedostatečné, jak bude podrobněji uvedeno níže."}\footnote{Tamtéž.}.\\

% --------------------------------------------------------------------------
% mít vůči majetku věcná práva \textit{erga omnes}

%I přes tuto zjevnou nekompatibilitu si však stále 

%https://www.epravo.cz/top/clanky/sverensky-fond-jako-vysledek-ceskeho-pokusu-o-pravni-transplantaci-trustu-zklamani-jako-dite-ocekavani-vybrana-zakonna-ustanoveni-z-pohledu-zahranicniho-pravnika-93493.html
% --------------------------------------------------------------------------

Možná i nejen proto Lepaulle později svoji teorii pozměnil a místo "depersonalizace" zvolil postup subjektivizace a potřebě na trust pohlížet jako na právnickou osobu\footfullcite{pierre_lepaulle_notion_1950}.\\

% --------------------------------------------------------------------------
%Doplnit zdroj Luboš Tichý.
% --------------------------------------------------------------------------

Obdobně jako Luboš Tichý proto považuji pojetí svěřenského fondu jako bezsubjektové vlastnictví, jako spolčeného jmenovatele velké spousty problému spojených se svěřenským fondem, problematiku nepominutelného dědice nevyjímaje. Návrhy na vyřešení těchto problémů v závěrečné části práce, tak do značné míry souvisí právě s pojetím bezsubjektového vlastnictví.\\

% --------------------------------------------------------------------------
%PŘEPSAT
%Lepaulle, ve snaze přiblížit trust kontinentální právní kultuře a jejím právníkům, šel "až na dřeň" a vymezil základní prvek trustu. Argumentuje, že po založení trustu není takový vůbec závislý na svém zakladateli, ani na správci, který může být vyměněn (srov. existence "bare" trustu, viz. kap. 3.4.). Obmyšleného pak jako esenciální prvek odmítá úplně. "Jedině, jak se dá k této otázce přistoupit, je hledat ty prvky, které jsou pro existenci trustu nezbytné", říká, a definuje trust jako "právní institut sestávající se z majetku nezávislého na právním subjektu, jehož jednotnost je definována učelovým určením v mezích zákona a veřejného pořádku". POPOVICI ALEXANDRA Nutným závěrem je pak vznik koncepce "patrimoine d`affectation", nebo též "patrimony by appropriation", konečně pak v OZ "účelově určený majetek", nebo též jak rozvádí Kocí v komentáři k OZ "majetek bez vlastníka vyčleněný zakladatelem k naplňování určitého účelu a spravovaný svěřenským správcem podle instrukcí vtělených zakladatelem do zakladatelského právního jednání, resp. do zvláštního statutu". KOCÍ KOMENTÁŘ 

%Doplnit reivindikační žalobu, dle RPTX_2011_1__0_370222_0_120694 staženého souboru na stránce 89 - HOTOVO

%Doplnit, že si myslím, že všechny problémy jsou zapříčiněny právě rozdílným pojetím vlastnictví.

%\subsubsection{Ležící pozůstalost}

%Za povšimnutí také stojí ležící pozůstalost, kterou v souvislosti s pojetí vlastnictví a svěřenskými fondy některé texty věnující se této problematice zmiňují. Zde bych chtěl pouze dodat, že ta zde není v zásadě aplikovatelná. 

%Zdroje 2 Pihera Nejpodivnější zvíře v lese a Luboš Tichý článek na advokátním děníku.

%\subsection{Zákonné nejasnosti změn svěřenských fondů před a za jejich trvání}

%https://www.ilaw.cas.cz/casopisy-a-knihy/casopisy/casopis-pravnik/hledat-v-archivu/detail-clanku.html?id=40182&r=%252Fcasopisy-a-knihy%252Fcasopisy%252Fcasopis-pravnik%252Fhledat-v-archivu.html%253Fnaki_search%253D1%2526form_state%253D%2526query%253DFunda%2525C4%25258Dn%2525C3%2525AD%252Bprincipy%252Ba%252Bhranice%252B%2526year_start%253D%2526year_end%253D%2526volume%253D%2526issue%253D%2526search_article_title%253D1%2526search_article_annotation%253D1%2526search_article_content%253D1%2526search_author%253D1%2526search_keywords%253D1%2526search_match_words%253Dall

%\subsection{Nezávislost správce}

%Kontrola správce, dohled nad správcem, odpovědnost správce, potenciální problémy, možná dopsat pokud zbyde čas a prostor.

%Diskrece posiluje korupční motivaci správce.

%Tady se jedná o něco co je ve veřejném zájmu.

%Kro samotné kontroly správce je ještě vhodné uvažovat o oprávnění zakladatele ve vztahu k činnosti správce. Jak poznamenává Lucie Josková, takováto kontrola je zřejmě možná i preventivně, je tedy možné vymezit určitý okruh právních jednání, s nimiž musí vyslovit zakladatel souhlas\footfullcite[Postavení svěřenského správce, strana 180]{trust_2014}. 

%Najít na internetu.

%Doplnit odkaz na smlouvu a doplnit ve smlouvě nějakou část dbající na kontrolu svěřenského fondu, říct že samozřejmě ne vždy je to možné, například pokud není nikdo, kdo by tuto kontrolu vykonával.

%Častečně už jsem popsal v předcházející kapitole o rozdílném pojetí vlastnictví. Je hrozně důležité kontrolovat a být si jistý osobu správce.
% --------------------------------------------------------------------------

\subsection{Návrhy na změny České právní úpravy}

Jako nejzásadnější otázku, která je zároveň hypotézou, považuji nedostatky právní úpravy zaměřující se na řešení otázek vznikajících při kolizi dědického práva a práva svěřenských fondů, blíže se tedy budu věnovat návrhu řešení tohoto problému. V samotném závěru této části práce také uvedu svůj pohled na ostatní mnou výše popsané nedostatky, které spatřuji v právní úpravě svěřenských fondů a sním souvisejících institutů.\\

Z jednotlivých ustanovení dědického práva a svěřenského fondu je nepochybná vůle zákonodárce chránit na jedné straně nepominutelné dědice a na straně druhé zřídit fiduciární institut, který umožní zůstavitelovi téměř neomezeně nakládat se svým majetkem. Záleží samozřejmě na optice, kterou pohlédneme na jednotlivá ustanovení svěřenského fondu a pořízení pro případ smrti, úmyslně zmiňuji jen pořízení pro případ smrti a s tím související založení svěřenského fondu pro případ smrti, neboť ostatní formy založení nejsou a ani nemohou být omezovány dědickým právem, je však nezpochybnitelné, že zavedení neomezeného fiduciárního institutu je neslučitelné s právem nepominutelných dědiců. Způsobů k vyřešení tohoto nedostatku právní úpravy se nabízí hned několik.\\

\subsubsection{Zrušení svěřenského fondu}

První možností jak řešit nedostatky spojené se svěřenským fondem je v zásadě legislativně, nikoli však politicky, nejjednodušší. Zrušení svěřenského fondu by dozajista vyřešilo všechny problémy s ním spojené, zapříčinilo by však mnoho jiných problémů, které by se objevili při zániku jednotlivých svěřenských fondů.\\

S návrhem na zrušení svěřenského fondu tak například přišla Pirátská strana, jedná se o návrh již z roku 2014, jednotlivé problematické body v něm zmíněné tak byly odstraněny novelou Občanského zákoníku účinné od 1.1.2018. Při snaze o objektivní posuzování dalších jednotlivých bodů se však nelze ubránit pocitu, že jsou věcné a mají, minimálně částečně, pravdu. S ohledem na ochranu věřitelů se dá souhlasit s bodem učiněným v závěru: "Do té doby lze očekávat obrovské problémy s vyhýbáním se daňové povinnosti, zneužíváním při zastření vlastnictví a absenci právní jistoty při obchodním styku s účastníky svěřenského fondu."\footfullcite{noauthor_stanovisko_nodate}. Z této citace je dle mého názoru aktuální její poslední část a to nejistota při obchodním styku s účastníky svěřenského fondu, zejména tedy se svěřenským správcem. Tato nejistota působí zejména ve chvíli zániku svěřenského fondu kvůli absenci likvidace svěřenského fondu. Uspokojení věřitelů tak může být problematické, viz kapitola 7.7 Možnost použití svěřenského fondu k obcházení věřitelů při zániku správy, ale do jisté míry může působit i při jeho vzniku, například při zvýšení majetku ve svěřenském fondu, přestože zákon předpokládá vložení dokladu o zvýšení majetku svěřenského fondu zřízeného za soukromým účelem, vzhledem k nízkým požadavkům na formu takovéhoto dokladu mu však nelze připisovat vysokou průkaznost\footfullcite{j_hollman_co_2016}, či ve vztahu k osobě blízké a neúčinnosti právních jednání v insolvenčním zákoně.\\ 

Nemyslím si však, že takovýto krok je vhodný, už jen kvůli tomu, že institut svěřenského fondu byl přijat relativně nedávno a jeho odstranění by nepřispělo k již tak, dle mého názoru relativně nestabilní, právní jistotě občanů Českého státu.\\

% --------------------------------------------------------------------------
%Napsat i o politických problémech, které by se tímto vyřešili?
% --------------------------------------------------------------------------

\subsubsection{Zrušení nebo omezení institutu nepominutelného dědice}

Další možností je pak úplné zrušení institutu nepominutelného dědice. Toto řešení by dozajista vyřešilo kolizní problematiku spojenou se svěřenským fondem, ale nepovažuji ho za zcela šťastné. Nejen proto, že institut nepominutelného dědice je historicky spjat s naším právním řádem, ale i kvůli jisté společenské potřebnosti tohoto institutu, kdy není dle mého názoru zcela vhodné, aby zůstavitelovi děti zůstaly po jeho smrti potenciálně zcela bez prostředků\footnote{Čistě prakticky na toto lze pohlédnout i optikou výživného mezi rodiči a dětmi a dědictví jako potenciálního pokračovatele této povinnosti.}.\\

Instituty obdobné nepominutelnému dědice se nacházejí ve většině právních řádech zemí jež stojí na systému civil law a také v mnoha islámských zemích. Dokonce i v oblastech, kde dochází k velkému promíchávání právních systémů civil law a common law, jako je například kanadská provincie Quebec, nebo americký stát Louisiana, existuje buď tento institut, nebo při nejmenším instituty jemu podobné.\\

Takovouto deviaci od normy bych proto nepovažoval za šťastnou. Možnost by představovalo samozřejmě omezení práv nepominutelných dědiců a přiznání ochrany například pouze v obdobných případech, v jakých je přiznává Občanský zákoník provincie Quebec.

% --------------------------------------------------------------------------
%Se zachováním určité ochrany, inspirovat se můžeme z Kanadské právní úpravy.
% --------------------------------------------------------------------------

\subsubsection{Změna institutu nepominutelného dědice, absolutní pořizovací vůle zůstavitele ve vztahu ke svěřenskému fondu}

Svěřenský fond by v tomto případě mohl být určen jako možný způsob vydědění, popřípadě by se na svěřenský fond založený \textit{mortis causa} nepohlíželo vůbec optikou dědického práva, jehož ustanovením by tedy za žádných okolností nepodléhal, a ani pořízení (v tomto případě by se muselo jednat pouze o smlouvu mimo zákonem stanovená pořízení) pro případ smrti by tedy v tomto ohledu nemuselo splňovat formální nároky na něj kladené, bylo by ale samozřejmě potřeba toto v zákoně jasně uvést a nedat prostor pro diskuzi.\\

V případě určení svěřenského fondu jako dalšího "důvodu" vydědění je pak otázkou proč by v rámci transferu majetku za pomocí instrumentů dědického práva měly být taxativně určeny důvody vydědění, pro které a pouze pro které zůstavitel může dědice vydědit, přičemž splnění těchto důvodů by se neaplikovalo na formu vydědění pomocí svěřenského fondu. Takto kladené nároky by dle mého názoru byly vůči sobě značně disproporční bez jakékoliv důvodové opory v zákoně či dokonce v samotné logice. Ochrana pro nepominutelné dědice by se v případě zvolení této cesty nevztahovala na případy, kdy je mezigeneračního převodu majetku činěno za pomoci svěřenského fondu. Troufám si tvrdit, že tomuto je současná právní úprava nejblíže, pokud by to však takto zamýšlel i zákonodárce, je stále třeba výslovného ustanovení toto upravující, aby se zabránilo disputacím a nejasnostem.\\

% --------------------------------------------------------------------------
%oddělení institutu nepominutelného dědice od institutu svěřenského fondu, ochrana by se tedy nevztahovala na případy, kdy je mezigeneračního převodu majetku docíleno pomocí svěřenských fondů. 

%Tomuto je, dle mého názoru, současná právní úprava nejblíže, je ale nicméně stále potřeba výslovné úpravy.
% --------------------------------------------------------------------------

\subsubsection{Změna institutu svěřenského fondu, ochrana nepominutelných dědiců}

% --------------------------------------------------------------------------
%Součást pozůstalosti, nebo ochrana nepominutelných dědiců již v rámci pořízení.

%Návrhy jsem vymezil v rámci obrázku na ploše obrazovky.

%Znovu zmínit tento článek: https://advokatnidenik.cz/2020/04/04/vlastnictvi-bez-vlastnika-v-ceskem-sverenskem-fondu-skica-na-margo-trustu-v-ceskem-pravu/

%Jako příklad by mohly sloužit jednotlivé ochrany nároků již existující v Quebeckém právu, viz kapitola Jiné obdobné zahraniční nároky připomínající.
% --------------------------------------------------------------------------

Otázkou je, zda je tedy konstrukce majetku vyčleněného za určitým účelem funkční, respektive do jaké míry je funkční v systému civil law. Pihera ostatně označuje svěřenský fond jako "nejpodivnější zvíře v lese", kterým beze sporu je, avšak taktéž dodává, že myšlenka takto vymezeného vlastnictví je ve své podstatě docela funkční\footfullcite{pihera_v_nejpodivnejsi_2012}. Já musím uznat, že funkční je, to že je ale něco funkční neznamená, že je to nejlepší možné řešení, problematika zavedení tohoto pojetí je zřejmá právě například ve vztahu k dědickému právo, ostatně je vidět že i sám zákonodárce má tendenci doktrinálně přibližovat svěřenský fond k právnické osobě\footfullcite[HAVEL B. Právní osobnost svěřenského fondu ve světle českého občanského zákoníku, strana 203]{trust_2014}. Zda by toto mohlo být dobré řešení rozebírá například Luboš Tichý ve svém článku\footfullcite{noauthor_vlastnictvi_2020}.\\

První řešením by mohlo být \textbf{stanovení svěřenského správce, jakožto nositele subjektivních vlastnických práv}\footnote{Přičemž osobní majetek správce by byl stále oddělen. Obdobně jako například v případě společného jmění manželů.}\textsuperscript{,}\footnote{Takto učinil například Supreme Court of Canada v případu Curran v. Davis.}, který by však ve výkonu vlastnického práva byl omezen statutem, tato omezení by však neměla účiny vůči třetím osobám\footnote{Toto řešení považoval za jednu z možností i Lepaulle, srovnání: An Outsider’s View Point of the Nature of Trusts, dostupné z: \detokenize{https://scholarship.law.cornell.edu/cgi/viewcontent.cgi?article=1276&context=clr}, strana 55.}. V případě vydání se touto cestou by se tak například vyřešil problém povinného dílu, neboť za dědice by se zcela jasně dal označit svěřenský správce. Svěřenský správce by měl rovněž věcná práva, přičemž vztah obmyšleného by byl založen na obligačněprávním nároku\footfullcite{noauthor_vlastnictvi_2020}.\\

Tento způsob ve formě faktického vlastnictví správce \textit{sui generis} byl i díky reálné subrogaci\footfullcite[Strana 71]{claxton_john_b_studies_2005} recipován v \textit{mixed legal system} v JAR, či Skotsku.\\

Druhým možným řešení by dle Tichého mohlo být \textbf{přizpůsobení vztahům uvnitř svěřenského fondu poměrům v trustu}. Pokud se však ohlédneme zpět za čistě \textit{common law} trustem, vyjevuje se problematika \textit{equitable ownership}, tedy duálního vlastnictví, tento způsob fungování trustu sám o sobě nelze přijmout neboť v českém právní řádu neexistuje nic jako \textit{equity}, možnost poskytovat věcněprávní účinky založené ne \textit{equity} tak systémy \textit{civil law} zkrátka nemohou. Je tedy logické, že kontinentální právní řády nemohou přijmout trust v té nejčistší \textit{common law} podobě. V cestě pojetí svěřenského trustu formou duálního vlastnictví stojí v cestě rovněž princip \textit{numerus clausus}\footnote{Viz kapitola 7.6 Rozdílné pojetí vlastnictví.}, tato koncepce byla rovněž neslučitelná s principem publicity, avšak tuto problematiku vyřešil zákonodárce povinností majetek v trustu registrovat v registrech s odpovídající poznámkou k němu zapsanou\footfullcite[Strana 136]{skuhravy_jan_trust_2010}. Dle mnohých autorů je ale odmítání přijetí této formy spíše dogmatického charakteru a v reálu mnoho nebrání tomu, aby takto přijat byl\footfullcite{Tamtéž, strana 128.}\textsuperscript{,}\footfullcite[Strana 206]{vera_bolgar_why_1953}\textsuperscript{,}\footnote{Příkladem budiž i dělené vlastnictví v § 357 ABGB.}\textsuperscript{,}\footfullcite[Strana 601]{gretton_g_trusts_2000}, samotná polemika o přijetí této formy však nebyla obsahem práce.\\

Třetí možností by mohlo být \textbf{přiznání vlastnických práv obmyšlenému}. V takovém případě by správce jednal jako zmocněnec obmyšleného, nebo by k majetku měl věcněprávní oprávnění a obmyšlenému by tak zbylo pouze holé vlastnictví. Je však otázkou, zda by se při takovémto pojetí zachoval insolvenční efekt. Rovněž se domnívám, že takovouto správu by šlo podřadit pod obecnou úpravu správy cizího majetku, úprava svěřenského fondu by tak byla dle mého názoru zbytečná.\\

% --------------------------------------------------------------------------
% Další možnosti jsou pak uvedeny například v této práci: https://webcache.googleusercontent.com/search?q=cache:j9oOk5qBAwYJ:https://is.cuni.cz/webapps/zzp/download/120175116+&cd=5&hl=en&ct=clnk&gl=cz#21
% --------------------------------------------------------------------------

Čtvrtý a poslední způsob, ke kterému ve svém výkladu Tichý nejvíce inklinuje, je \textbf{pohled na svěřenský fond jako na právnickou osobu}. Jak již jsem zmínil v kapitole 6.8, pro tuto možnost ji připraven poměrně solidní základ, neboť svěřenský fond vykazuje mnoho znaků právnické osoby a sám zákonodárce svěřenský fond nezřídkakdy jako právnickou osobu finguje. Toto pojetí by vyřešilo problém nepominutelného dědice, subjekt, který předpokládán vlastnictví, by byl samotný svěřenský fond a ten by tak majetek standardně vlastnil, nejednalo by se tak nadále o samostatně určený majetek.\\

Závěrem je tedy možné říci, že všechna výše uvedená řešení by problém s nepominutelným dědicem vyřešilo, neboť by byla vždy osoba s vlastnickým titulem, duální vlastnictví však není v českém právním prostředí myslitelný a přiznání vlastnických práv obmyšlenému by dle mého názoru vedlo k nadbytečnosti svěřenského fondu v českém právním řádu. Domnívám se proto, že možná řešení jsou ustanovit svěřenského správce jako vlastníka majetku \textit{"sui generis"}, nebo přiznání právní subjektivity svěřenskému fondu jak navrhuje Tichý. Otázkou je však, jaký dopad by mělo přiznání právní subjektivity na flexibilitu a jednoduchost svěřenského fondu a tím i na jeho atraktivitu.\\

S ohledem na problematiku zásadních koncepčních změn svěřenského fondu se tak jeví jako nejjednodušším a nejpraktičtějším řešením problematiky svěřenského fondu a nepominutelného dědice stanovení kolizní úpravy, která by jasně stanovovala, zda má přednost ochrana nepominutelného dědice, nebo svobodné projevení vůle zůstavitele a její následný respekt. Dalším možným řešením je pak zcela vyjmout způsob konstituování svěřenského fondu \textit{mortis causa} z právní úpravy pořízení pro případ smrti a stanovení možnosti zřízení této formy za pomocí smlouvy pro případ smrti\footnote{Viz kapitola 6.2.3 Vznik svěřenského fondu \textit{inter vivos} s odkládací podmínkou.}.\\

\subsubsection{Návrhy na změny vedoucí k vyřešení ostatních nejasností svěřenského fondu}

Je zřejmé, že výše uvedené problémy uvedené nad rozsah hypotetické otázky zdaleka nepokrývají celou další škálu potenciálních problémů vzniklých v souvislosti s implementací právní úpravy svěřenského fondu do našeho právního řádu a poskytují tak spíše pouze demonstrativní výčet dalších možných problémů, jímž by se měl, dle mého názoru, zákonodárce zabývat.\\

Jak již jsem stanovil výše, považuji samotnou konstrukci svěřenského fondu za společného jmenovatele mnoha problémů, mé doporučení v závěru předchozí kapitoly k zavedení kolizní úpravy, nebo vyjmutí způsobu zakládání svěřenského fondu \textit{mortis causa} z právní úpravy pořízení pro případ smrti tak zdaleka neřeší další širokou škálu problémů a nejasností, které se svěřenským fondem souvisí.\\

Je vhodné poznamenat, že tyto problémy\footnote{Viz kapitola 7.5 Další systémové nedostatky.} jsou často vhodně a proporcionálně řešeny ve vztahu k jiným, nejobecněji řečeno, entitám, tedy fyzickým a právnickým osobám. Takováto úprava však ne vždy poskytuje i svěřenský fond, jakožto účelově určený majetek bez právní subjektivity. Změna pojetí konstrukce svěřenského fondu ve smyslu přiznání vlastnického práva některé ze stran zainteresovaných na založeném svěřenském fondu, nebo jeho přiblížení k právnickým osobám by mnoho problému vyřešilo, ale potenciálně svěřenský fond i znevýhodnilo a zneatraktivnělo. Pokud bychom tak chtěli ponechat stávající pojetí svěřenského fondu je nezbytné, aby začal zákonodárce dané problémy objevovat a především řešit. Je však otázka, zda přemíra právních norem a kazuistická úprava je to, co český právní řád potřebuje, další možností je tak ponechat mnoho otázek na soudech a spolehnout se na soudní praxi.\\

\newpage

\section{Závěr}

% --------------------------------------------------------------------------
%Shrnutí toho, co práce obsahovala. Následovat bude shrnutí praktické části práce, tedy že v tato úprava chybí a zároveň není ustálená rozhodovací praxe soudů. Zhodnotit návrhy řešení a nastínit, zda by mohli pomoci, nebo ne.
% --------------------------------------------------------------------------

\textit{"Různé podoby svěřenství či trustů jsou běžně používány už více než dva tísíce let od římského fideikomisu pro správu majetku ve většině západních zemí. Jejich návrat do českého právního řádu v podobě svěřenského fondu není ničím náhodným nebo překvapivým, nýbrž je logickým důsledkem návratu naší země do civilizačního okruhu, ke kterému po staletí patřila. (...) svěřenství bylo ve svých dobových a místních podobách oblíbeným a velmi užitečným nástrojem správy majetku (...) a stejně tak bylo velmi často používáno pro udržení celistvosti a nezcizitelnosti zejména rodinného majetku při jeho předávání následujícím generacím."}\footfullcite[Strana 173]{bednarikova_barbora_sverenske_2014}\\

Bednaříková takto shrnula historickou hodnotu svěřenství a trustů a vyslovila jejich užitečnost. Je nepochybné, že tyto instituty mají v moderních společnostech důležitou roli a mnoho využití, přičemž jedno z nich je použití za účelem mezigeneračního převodu majetku. Tato myšlenka ostatně prostupuje celou prací. Otázka však zní, zda jejich implemetace ve formě svěřenského fondu je k mezigeneračním převodům majetku, s ohledem na součanou právní úpravu, stejně dobře využitelná.\\

Základní idea dědického práva v novém občanském zákoníku je, že je třeba co nejvíce respektovat vůli zůstavitele. Tato idea však není omezena pouze na dědické právo, ale je obecně reflektována v celém občanském zákoníku. Ke snaze o vytvoření co největší volnosti občana, ve smyslu mezigeneračního převodu majetku, nově přispívá rovněž i možnost použití svěřenského fondu k tomuto účelu. Problémem však i nadále zůstává, že tato idea mnohdy potlačuje do pozadí oprávněné zájmy nepominutelných dědiců.\\

Za základní problém u svěřenských fondů považuji právě nedostatečnou úpravu ve vztahu k dědickému právu, a ať již pohlížíme na institut svěřenského fondu jako na posílení možností pořízení o svém majetku ze strany zůstavitele, je nepochybné, že by tuto možnost měla "podpírat" solidní a jasná právní úprava. Ta žal ke dni vyhotovení této práce chybí a zapříčiňuje tak ve svém důsledku právní nejistotu občana, který by mohl mít zájem o založení svěřenského fondu a v souvislosti s tímto i obmyšlených a dědiců.\\

Samotný fakt, že na svěřenský fond lze pohlédnout z mnoha perspektiv a dle toho dospět k patřičnému závěru ve vztahu k tomu, jak bude vyřešena kolize s dědickým právem ukazuje, že právní úprava svěřenského fondu není finalizována a že je naopak třeba značných změn, aby byla právní nejistota všech zúčastněných stran odstraněna. Je tedy zřejmé, že k plnému pochopení a integrování tohoto institutu a jeho běžnějšímu využívání tak povede ještě dlouhá a dost možná trnitá cesta.\\

Není pochyb, že zavedení institutu svěřenského fondu do našeho právního řádu je jevem pozitivním a vkládá zůstaviteli do rukou možnosti, kterými lze dosáhnout mnoha rozmanitých uspořádání. Je tedy nepochybné, že OZ klade důraz nejen na posílení testovací vůle zakladatele, ale i na možnost dosáhnout uspořádání svého majetku i jinými způsoby, mezi které lze zařadit právě svěřenský fond.\\

Jeho samotné zavedení a formy založení mezi živými a mezi živými s odkládací podmínkou smrti nepředstavují žádný zásadní rozpor mezi právem dědickým a právem zůstavitele nakládat se svým majetkem. Při aplikování zásady autonomie vůle na založení svěřenského fondu pro případ smrti však dochází k zásadní kolizi mezi vůlí zakladatele o uspořádání vlastního majetku a rigidní úpravou dědického práva chránící zájmy nepominutelných dědiců. Respektování vůle zůstavitele je velice důležité, neboť její její podcenění by mohlo vést k popření základní funkce občanského práva a nového občanského zákoníku, kterou zákonodárce popsal v důvodové zprávě. Tato funkce umožňuje svobodné rozvíjení soukromého života a soukromoprávních vztahů a zahrnuje respektování osobnosti člověka jako svobodného individua způsobilého žít podle svého a rozhodovat o svých soukromých věcech samostatně dle vlastní svobodné vůle. Rozhodování o svých soukromých věcech zahrnuje i určení toho, co se stane se jměním osoby po její smrti. Je však nepochybné, že v případě aplikování ustanovení dědického práva na toto rozhodování je zůstavitel povinen tyto ustanovení respektovat.\\

Zde vyvstává základní problém úpravy svěřenského fondu založeného pro případ smrti, kdy není ze zákona jasné, jakým způsobem je zajištěno zachování práva nepominutelného dědice. K tomuto přispívá i neexistence soudních rozhodnutí, které by se tomuto tématu věnovali. Samotná právní obec je v pohledu na tento problém rozdělena a lze se setkat s různými pohledy na problematiku potenciálního obcházení práva nepominutelných dědiců pomocí svěřenského fondu.\\

S ohledem na momentální stav právní úpravy a názorů jednotlivých autorů, z jejichž prací jsem v průběhu své bakalářské práce čerpal nelze se stoprocentní jistotou zodpovědět moje hypotéza stanovená v úvodu práce a dochází tak k právní nejistotě u zakladatelů svěřenských fondů zda mohou použít svěřenský fond k faktickému vydědění jejich nepominutelných dědiců a obdobně dochází k nejistotě na straně dědiců, kterým neposkytuje momentální právní úprava jasné stanovisko, zda by se měli o své právo na povinný díl brát, nebo nikoli. Lze tak pouze doporučit, aby se zákonodárce na tento problém co nejdříve zaměřil a provedl patřičné novelizace, které tyto otázky zodpoví a vyřeší problémy, ke kterým neodvratně dojde, pokud tak neučiní.\\

Jak již jsem nastínil v poslední kapitole Systémové nedostatky a návrhy na řešení, problémy spojené s přijetí právní úpravy trustu do našeho Českého právního řádu ve formě svěřenského fondu přineslo krom nejasností ve vztahu k dědickému právu mnoho dalších problémů a krom toho tedy zůstává otevřeno mnoho otázek, jejichž řešení může pomocí novelizace momentální nedostačující právní úpravy zodpovědět pouze zákonodárce.\\

Lze tedy uzavřít s tím, že aplikace právní úpravy trustu byla v našem právním řádu a právním prostředí potřebná, nelze se však ztotožnit a spokojit s momentální právní úpravou, která není sto schopná unést otázky na ní padající ve vztahu k tomuto nově implementovanému institutu.\\

Svěřenský fond jako takový tedy přináší příjemné osvěžení k možnostem převodu rodinného majetku, přináší ale zároveň spoustu potenciálních problému, na které momentální právní úprava nenabízí řešení. Pokud by tedy otázka stála ve smyslu: "Je vhodné užití svěřenského fondu na mezigenerační převody majetku?", tak bych odpověděl že s ohledem na momentální právní úpravu pouze ve formě založení mezi živými, k práci je proto přiložen hrubý návrh smlouvy s odkládací podmínkou smrti, u které se domnívám, že může do značné míry nahradit způsob zřízení \textit{mortis causa}, avšak bez problému ochrany nepominutelných dědiců\footnote{Ten je však pouze vzorem se snahou postihnout po obsahové stránce základní požadavky na něj kladené, úmysl nebyl koncipovat ho jako jasný návod na to, jak má daná smlouva vypadat, ostatně vzhledem k šíři možných účelů svěřenského fondu toto není ani dobře možné.}. Založení pro případ smrti s sebou nese tolik potenciálních problémů, které budou muset zůstavitelovi dědici, obmyšlení a svěřenský správce řešit, že dle mého názoru není jeho použití vhodné a tímto bych ho také ani, mimo případy, kdy problematika nepominutelných dědiců nepředstavuje potenciální problém, tedy například při absenci nepominutelných dědiců, nedoporučil\footnote{Diskuze o problematice zřízení svěřenského fondu \textit{mortis causa} a jeho kolizi s úpravou dědického práva je poměrně arbitrární v případě, že nepominutelných dědiců není. V takovém případě, i přes existenci dalších potenciálních problémů, představuje i tato forma zřízení svěřenského fondu vhodný způsob vyjádření vůle zakladatele a způsob správy nebo transferu majetku.}.\\

Přes všechna pozitiva, kterých je mnoho, tak lze pouze doufat, že problémy spojené s právní úpravou svěřenských fondů budou co nejdříve odstraněny, aby česká společnost mohla tento velice užitečný institut začít požívat plnými doušky a se zajištěnou právní jistotou. Doufejme tedy, že s dopřáním více času se zákonodárci povede svěřenský fond opravit a dostat na výsluní tak, aby byl atraktivní a použitelný i pro běžné lidi. Pak se možná zvýší i počet založených svěřenských fondů, kterému je ke dnešnímu dni jen něco okolo dvou tisíc\footfullcite[Zde lze zjistit momentální číslo, evidence svěřenských fondů]{noauthor_isesf_nodate}.\\

% --------------------------------------------------------------------------
%Jak je to v případě, že zůstavitel napíše dědickou smlouvu a majetek, který je součástí takovézo smlouvy pak za svého života zcizí? Je to v pořádku, nebo je to zakázané, musím se podívat do úpravy dědické smlouvy. V případě neshod mezi dědickou smlouvou a závětí má přednost dědická smlouva, takže tady o tomto není pochyb. Co však kdyby proti sobě stálo několik dědických nároků, například pokud by zůstavitel pořídil o větším majetku, než který by poté byl souřástí dědického řízení? Co když pořídí více závětí, je pak platná ta naposledny pořízená závěť?

%Podívat se ještě na google na tokové ty články týkající se svěřenského fondu, často ještě pojednávali o dalších problémech, které bych tady v závěru mohl zmínit a podpořit jím tak stanovisko, že použití svěřenského fondu pro případ smrti není vhodné a rozhodně bych ho nedoporučoval.
% --------------------------------------------------------------------------

% --------------------------------------------------------------------------
% Konec hlavní části práce, následují zdroje
% --------------------------------------------------------------------------
\newpage
{\pagestyle{Contents}
\section*{Zdroje}\markright{ZDROJE}

% --------------------------------------------------------------------------
% Counter definition pro přidání čísla ke zdrojům v obsahu práce, možná číslo odeberu
% --------------------------------------------------------------------------

\newcounter{SecZdroje}
\setcounter{SecZdroje}{\thesection}
\addtocounter{SecZdroje}{1}
\addcontentsline{toc}{section}{\theSecZdroje \hspace{1,7 mm} Zdroje}

% --------------------------------------------------------------------------
% Printi bibliografie, která není přímo odcitována v textu
% --------------------------------------------------------------------------

\nocite{lepaulle_p_outsiders_1928}

% --------------------------------------------------------------------------
% Zde je definováno jak budou vypsány zdroje, přesná specifikace je obsažena v mappingu
% --------------------------------------------------------------------------

\printbibliography[omitnumbers=true,type=misc,heading=subbibliography,title={Online zdroje}]
\printbibliography[omitnumbers=true,type=book,heading=subbibliography,title={Knižní zdroje}]
\printbibliography[omitnumbers=true,type=article,heading=subbibliography,title={Články}]
\printbibliography[omitnumbers=true,type=proceedings,heading=subbibliography,title={Zákony}]
\printbibliography[omitnumbers=true,type=thesis,heading=subbibliography,title={Kvalifikační práce}]	
\clearpage
}

\newpage

\pagenumbering{Roman}
\addtocounter{section}{1}
\section{Seznam příloh}

\textbf{\normalsize{1\hspace*{1.8em}Vzor smlouvy o založení svěřenského fondu}} . . . . . . . . . . . . . . . . . . . . . . . . . . . . . . . \rom{2}\\
\noindent\textbf{\normalsize{2\hspace*{1.8em}Vzor listiny přání zakladatele o určení obmyšlených svěřenského fondu}} . . . . . . . . . . . . \rom{9}

% --------------------------------------------------------------------------
%\includepdf[pages=-]{Bakalarska_prace_navrh_smlouvy.pdf}
% --------------------------------------------------------------------------

\includepdf[scale=1,pages=-]{\detokenize{Bakalářská_práce_návrh_smlouvy}.pdf}

% --------------------------------------------------------------------------
%\includepdf[scale=0.8,pages=2-,pagecommand={\pagestyle{fancy}}]{\detokenize{Bakalářská_práce_návrh_smlouvy}.pdf}
% --------------------------------------------------------------------------

% --------------------------------------------------------------------------
% Vložení zadání bakalářské práce
% --------------------------------------------------------------------------

\newpage

\includepdf[scale=1,pages=-]{\detokenize{Zadání_závěrečné_práce}.pdf}

\end{document}
